%% Generated by Sphinx.
\def\sphinxdocclass{jsbook}
\documentclass[letterpaper,10pt,dvipdfmx]{sphinxmanual}
\ifdefined\pdfpxdimen
   \let\sphinxpxdimen\pdfpxdimen\else\newdimen\sphinxpxdimen
\fi \sphinxpxdimen=.75bp\relax
\ifdefined\pdfimageresolution
    \pdfimageresolution= \numexpr \dimexpr1in\relax/\sphinxpxdimen\relax
\fi
%% let collapsible pdf bookmarks panel have high depth per default
\PassOptionsToPackage{bookmarksdepth=5}{hyperref}


\PassOptionsToPackage{warn}{textcomp}


\usepackage{cmap}
\usepackage[T1]{fontenc}
\usepackage{amsmath,amssymb,amstext}




\usepackage{tgtermes}
\usepackage{tgheros}
\renewcommand{\ttdefault}{txtt}




\usepackage[,numfigreset=1,mathnumfig]{sphinx}

\fvset{fontsize=auto}
\usepackage[dvipdfm]{geometry}


% Include hyperref last.
\usepackage{hyperref}
% Fix anchor placement for figures with captions.
\usepackage{hypcap}% it must be loaded after hyperref.
% Set up styles of URL: it should be placed after hyperref.
\urlstyle{same}
\usepackage{pxjahyper}


\usepackage{sphinxmessages}
\setcounter{tocdepth}{2}


\makeindex
\makeatletter
%% 「\spxentry キーワード … ページ」をそのまま「\item キーワード … ページ」に置き換える
\renewcommand{\spxentry}{\item}
\renewcommand\sphinxlineitem[2]{%
  %% safe test of whether #2 is \sphinxlineitem
  \sphinx@gobto@sphinxlineitem#2\@gobbletwo\sphinxlineitem\unless
  \iftrue
    %% Accumulate successive terms until actual definition or sub-list is found
    \spx@lineitemlabel\expandafter{\the\spx@lineitemlabel\strut#1\\}%
  \else
    %% Issue the \item command with possibly multi-line contents
    \item[\kern\labelwidth\kern-\itemindent\kern-\leftmargin
          {\parbox[t]{1.4\linewidth}{%% <- ここで幅を調整
          \raggedright
          \the\spx@lineitemlabel%% Accumulated terms before this one, CR separated
          \strut#1}}%% No \par token allowed here, but the \parbox will insert one tacitly at end
          \kern-\labelsep]%
    \spx@lineitemlabel{}%
    %% This causes the label to be typeset (filling up the line), clearing up
    %% things in case a nested list follows.
    \leavevmode
  \fi #2%
}
\makeatother


\title{BlackPoker}
\date{2025年08月21日}
\release{2025/6/30}
\author{もりくま}
\newcommand{\sphinxlogo}{\vbox{}}
\renewcommand{\releasename}{リリース}
\makeindex
\begin{document}

\pagestyle{empty}

\begin{titlepage}
\begin{center}
\vspace*{25mm}

% ロゴを PDF で中央配置
\includegraphics[width=4cm]{blackpoker_logo.pdf}  % ファイルパスは latex_additional_files に合わせる

\vspace{20mm}
{\Huge \textbf{BlackPoker 公式ルール}}\\[10mm]
{\Large トランプだけでトレーディングカードゲームみたいに遊ぶ方法}\\[30mm]

% バージョン・リリース情報を埋め込む
{\huge 第8.1版}\\[20mm]

{\Large もりくま 著} \\[5mm]
{\Large 2025/6/30 発行}\\[10mm]

{\small © 2013 BlackPoker}

\end{center}
\end{titlepage}

% 表紙の次に空白ページを挟む
\clearpage
\thispagestyle{empty}  % ページ番号などを非表示にする
\phantomsection
\null                  % 何もないページを出力
\clearpage             % さらに改ページして本文へ

\pagestyle{plain}
\sphinxtableofcontents
\pagestyle{normal}
\phantomsection\label{\detokenize{index::doc}}


\sphinxAtStartPar
release: 2025/6/30

\sphinxstepscope


\chapter{はじめに}
\label{\detokenize{init/init:init-rst}}\label{\detokenize{init/init:id1}}\label{\detokenize{init/init::doc}}
\sphinxAtStartPar
この文書は、トランプゲーム「BlackPoker」の公式ルールをまとめたものです。

\sphinxAtStartPar
BlackPokerは、誰でも持っているトランプで遊べる、
「トレーディングカードゲームのようなプレイ感」を目指した完全オリジナルのゲームです。

\sphinxAtStartPar
ルール文書には詳細な内容が含まれており、最初はその量に圧倒されるかもしれません。
ゲームを始める際にすべてを暗記する必要はありませんが、ルールの仕組みや拡張性を深く知りたい方には
ぜひ一度じっくり読み込んでいただくことをおすすめします。


\section{BlackPokerとは}
\label{\detokenize{init/init:blackpoker}}
\sphinxAtStartPar
当サークルが考案した、1人1セットのトランプを使って遊ぶ、
ターン制かつ割込み可能な独自ルールのカードゲームです。
自分だけのオリジナルトランプを使って、友達と本格的なカードバトルを楽しむことができます。

\sphinxAtStartPar
プレイしているときのイメージは次のような感じです。(\hyperref[\detokenize{init/init:play-image}]{Fig.\@ \ref{\detokenize{init/init:play-image}}})

\begin{figure}[htbp]
\centering
\capstart

\noindent\sphinxincludegraphics{{play-image}.pdf}
\caption{プレイ風景}\label{\detokenize{init/init:id7}}\label{\detokenize{init/init:play-image}}\end{figure}


\section{きっかけ}
\label{\detokenize{init/init:id2}}
\begin{sphinxVerbatim}[commandchars=\\\{\}]
「また昔みたいに友達とカードゲームがしたい」
\end{sphinxVerbatim}

\sphinxAtStartPar
けど、仕事に追われ時間もなく、お金もかけられないのでカード資産も抱えられない。
昔のデッキを引っ張り出してもカードパワーが違って平等に遊べないし、コミュニティも狭い。
そんな悩みを解消しようと、誰でも持っているトランプでカードゲーム風のルールを作ろうと思いました。


\section{ゲームのストーリー}
\label{\detokenize{init/init:id3}}
\sphinxAtStartPar
あなたのトランプは、ただのトランプではありません。
かつて大地を支配していた偉大な国々や、強固な絆で結ばれた組織の魂が宿る遺産です。

\sphinxAtStartPar
これらのカードには、世界の歴史を彩る多彩な物語が織り込まれており、
その物語たちはまるで時間を超えたタペストリーのように絡み合っています。

\sphinxAtStartPar
プレイヤーは歴史の再現者、リアニメーターとして、
デッキに宿る過去の栄光を再現し、戦いを通じてその魅力を相手に示してください。

\sphinxAtStartPar
え？そんな国知らない？でしたら、あなたのカードから創造してみてはいかがでしょうか。


\section{読み方}
\label{\detokenize{init/init:id4}}
\sphinxAtStartPar
各章は次のような内容を説明しています。
\begin{description}
\sphinxlineitem{\hyperref[\detokenize{init/init:init-rst}]{\ref{\detokenize{init/init:init-rst}} \nameref{\detokenize{init/init:init-rst}}}}
\sphinxAtStartPar
ルールの背景や構成、全体像を説明（この章）

\sphinxlineitem{\hyperref[\detokenize{match-regulations/match-regulations:match-regulations-rst}]{\ref{\detokenize{match-regulations/match-regulations:match-regulations-rst}} \nameref{\detokenize{match-regulations/match-regulations:match-regulations-rst}}}}
\sphinxAtStartPar
対戦前にプレイヤー間で取り決める「遊び方のルール」（対戦レギュレーション）の選び方を説明

\sphinxlineitem{\hyperref[\detokenize{common/common:common-rst}]{\ref{\detokenize{common/common:common-rst}} \nameref{\detokenize{common/common:common-rst}}}}
\sphinxAtStartPar
BlackPokerの全ルールで共通して使う基本ルールを説明

\sphinxlineitem{\hyperref[\detokenize{format/format:format-rst}]{\ref{\detokenize{format/format:format-rst}} \nameref{\detokenize{format/format:format-rst}}}}
\sphinxAtStartPar
難易度に応じたルール制限（フォーマット）の詳細を説明

\sphinxlineitem{\hyperref[\detokenize{frame/frame:frame-rst}]{\ref{\detokenize{frame/frame:frame-rst}} \nameref{\detokenize{frame/frame:frame-rst}}}}
\sphinxAtStartPar
デッキ構成や初期配置など遊び方のパターン（フレーム）を説明

\sphinxlineitem{\hyperref[\detokenize{core/core:core-rst}]{\ref{\detokenize{core/core:core-rst}} \nameref{\detokenize{core/core:core-rst}}}}
\sphinxAtStartPar
アクションの処理や割込み、ステージの仕組みなどのコアシステムを説明

\end{description}


\subsection{用途別読み方}
\label{\detokenize{init/init:id5}}
\sphinxAtStartPar
ルール文書はやや複雑な構成になっているため、用途に合わせて読むことをおすすめします。
\begin{description}
\sphinxlineitem{ゲームを始めたい方（初心者）}
\sphinxAtStartPar
\hyperref[\detokenize{match-regulations/match-regulations:match-regulations-rst}]{\ref{\detokenize{match-regulations/match-regulations:match-regulations-rst}} \nameref{\detokenize{match-regulations/match-regulations:match-regulations-rst}}} \(\rightarrow\) \hyperref[\detokenize{common/common:common-rst}]{\ref{\detokenize{common/common:common-rst}} \nameref{\detokenize{common/common:common-rst}}} の順に読むのがおすすめです。

\sphinxAtStartPar
\hyperref[\detokenize{core/core:core-rst}]{\ref{\detokenize{core/core:core-rst}} \nameref{\detokenize{core/core:core-rst}}} （高度なルール）は飛ばしても構いません。

\sphinxlineitem{ゲーム中の処理で悩んでいる方}
\sphinxAtStartPar
アクションの起こし方や割込み、解決順などを \hyperref[\detokenize{core/core:core-rst}]{\ref{\detokenize{core/core:core-rst}} \nameref{\detokenize{core/core:core-rst}}} で確認してください。

\end{description}


\section{ルール指針}
\label{\detokenize{init/init:id6}}
\sphinxAtStartPar
BlackPokerのルールは、\sphinxstylestrong{できるだけシンプルに、わかりやすく} を基本方針としています。
そのうえで、より多くの人が楽しめるよう、以下の指針に基づいてルールを設計・調整しています。
\begin{description}
\sphinxlineitem{\sphinxstylestrong{誰とでも戦える \textasciitilde{}目指すは老若男女\textasciitilde{}}}
\sphinxAtStartPar
ルールを知りトランプを持っていれば誰とでも遊べるゲームを目指します。

\sphinxlineitem{\sphinxstylestrong{個性が出せる \textasciitilde{}オリジナルトランプ・デッキ構築\textasciitilde{}}}
\sphinxAtStartPar
さまざまなトランプが使え見た目で個性を出せるのはもちろんのこと、
デッキ構築の面でも自分のしたい戦い方が表現できることを目指します。

\sphinxlineitem{\sphinxstylestrong{短く終わる \textasciitilde{}1戦15分\textasciitilde{}}}
\sphinxAtStartPar
時間をかけずさっと遊べることを目指します。

\sphinxlineitem{\sphinxstylestrong{ずっと使えるデッキ}}
\sphinxAtStartPar
愛着のあるカードがずっと使えるようなルールとします。

\sphinxlineitem{\sphinxstylestrong{必要な物は最小限 \textasciitilde{}トランプのみ\textasciitilde{}}}
\sphinxAtStartPar
用意するものはトランプのみ。それ以外の道具は必要ないルールとします。

\sphinxlineitem{\sphinxstylestrong{プレイング重視 \textasciitilde{}5：3：2＝技：運：構築\textasciitilde{}}}
\sphinxAtStartPar
運やデッキ構築より技量を重視したルールを目指します。

\sphinxlineitem{\sphinxstylestrong{ベースルールはトレーディングカードゲーム}}
\sphinxAtStartPar
カードゲームプレイヤーが覚えやすいルールを目指します。

\sphinxlineitem{\sphinxstylestrong{カスタマイズ可能 \textasciitilde{}基本と拡張の分離\textasciitilde{}}}
\sphinxAtStartPar
基本ルールと拡張ルールを分離し、大富豪のようにローカルルールが作成できることを目指します。

\sphinxlineitem{\sphinxstylestrong{ルールの更新 \textasciitilde{}飽き防止＆不備改善\textasciitilde{}}}
\sphinxAtStartPar
新たなルールを度々公開し、飽きを防止します。またルールに不備がある場合、随時改善します。

\sphinxlineitem{\sphinxstylestrong{相手のカードに触らない}}
\sphinxAtStartPar
盗難防止とネット対戦対応に努めます。

\end{description}

\sphinxstepscope


\chapter{対戦レギュレーション}
\label{\detokenize{match-regulations/match-regulations:match-regulations-rst}}\label{\detokenize{match-regulations/match-regulations:id1}}\label{\detokenize{match-regulations/match-regulations::doc}}
\index{た\textbar{}対戦レギュレーション@\spxentry{た\textbar{}対戦レギュレーション}}\ignorespaces 

\section{対戦レギュレーションとは}
\label{\detokenize{match-regulations/match-regulations:index-0}}\label{\detokenize{match-regulations/match-regulations:id2}}
\sphinxAtStartPar
BlackPoker では、ゲームを始める前に「どのようなルールで遊ぶか」をプレイヤー同士で決める必要があります。

\sphinxAtStartPar
この事前に決めるルールの組み合わせを \sphinxstylestrong{対戦レギュレーション} と呼びます。

\sphinxAtStartPar
BlackPoker は、標準的なトランプだけを使って遊ぶゲームです。
カードは同じでも、ルール（対戦レギュレーション）を変えることで、
初心者から上級者まで、さまざまなプレイヤーが自分に合った難易度・遊び方で楽しめるようになっています。

\sphinxAtStartPar
この章では、対戦レギュレーションの構成や決め方について説明します。


\section{定義項目}
\label{\detokenize{match-regulations/match-regulations:id3}}
\sphinxAtStartPar
対戦レギュレーションは、次の2つの要素を組み合わせて構成されます。

\index{フォーマット(対戦レギュレーション)@\spxentry{フォーマット(対戦レギュレーション)}}\ignorespaces \begin{description}
\sphinxlineitem{フォーマット}
\sphinxAtStartPar
\sphinxstylestrong{プレイヤーの熟練度に応じて、使える行動（アクション）や登場キャラクターなどを制限した難易度の設定です。}

\sphinxAtStartPar
たとえば「ライト」は初心者向けに覚える行動を最小限に絞ったフォーマット、
「スタンダード」は基本的な要素を一通り含んだ標準的なフォーマットです。

\sphinxAtStartPar
カードゲームでいう「カードプール」（どのカードを使ってよいか）に近い考え方です。

\sphinxAtStartPar
詳しくは \hyperref[\detokenize{format/format:format-rst}]{\ref{\detokenize{format/format:format-rst}} \nameref{\detokenize{format/format:format-rst}}} を参照してください。

\end{description}

\index{フレーム(対戦レギュレーション)@\spxentry{フレーム(対戦レギュレーション)}}\ignorespaces \begin{description}
\sphinxlineitem{フレーム}
\sphinxAtStartPar
\sphinxstylestrong{デッキの構成やカードの初期配置、場の使い方などを定めた、遊び方のパターンです。}

\sphinxAtStartPar
たとえば「エントリー20」は、決められた20枚のカードを使って手軽に遊べる入門向けフレームです。
一方「パック」や「ストラテジー」では、構築や初期戦略の自由度が高く、プレイヤーの個性がより強く表れます。

\sphinxAtStartPar
フォーマットは「どのくらい覚えるか」で選び、
フレームは「どのように遊びたいか」で選ぶイメージです。

\sphinxAtStartPar
詳しくは \hyperref[\detokenize{frame/frame:frame-rst}]{\ref{\detokenize{frame/frame:frame-rst}} \nameref{\detokenize{frame/frame:frame-rst}}} を参照してください。

\end{description}


\subsection{対戦レギュレーションの表記}
\label{\detokenize{match-regulations/match-regulations:id4}}
\sphinxAtStartPar
対戦レギュレーションは、 \sphinxstylestrong{フォーマット + フレーム} の形式で表記します。

\begin{sphinxVerbatim}[commandchars=\\\{\}]
\PYG{n}{フォーマット}\PYG{o}{+}\PYG{n}{フレーム}
\end{sphinxVerbatim}

\sphinxAtStartPar
たとえば、次のように記述されます。

\begin{sphinxVerbatim}[commandchars=\\\{\}]
\PYG{n}{ライト}\PYG{o}{+}\PYG{n}{エントリー}
\end{sphinxVerbatim}

\sphinxAtStartPar
これは「ライトフォーマットで、エントリーフレームを使用して対戦する」ことを意味します。


\section{対戦レギュレーションの決め方}
\label{\detokenize{match-regulations/match-regulations:id5}}
\sphinxAtStartPar
以下の手順で、プレイヤー同士が合意して対戦レギュレーションを決めてください。

\sphinxAtStartPar
一部の組み合わせは公式に対応していないため、必ず以下の手順に従って選択してください。
\begin{enumerate}
\sphinxsetlistlabels{\arabic}{enumi}{enumii}{}{.}%
\item {} 
\sphinxAtStartPar
\sphinxstylestrong{フォーマットを決める}

\sphinxAtStartPar
「ライト」「スタンダード」「プロ」などのフォーマットから選びます。

\sphinxAtStartPar
フォーマットが変わると、使用可能なアクションの種類や難易度が大きく変わります。
プレイヤーの経験や慣れ具合をもとに、無理のないフォーマットを選びましょう。

\item {} 
\sphinxAtStartPar
\sphinxstylestrong{フレームを決める}

\sphinxAtStartPar
フォーマットを決めたら、次に使用するフレームを選びます。

\sphinxAtStartPar
各フレームが対応しているフォーマットの一覧は以下の表をご覧ください。
「◯」と表記されている組み合わせが選択可能です。


\begin{savenotes}\sphinxattablestart
\sphinxthistablewithglobalstyle
\centering
\sphinxcapstartof{table}
\sphinxthecaptionisattop
\sphinxcaption{フレーム対応一覧}\label{\detokenize{match-regulations/match-regulations:id6}}\label{\detokenize{match-regulations/match-regulations:frame-format}}
\sphinxaftertopcaption
\begin{tabulary}{\linewidth}[t]{|T|T|T|T|T|}
\sphinxtoprule
\sphinxtableatstartofbodyhook
\sphinxAtStartPar
【フレーム】
&
\sphinxAtStartPar
ライト
&
\sphinxAtStartPar
スタンダード
&
\sphinxAtStartPar
プロ
&
\sphinxAtStartPar
マスター
\\
\sphinxhline
\sphinxAtStartPar
エントリー20
&
\sphinxAtStartPar
◯
&
\sphinxAtStartPar
x
&
\sphinxAtStartPar
x
&
\sphinxAtStartPar
x
\\
\sphinxhline
\sphinxAtStartPar
パック
&
\sphinxAtStartPar
x
&
\sphinxAtStartPar
◯
&
\sphinxAtStartPar
◯
&
\sphinxAtStartPar
◯
\\
\sphinxhline
\sphinxAtStartPar
レアパック
&
\sphinxAtStartPar
x
&
\sphinxAtStartPar
◯
&
\sphinxAtStartPar
◯
&
\sphinxAtStartPar
◯
\\
\sphinxhline
\sphinxAtStartPar
ストラテジー
&
\sphinxAtStartPar
x
&
\sphinxAtStartPar
x
&
\sphinxAtStartPar
◯
&
\sphinxAtStartPar
◯
\\
\sphinxhline
\sphinxAtStartPar
エクストラ
&
\sphinxAtStartPar
x
&
\sphinxAtStartPar
x
&
\sphinxAtStartPar
x
&
\sphinxAtStartPar
◯
\\
\sphinxbottomrule
\end{tabulary}
\sphinxtableafterendhook\par
\sphinxattableend\end{savenotes}

\end{enumerate}

\sphinxAtStartPar
\sphinxstylestrong{補足：} 初めての方には「ライト+エントリー20」がおすすめです。
決められた20枚のカードを使って、最も手軽にBlackPokerを体験できます。


\bigskip\hrule\bigskip


\sphinxAtStartPar
このあと、　\sphinxstylestrong{フォーマットやフレームに関係なく適用される「共通ルール」} を説明します。
まずは次の \hyperref[\detokenize{common/common:common-rst}]{\ref{\detokenize{common/common:common-rst}} \nameref{\detokenize{common/common:common-rst}}} を読み進めてください。

\sphinxAtStartPar
フォーマットやフレームの詳細を知りたい方は、 \hyperref[\detokenize{common/common:common-rst}]{\ref{\detokenize{common/common:common-rst}} \nameref{\detokenize{common/common:common-rst}}} を読んだあとに
\hyperref[\detokenize{format/format:format-rst}]{\ref{\detokenize{format/format:format-rst}} \nameref{\detokenize{format/format:format-rst}}} および \hyperref[\detokenize{frame/frame:frame-rst}]{\ref{\detokenize{frame/frame:frame-rst}} \nameref{\detokenize{frame/frame:frame-rst}}} を参照するのがおすすめです。

\sphinxstepscope


\chapter{共通ルール}
\label{\detokenize{common/common:common-rst}}\label{\detokenize{common/common:id1}}\label{\detokenize{common/common::doc}}
\sphinxAtStartPar
この章では、BlackPokerのすべての対戦で共通して使用される共通ルールを説明します。
フォーマットやフレームに関係なく、BlackPokerを遊ぶうえでの土台となるルールです。

\sphinxAtStartPar
プレイ人数、カードの種類、場の構成、アクションの処理など、ゲームの進行に必要な共通事項をまとめています。

\sphinxAtStartPar
※ フォーマットやフレームにより細かな違いがある場合は、フォーマットやフレーム側で補足されます。


\section{プレイ人数}
\label{\detokenize{common/common:id2}}
\sphinxAtStartPar
2人です。


\section{用意するもの}
\label{\detokenize{common/common:id3}}\begin{itemize}
\item {} 
\sphinxAtStartPar
1人1セットのトランプが必要です。

\item {} 
\sphinxAtStartPar
覚えていない場合、対戦レギュレーションに応じてアクションリストやキャラクターリストなどがあると便利です。

\end{itemize}


\section{使用できるトランプ}
\label{\detokenize{common/common:id4}}
\sphinxAtStartPar
BlackPokerでは次の条件を満たしたトランプを使うことができます。
一般的なトランプなら満たす条件となっています。
\begin{itemize}
\item {} 
\sphinxAtStartPar
スートと数字が分かる

\item {} 
\sphinxAtStartPar
スートの{\normalsize $\spadesuit$} {\normalsize $\heartsuit$} {\normalsize $\diamondsuit$} {\normalsize $\clubsuit$} が判断できる

\item {} 
\sphinxAtStartPar
数字のA\sphinxhyphen{}K（1\sphinxhyphen{}13）が判断できる

\item {} 
\sphinxAtStartPar
スートと数字の組合せが重複していない

\item {} 
\sphinxAtStartPar
裏から表がわからない

\item {} 
\sphinxAtStartPar
縦向き横向きが判断できる

\item {} 
\sphinxAtStartPar
Jokerは2枚まで入れられる

\item {} 
\sphinxAtStartPar
54枚無くてもよい

\end{itemize}

\sphinxAtStartPar
対戦レギュレーションにより使用できるトランプの枚数など異なる場合があるため、
対戦する際に、対戦レギュレーションを確認してください。


\section{トランプの数字}
\label{\detokenize{common/common:id5}}
\sphinxAtStartPar
ゲーム全体を通してトランプの数字は次のような数値として扱います。(\hyperref[\detokenize{common/common:cardrank}]{Table \ref{\detokenize{common/common:cardrank}}})


\begin{savenotes}\sphinxattablestart
\sphinxthistablewithglobalstyle
\centering
\sphinxcapstartof{table}
\sphinxthecaptionisattop
\sphinxcaption{トランプの数字}\label{\detokenize{common/common:id43}}\label{\detokenize{common/common:cardrank}}
\sphinxaftertopcaption
\begin{tabulary}{\linewidth}[t]{|T|T|}
\sphinxtoprule
\sphinxstyletheadfamily 
\sphinxAtStartPar
カード
&\sphinxstyletheadfamily 
\sphinxAtStartPar
数字
\\
\sphinxmidrule
\sphinxtableatstartofbodyhook
\sphinxAtStartPar
A
&
\sphinxAtStartPar
1
\\
\sphinxhline
\sphinxAtStartPar
2〜10
&
\sphinxAtStartPar
表記どおり
\\
\sphinxhline
\sphinxAtStartPar
J
&
\sphinxAtStartPar
11
\\
\sphinxhline
\sphinxAtStartPar
Q
&
\sphinxAtStartPar
12
\\
\sphinxhline
\sphinxAtStartPar
K
&
\sphinxAtStartPar
13
\\
\sphinxhline
\sphinxAtStartPar
Joker
&
\sphinxAtStartPar
0
\\
\sphinxbottomrule
\end{tabulary}
\sphinxtableafterendhook\par
\sphinxattableend\end{savenotes}

\index{り\textbar{}領域@\spxentry{り\textbar{}領域}}\index{ゾーン@\spxentry{ゾーン}}\ignorespaces 

\section{領域（ゾーン）}
\label{\detokenize{common/common:zone-section}}\label{\detokenize{common/common:index-0}}\label{\detokenize{common/common:id6}}
\sphinxAtStartPar
BlackPokerには、「領域（ゾーン）」と呼ぶ、カードを置く場所がいくつかあります。
各領域には配置可能なカード種類（コンポーネント）が定義されています。

\index{コンポーネントキーカード@\spxentry{コンポーネントキーカード}}\index{ユニットカード@\spxentry{ユニットカード}}\index{ユニット@\spxentry{ユニット}}\ignorespaces 
\sphinxAtStartPar
ゲーム中に配置されたカードのまとまりを「ユニット」と呼びます。
ユニットを構成する各カードを「ユニットカード」と呼びます。
また、そのユニットがどのコンポーネント定義に対応するかを示すカードを
「コンポーネントキーカード」と呼びます。

\sphinxAtStartPar
領域ごとに存在するコンポーネントの一覧です。(\hyperref[\detokenize{common/common:zone-component}]{Table \ref{\detokenize{common/common:zone-component}}})


\begin{savenotes}\sphinxattablestart
\sphinxthistablewithglobalstyle
\centering
\sphinxcapstartof{table}
\sphinxthecaptionisattop
\sphinxcaption{領域ごとに存在するコンポーネント一覧}\label{\detokenize{common/common:id44}}\label{\detokenize{common/common:zone-component}}
\sphinxaftertopcaption
\begin{tabular}[t]{|\X{1}{3}|\X{1}{3}|\X{1}{3}|}
\sphinxtoprule
\sphinxstyletheadfamily 
\sphinxAtStartPar
領域名
&\sphinxstyletheadfamily 
\sphinxAtStartPar
存在するコンポーネント
&\sphinxstyletheadfamily 
\sphinxAtStartPar
コンポーネントキーカードの名称
\\
\sphinxmidrule
\sphinxtableatstartofbodyhook
\sphinxAtStartPar
場
&
\sphinxAtStartPar
キャラクター
&
\sphinxAtStartPar
キャラクターキーカード
\\
\sphinxhline
\sphinxAtStartPar
ライフ
&
\sphinxAtStartPar
ライフ
&
\sphinxAtStartPar
\sphinxhyphen{}
\\
\sphinxhline
\sphinxAtStartPar
墓地
&
\sphinxAtStartPar
墓地
&
\sphinxAtStartPar
\sphinxhyphen{}
\\
\sphinxhline
\sphinxAtStartPar
フォグ
&
\sphinxAtStartPar
フォグ
&
\sphinxAtStartPar
フォグキーカード
\\
\sphinxhline
\sphinxAtStartPar
手札
&
\sphinxAtStartPar
手札
&
\sphinxAtStartPar
\sphinxhyphen{}
\\
\sphinxhline
\sphinxAtStartPar
パック
&
\sphinxAtStartPar
パック
&
\sphinxAtStartPar
\sphinxhyphen{}
\\
\sphinxhline
\sphinxAtStartPar
レアカード
&
\sphinxAtStartPar
レアカード
&
\sphinxAtStartPar
\sphinxhyphen{}
\\
\sphinxhline
\sphinxAtStartPar
切札
&
\sphinxAtStartPar
切札
&
\sphinxAtStartPar
切札キーカード
\\
\sphinxbottomrule
\end{tabular}
\sphinxtableafterendhook\par
\sphinxattableend\end{savenotes}

\sphinxAtStartPar
※ 表中では「ライフ」「墓地」「手札」などと記載しますが、「ライフ領域」「墓地領域」「手札領域」と呼んでも構いません。

\sphinxAtStartPar
各コンポーネントの詳細な定義や項目については、\hyperref[\detokenize{component/component:component-rst}]{\ref{\detokenize{component/component:component-rst}} \nameref{\detokenize{component/component:component-rst}}} で詳しく説明します。

\sphinxAtStartPar
ここまでで「どの領域に」「どんなコンポーネントが置かれるか」を整理しました。
次節では、実際のプレイ中にそれらがどのように配置されるかを図示して解説します。


\section{カードの配置}
\label{\detokenize{common/common:id7}}
\sphinxAtStartPar
カードの配置には次のような場所があります。(\hyperref[\detokenize{common/common:field}]{Fig.\@ \ref{\detokenize{common/common:field}}})

\sphinxAtStartPar
※これらの領域は、フレームによって使用する領域が異なります。

\begin{figure}[htbp]
\centering
\capstart

\noindent\sphinxincludegraphics{{extra}.pdf}
\caption{プレイ中のカードの配置}\label{\detokenize{common/common:id45}}\label{\detokenize{common/common:field}}\end{figure}

\sphinxAtStartPar
各領域で定義されたコンポーネントは、この配置図の対応する位置に置かれます。
以下では主要な領域について説明します。

\index{ライフ@\spxentry{ライフ}}\ignorespaces \begin{description}
\sphinxlineitem{ライフ}
\sphinxAtStartPar
山札。ゲームを始める時に自分のデッキを裏向きで置きます。
ダメージを受けるとライフの一番上から墓地にカードを移します。

\end{description}

\index{ぼ\textbar{}墓地@\spxentry{ぼ\textbar{}墓地}}\ignorespaces \begin{description}
\sphinxlineitem{墓地}
\sphinxAtStartPar
捨て札置き場。ダメージを受けた時などに表向きでカードを重ねて置きます。

\end{description}

\index{ば\textbar{}場@\spxentry{ば\textbar{}場}}\ignorespaces \begin{description}
\sphinxlineitem{場}
\sphinxAtStartPar
兵士や防壁などのキャラクターを置きます。

\end{description}

\index{て\textbar{}手札@\spxentry{て\textbar{}手札}}\ignorespaces \begin{description}
\sphinxlineitem{手札}
\sphinxAtStartPar
ライフから引いたカードを持っておく場所です。相手から見えないようにしましょう。

\end{description}

\index{フォグ@\spxentry{フォグ}}\ignorespaces \begin{description}
\sphinxlineitem{フォグ}
\sphinxAtStartPar
このターンのみ影響を与えるカードを置きます。
領域定義で示したフォグキーカードは、すべてここに配置されます。

\end{description}

\index{パック@\spxentry{パック}}\ignorespaces \begin{description}
\sphinxlineitem{パック}
\sphinxAtStartPar
一部のフレームで使用される領域です。デッキから一部のカードを取り分けて伏せて置きます。
アクション（例：パック開封）によって順に使用されることがあります。

\end{description}

\index{レアカード@\spxentry{レアカード}}\ignorespaces \begin{description}
\sphinxlineitem{レアカード}
\sphinxAtStartPar
特別な1枚を配置するための領域です。主にゲーム開始時に選ばれ、フレームによって使い方が異なります。
特定のアクション（例：レアドロー、レア召喚）と組み合わせて使われます。

\end{description}

\index{き\textbar{}切札@\spxentry{き\textbar{}切札}}\ignorespaces \begin{description}
\sphinxlineitem{切札}
\sphinxAtStartPar
切札領域に配置するコンポーネントです。各カードに割り当てられた能力を活用でき、使用可能なフレームは個別に定義されます。

\end{description}


\subsection{デッキとライフ}
\label{\detokenize{common/common:id8}}
\sphinxAtStartPar
対戦レギュレーションなどでデッキという表現が出てきます。

\index{デッキ@\spxentry{デッキ}}\ignorespaces \begin{description}
\sphinxlineitem{デッキ}
\sphinxAtStartPar
ゲーム開始前にゲームで使用するカードの束（カード構成）。
フレームに応じて使用できるデッキの条件が異なります。

\end{description}

\sphinxAtStartPar
ゲームの始め方を経てデッキはライフとなります。詳細は \hyperref[\detokenize{common/common:common-gamestart}]{\ref{\detokenize{common/common:common-gamestart}} \nameref{\detokenize{common/common:common-gamestart}}} で説明します。


\section{勝利条件}
\label{\detokenize{common/common:id9}}
\sphinxAtStartPar
プレイヤーは順に対戦相手に対し攻撃を行い、ダメージを与え先に相手のライフを0枚にした方が勝ちです。ダメージは1点につき、ライフが1枚減ります。

\index{ダメージ@\spxentry{ダメージ}}\ignorespaces 

\section{ダメージ}
\label{\detokenize{common/common:index-11}}\label{\detokenize{common/common:id10}}
\sphinxAtStartPar
プレイヤーがダメージを受けた場合、ライフの一番上から受けた点数分墓地にカードを表向きで移動します。移動する際は、カードの表を対戦相手に見せる必要はありません。

\index{キャラクター@\spxentry{キャラクター}}\ignorespaces 

\section{キャラクター}
\label{\detokenize{common/common:index-12}}\label{\detokenize{common/common:id11}}
\sphinxAtStartPar
キャラクターとは、場に存在する兵士や防壁のことを指し、TCGでいう「クリーチャー」や「モンスター」に相当します。

\sphinxAtStartPar
キャラクターは1枚のカードで1体を表すこともあれば、
複数枚で1体を表すこともあります。(\hyperref[\detokenize{common/common:character}]{Fig.\@ \ref{\detokenize{common/common:character}}})

\begin{figure}[htbp]
\centering
\capstart

\noindent\sphinxincludegraphics{{character}.pdf}
\caption{キャラクターの例}\label{\detokenize{common/common:id46}}\label{\detokenize{common/common:character}}\end{figure}

\sphinxAtStartPar
具体的なキャラクターの持つ項目やサイズ、能力などは、 \hyperref[\detokenize{component/component:character-component}]{\ref{\detokenize{component/component:character-component}} \nameref{\detokenize{component/component:character-component}}} を参照してください。

\index{チャージ@\spxentry{チャージ}}\index{ドライブ@\spxentry{ドライブ}}\ignorespaces 

\subsection{チャージとドライブ}
\label{\detokenize{common/common:index-13}}\label{\detokenize{common/common:id12}}
\sphinxAtStartPar
キャラクターには、チャージ状態とドライブ状態が存在します。
チャージ状態は未使用状態を示し、ドライブ状態は使用済み状態を示しています。
また、キャラクターを横向きにすることを「ドライブ」、縦向きにすることを「チャージ」と言います。(\hyperref[\detokenize{common/common:chargedrive}]{Fig.\@ \ref{\detokenize{common/common:chargedrive}}})

\begin{figure}[htbp]
\centering
\capstart

\noindent\sphinxincludegraphics{{charge&drive}.pdf}
\caption{チャージとドライブ}\label{\detokenize{common/common:id47}}\label{\detokenize{common/common:chargedrive}}\end{figure}


\section{ゲームの始め方}
\label{\detokenize{common/common:common-gamestart}}\label{\detokenize{common/common:id13}}
\sphinxAtStartPar
BlackPokerでは、選択した「フレーム」に応じてゲームの開始手順が異なります。
各フレームでは、「デッキ構築」や「初期配置」、「防壁・兵士のプリセット状態」などが個別に定義されています。

\sphinxAtStartPar
そのため、ゲームを始める際は、まずフレームごとの「ゲームの始め方」に従って準備を行ってください。

\sphinxAtStartPar
各フレームに共通する手順を次に示します。ただし、エクストラではプリセット配置の手順が異なります。


\subsection{共通手順：プリセット配置}
\label{\detokenize{common/common:common-gamestart-preset}}\label{\detokenize{common/common:id14}}\begin{enumerate}
\sphinxsetlistlabels{\arabic}{enumi}{enumii}{}{.}%
\item {} 
\sphinxAtStartPar
ライフの一番上のカードを表向きにして「防壁」として場に出す。

\item {} 
\sphinxAtStartPar
続けてライフの次のカードを表向きにして「兵士」として場に出す。

\end{enumerate}

\sphinxAtStartPar
この手順は「プリセット」と呼ばれています。


\subsection{共通手順：先攻決定}
\label{\detokenize{common/common:common-gamestart-first}}\label{\detokenize{common/common:id15}}\begin{enumerate}
\sphinxsetlistlabels{\arabic}{enumi}{enumii}{}{.}%
\item {} 
\sphinxAtStartPar
両者ライフの一番上のカードを表にする。

\item {} 
\sphinxAtStartPar
数字の大きいプレイヤーが先攻（数値は \hyperref[\detokenize{common/common:cardrank}]{Table \ref{\detokenize{common/common:cardrank}}} を参照）。

\item {} 
\sphinxAtStartPar
同じ数値の場合、さらにライフの一番上をめくって同様に比較。

\item {} 
\sphinxAtStartPar
表にしたカードはすべて墓地に移す。

\end{enumerate}


\subsection{共通手順：ゲーム開始}
\label{\detokenize{common/common:common-gamestart-start}}\label{\detokenize{common/common:id16}}\begin{enumerate}
\sphinxsetlistlabels{\arabic}{enumi}{enumii}{}{.}%
\item {} 
\sphinxAtStartPar
先攻プレイヤーはライフから1枚引いて手札に加える。

\item {} 
\sphinxAtStartPar
先攻プレイヤーがターンとチャンスを持ち、ゲームが開始される。

\end{enumerate}

\begin{sphinxadmonition}{note}{注釈:}
\sphinxAtStartPar
ゲームの始め方の補足
\begin{itemize}
\item {} 
\sphinxAtStartPar
フレームによっては、上記の「プリセット」や「先攻決定」「ゲーム開始」に入る前に、特別な手順（例：ストラテジー選択や切札設定など）があります。

\item {} 
\sphinxAtStartPar
必ず対戦前に、使用するフレームの内容を確認してください。

\end{itemize}
\end{sphinxadmonition}

\index{アクション@\spxentry{アクション}}\ignorespaces 

\section{アクション}
\label{\detokenize{common/common:index-14}}\label{\detokenize{common/common:id17}}
\sphinxAtStartPar
BlackPokerは割込み可能なターン制ゲームです。

\sphinxAtStartPar
例えば次の状況をイメージしてください。
\begin{itemize}
\item {} 
\sphinxAtStartPar
\sphinxstylestrong{Aくん}：「この兵士アップします。」

\item {} 
\sphinxAtStartPar
\sphinxstylestrong{Bさん}：「その前にこの兵士ダウンします。」

\item {} 
\sphinxAtStartPar
\sphinxstylestrong{Aくん}：「じゃあそのダウンをカウンターします。」

\item {} 
\sphinxAtStartPar
\sphinxstylestrong{Bさん}：「それをさらにカウンターします。」

\item {} 
\sphinxAtStartPar
\sphinxstylestrong{Aくん}：「・・・（泣）」

\item {} 
\sphinxAtStartPar
\sphinxstylestrong{Bさん}：「（どやっ！）」

\end{itemize}

\sphinxAtStartPar
このやり取りの中で「アップします」や「ダウンします」などの1行1行がアクションになります。

\sphinxAtStartPar
割込み可能なターン制ゲームは、見方を変えると“許可制のゲーム”とも表現できます。

\sphinxAtStartPar
このアクションを実行したいとルールシステムに要求(リクエスト)し、相手に許可を得てリクエストが実行されます。もちろん相手はリクエストに対して割り込んでリクエストすることもできます。

\sphinxAtStartPar
アクションには、プレイヤーのすべての行動を定義しており、従来のTCGでいう「魔法」や「ターン制御」が含まれています。

\sphinxAtStartPar
参考: \hyperref[\detokenize{core/core:core-rst}]{\ref{\detokenize{core/core:core-rst}} \nameref{\detokenize{core/core:core-rst}}}


\subsection{アクションが持つ項目}
\label{\detokenize{common/common:id18}}
\sphinxAtStartPar
アクションが持つ項目について説明します。

\index{アクション名@\spxentry{アクション名}}\ignorespaces \begin{description}
\sphinxlineitem{アクション名}
\sphinxAtStartPar
アクションの名称を示します。

\end{description}

\index{タイプ(アクション)@\spxentry{タイプ(アクション)}}\ignorespaces \begin{description}
\sphinxlineitem{タイプ}
\sphinxAtStartPar
アクションの種類を表します。アクション名の後に括弧書きで記載します。

\end{description}

\index{トリガー@\spxentry{トリガー}}\ignorespaces \begin{description}
\sphinxlineitem{トリガー}
\sphinxAtStartPar
アクションには自分で起こせるアクションと誘発するアクションがあります。
トリガー項目では「直接」か「誘発」が設定されています。

\sphinxAtStartPar
参考: \hyperref[\detokenize{core/core:trigger-core}]{\ref{\detokenize{core/core:trigger-core}} \nameref{\detokenize{core/core:trigger-core}}}

\end{description}

\index{スピード@\spxentry{スピード}}\ignorespaces \begin{description}
\sphinxlineitem{スピード}
\sphinxAtStartPar
アクションはすぐに効果が解決されるものとそうでないものがあります。
スピード項目では「即時」か「通常」が設定されています。

\sphinxAtStartPar
参考: \hyperref[\detokenize{core/core:speed-core}]{\ref{\detokenize{core/core:speed-core}} \nameref{\detokenize{core/core:speed-core}}}

\end{description}

\index{タイミング@\spxentry{タイミング}}\ignorespaces \begin{description}
\sphinxlineitem{タイミング}
\sphinxAtStartPar
アクションは起こせるタイミングが2種類あります。「メイン」は自分のターンかつステージが空のときに起こせます。
「クイック」はいつでも起こすことができます。

\sphinxAtStartPar
参考: \hyperref[\detokenize{core/core:timing}]{\ref{\detokenize{core/core:timing}} \nameref{\detokenize{core/core:timing}}}

\end{description}

\index{キーカード@\spxentry{キーカード}}\ignorespaces \begin{description}
\sphinxlineitem{キーカード}
\sphinxAtStartPar
アクションの核となるカードを示します。
キーカードは★で表記します。
特に指定がない場合は、手札から選びます。
兵士など場のカードを指定する場合は、その指定に従って選択してください。

\end{description}

\index{コスト@\spxentry{コスト}}\ignorespaces \begin{description}
\sphinxlineitem{コスト}
\sphinxAtStartPar
アクションを起こすのに必要な対価です。
コストは＄を使って表記し、コストの支払いはアクションを起こすプレイヤーが行います。コストの種類は \hyperref[\detokenize{common/common:cost}]{\ref{\detokenize{common/common:cost}} \nameref{\detokenize{common/common:cost}}} で説明します。

\end{description}

\index{と\textbar{}特記事項@\spxentry{と\textbar{}特記事項}}\ignorespaces \begin{description}
\sphinxlineitem{特記事項}
\sphinxAtStartPar
特記事項は※を使って表記し、その他の項目では書き表せない条件を示します。

\end{description}

\index{き\textbar{}起動条件@\spxentry{き\textbar{}起動条件}}\ignorespaces \begin{description}
\sphinxlineitem{起動条件}
\sphinxAtStartPar
起動するための条件を示します。

\end{description}

\index{ゆ\textbar{}誘発条件@\spxentry{ゆ\textbar{}誘発条件}}\ignorespaces \begin{description}
\sphinxlineitem{誘発条件}
\sphinxAtStartPar
誘発する条件を示します。

\end{description}

\index{た\textbar{}対象@\spxentry{た\textbar{}対象}}\ignorespaces \begin{description}
\sphinxlineitem{対象}
\sphinxAtStartPar
効果の対象を示します。

\end{description}

\index{つ\textbar{}通常効果@\spxentry{つ\textbar{}通常効果}!そ\textbar{}即時効果@\spxentry{そ\textbar{}即時効果}}\index{そ\textbar{}即時効果@\spxentry{そ\textbar{}即時効果}!つ\textbar{}通常効果@\spxentry{つ\textbar{}通常効果}}\ignorespaces \begin{description}
\sphinxlineitem{効果}
\sphinxAtStartPar
効果の内容を示します。

\end{description}

\begin{sphinxadmonition}{note}{注釈:}
\sphinxAtStartPar
トリガー,スピード,タイミングの表記

\sphinxAtStartPar
トリガー,スピード,タイミングは@を使って次のように表記されます。

\sphinxAtStartPar
@{[}トリガー{]}\sphinxhyphen{}{[}スピード{]}\sphinxhyphen{}{[}タイミング{]}

\sphinxAtStartPar
例えば次のようになります。

\sphinxAtStartPar
@誘発\sphinxhyphen{}即時\sphinxhyphen{}クイック
\end{sphinxadmonition}


\subsubsection{記載されていないアクションの項目}
\label{\detokenize{common/common:id19}}
\sphinxAtStartPar
アクションによっては記載されていない項目もあります。
記載されていない項目は無視して構いません。
たとえばコスト項目がなければコストを支払う必要はありません。


\subsection{コストの種類}
\label{\detokenize{common/common:cost}}\label{\detokenize{common/common:id20}}
\sphinxAtStartPar
アクションによって支払うコストが異なります。
コストには次の種類があり、それぞれ支払い方が異なります。(\hyperref[\detokenize{common/common:table-cost}]{Table \ref{\detokenize{common/common:table-cost}}})


\begin{savenotes}\sphinxattablestart
\sphinxthistablewithglobalstyle
\centering
\sphinxcapstartof{table}
\sphinxthecaptionisattop
\sphinxcaption{コストの種類}\label{\detokenize{common/common:id48}}\label{\detokenize{common/common:table-cost}}
\sphinxaftertopcaption
\begin{tabulary}{\linewidth}[t]{|T|T|}
\sphinxtoprule
\sphinxstyletheadfamily 
\sphinxAtStartPar
表記(名称)
&\sphinxstyletheadfamily 
\sphinxAtStartPar
対価
\\
\sphinxmidrule
\sphinxtableatstartofbodyhook
\sphinxAtStartPar
B (Bulwark)
&
\sphinxAtStartPar
防壁をドライブする
\\
\sphinxhline
\sphinxAtStartPar
L (Life)
&
\sphinxAtStartPar
1点ダメージを受ける
\\
\sphinxhline
\sphinxAtStartPar
D (Discard)
&
\sphinxAtStartPar
手札を1枚捨てる
\\
\sphinxhline
\sphinxAtStartPar
S (Sacrifice)
&
\sphinxAtStartPar
キャラクター1体を墓地に移す
\\
\sphinxhline
\sphinxAtStartPar
K (KeyDrive)
&
\sphinxAtStartPar
キーカードをドライブする
\\
\sphinxbottomrule
\end{tabulary}
\sphinxtableafterendhook\par
\sphinxattableend\end{savenotes}

\sphinxAtStartPar
たとえばコストが \sphinxstylestrong{「\$BL」} の場合、自分の場にいるチャージ状態の防壁を1体ドライブし、1点ダメージを受けることでコストが支払われたことになります。


\subsection{アクションの起こし方(リクエスト)}
\label{\detokenize{common/common:id21}}
\sphinxAtStartPar
BlackPokerは実行したいアクションを要求(リクエスト)し、進める形式のゲームです。

\sphinxAtStartPar
アクションを要求することを「アクションを起こす」または「アクションをリクエストする」といいます。

\sphinxAtStartPar
次の手順でアクションをリクエストすることができます。
\begin{enumerate}
\sphinxsetlistlabels{\arabic}{enumi}{enumii}{}{.}%
\item {} 
\sphinxAtStartPar
起こすアクションを対戦相手に伝える。

\item {} 
\sphinxAtStartPar
アクションに応じたコストを支払う。

\item {} 
\sphinxAtStartPar
必要なら手札からキーカードを出す。

\item {} 
\sphinxAtStartPar
対象の指定が必要な場合、対象を指定する。

\end{enumerate}


\subsubsection{アクションを起こすときの注意点}
\label{\detokenize{common/common:id22}}

\paragraph{対象を指定しないでアクションを起こせるか？}
\label{\detokenize{common/common:id23}}
\sphinxAtStartPar
「対象」項目がある場合、記載された条件を満たした対象を指定できなければ、
そのアクションを起こすことはできません。


\paragraph{アクションを対象とするアクションは自身を対象にできるか？}
\label{\detokenize{common/common:id24}}
\sphinxAtStartPar
アクションは、自分自身を対象とすることはできません。
そのため、「カウンター」アクションのようにアクションを対象とするアクションは
自身を対象とすることはできません。


\subsection{アクションの解決}
\label{\detokenize{common/common:id25}}
\sphinxAtStartPar
リクエストされたアクションを実行済みにすることを「アクションを解決する」といいます。

\sphinxAtStartPar
実際にアクションが解決される流れを見ていきましょう。

\sphinxAtStartPar
\sphinxstylestrong{用語}

\sphinxAtStartPar
登場する用語を説明します。
\begin{description}
\sphinxlineitem{\sphinxstylestrong{チャンス}}
\sphinxAtStartPar
アクションをリクエストする(ステージに積む)権利（他TCGの優先権）

\sphinxlineitem{\sphinxstylestrong{ターン}}
\sphinxAtStartPar
手番を示す印

\sphinxlineitem{\sphinxstylestrong{ステージ}}
\sphinxAtStartPar
アクションが蓄積される場所（他TCGのスタック）

\end{description}
\begin{enumerate}
\sphinxsetlistlabels{\arabic}{enumi}{enumii}{}{.}%
\item {} 
\sphinxAtStartPar
プレイヤーAがアクションを積む

\end{enumerate}

\begin{figure}[htbp]
\centering
\capstart

\noindent\sphinxincludegraphics{{action-request1}.pdf}
\caption{アクションの解決1}\label{\detokenize{common/common:id49}}\label{\detokenize{common/common:action-request1-image}}\end{figure}
\begin{enumerate}
\sphinxsetlistlabels{\arabic}{enumi}{enumii}{}{.}%
\setcounter{enumi}{1}
\item {} 
\sphinxAtStartPar
プレイヤーBがアクションを積む

\end{enumerate}

\begin{figure}[htbp]
\centering
\capstart

\noindent\sphinxincludegraphics{{action-request2}.pdf}
\caption{アクションの解決2}\label{\detokenize{common/common:id50}}\label{\detokenize{common/common:action-request2-image}}\end{figure}
\begin{enumerate}
\sphinxsetlistlabels{\arabic}{enumi}{enumii}{}{.}%
\setcounter{enumi}{2}
\item {} 
\sphinxAtStartPar
アクションを実行（解決）

\end{enumerate}

\begin{figure}[htbp]
\centering
\capstart

\noindent\sphinxincludegraphics{{action-request3}.pdf}
\caption{アクションの解決3}\label{\detokenize{common/common:id51}}\label{\detokenize{common/common:action-request3-image}}\end{figure}

\sphinxAtStartPar
大まかな流れは図の通りとなります。

\sphinxAtStartPar
更に厳密な処理は \hyperref[\detokenize{core/core:coreflowsec}]{\ref{\detokenize{core/core:coreflowsec}} \nameref{\detokenize{core/core:coreflowsec}}} を参照してください。

\sphinxAtStartPar
次に細かな部分を説明します。

\sphinxAtStartPar
補足ですが次に説明する部分は、 \hyperref[\detokenize{core/core:coreflowsec}]{\ref{\detokenize{core/core:coreflowsec}} \nameref{\detokenize{core/core:coreflowsec}}} の {\hyperref[\detokenize{core/core:actresolve}]{\sphinxcrossref{\DUrole{std,std-ref}{{[}8{]} リクエストの解決}}}} で順に行われます。


\subsubsection{対象条件を確認}
\label{\detokenize{common/common:id26}}
\sphinxAtStartPar
対象を指定するアクションが効果を発揮しようとした時に次の条件に該当する場合、効果を発揮する対象を失うため効果が発揮されず
アクションが解決されます。
\begin{itemize}
\item {} 
\sphinxAtStartPar
対象が存在していない場合

\item {} 
\sphinxAtStartPar
対象が分裂した場合

\end{itemize}

\sphinxAtStartPar
たとえば兵士に対して「アップ」アクションを起こし、対応して「ダウン」
アクションを起こされました。
「ダウン」の方が先に解決されるため、「アップ」を解決する時には
兵士が墓地に移っていたとします。その場合、「アップ」アクションは効果を発揮せず解決されます。

\sphinxAtStartPar
「リバース」による対象分裂も同様です。
たとえば装備兵に対して「ツイスト」アクションを起こし、対応して「リバース」アクションを起こしたとします。
この場合、「リバース」が先に解決され、装備兵が分裂します。
その場合、「ツイスト」は対象を失いアクションの効果を発揮せず解決されます。


\subsubsection{効果を発揮}
\label{\detokenize{common/common:id27}}
\sphinxAtStartPar
リクエストが解決する際に、アクションの効果に定義されている内容を実行します。
効果の中に実行不可能な部分がある場合、可能な部分のみ実行します。

\sphinxAtStartPar
たとえば、ライフの枚数が残1枚の時に5点のダメージを受けたとします。
ライフは1枚しかないので5点ダメージを受けることはできませんが、
1点までなら受けることが可能なため、
この場合1点のダメージを受けることになります。


\subsubsection{解決は墓地移動までを含む}
\label{\detokenize{common/common:id28}}
\sphinxAtStartPar
「リクエストを解決する」という文言には、キーカードを墓地に移動し終えるまでが含まれています。

\sphinxAtStartPar
「リクエストを解決する」をまとめると次のようになります。
\begin{enumerate}
\sphinxsetlistlabels{\arabic}{enumi}{enumii}{}{.}%
\item {} 
\sphinxAtStartPar
ステージの一番上にあるリクエストを特定する

\item {} 
\sphinxAtStartPar
リクエストの対象が正しいか確認する

\item {} 
\sphinxAtStartPar
正しい場合、リクエストされたアクションの効果を可能な限り実行する

\item {} 
\sphinxAtStartPar
リクエストをステージから取り除く

\item {} 
\sphinxAtStartPar
キーカードを墓地に移す

\end{enumerate}


\subsubsection{ステージ上で効果を発揮}
\label{\detokenize{common/common:id29}}
\sphinxAtStartPar
アクションの効果を実行する際にリクエストはまだステージ上にあります。

\sphinxAtStartPar
効果の実行が完了した後、ステージ上から取り除かれます。


\subsubsection{キーカードを墓地に移す}
\label{\detokenize{common/common:keycard-gy}}\label{\detokenize{common/common:id30}}
\sphinxAtStartPar
効果を発揮した後、そのアクションをステージから取り除き、キーカードを墓地に移します。
ただし、次の場合はキーカードを墓地に移しません。
\begin{itemize}
\item {} 
\sphinxAtStartPar
効果の中でキーカードを場に出した場合

\item {} 
\sphinxAtStartPar
効果の中でキーカードを手札に戻した場合

\item {} 
\sphinxAtStartPar
支払いコストに K (KeyDrive) が含まれる場合は、キーカードをドライブ状態のままにする

\end{itemize}


\subsection{勝敗判定}
\label{\detokenize{common/common:id31}}
\sphinxAtStartPar
アクションを解決するたびに勝敗判定が行われます。

\sphinxAtStartPar
勝敗はライフを確認し0枚の場合そのプレイヤーは敗北となります。

\sphinxAtStartPar
勝敗判定はターンプレイヤーから行われます。
もし、両プレイヤーのライフが0枚の場合、ターンプレイヤーの負けとなります。

\sphinxAtStartPar
補足ですがこの勝敗判定は、 \hyperref[\detokenize{core/core:coreflowsec}]{\ref{\detokenize{core/core:coreflowsec}} \nameref{\detokenize{core/core:coreflowsec}}} の {\hyperref[\detokenize{core/core:winlose}]{\sphinxcrossref{\DUrole{std,std-ref}{{[}9{]} 勝敗判定}}}} で確認されます。


\subsection{誘発チェック}
\label{\detokenize{common/common:id32}}
\sphinxAtStartPar
アクションが解決された際に、
アクションの誘発条件に該当するとアクションが誘発されることがあります。

\sphinxAtStartPar
誘発とは、自動的にアクションがリクエストされることです。

\sphinxAtStartPar
BlackPokerの基本的なルールでは、次の2つの誘発パターンがほとんどです。

\sphinxAtStartPar
それ以外の場合は、\hyperref[\detokenize{core/core:coreflowsec}]{\ref{\detokenize{core/core:coreflowsec}} \nameref{\detokenize{core/core:coreflowsec}}} を参照してください。


\subsubsection{パターン1: フェイズ系}
\label{\detokenize{common/common:id33}}
\sphinxAtStartPar
アタックアクションでは対象のアクションが解決すると次のアクションが誘発します。
\begin{enumerate}
\sphinxsetlistlabels{\arabic}{enumi}{enumii}{}{.}%
\item {} 
\sphinxAtStartPar
アタック

\item {} 
\sphinxAtStartPar
ブロック

\item {} 
\sphinxAtStartPar
ダメージ判定

\end{enumerate}

\sphinxAtStartPar
アタックアクションが解決すると、ブロックアクションが誘発します。

\sphinxAtStartPar
ブロックアクションが解決するとダメージ判定アクションが誘発します。

\sphinxAtStartPar
BlackPokerでは一般的なTCGではフェイズとして扱われるものも全てアクションとして定義されています。

\sphinxAtStartPar
同様に、エンド、チャージ、ドローアクションも定義されています。


\subsubsection{パターン2: 世代交代}
\label{\detokenize{common/common:id34}}
\sphinxAtStartPar
BlackPokerでは、Joker,A,J,Q,Kのカードが場から墓地に移った場合、世代交代というアクションが誘発します。

\sphinxAtStartPar
効果の内容はアクションリストを参照してください。

\sphinxAtStartPar
\hyperref[\detokenize{auto/actionlist:act-nextgeneration}]{\ref{\detokenize{auto/actionlist:act-nextgeneration}} \nameref{\detokenize{auto/actionlist:act-nextgeneration}}}


\section{その他のルール}
\label{\detokenize{common/common:id35}}
\sphinxAtStartPar
この章では、
公開レベルやシャッフルの仕方といった
細かな決まりごとを説明します。


\subsection{公開レベル}
\label{\detokenize{common/common:id36}}
\sphinxAtStartPar
配置されているカードには、アクションの効果
を使わなくても中身や枚数を知れるものがあります。
知れる度合いには次の種類があります。
\begin{description}
\sphinxlineitem{完全公開}
\sphinxAtStartPar
全てのプレイヤーが知ることができ、
聞かれたプレイヤーは正しく答える必要がある

\sphinxlineitem{個人公開}
\sphinxAtStartPar
ライフの持ち主のみ知ることができる

\sphinxlineitem{非公開}
\sphinxAtStartPar
全てのプレイヤーは知ることができない

\end{description}

\sphinxAtStartPar
完全公開の情報であれば、ゲーム中いつでも対戦相手に聞くことができます。
各カードの配置と公開・非公開の度合いは次のとおりです。
\begin{description}
\sphinxlineitem{ライフ}
\begin{DUlineblock}{0em}
\item[] 完全公開：10枚未満のライフ枚数
\item[] 個人公開：ライフの枚数
\item[] 非公開：ライフの中身
\end{DUlineblock}

\sphinxlineitem{墓地}
\begin{DUlineblock}{0em}
\item[] 完全公開：墓地の一番上のカード
\item[] 個人公開：墓地の中身
\item[] 非公開：なし
\end{DUlineblock}

\sphinxlineitem{場}
\begin{DUlineblock}{0em}
\item[] 完全公開：表裏を変えずに見えるカード
\item[] 個人公開：伏せてあるカード
\item[] 非公開：なし
\end{DUlineblock}

\sphinxlineitem{手札}
\begin{DUlineblock}{0em}
\item[] 完全公開：手札の枚数
\item[] 個人公開：手札の中身
\item[] 非公開：なし
\end{DUlineblock}

\sphinxlineitem{フォグ}
\begin{DUlineblock}{0em}
\item[] 完全公開：表裏を変えずに見えるカード
\item[] 個人公開：伏せてあるカード
\item[] 非公開：なし
\end{DUlineblock}

\end{description}


\subsubsection{残りライフを聞かれたらどうしたらいいの？}
\label{\detokenize{common/common:id37}}
\sphinxAtStartPar
対戦相手から残りのライフを聞かれた場合、自分のライフの枚数を10枚まで数えます。
10枚未満であれば枚数を答え、10枚以上の場合「10枚以上です」と答えて下さい。
10枚以上の場合、正確な枚数を答える必要はありません。


\subsubsection{墓地の一番上のカードはいつ決まるのか？}
\label{\detokenize{common/common:id38}}
\sphinxAtStartPar
カードを墓地に移す際に移すカードの中から1枚を公開してください。
すでに墓地にあるカードを改めて公開しないでください。


\subsection{デッキのシャッフルについて}
\label{\detokenize{common/common:id39}}
\sphinxAtStartPar
BlackPokerでは
コンセプトの一つに「相手のカードに触らない」があるため、
対戦相手にデッキのシャッフルをお願いする必要はありません。

\sphinxAtStartPar
ただし、シャッフルしてほしい場合は、対戦相手にお願いしても構いません。
逆に、対戦相手があまりシャッフルしていない場合は、
さらにシャッフルをお願いできます。


\subsection{防壁の置き方}
\label{\detokenize{common/common:id40}}
\sphinxAtStartPar
防壁を場に出す際は、次のルールに従って配置してください。(\hyperref[\detokenize{common/common:set-bulwork}]{Fig.\@ \ref{\detokenize{common/common:set-bulwork}}})
\begin{itemize}
\item {} 
\sphinxAtStartPar
防壁を置く際は、ライフ側に寄せて配置してください。

\item {} 
\sphinxAtStartPar
防壁の左右の入れ替えは行わないでください。

\end{itemize}

\begin{figure}[htbp]
\centering
\capstart

\noindent\sphinxincludegraphics{{set-bulwork}.pdf}
\caption{防壁の置き方}\label{\detokenize{common/common:id52}}\label{\detokenize{common/common:set-bulwork}}\end{figure}


\subsection{フォグの置き方}
\label{\detokenize{common/common:id41}}
\sphinxAtStartPar
フォグにアップなどのカードを置く場合は、次のルールに従って配置してください。(\hyperref[\detokenize{common/common:set-fog}]{Fig.\@ \ref{\detokenize{common/common:set-fog}}})
\begin{itemize}
\item {} 
\sphinxAtStartPar
対象の向きにカードを傾けて置いてください。

\item {} 
\sphinxAtStartPar
ダウンなど対戦相手のカードを対象とする場合も同様に置いてください。

\item {} 
\sphinxAtStartPar
フォースなど対象を取らない場合は、対象がないため、傾ける必要はありません。

\item {} 
\sphinxAtStartPar
フォグのカードは、場のカードと重ならないように間隔を空けて配置してください。

\end{itemize}

\begin{figure}[htbp]
\centering
\capstart

\noindent\sphinxincludegraphics{{set-fog}.pdf}
\caption{フォグの置き方}\label{\detokenize{common/common:id53}}\label{\detokenize{common/common:set-fog}}\end{figure}


\subsection{ステージの置き方}
\label{\detokenize{common/common:id42}}
\sphinxAtStartPar
ステージ上にあるリクエストのキーカードは、場のカードと区別できるように、次のルールに従って配置してください。(\hyperref[\detokenize{common/common:set-stage}]{Fig.\@ \ref{\detokenize{common/common:set-stage}}})
\begin{itemize}
\item {} 
\sphinxAtStartPar
ステージ上にあるリクエストのキーカードは傾けて置いてください。

\item {} 
\sphinxAtStartPar
フォグと区別できるように置いてください。

\item {} 
\sphinxAtStartPar
対象がある場合は、その方向に傾けることを推奨します。

\end{itemize}

\begin{figure}[htbp]
\centering
\capstart

\noindent\sphinxincludegraphics{{set-stage}.pdf}
\caption{ステージの置き方}\label{\detokenize{common/common:id54}}\label{\detokenize{common/common:set-stage}}\end{figure}

\sphinxstepscope


\chapter{コンポーネント}
\label{\detokenize{component/component:component-rst}}\label{\detokenize{component/component:id1}}\label{\detokenize{component/component::doc}}
\sphinxAtStartPar
この章では、各コンポーネント（ライフカード、キャラクターキーカードなど）の詳細な構造や振る舞いについて説明します。

\sphinxAtStartPar
各コンポーネントは、配置先の「領域（ゾーン）」に置かれて使用されます。
詳細は共通ルール \hyperref[\detokenize{common/common:zone-section}]{\ref{\detokenize{common/common:zone-section}} \nameref{\detokenize{common/common:zone-section}}} を参照してください。

\begin{sphinxShadowBox}
\sphinxstyletopictitle{コンポーネント一覧}
\begin{itemize}
\item {} 
\sphinxAtStartPar
\phantomsection\label{\detokenize{component/component:id20}}{\hyperref[\detokenize{component/component:id3}]{\sphinxcrossref{ライフ}}}

\item {} 
\sphinxAtStartPar
\phantomsection\label{\detokenize{component/component:id21}}{\hyperref[\detokenize{component/component:id4}]{\sphinxcrossref{墓地}}}

\item {} 
\sphinxAtStartPar
\phantomsection\label{\detokenize{component/component:id22}}{\hyperref[\detokenize{component/component:id5}]{\sphinxcrossref{手札}}}

\item {} 
\sphinxAtStartPar
\phantomsection\label{\detokenize{component/component:id23}}{\hyperref[\detokenize{component/component:character-component}]{\sphinxcrossref{キャラクター}}}
\begin{itemize}
\item {} 
\sphinxAtStartPar
\phantomsection\label{\detokenize{component/component:id24}}{\hyperref[\detokenize{component/component:id7}]{\sphinxcrossref{定義リスト}}}

\item {} 
\sphinxAtStartPar
\phantomsection\label{\detokenize{component/component:id25}}{\hyperref[\detokenize{component/component:id8}]{\sphinxcrossref{キャラクターの持つ項目}}}

\item {} 
\sphinxAtStartPar
\phantomsection\label{\detokenize{component/component:id26}}{\hyperref[\detokenize{component/component:id10}]{\sphinxcrossref{キャラクターの注意点}}}

\end{itemize}

\item {} 
\sphinxAtStartPar
\phantomsection\label{\detokenize{component/component:id27}}{\hyperref[\detokenize{component/component:id12}]{\sphinxcrossref{フォグ}}}
\begin{itemize}
\item {} 
\sphinxAtStartPar
\phantomsection\label{\detokenize{component/component:id28}}{\hyperref[\detokenize{component/component:id13}]{\sphinxcrossref{定義リスト}}}

\item {} 
\sphinxAtStartPar
\phantomsection\label{\detokenize{component/component:id29}}{\hyperref[\detokenize{component/component:id14}]{\sphinxcrossref{フォグの持つ項目}}}

\end{itemize}

\item {} 
\sphinxAtStartPar
\phantomsection\label{\detokenize{component/component:id30}}{\hyperref[\detokenize{component/component:id15}]{\sphinxcrossref{パック}}}

\item {} 
\sphinxAtStartPar
\phantomsection\label{\detokenize{component/component:id31}}{\hyperref[\detokenize{component/component:id16}]{\sphinxcrossref{レアカード}}}

\item {} 
\sphinxAtStartPar
\phantomsection\label{\detokenize{component/component:id32}}{\hyperref[\detokenize{component/component:id17}]{\sphinxcrossref{切札}}}
\begin{itemize}
\item {} 
\sphinxAtStartPar
\phantomsection\label{\detokenize{component/component:id33}}{\hyperref[\detokenize{component/component:id18}]{\sphinxcrossref{定義リスト}}}

\item {} 
\sphinxAtStartPar
\phantomsection\label{\detokenize{component/component:id34}}{\hyperref[\detokenize{component/component:id19}]{\sphinxcrossref{切札の持つ項目}}}

\end{itemize}

\end{itemize}
\end{sphinxShadowBox}


\section{ライフ}
\label{\detokenize{component/component:id3}}
\index{ライフ@\spxentry{ライフ}}\ignorespaces 
\sphinxAtStartPar
ライフカードは、ゲーム開始時に「ライフ領域」に裏向きで配置されるカードです。
プレイヤーの体力のようなもので、ダメージを受けるとライフが減っていきます。
\begin{itemize}
\item {} 
\sphinxAtStartPar
中身は非公開です。

\item {} 
\sphinxAtStartPar
上から順に墓地へ移動します。

\item {} 
\sphinxAtStartPar
定義リスト：なし

\end{itemize}


\section{墓地}
\label{\detokenize{component/component:id4}}
\index{ぼ\textbar{}墓地@\spxentry{ぼ\textbar{}墓地}}\ignorespaces 
\sphinxAtStartPar
墓地カードは、使用済みや破壊されたカードを表向きに重ねて配置するコンポーネントです。
\begin{itemize}
\item {} 
\sphinxAtStartPar
ライフや場などから移動したカードが置かれます。

\item {} 
\sphinxAtStartPar
墓地の一番上のカードは完全公開です。

\item {} 
\sphinxAtStartPar
領域としては「表向きの捨て札置き場」です。

\item {} 
\sphinxAtStartPar
定義リスト：なし

\end{itemize}


\section{手札}
\label{\detokenize{component/component:id5}}
\index{て\textbar{}手札@\spxentry{て\textbar{}手札}}\ignorespaces 
\sphinxAtStartPar
手札カードは、ライフから引いてプレイヤーの手元に保持されるカードです。
\begin{itemize}
\item {} 
\sphinxAtStartPar
枚数は完全公開、内容は個人公開です。

\item {} 
\sphinxAtStartPar
対象の選択やアクションの起点としてよく使われます。

\item {} 
\sphinxAtStartPar
定義リスト：なし

\end{itemize}


\section{キャラクター}
\label{\detokenize{component/component:character-component}}\label{\detokenize{component/component:id6}}
\index{キャラクター@\spxentry{キャラクター}}\ignorespaces 
\sphinxAtStartPar
「場」領域に置かれるコンポーネントです。兵士や防壁を表します。
\begin{itemize}
\item {} 
\sphinxAtStartPar
1枚または複数枚で構成

\item {} 
\sphinxAtStartPar
ドライブ／チャージ状態を持つ

\end{itemize}


\subsection{定義リスト}
\label{\detokenize{component/component:id7}}
\sphinxAtStartPar
次のリストで定義されています。
\begin{itemize}
\item {} 
\sphinxAtStartPar
\hyperref[\detokenize{auto/actionlist:characterlist-act}]{\ref{\detokenize{auto/actionlist:characterlist-act}} \nameref{\detokenize{auto/actionlist:characterlist-act}}}

\end{itemize}


\subsection{キャラクターの持つ項目}
\label{\detokenize{component/component:id8}}
\sphinxAtStartPar
キャラクターの持つ項目について説明します。

\index{キャラクター名@\spxentry{キャラクター名}}\ignorespaces \begin{description}
\sphinxlineitem{キャラクター名}
\sphinxAtStartPar
キャラクターの名称を示します。

\end{description}

\index{タイプ(キャラクター)@\spxentry{タイプ(キャラクター)}}\ignorespaces \begin{description}
\sphinxlineitem{タイプ}
\sphinxAtStartPar
キャラクターのタイプを示します。タイプは兵士と防壁の2種類が存在します。

\end{description}

\index{キャラクターキーカード@\spxentry{キャラクターキーカード}}\ignorespaces \begin{description}
\sphinxlineitem{キャラクターキーカード}
\sphinxAtStartPar
キャラクターを示すカードです。複数枚で１体を表現する場合もあります。

\end{description}

\index{ラベル(キャラクター)@\spxentry{ラベル(キャラクター)}}\ignorespaces \begin{description}
\sphinxlineitem{ラベル}
\sphinxAtStartPar
キャラクターに付与された属性を示します。\textless{}速攻\textgreater{} や \textless{}アタッカー\textgreater{} など様々なラベルがあります。
ラベルは\textless{}\textgreater{}で括り記載されます。

\end{description}

\index{サイズ@\spxentry{サイズ}}\ignorespaces \begin{description}
\sphinxlineitem{サイズ}
\sphinxAtStartPar
兵士の持つ大きさを示します。

\end{description}

\index{の\textbar{}能力(キャラクター)@\spxentry{の\textbar{}能力(キャラクター)}}\ignorespaces \begin{description}
\sphinxlineitem{能力}
\sphinxAtStartPar
キャラクターが持っている能力を記載しています。

\end{description}


\subsubsection{キャラクターのサイズ}
\label{\detokenize{component/component:id9}}
\sphinxAtStartPar
トランプの数字は、キャラクターの強さを示します。
基本はカードに記載された数字に準じますが、魔法などのアクションを使うことで
加算や減算されることがあります。


\subsection{キャラクターの注意点}
\label{\detokenize{component/component:id10}}

\subsubsection{複数枚で1体となるキャラクターが防壁になったら？}
\label{\detokenize{component/component:id11}}
\sphinxAtStartPar
アクションの効果で兵士を防壁にすることがあります。
防壁は1枚で1体のキャラクターであるため、
複数枚からなるキャラクターが防壁となった場合、
複数体の防壁となります。

\sphinxAtStartPar
なお、複数枚からなるキャラクターが
墓地や手札に移った場合、
1体のキャラクターとして
扱うため複数枚合わせて移します。
チャージ状態、ドライブ状態となった場合も同様に1体のキャラクター
として扱います。


\section{フォグ}
\label{\detokenize{component/component:id12}}
\index{フォグ@\spxentry{フォグ}}\ignorespaces 
\sphinxAtStartPar
フォグキーカードは、フォグ領域に一時的に置かれるコンポーネントです。
ターン中のみ効果を発揮する一時的な魔法のような存在です。
\begin{itemize}
\item {} 
\sphinxAtStartPar
置かれるカードとタイミングは、アクションによって定義されます。

\item {} 
\sphinxAtStartPar
効果はターン終了時に自然に失われます。

\end{itemize}


\subsection{定義リスト}
\label{\detokenize{component/component:id13}}
\sphinxAtStartPar
次のリストで定義されています。
\begin{itemize}
\item {} 
\sphinxAtStartPar
\hyperref[\detokenize{auto/actionlist:foglist-act}]{\ref{\detokenize{auto/actionlist:foglist-act}} \nameref{\detokenize{auto/actionlist:foglist-act}}}

\end{itemize}


\subsection{フォグの持つ項目}
\label{\detokenize{component/component:id14}}
\sphinxAtStartPar
フォグの持つ項目について説明します。

\index{フォグキーカード@\spxentry{フォグキーカード}}\ignorespaces \begin{description}
\sphinxlineitem{フォグキーカード}
\sphinxAtStartPar
フォグを示すカードです。複数枚で１つを表現する場合もあります。

\end{description}

\index{の\textbar{}能力(フォグ)@\spxentry{の\textbar{}能力(フォグ)}}\ignorespaces \begin{description}
\sphinxlineitem{能力}
\sphinxAtStartPar
フォグが持っている能力を記載しています。

\end{description}


\section{パック}
\label{\detokenize{component/component:id15}}
\index{パック@\spxentry{パック}}\ignorespaces 
\sphinxAtStartPar
パックカードは、ゲーム開始時にデッキから一部を抜き出してパック領域に置かれるカードです。
特定のフレームでのみ使用されます。
\begin{itemize}
\item {} 
\sphinxAtStartPar
通常、非公開で裏向きに配置されます。

\item {} 
\sphinxAtStartPar
アクション（例：パック開封）によって使用されます。

\item {} 
\sphinxAtStartPar
定義リスト：なし

\end{itemize}


\section{レアカード}
\label{\detokenize{component/component:id16}}
\index{レアカード@\spxentry{レアカード}}\ignorespaces 
\sphinxAtStartPar
レアカードは、ゲーム開始時に1枚選ばれ、特別な領域に伏せて配置されるカードです。
アクションで呼び出されるまで公開されないことが多いです。
\begin{itemize}
\item {} 
\sphinxAtStartPar
対象となるカードは1枚だけです。

\item {} 
\sphinxAtStartPar
レアドローやレア召喚など、専用のアクションで使用されます。

\item {} 
\sphinxAtStartPar
定義リスト：なし

\end{itemize}


\section{切札}
\label{\detokenize{component/component:id17}}
\index{き\textbar{}切札@\spxentry{き\textbar{}切札}}\ignorespaces 
\sphinxAtStartPar
切札は、切札領域に置かれたカードを示します。
各切札には固有の能力が割り当てられており、表にするとその能力が発動します。
一度表になった切札は、裏返すか破壊されるまで能力が持続します。


\subsection{定義リスト}
\label{\detokenize{component/component:id18}}
\sphinxAtStartPar
次のリストで定義されています。
\begin{itemize}
\item {} 
\sphinxAtStartPar
\hyperref[\detokenize{auto/frameActionlist:trumplist-act-frame}]{\ref{\detokenize{auto/frameActionlist:trumplist-act-frame}} \nameref{\detokenize{auto/frameActionlist:trumplist-act-frame}}}

\end{itemize}


\subsection{切札の持つ項目}
\label{\detokenize{component/component:id19}}
\sphinxAtStartPar
切札の持つ項目について説明します。

\index{き\textbar{}切札キーカード@\spxentry{き\textbar{}切札キーカード}}\ignorespaces \begin{description}
\sphinxlineitem{切札キーカード}
\sphinxAtStartPar
切札を示すカードです。スートと数字によって能力が異なります。
裏向きで配置し、「オープン」アクションによって表にします。

\end{description}

\index{ラベル(切札)@\spxentry{ラベル(切札)}}\ignorespaces \begin{description}
\sphinxlineitem{ラベル}
\sphinxAtStartPar
切札に付与された属性を示します。

\end{description}

\index{の\textbar{}能力(切札)@\spxentry{の\textbar{}能力(切札)}}\ignorespaces \begin{description}
\sphinxlineitem{能力}
\sphinxAtStartPar
切札が表になった際に発揮する特殊効果を記載します。
能力の内容や発動条件は、エクストラリスト（切札リスト）に定義されています。

\end{description}

\sphinxstepscope


\chapter{フォーマット}
\label{\detokenize{format/format:format-rst}}\label{\detokenize{format/format:id1}}\label{\detokenize{format/format::doc}}
\index{フォーマット@\spxentry{フォーマット}}\ignorespaces 

\section{フォーマットとは}
\label{\detokenize{format/format:index-0}}\label{\detokenize{format/format:id2}}
\sphinxAtStartPar
BlackPokerにはいくつかのフォーマットがあり、
フォーマットによりゲーム内でできる行動が異なります。
同じトランプでもフォーマットを変えることで様々な遊び方ができます。


\section{定義項目}
\label{\detokenize{format/format:id3}}
\sphinxAtStartPar
フォーマットには次の項目が定義されています。
\begin{description}
\sphinxlineitem{アクションリスト}
\sphinxAtStartPar
リクエスト可能なアクションの一覧

\sphinxlineitem{キャラクターリスト}
\sphinxAtStartPar
場に登場可能なキャラクターの一覧

\sphinxlineitem{フォグリスト}
\sphinxAtStartPar
アクションなどによってフォグに置かれるコンポーネントの一覧

\end{description}


\section{フォーマット定義}
\label{\detokenize{format/format:id4}}
\sphinxAtStartPar
公式として次のフォーマットを定義しています。
\begin{description}
\sphinxlineitem{ライト}
\sphinxAtStartPar
初心者向けフォーマット。覚えるアクションを最小限にしています。

\sphinxlineitem{スタンダード}
\sphinxAtStartPar
中級者向けフォーマット。ライトより幅広い戦略が楽しめます。

\sphinxlineitem{プロ}
\sphinxAtStartPar
上級者向けフォーマット。高度なアクションが多数含まれています。

\sphinxlineitem{マスター}
\sphinxAtStartPar
達人向けフォーマット。全てのアクションが含まれています。

\end{description}

\sphinxAtStartPar
各フォーマットの詳細は次の公式フォーマットを参照してください。

\sphinxstepscope


\subsection{アクションリスト}
\label{\detokenize{auto/actionlist:actionlist-act}}\label{\detokenize{auto/actionlist:act-act}}\label{\detokenize{auto/actionlist:id1}}\label{\detokenize{auto/actionlist::doc}}
\begin{sphinxShadowBox}
\sphinxstyletopictitle{目次}
\begin{itemize}
\item {} 
\sphinxAtStartPar
\phantomsection\label{\detokenize{auto/actionlist:id58}}{\hyperref[\detokenize{auto/actionlist:id3}]{\sphinxcrossref{基本}}}
\begin{itemize}
\item {} 
\sphinxAtStartPar
\phantomsection\label{\detokenize{auto/actionlist:id59}}{\hyperref[\detokenize{auto/actionlist:act-end}]{\sphinxcrossref{エンド (基本)}}}

\item {} 
\sphinxAtStartPar
\phantomsection\label{\detokenize{auto/actionlist:id60}}{\hyperref[\detokenize{auto/actionlist:act-charge}]{\sphinxcrossref{チャージ (基本)}}}

\item {} 
\sphinxAtStartPar
\phantomsection\label{\detokenize{auto/actionlist:id61}}{\hyperref[\detokenize{auto/actionlist:act-draw}]{\sphinxcrossref{ドロー (基本)}}}

\item {} 
\sphinxAtStartPar
\phantomsection\label{\detokenize{auto/actionlist:id62}}{\hyperref[\detokenize{auto/actionlist:act-attack}]{\sphinxcrossref{アタック (基本)}}}

\item {} 
\sphinxAtStartPar
\phantomsection\label{\detokenize{auto/actionlist:id63}}{\hyperref[\detokenize{auto/actionlist:act-block}]{\sphinxcrossref{ブロック (基本)}}}

\item {} 
\sphinxAtStartPar
\phantomsection\label{\detokenize{auto/actionlist:id64}}{\hyperref[\detokenize{auto/actionlist:act-damagejudge}]{\sphinxcrossref{ダメージ判定 (基本)}}}

\item {} 
\sphinxAtStartPar
\phantomsection\label{\detokenize{auto/actionlist:id65}}{\hyperref[\detokenize{auto/actionlist:act-nextgeneration}]{\sphinxcrossref{世代交代 (誘発魔法)}}}

\end{itemize}

\item {} 
\sphinxAtStartPar
\phantomsection\label{\detokenize{auto/actionlist:id66}}{\hyperref[\detokenize{auto/actionlist:id11}]{\sphinxcrossref{召喚}}}
\begin{itemize}
\item {} 
\sphinxAtStartPar
\phantomsection\label{\detokenize{auto/actionlist:id67}}{\hyperref[\detokenize{auto/actionlist:act-setbulwark}]{\sphinxcrossref{防壁設置 (召喚)}}}

\item {} 
\sphinxAtStartPar
\phantomsection\label{\detokenize{auto/actionlist:id68}}{\hyperref[\detokenize{auto/actionlist:act-summonssoldier}]{\sphinxcrossref{兵士召喚 (召喚)}}}

\item {} 
\sphinxAtStartPar
\phantomsection\label{\detokenize{auto/actionlist:id69}}{\hyperref[\detokenize{auto/actionlist:act-summonshero}]{\sphinxcrossref{英雄召喚 (召喚)}}}

\item {} 
\sphinxAtStartPar
\phantomsection\label{\detokenize{auto/actionlist:id70}}{\hyperref[\detokenize{auto/actionlist:act-summonsace}]{\sphinxcrossref{エース召喚 (召喚)}}}

\item {} 
\sphinxAtStartPar
\phantomsection\label{\detokenize{auto/actionlist:id71}}{\hyperref[\detokenize{auto/actionlist:act-quicksummonsace}]{\sphinxcrossref{クイック召喚 (召喚)}}}

\item {} 
\sphinxAtStartPar
\phantomsection\label{\detokenize{auto/actionlist:id72}}{\hyperref[\detokenize{auto/actionlist:act-summonsmagic}]{\sphinxcrossref{魔術士召喚 (召喚)}}}

\item {} 
\sphinxAtStartPar
\phantomsection\label{\detokenize{auto/actionlist:id73}}{\hyperref[\detokenize{auto/actionlist:act-mountsoldier}]{\sphinxcrossref{装備 (召喚)}}}

\end{itemize}

\item {} 
\sphinxAtStartPar
\phantomsection\label{\detokenize{auto/actionlist:id74}}{\hyperref[\detokenize{auto/actionlist:id19}]{\sphinxcrossref{速攻魔法}}}
\begin{itemize}
\item {} 
\sphinxAtStartPar
\phantomsection\label{\detokenize{auto/actionlist:id75}}{\hyperref[\detokenize{auto/actionlist:act-up}]{\sphinxcrossref{アップ (速攻魔法)}}}

\item {} 
\sphinxAtStartPar
\phantomsection\label{\detokenize{auto/actionlist:id76}}{\hyperref[\detokenize{auto/actionlist:act-down}]{\sphinxcrossref{ダウン (速攻魔法)}}}

\item {} 
\sphinxAtStartPar
\phantomsection\label{\detokenize{auto/actionlist:id77}}{\hyperref[\detokenize{auto/actionlist:act-twist}]{\sphinxcrossref{ツイスト (速攻魔法)}}}

\item {} 
\sphinxAtStartPar
\phantomsection\label{\detokenize{auto/actionlist:id78}}{\hyperref[\detokenize{auto/actionlist:act-counter}]{\sphinxcrossref{カウンター (速攻魔法)}}}

\item {} 
\sphinxAtStartPar
\phantomsection\label{\detokenize{auto/actionlist:id79}}{\hyperref[\detokenize{auto/actionlist:act-reunion}]{\sphinxcrossref{再会 (速攻魔法)}}}

\item {} 
\sphinxAtStartPar
\phantomsection\label{\detokenize{auto/actionlist:id80}}{\hyperref[\detokenize{auto/actionlist:act-kill}]{\sphinxcrossref{キル (速攻魔法)}}}

\item {} 
\sphinxAtStartPar
\phantomsection\label{\detokenize{auto/actionlist:id81}}{\hyperref[\detokenize{auto/actionlist:act-truce}]{\sphinxcrossref{停戦 (速攻魔法)}}}

\item {} 
\sphinxAtStartPar
\phantomsection\label{\detokenize{auto/actionlist:id82}}{\hyperref[\detokenize{auto/actionlist:act-changetarget}]{\sphinxcrossref{対象変更 (速攻魔法)}}}

\item {} 
\sphinxAtStartPar
\phantomsection\label{\detokenize{auto/actionlist:id83}}{\hyperref[\detokenize{auto/actionlist:act-search}]{\sphinxcrossref{サーチ (速攻魔法)}}}

\item {} 
\sphinxAtStartPar
\phantomsection\label{\detokenize{auto/actionlist:id84}}{\hyperref[\detokenize{auto/actionlist:bj}]{\sphinxcrossref{B・J (速攻魔法)}}}

\item {} 
\sphinxAtStartPar
\phantomsection\label{\detokenize{auto/actionlist:id85}}{\hyperref[\detokenize{auto/actionlist:rsf}]{\sphinxcrossref{R・S・F (速攻魔法)}}}

\item {} 
\sphinxAtStartPar
\phantomsection\label{\detokenize{auto/actionlist:id86}}{\hyperref[\detokenize{auto/actionlist:act-reverse}]{\sphinxcrossref{リバース (速攻魔法)}}}

\item {} 
\sphinxAtStartPar
\phantomsection\label{\detokenize{auto/actionlist:id87}}{\hyperref[\detokenize{auto/actionlist:act-unsummons}]{\sphinxcrossref{帰還 (速攻魔法)}}}

\end{itemize}

\item {} 
\sphinxAtStartPar
\phantomsection\label{\detokenize{auto/actionlist:id88}}{\hyperref[\detokenize{auto/actionlist:id31}]{\sphinxcrossref{通常魔法}}}
\begin{itemize}
\item {} 
\sphinxAtStartPar
\phantomsection\label{\detokenize{auto/actionlist:id89}}{\hyperref[\detokenize{auto/actionlist:act-destroybulwark}]{\sphinxcrossref{防壁破壊 (通常魔法)}}}

\item {} 
\sphinxAtStartPar
\phantomsection\label{\detokenize{auto/actionlist:id90}}{\hyperref[\detokenize{auto/actionlist:act-throwing}]{\sphinxcrossref{投擲 (通常魔法)}}}

\item {} 
\sphinxAtStartPar
\phantomsection\label{\detokenize{auto/actionlist:id91}}{\hyperref[\detokenize{auto/actionlist:act-deathlance}]{\sphinxcrossref{死の槍 (通常魔法)}}}

\item {} 
\sphinxAtStartPar
\phantomsection\label{\detokenize{auto/actionlist:id92}}{\hyperref[\detokenize{auto/actionlist:act-addbulwark}]{\sphinxcrossref{防壁補充 (通常魔法)}}}

\item {} 
\sphinxAtStartPar
\phantomsection\label{\detokenize{auto/actionlist:id93}}{\hyperref[\detokenize{auto/actionlist:act-reanimate}]{\sphinxcrossref{リアニメイト (通常魔法)}}}

\item {} 
\sphinxAtStartPar
\phantomsection\label{\detokenize{auto/actionlist:id94}}{\hyperref[\detokenize{auto/actionlist:act-handeth}]{\sphinxcrossref{ハンデス (通常魔法)}}}

\item {} 
\sphinxAtStartPar
\phantomsection\label{\detokenize{auto/actionlist:id95}}{\hyperref[\detokenize{auto/actionlist:act-force}]{\sphinxcrossref{フォース (通常魔法)}}}

\item {} 
\sphinxAtStartPar
\phantomsection\label{\detokenize{auto/actionlist:id96}}{\hyperref[\detokenize{auto/actionlist:act-swordrain}]{\sphinxcrossref{剣の雨 (通常魔法)}}}

\item {} 
\sphinxAtStartPar
\phantomsection\label{\detokenize{auto/actionlist:id97}}{\hyperref[\detokenize{auto/actionlist:act-recruit}]{\sphinxcrossref{徴募 (通常魔法)}}}

\item {} 
\sphinxAtStartPar
\phantomsection\label{\detokenize{auto/actionlist:id98}}{\hyperref[\detokenize{auto/actionlist:act-surprise}]{\sphinxcrossref{奇襲 (通常魔法)}}}

\end{itemize}

\end{itemize}
\end{sphinxShadowBox}


\bigskip\hrule\bigskip



\subsubsection{基本}
\label{\detokenize{auto/actionlist:id3}}

\paragraph{エンド (基本)}
\label{\detokenize{auto/actionlist:act-end}}\label{\detokenize{auto/actionlist:id4}}
\sphinxAtStartPar
\sphinxstylestrong{フォーマット:}


\begin{savenotes}\sphinxattablestart
\sphinxthistablewithglobalstyle
\centering
\begin{tabular}[t]{|\X{1}{4}|\X{1}{4}|\X{1}{4}|\X{1}{4}|}
\sphinxtoprule
\sphinxstyletheadfamily 
\sphinxAtStartPar
ライト
&\sphinxstyletheadfamily 
\sphinxAtStartPar
スタンダード
&\sphinxstyletheadfamily 
\sphinxAtStartPar
プロ
&\sphinxstyletheadfamily 
\sphinxAtStartPar
マスター
\\
\sphinxmidrule
\sphinxtableatstartofbodyhook
\sphinxAtStartPar
◯
&
\sphinxAtStartPar
◯
&
\sphinxAtStartPar
◯
&
\sphinxAtStartPar
◯
\\
\sphinxbottomrule
\end{tabular}
\sphinxtableafterendhook\par
\sphinxattableend\end{savenotes}

\sphinxAtStartPar
\sphinxstylestrong{トリガー:} 直接

\sphinxAtStartPar
\sphinxstylestrong{スピード:} 通常

\sphinxAtStartPar
\sphinxstylestrong{タイミング:} メイン

\sphinxAtStartPar
\sphinxstylestrong{効果:}
\begin{enumerate}
\sphinxsetlistlabels{\arabic}{enumi}{enumii}{}{.}%
\item {} 
\sphinxAtStartPar
手札が7枚を越えた場合、7枚になるよう手札を捨てる。

\item {} 
\sphinxAtStartPar
自分のフォグにあるカードを全て墓地に移す。

\item {} 
\sphinxAtStartPar
自分のターンを終了し、対戦相手にターンを渡す。

\end{enumerate}

\sphinxAtStartPar
\sphinxstylestrong{導入バージョン:}  v0.2

\sphinxAtStartPar
\sphinxstylestrong{更新バージョン:}  v8.0


\bigskip\hrule\bigskip



\paragraph{チャージ (基本)}
\label{\detokenize{auto/actionlist:act-charge}}\label{\detokenize{auto/actionlist:id5}}
\sphinxAtStartPar
\sphinxstylestrong{フォーマット:}


\begin{savenotes}\sphinxattablestart
\sphinxthistablewithglobalstyle
\centering
\begin{tabular}[t]{|\X{1}{4}|\X{1}{4}|\X{1}{4}|\X{1}{4}|}
\sphinxtoprule
\sphinxstyletheadfamily 
\sphinxAtStartPar
ライト
&\sphinxstyletheadfamily 
\sphinxAtStartPar
スタンダード
&\sphinxstyletheadfamily 
\sphinxAtStartPar
プロ
&\sphinxstyletheadfamily 
\sphinxAtStartPar
マスター
\\
\sphinxmidrule
\sphinxtableatstartofbodyhook
\sphinxAtStartPar
◯
&
\sphinxAtStartPar
◯
&
\sphinxAtStartPar
◯
&
\sphinxAtStartPar
◯
\\
\sphinxbottomrule
\end{tabular}
\sphinxtableafterendhook\par
\sphinxattableend\end{savenotes}

\sphinxAtStartPar
\sphinxstylestrong{トリガー:} 誘発

\sphinxAtStartPar
\sphinxstylestrong{スピード:} 即時

\sphinxAtStartPar
\sphinxstylestrong{タイミング:} メイン

\sphinxAtStartPar
\sphinxstylestrong{誘発条件:}

\sphinxAtStartPar
ターンプレイヤーかつ、エンドアクションが解決した時に誘発する。

\sphinxAtStartPar
\sphinxstylestrong{効果:}

\sphinxAtStartPar
ターンを持っているプレイヤーの場にいるキャラクターを全てチャージ状態にする。

\sphinxAtStartPar
\sphinxstylestrong{導入バージョン:}  v4.0

\sphinxAtStartPar
\sphinxstylestrong{更新バージョン:}  v8.0


\bigskip\hrule\bigskip



\paragraph{ドロー (基本)}
\label{\detokenize{auto/actionlist:act-draw}}\label{\detokenize{auto/actionlist:id6}}
\sphinxAtStartPar
\sphinxstylestrong{フォーマット:}


\begin{savenotes}\sphinxattablestart
\sphinxthistablewithglobalstyle
\centering
\begin{tabular}[t]{|\X{1}{4}|\X{1}{4}|\X{1}{4}|\X{1}{4}|}
\sphinxtoprule
\sphinxstyletheadfamily 
\sphinxAtStartPar
ライト
&\sphinxstyletheadfamily 
\sphinxAtStartPar
スタンダード
&\sphinxstyletheadfamily 
\sphinxAtStartPar
プロ
&\sphinxstyletheadfamily 
\sphinxAtStartPar
マスター
\\
\sphinxmidrule
\sphinxtableatstartofbodyhook
\sphinxAtStartPar
◯
&
\sphinxAtStartPar
◯
&
\sphinxAtStartPar
◯
&
\sphinxAtStartPar
◯
\\
\sphinxbottomrule
\end{tabular}
\sphinxtableafterendhook\par
\sphinxattableend\end{savenotes}

\sphinxAtStartPar
\sphinxstylestrong{トリガー:} 誘発

\sphinxAtStartPar
\sphinxstylestrong{スピード:} 通常

\sphinxAtStartPar
\sphinxstylestrong{タイミング:} メイン

\sphinxAtStartPar
\sphinxstylestrong{誘発条件:}

\sphinxAtStartPar
ターンプレイヤーかつ、チャージアクションが解決した時に誘発する。

\sphinxAtStartPar
\sphinxstylestrong{効果:}
\begin{enumerate}
\sphinxsetlistlabels{\arabic}{enumi}{enumii}{}{.}%
\item {} 
\sphinxAtStartPar
ライフの一番上からカードを1枚引き、手札に加える。

\item {} 
\sphinxAtStartPar
必要であれば、更にライフの一番上からカードを1枚引き、手札に加える。

\end{enumerate}

\sphinxAtStartPar
\sphinxstylestrong{導入バージョン:}  v4.0

\sphinxAtStartPar
\sphinxstylestrong{更新バージョン:}  v8.0


\bigskip\hrule\bigskip



\paragraph{アタック (基本)}
\label{\detokenize{auto/actionlist:act-attack}}\label{\detokenize{auto/actionlist:id7}}
\sphinxAtStartPar
\sphinxstylestrong{フォーマット:}


\begin{savenotes}\sphinxattablestart
\sphinxthistablewithglobalstyle
\centering
\begin{tabular}[t]{|\X{1}{4}|\X{1}{4}|\X{1}{4}|\X{1}{4}|}
\sphinxtoprule
\sphinxstyletheadfamily 
\sphinxAtStartPar
ライト
&\sphinxstyletheadfamily 
\sphinxAtStartPar
スタンダード
&\sphinxstyletheadfamily 
\sphinxAtStartPar
プロ
&\sphinxstyletheadfamily 
\sphinxAtStartPar
マスター
\\
\sphinxmidrule
\sphinxtableatstartofbodyhook
\sphinxAtStartPar
◯
&
\sphinxAtStartPar
◯
&
\sphinxAtStartPar
◯
&
\sphinxAtStartPar
◯
\\
\sphinxbottomrule
\end{tabular}
\sphinxtableafterendhook\par
\sphinxattableend\end{savenotes}

\sphinxAtStartPar
\sphinxstylestrong{トリガー:} 直接

\sphinxAtStartPar
\sphinxstylestrong{スピード:} 通常

\sphinxAtStartPar
\sphinxstylestrong{タイミング:} メイン

\sphinxAtStartPar
\sphinxstylestrong{起動条件:}

\sphinxAtStartPar
プレイヤーは、1ターンに1回までこのアクションを起こすことができる。

\sphinxAtStartPar
\sphinxstylestrong{効果:}

\sphinxAtStartPar
対戦相手を攻撃するアタッカーをドライブして指定する。

\sphinxAtStartPar
・アタッカーは \textless{}アタッカー\textgreater{} を持っているキャラクターを指定できる。

\sphinxAtStartPar
・アタッカーは複数指定可能。ドライブ状態の兵士は指定できない。

\sphinxAtStartPar
・このターン場に出たキャラクターは指定できない、ただし \textless{}速攻\textgreater{} があるキャラクターは指定できる。

\sphinxAtStartPar
\sphinxstylestrong{導入バージョン:}  v0.2

\sphinxAtStartPar
\sphinxstylestrong{更新バージョン:}  v8.0,v8.1


\bigskip\hrule\bigskip



\paragraph{ブロック (基本)}
\label{\detokenize{auto/actionlist:act-block}}\label{\detokenize{auto/actionlist:id8}}
\sphinxAtStartPar
\sphinxstylestrong{フォーマット:}


\begin{savenotes}\sphinxattablestart
\sphinxthistablewithglobalstyle
\centering
\begin{tabular}[t]{|\X{1}{4}|\X{1}{4}|\X{1}{4}|\X{1}{4}|}
\sphinxtoprule
\sphinxstyletheadfamily 
\sphinxAtStartPar
ライト
&\sphinxstyletheadfamily 
\sphinxAtStartPar
スタンダード
&\sphinxstyletheadfamily 
\sphinxAtStartPar
プロ
&\sphinxstyletheadfamily 
\sphinxAtStartPar
マスター
\\
\sphinxmidrule
\sphinxtableatstartofbodyhook
\sphinxAtStartPar
◯
&
\sphinxAtStartPar
◯
&
\sphinxAtStartPar
◯
&
\sphinxAtStartPar
◯
\\
\sphinxbottomrule
\end{tabular}
\sphinxtableafterendhook\par
\sphinxattableend\end{savenotes}

\sphinxAtStartPar
\sphinxstylestrong{トリガー:} 誘発

\sphinxAtStartPar
\sphinxstylestrong{スピード:} 通常

\sphinxAtStartPar
\sphinxstylestrong{タイミング:} メイン

\sphinxAtStartPar
\sphinxstylestrong{誘発条件:}

\sphinxAtStartPar
ターンプレイヤーかつ、アタックアクションが解決した時にアタッカーが1体以上いる場合、誘発する。

\sphinxAtStartPar
\sphinxstylestrong{効果:}

\sphinxAtStartPar
対戦相手はアタックアクションにて指定されたアタッカー毎にそれをブロックするキャラクター（ブロッカー）を指定する。

\sphinxAtStartPar
・ブロッカーは \textless{}ブロッカー\textgreater{} を持っているキャラクターを指定できる。

\sphinxAtStartPar
・兵士でブロックする場合、1アタッカーに対して複数の兵士を指定できる。ドライブ状態のキャラクターは指定できない。

\sphinxAtStartPar
\sphinxstylestrong{導入バージョン:}  v0.2

\sphinxAtStartPar
\sphinxstylestrong{更新バージョン:}  v8.0,v8.1


\bigskip\hrule\bigskip



\paragraph{ダメージ判定 (基本)}
\label{\detokenize{auto/actionlist:act-damagejudge}}\label{\detokenize{auto/actionlist:id9}}
\sphinxAtStartPar
\sphinxstylestrong{フォーマット:}


\begin{savenotes}\sphinxattablestart
\sphinxthistablewithglobalstyle
\centering
\begin{tabular}[t]{|\X{1}{4}|\X{1}{4}|\X{1}{4}|\X{1}{4}|}
\sphinxtoprule
\sphinxstyletheadfamily 
\sphinxAtStartPar
ライト
&\sphinxstyletheadfamily 
\sphinxAtStartPar
スタンダード
&\sphinxstyletheadfamily 
\sphinxAtStartPar
プロ
&\sphinxstyletheadfamily 
\sphinxAtStartPar
マスター
\\
\sphinxmidrule
\sphinxtableatstartofbodyhook
\sphinxAtStartPar
◯
&
\sphinxAtStartPar
◯
&
\sphinxAtStartPar
◯
&
\sphinxAtStartPar
◯
\\
\sphinxbottomrule
\end{tabular}
\sphinxtableafterendhook\par
\sphinxattableend\end{savenotes}

\sphinxAtStartPar
\sphinxstylestrong{トリガー:} 誘発

\sphinxAtStartPar
\sphinxstylestrong{スピード:} 通常

\sphinxAtStartPar
\sphinxstylestrong{タイミング:} メイン

\sphinxAtStartPar
\sphinxstylestrong{誘発条件:}

\sphinxAtStartPar
ターンプレイヤーかつ、ブロックアクションが解決した時に誘発する。

\sphinxAtStartPar
\sphinxstylestrong{効果:}

\sphinxAtStartPar
アタッカーとブロッカーを比較する
\begin{enumerate}
\sphinxsetlistlabels{\arabic}{enumi}{enumii}{}{.}%
\item {} 
\sphinxAtStartPar
兵士（アタッカー）と兵士（ブロッカー）の場合、アタッカーとブロッカーで数字を比較し、少ない方を墓地に移す。同じ場合は両方を墓地に移動する。1アタッカーに対して複数ブロッカーがいる場合、ブロッカーの合計数字と比較する。

\item {} 
\sphinxAtStartPar
兵士（アタッカー）と防壁（ブロッカー）の場合、次を行う。
\begin{enumerate}
\sphinxsetlistlabels{\arabic}{enumii}{enumiii}{}{.}%
\item {} 
\sphinxAtStartPar
防壁を表にし、防壁が次の条件に当てはまる場合、アタッカーを墓地に移す。
\begin{enumerate}
\sphinxsetlistlabels{\Alph}{enumiii}{enumiv}{}{.}%
\item {} 
\sphinxAtStartPar
防壁がJokerの場合

\item {} 
\sphinxAtStartPar
防壁のカードに記載されている数字と同じ数字がアタッカーのカードに含まれている場合

\end{enumerate}

\item {} 
\sphinxAtStartPar
防壁を墓地に移す。

\end{enumerate}

\item {} 
\sphinxAtStartPar
アタッカーをブロックするブロッカーが場に存在しない場合、アタッカーの数字だけ対戦相手にダメージを与える。

\end{enumerate}

\sphinxAtStartPar
\sphinxstylestrong{導入バージョン:}  v0.2

\sphinxAtStartPar
\sphinxstylestrong{更新バージョン:}  v8.0


\bigskip\hrule\bigskip



\paragraph{世代交代 (誘発魔法)}
\label{\detokenize{auto/actionlist:act-nextgeneration}}\label{\detokenize{auto/actionlist:id10}}
\sphinxAtStartPar
\sphinxstylestrong{フォーマット:}


\begin{savenotes}\sphinxattablestart
\sphinxthistablewithglobalstyle
\centering
\begin{tabular}[t]{|\X{1}{4}|\X{1}{4}|\X{1}{4}|\X{1}{4}|}
\sphinxtoprule
\sphinxstyletheadfamily 
\sphinxAtStartPar
ライト
&\sphinxstyletheadfamily 
\sphinxAtStartPar
スタンダード
&\sphinxstyletheadfamily 
\sphinxAtStartPar
プロ
&\sphinxstyletheadfamily 
\sphinxAtStartPar
マスター
\\
\sphinxmidrule
\sphinxtableatstartofbodyhook
\sphinxAtStartPar
◯
&
\sphinxAtStartPar
◯
&
\sphinxAtStartPar
◯
&
\sphinxAtStartPar
◯
\\
\sphinxbottomrule
\end{tabular}
\sphinxtableafterendhook\par
\sphinxattableend\end{savenotes}

\sphinxAtStartPar
\sphinxstylestrong{トリガー:} 誘発

\sphinxAtStartPar
\sphinxstylestrong{スピード:} 即時

\sphinxAtStartPar
\sphinxstylestrong{タイミング:} クイック

\sphinxAtStartPar
\sphinxstylestrong{誘発条件:}

\sphinxAtStartPar
Joker,A,J,Q,Kのカードが自分の場から墓地に移るたびに1枚につき1回誘発する。

\sphinxAtStartPar
\sphinxstylestrong{効果:}

\sphinxAtStartPar
ライフの一番上からJoker,A,J,Q,Kのいずれかが出るまで墓地にカードを移動し、出たら手札に加える。

\sphinxAtStartPar
\sphinxstylestrong{導入バージョン:}  v1.0

\sphinxAtStartPar
\sphinxstylestrong{更新バージョン:}  v8.0


\bigskip\hrule\bigskip



\subsubsection{召喚}
\label{\detokenize{auto/actionlist:id11}}

\paragraph{防壁設置 (召喚)}
\label{\detokenize{auto/actionlist:act-setbulwark}}\label{\detokenize{auto/actionlist:id12}}
\sphinxAtStartPar
\sphinxstylestrong{フォーマット:}


\begin{savenotes}\sphinxattablestart
\sphinxthistablewithglobalstyle
\centering
\begin{tabular}[t]{|\X{1}{4}|\X{1}{4}|\X{1}{4}|\X{1}{4}|}
\sphinxtoprule
\sphinxstyletheadfamily 
\sphinxAtStartPar
ライト
&\sphinxstyletheadfamily 
\sphinxAtStartPar
スタンダード
&\sphinxstyletheadfamily 
\sphinxAtStartPar
プロ
&\sphinxstyletheadfamily 
\sphinxAtStartPar
マスター
\\
\sphinxmidrule
\sphinxtableatstartofbodyhook
\sphinxAtStartPar
◯
&
\sphinxAtStartPar
◯
&
\sphinxAtStartPar
◯
&
\sphinxAtStartPar
◯
\\
\sphinxbottomrule
\end{tabular}
\sphinxtableafterendhook\par
\sphinxattableend\end{savenotes}

\sphinxAtStartPar
\sphinxstylestrong{トリガー:} 直接

\sphinxAtStartPar
\sphinxstylestrong{スピード:} 即時

\sphinxAtStartPar
\sphinxstylestrong{タイミング:} メイン

\sphinxAtStartPar
\sphinxstylestrong{コスト:} L

\sphinxAtStartPar
\sphinxstylestrong{起動条件:}

\sphinxAtStartPar
プレイヤーは、1ターンに1回までこのアクションを起こすことができる。

\sphinxAtStartPar
\sphinxstylestrong{効果:}

\sphinxAtStartPar
手札からカード1枚を防壁として裏向きかつチャージ状態で場に出す。防壁の置き方は「防壁の置き方」参照。

\sphinxAtStartPar
\sphinxstylestrong{ストーリー:}

\sphinxAtStartPar
どの国もいきなり攻撃はしません。自国を守るため、防衛拠点（防壁：bulwark）を作り、徐々に相手に攻め入ります。

\sphinxAtStartPar
\sphinxstylestrong{フレーバー:}

\sphinxAtStartPar
\sphinxstyleemphasis{「まずは、防壁を出す。話はそれからだ。」}

\sphinxAtStartPar
\textasciitilde{} 再現の師より \textasciitilde{}

\sphinxAtStartPar
\sphinxstylestrong{導入バージョン:}  v0.1

\sphinxAtStartPar
\sphinxstylestrong{更新バージョン:}  v8.0


\bigskip\hrule\bigskip



\paragraph{兵士召喚 (召喚)}
\label{\detokenize{auto/actionlist:act-summonssoldier}}\label{\detokenize{auto/actionlist:id13}}
\sphinxAtStartPar
\sphinxstylestrong{フォーマット:}


\begin{savenotes}\sphinxattablestart
\sphinxthistablewithglobalstyle
\centering
\begin{tabular}[t]{|\X{1}{4}|\X{1}{4}|\X{1}{4}|\X{1}{4}|}
\sphinxtoprule
\sphinxstyletheadfamily 
\sphinxAtStartPar
ライト
&\sphinxstyletheadfamily 
\sphinxAtStartPar
スタンダード
&\sphinxstyletheadfamily 
\sphinxAtStartPar
プロ
&\sphinxstyletheadfamily 
\sphinxAtStartPar
マスター
\\
\sphinxmidrule
\sphinxtableatstartofbodyhook
\sphinxAtStartPar
◯
&
\sphinxAtStartPar
◯
&
\sphinxAtStartPar
◯
&
\sphinxAtStartPar
◯
\\
\sphinxbottomrule
\end{tabular}
\sphinxtableafterendhook\par
\sphinxattableend\end{savenotes}

\sphinxAtStartPar
\sphinxstylestrong{トリガー:} 直接

\sphinxAtStartPar
\sphinxstylestrong{スピード:} 通常

\sphinxAtStartPar
\sphinxstylestrong{タイミング:} メイン

\sphinxAtStartPar
\sphinxstylestrong{キーカード:} 2〜10

\sphinxAtStartPar
\sphinxstylestrong{コスト:} BL

\sphinxAtStartPar
\sphinxstylestrong{効果:}

\sphinxAtStartPar
キーカードを一般兵として表向きかつチャージ状態で場に出す。

\sphinxAtStartPar
\sphinxstylestrong{ストーリー:}

\sphinxAtStartPar
兵士をいきない戦場に出すことはできません。防衛拠点に受け入れてもらい、休息をとってから戦場に出ます。当然民の支援も必要です。

\sphinxAtStartPar
\sphinxstylestrong{フレーバー:}

\sphinxAtStartPar
\sphinxstyleemphasis{「防壁の次は兵士を出します。これが戦術の基本です。」}

\sphinxAtStartPar
\textasciitilde{} 再現者の心得第二章より \textasciitilde{}

\sphinxAtStartPar
\sphinxstylestrong{導入バージョン:}  v3.0

\sphinxAtStartPar
\sphinxstylestrong{更新バージョン:}  v8.0


\bigskip\hrule\bigskip



\paragraph{英雄召喚 (召喚)}
\label{\detokenize{auto/actionlist:act-summonshero}}\label{\detokenize{auto/actionlist:id14}}
\sphinxAtStartPar
\sphinxstylestrong{フォーマット:}


\begin{savenotes}\sphinxattablestart
\sphinxthistablewithglobalstyle
\centering
\begin{tabular}[t]{|\X{1}{4}|\X{1}{4}|\X{1}{4}|\X{1}{4}|}
\sphinxtoprule
\sphinxstyletheadfamily 
\sphinxAtStartPar
ライト
&\sphinxstyletheadfamily 
\sphinxAtStartPar
スタンダード
&\sphinxstyletheadfamily 
\sphinxAtStartPar
プロ
&\sphinxstyletheadfamily 
\sphinxAtStartPar
マスター
\\
\sphinxmidrule
\sphinxtableatstartofbodyhook
\sphinxAtStartPar
◯
&
\sphinxAtStartPar
◯
&
\sphinxAtStartPar
◯
&
\sphinxAtStartPar
◯
\\
\sphinxbottomrule
\end{tabular}
\sphinxtableafterendhook\par
\sphinxattableend\end{savenotes}

\sphinxAtStartPar
\sphinxstylestrong{トリガー:} 直接

\sphinxAtStartPar
\sphinxstylestrong{スピード:} 通常

\sphinxAtStartPar
\sphinxstylestrong{タイミング:} メイン

\sphinxAtStartPar
\sphinxstylestrong{キーカード:} J〜K

\sphinxAtStartPar
\sphinxstylestrong{コスト:} BBL

\sphinxAtStartPar
\sphinxstylestrong{効果:}

\sphinxAtStartPar
キーカードを英雄として表向きかつチャージ状態で場に出す。

\sphinxAtStartPar
\sphinxstylestrong{ストーリー:}

\sphinxAtStartPar
英雄の召喚には一般兵よりコストがかかります。防壁による厳重な警護。十分な休息。そして民の支援です。

\sphinxAtStartPar
\sphinxstylestrong{フレーバー:}

\sphinxAtStartPar
\sphinxstyleemphasis{「実力はわからんが一応英雄様の出陣じゃ。粗相のないようしっかりやれよ。」}

\sphinxAtStartPar
\textasciitilde{} 古参兵より \textasciitilde{}

\sphinxAtStartPar
\sphinxstylestrong{導入バージョン:}  v3.0

\sphinxAtStartPar
\sphinxstylestrong{更新バージョン:}  v8.0


\bigskip\hrule\bigskip



\paragraph{エース召喚 (召喚)}
\label{\detokenize{auto/actionlist:act-summonsace}}\label{\detokenize{auto/actionlist:id15}}
\sphinxAtStartPar
\sphinxstylestrong{フォーマット:}


\begin{savenotes}\sphinxattablestart
\sphinxthistablewithglobalstyle
\centering
\begin{tabular}[t]{|\X{1}{4}|\X{1}{4}|\X{1}{4}|\X{1}{4}|}
\sphinxtoprule
\sphinxstyletheadfamily 
\sphinxAtStartPar
ライト
&\sphinxstyletheadfamily 
\sphinxAtStartPar
スタンダード
&\sphinxstyletheadfamily 
\sphinxAtStartPar
プロ
&\sphinxstyletheadfamily 
\sphinxAtStartPar
マスター
\\
\sphinxmidrule
\sphinxtableatstartofbodyhook
\sphinxAtStartPar
◯
&
\sphinxAtStartPar
◯
&
\sphinxAtStartPar
◯
&
\sphinxAtStartPar
◯
\\
\sphinxbottomrule
\end{tabular}
\sphinxtableafterendhook\par
\sphinxattableend\end{savenotes}

\sphinxAtStartPar
\sphinxstylestrong{トリガー:} 直接

\sphinxAtStartPar
\sphinxstylestrong{スピード:} 通常

\sphinxAtStartPar
\sphinxstylestrong{タイミング:} メイン

\sphinxAtStartPar
\sphinxstylestrong{キーカード:} A

\sphinxAtStartPar
\sphinxstylestrong{コスト:} L

\sphinxAtStartPar
\sphinxstylestrong{効果:}

\sphinxAtStartPar
キーカードをエースとして表向きかつチャージ状態で場に出す。

\sphinxAtStartPar
\sphinxstylestrong{ストーリー:}

\sphinxAtStartPar
エースは素早く行動を起こします。いきなり戦場に登場し切り込んで行きます。

\sphinxAtStartPar
\sphinxstylestrong{フレーバー:}

\sphinxAtStartPar
\sphinxstyleemphasis{「あいつ！何も持たずに飛び出して行きやがった！」}

\sphinxAtStartPar
\textasciitilde{} 親方より \textasciitilde{}

\sphinxAtStartPar
\sphinxstylestrong{導入バージョン:}  v3.0

\sphinxAtStartPar
\sphinxstylestrong{更新バージョン:}  v8.0


\bigskip\hrule\bigskip



\paragraph{クイック召喚 (召喚)}
\label{\detokenize{auto/actionlist:act-quicksummonsace}}\label{\detokenize{auto/actionlist:id16}}
\sphinxAtStartPar
\sphinxstylestrong{フォーマット:}


\begin{savenotes}\sphinxattablestart
\sphinxthistablewithglobalstyle
\centering
\begin{tabular}[t]{|\X{1}{4}|\X{1}{4}|\X{1}{4}|\X{1}{4}|}
\sphinxtoprule
\sphinxstyletheadfamily 
\sphinxAtStartPar
ライト
&\sphinxstyletheadfamily 
\sphinxAtStartPar
スタンダード
&\sphinxstyletheadfamily 
\sphinxAtStartPar
プロ
&\sphinxstyletheadfamily 
\sphinxAtStartPar
マスター
\\
\sphinxmidrule
\sphinxtableatstartofbodyhook
\sphinxAtStartPar
x
&
\sphinxAtStartPar
x
&
\sphinxAtStartPar
◯
&
\sphinxAtStartPar
◯
\\
\sphinxbottomrule
\end{tabular}
\sphinxtableafterendhook\par
\sphinxattableend\end{savenotes}

\sphinxAtStartPar
\sphinxstylestrong{トリガー:} 直接

\sphinxAtStartPar
\sphinxstylestrong{スピード:} 通常

\sphinxAtStartPar
\sphinxstylestrong{タイミング:} クイック

\sphinxAtStartPar
\sphinxstylestrong{キーカード:} A

\sphinxAtStartPar
\sphinxstylestrong{コスト:} D

\sphinxAtStartPar
\sphinxstylestrong{起動条件:}

\sphinxAtStartPar
ターンを持っていない時しかこのアクションを起こすことができない。

\sphinxAtStartPar
\sphinxstylestrong{効果:}

\sphinxAtStartPar
キーカードをエースとして表向きかつチャージ状態で場に出す。

\sphinxAtStartPar
もしくは、キーカードを防壁として裏向きかつチャージ状態で場に出す。

\sphinxAtStartPar
防壁の置き方は「防壁の置き方」参照。

\sphinxAtStartPar
\sphinxstylestrong{導入バージョン:}  v6.0

\sphinxAtStartPar
\sphinxstylestrong{更新バージョン:}  v8.0


\bigskip\hrule\bigskip



\paragraph{魔術士召喚 (召喚)}
\label{\detokenize{auto/actionlist:act-summonsmagic}}\label{\detokenize{auto/actionlist:id17}}
\sphinxAtStartPar
\sphinxstylestrong{フォーマット:}


\begin{savenotes}\sphinxattablestart
\sphinxthistablewithglobalstyle
\centering
\begin{tabular}[t]{|\X{1}{4}|\X{1}{4}|\X{1}{4}|\X{1}{4}|}
\sphinxtoprule
\sphinxstyletheadfamily 
\sphinxAtStartPar
ライト
&\sphinxstyletheadfamily 
\sphinxAtStartPar
スタンダード
&\sphinxstyletheadfamily 
\sphinxAtStartPar
プロ
&\sphinxstyletheadfamily 
\sphinxAtStartPar
マスター
\\
\sphinxmidrule
\sphinxtableatstartofbodyhook
\sphinxAtStartPar
x
&
\sphinxAtStartPar
◯
&
\sphinxAtStartPar
◯
&
\sphinxAtStartPar
◯
\\
\sphinxbottomrule
\end{tabular}
\sphinxtableafterendhook\par
\sphinxattableend\end{savenotes}

\sphinxAtStartPar
\sphinxstylestrong{トリガー:} 直接

\sphinxAtStartPar
\sphinxstylestrong{スピード:} 通常

\sphinxAtStartPar
\sphinxstylestrong{タイミング:} メイン

\sphinxAtStartPar
\sphinxstylestrong{キーカード:} Joker

\sphinxAtStartPar
\sphinxstylestrong{コスト:} BD

\sphinxAtStartPar
\sphinxstylestrong{効果:}

\sphinxAtStartPar
キーカードを魔術士として表向きかつチャージ状態で場に出す。魔術士の能力は、キャラクターリスト参照。

\sphinxAtStartPar
\sphinxstylestrong{ストーリー:}

\sphinxAtStartPar
魔法使いは戦場では魔術士として戦いに臨みます。
コストとして民の支援より手元の予定を狂わせることを要求するあたり捻くれ者であることが伺えます。

\sphinxAtStartPar
\sphinxstylestrong{フレーバー:}

\sphinxAtStartPar
\sphinxstyleemphasis{「場に出れば強いんだよなぁ。場に出れば。」}

\sphinxAtStartPar
\textasciitilde{} 再現者見習いより \textasciitilde{}

\sphinxAtStartPar
\sphinxstylestrong{導入バージョン:}  v3.0

\sphinxAtStartPar
\sphinxstylestrong{更新バージョン:}  v8.0


\bigskip\hrule\bigskip



\paragraph{装備 (召喚)}
\label{\detokenize{auto/actionlist:act-mountsoldier}}\label{\detokenize{auto/actionlist:id18}}
\sphinxAtStartPar
\sphinxstylestrong{フォーマット:}


\begin{savenotes}\sphinxattablestart
\sphinxthistablewithglobalstyle
\centering
\begin{tabular}[t]{|\X{1}{4}|\X{1}{4}|\X{1}{4}|\X{1}{4}|}
\sphinxtoprule
\sphinxstyletheadfamily 
\sphinxAtStartPar
ライト
&\sphinxstyletheadfamily 
\sphinxAtStartPar
スタンダード
&\sphinxstyletheadfamily 
\sphinxAtStartPar
プロ
&\sphinxstyletheadfamily 
\sphinxAtStartPar
マスター
\\
\sphinxmidrule
\sphinxtableatstartofbodyhook
\sphinxAtStartPar
◯
&
\sphinxAtStartPar
◯
&
\sphinxAtStartPar
◯
&
\sphinxAtStartPar
◯
\\
\sphinxbottomrule
\end{tabular}
\sphinxtableafterendhook\par
\sphinxattableend\end{savenotes}

\sphinxAtStartPar
\sphinxstylestrong{トリガー:} 直接

\sphinxAtStartPar
\sphinxstylestrong{スピード:} 通常

\sphinxAtStartPar
\sphinxstylestrong{タイミング:} メイン

\sphinxAtStartPar
\sphinxstylestrong{キーカード:} A〜K

\sphinxAtStartPar
\sphinxstylestrong{コスト:} BL

\sphinxAtStartPar
\sphinxstylestrong{対象:}

\sphinxAtStartPar
キーカードとスートが等しい自分の兵士1体

\sphinxAtStartPar
\sphinxstylestrong{効果:}

\sphinxAtStartPar
対象とした兵士の上にキーカードを置き装備兵とする。装備兵の能力は、キャラクターリスト参照。

\sphinxAtStartPar
\sphinxstylestrong{ストーリー:}

\sphinxAtStartPar
兵士を強化することができます。強化した分強くなるのですが、弱点も増えるので要注意です。

\sphinxAtStartPar
\sphinxstylestrong{フレーバー:}

\sphinxAtStartPar
\sphinxstyleemphasis{「こんなにゴテゴテしてなんになるつもりだ？結局一人だろ？」}

\sphinxAtStartPar
\textasciitilde{} ライフよりの防壁より \textasciitilde{}

\sphinxAtStartPar
\sphinxstylestrong{導入バージョン:}  v0.1

\sphinxAtStartPar
\sphinxstylestrong{更新バージョン:}  v8.0,v8.1


\bigskip\hrule\bigskip



\subsubsection{速攻魔法}
\label{\detokenize{auto/actionlist:id19}}

\paragraph{アップ (速攻魔法)}
\label{\detokenize{auto/actionlist:act-up}}\label{\detokenize{auto/actionlist:id20}}
\sphinxAtStartPar
\sphinxstylestrong{フォーマット:}


\begin{savenotes}\sphinxattablestart
\sphinxthistablewithglobalstyle
\centering
\begin{tabular}[t]{|\X{1}{4}|\X{1}{4}|\X{1}{4}|\X{1}{4}|}
\sphinxtoprule
\sphinxstyletheadfamily 
\sphinxAtStartPar
ライト
&\sphinxstyletheadfamily 
\sphinxAtStartPar
スタンダード
&\sphinxstyletheadfamily 
\sphinxAtStartPar
プロ
&\sphinxstyletheadfamily 
\sphinxAtStartPar
マスター
\\
\sphinxmidrule
\sphinxtableatstartofbodyhook
\sphinxAtStartPar
◯
&
\sphinxAtStartPar
◯
&
\sphinxAtStartPar
◯
&
\sphinxAtStartPar
◯
\\
\sphinxbottomrule
\end{tabular}
\sphinxtableafterendhook\par
\sphinxattableend\end{savenotes}

\sphinxAtStartPar
\sphinxstylestrong{トリガー:} 直接

\sphinxAtStartPar
\sphinxstylestrong{スピード:} 通常

\sphinxAtStartPar
\sphinxstylestrong{タイミング:} クイック

\sphinxAtStartPar
\sphinxstylestrong{キーカード:} {\normalsize $\heartsuit$} A〜10

\sphinxAtStartPar
\sphinxstylestrong{コスト:} D

\sphinxAtStartPar
\sphinxstylestrong{対象:}

\sphinxAtStartPar
兵士1体

\sphinxAtStartPar
\sphinxstylestrong{効果:}

\sphinxAtStartPar
対象とした兵士のサイズは、このターンが終わるまでキーカードの数字分加算される。その印としてフォグにキーカードを置き、アップとする。

\sphinxAtStartPar
\sphinxstylestrong{導入バージョン:}  v0.1

\sphinxAtStartPar
\sphinxstylestrong{更新バージョン:}  v8.0,v8.1


\bigskip\hrule\bigskip



\paragraph{ダウン (速攻魔法)}
\label{\detokenize{auto/actionlist:act-down}}\label{\detokenize{auto/actionlist:id21}}
\sphinxAtStartPar
\sphinxstylestrong{フォーマット:}


\begin{savenotes}\sphinxattablestart
\sphinxthistablewithglobalstyle
\centering
\begin{tabular}[t]{|\X{1}{4}|\X{1}{4}|\X{1}{4}|\X{1}{4}|}
\sphinxtoprule
\sphinxstyletheadfamily 
\sphinxAtStartPar
ライト
&\sphinxstyletheadfamily 
\sphinxAtStartPar
スタンダード
&\sphinxstyletheadfamily 
\sphinxAtStartPar
プロ
&\sphinxstyletheadfamily 
\sphinxAtStartPar
マスター
\\
\sphinxmidrule
\sphinxtableatstartofbodyhook
\sphinxAtStartPar
◯
&
\sphinxAtStartPar
◯
&
\sphinxAtStartPar
◯
&
\sphinxAtStartPar
◯
\\
\sphinxbottomrule
\end{tabular}
\sphinxtableafterendhook\par
\sphinxattableend\end{savenotes}

\sphinxAtStartPar
\sphinxstylestrong{トリガー:} 直接

\sphinxAtStartPar
\sphinxstylestrong{スピード:} 通常

\sphinxAtStartPar
\sphinxstylestrong{タイミング:} クイック

\sphinxAtStartPar
\sphinxstylestrong{キーカード:} {\normalsize $\spadesuit$} A〜10

\sphinxAtStartPar
\sphinxstylestrong{コスト:} D

\sphinxAtStartPar
\sphinxstylestrong{対象:}

\sphinxAtStartPar
兵士1体

\sphinxAtStartPar
\sphinxstylestrong{効果:}

\sphinxAtStartPar
対象とした兵士のサイズは、このターンが終わるまでキーカードの数字分減算される。

\sphinxAtStartPar
もし対象のサイズが0以下となった場合、対象を墓地に移す。

\sphinxAtStartPar
もし対象のサイズが0以下とならなかった場合、その印としてフォグにキーカードを置き、ダウンとする。

\sphinxAtStartPar
\sphinxstylestrong{導入バージョン:}  v0.1

\sphinxAtStartPar
\sphinxstylestrong{更新バージョン:}  v8.0,v8.1


\bigskip\hrule\bigskip



\paragraph{ツイスト (速攻魔法)}
\label{\detokenize{auto/actionlist:act-twist}}\label{\detokenize{auto/actionlist:id22}}
\sphinxAtStartPar
\sphinxstylestrong{フォーマット:}


\begin{savenotes}\sphinxattablestart
\sphinxthistablewithglobalstyle
\centering
\begin{tabular}[t]{|\X{1}{4}|\X{1}{4}|\X{1}{4}|\X{1}{4}|}
\sphinxtoprule
\sphinxstyletheadfamily 
\sphinxAtStartPar
ライト
&\sphinxstyletheadfamily 
\sphinxAtStartPar
スタンダード
&\sphinxstyletheadfamily 
\sphinxAtStartPar
プロ
&\sphinxstyletheadfamily 
\sphinxAtStartPar
マスター
\\
\sphinxmidrule
\sphinxtableatstartofbodyhook
\sphinxAtStartPar
◯
&
\sphinxAtStartPar
◯
&
\sphinxAtStartPar
◯
&
\sphinxAtStartPar
◯
\\
\sphinxbottomrule
\end{tabular}
\sphinxtableafterendhook\par
\sphinxattableend\end{savenotes}

\sphinxAtStartPar
\sphinxstylestrong{トリガー:} 直接

\sphinxAtStartPar
\sphinxstylestrong{スピード:} 通常

\sphinxAtStartPar
\sphinxstylestrong{タイミング:} クイック

\sphinxAtStartPar
\sphinxstylestrong{キーカード:} {\normalsize $\diamondsuit$} A〜10

\sphinxAtStartPar
\sphinxstylestrong{コスト:} D

\sphinxAtStartPar
\sphinxstylestrong{対象:}

\sphinxAtStartPar
キャラクター1体

\sphinxAtStartPar
\sphinxstylestrong{効果:}

\sphinxAtStartPar
対象のキャラクターをドライブ状態またはチャージ状態にする。

\sphinxAtStartPar
\sphinxstylestrong{導入バージョン:}  v0.2

\sphinxAtStartPar
\sphinxstylestrong{更新バージョン:}  v8.0,v8.1


\bigskip\hrule\bigskip



\paragraph{カウンター (速攻魔法)}
\label{\detokenize{auto/actionlist:act-counter}}\label{\detokenize{auto/actionlist:id23}}
\sphinxAtStartPar
\sphinxstylestrong{フォーマット:}


\begin{savenotes}\sphinxattablestart
\sphinxthistablewithglobalstyle
\centering
\begin{tabular}[t]{|\X{1}{4}|\X{1}{4}|\X{1}{4}|\X{1}{4}|}
\sphinxtoprule
\sphinxstyletheadfamily 
\sphinxAtStartPar
ライト
&\sphinxstyletheadfamily 
\sphinxAtStartPar
スタンダード
&\sphinxstyletheadfamily 
\sphinxAtStartPar
プロ
&\sphinxstyletheadfamily 
\sphinxAtStartPar
マスター
\\
\sphinxmidrule
\sphinxtableatstartofbodyhook
\sphinxAtStartPar
◯
&
\sphinxAtStartPar
◯
&
\sphinxAtStartPar
◯
&
\sphinxAtStartPar
◯
\\
\sphinxbottomrule
\end{tabular}
\sphinxtableafterendhook\par
\sphinxattableend\end{savenotes}

\sphinxAtStartPar
\sphinxstylestrong{トリガー:} 直接

\sphinxAtStartPar
\sphinxstylestrong{スピード:} 通常

\sphinxAtStartPar
\sphinxstylestrong{タイミング:} クイック

\sphinxAtStartPar
\sphinxstylestrong{キーカード:} {\normalsize $\clubsuit$} A〜10

\sphinxAtStartPar
\sphinxstylestrong{コスト:} D

\sphinxAtStartPar
\sphinxstylestrong{対象:}

\sphinxAtStartPar
キーカードが1枚または2枚のリクエスト

\sphinxAtStartPar
\sphinxstylestrong{効果:}

\sphinxAtStartPar
次のいずれかの場合、対象のリクエストを無効にする。その場合、対象リクエストをステージから取り除き、対象リクエストのキーカードを墓地に移す。

\sphinxAtStartPar
・対象リクエストのキーカードが1枚かつこのリクエストのキーカードの数字が対象リクエストのキーカードの数字以上

\sphinxAtStartPar
・対象リクエストのキーカードが2枚

\sphinxAtStartPar
\sphinxstylestrong{導入バージョン:}  v0.1

\sphinxAtStartPar
\sphinxstylestrong{更新バージョン:}  v8.0,v8.1


\bigskip\hrule\bigskip



\paragraph{再会 (速攻魔法)}
\label{\detokenize{auto/actionlist:act-reunion}}\label{\detokenize{auto/actionlist:id24}}
\sphinxAtStartPar
\sphinxstylestrong{フォーマット:}


\begin{savenotes}\sphinxattablestart
\sphinxthistablewithglobalstyle
\centering
\begin{tabular}[t]{|\X{1}{4}|\X{1}{4}|\X{1}{4}|\X{1}{4}|}
\sphinxtoprule
\sphinxstyletheadfamily 
\sphinxAtStartPar
ライト
&\sphinxstyletheadfamily 
\sphinxAtStartPar
スタンダード
&\sphinxstyletheadfamily 
\sphinxAtStartPar
プロ
&\sphinxstyletheadfamily 
\sphinxAtStartPar
マスター
\\
\sphinxmidrule
\sphinxtableatstartofbodyhook
\sphinxAtStartPar
x
&
\sphinxAtStartPar
x
&
\sphinxAtStartPar
◯
&
\sphinxAtStartPar
◯
\\
\sphinxbottomrule
\end{tabular}
\sphinxtableafterendhook\par
\sphinxattableend\end{savenotes}

\sphinxAtStartPar
\sphinxstylestrong{トリガー:} 直接

\sphinxAtStartPar
\sphinxstylestrong{スピード:} 通常

\sphinxAtStartPar
\sphinxstylestrong{タイミング:} クイック

\sphinxAtStartPar
\sphinxstylestrong{キーカード:} {\normalsize $\heartsuit$} A〜10 を 2枚

\sphinxAtStartPar
\sphinxstylestrong{効果:}

\sphinxAtStartPar
自分の墓地からカードを1枚選び対戦相手に見せ手札に加える。

\sphinxAtStartPar
\sphinxstylestrong{ストーリー:}

\sphinxAtStartPar
親愛な仲間との再会を強く強く願った。どんな手段も厭わないその願いが今叶った。

\sphinxAtStartPar
\sphinxstylestrong{フレーバー:}

\sphinxAtStartPar
\sphinxstyleemphasis{「ねぇねぇ知ってた。この後ヒロインの願いが届いて死んだはずの主人公が登場して、二人で敵を全員やっつけちゃうの！ありがちよねぇ。」}

\sphinxAtStartPar
\textasciitilde{} ネタバレ娘より \textasciitilde{}

\sphinxAtStartPar
\sphinxstylestrong{導入バージョン:}  v1.0

\sphinxAtStartPar
\sphinxstylestrong{更新バージョン:}  v8.0


\bigskip\hrule\bigskip



\paragraph{キル (速攻魔法)}
\label{\detokenize{auto/actionlist:act-kill}}\label{\detokenize{auto/actionlist:id25}}
\sphinxAtStartPar
\sphinxstylestrong{フォーマット:}


\begin{savenotes}\sphinxattablestart
\sphinxthistablewithglobalstyle
\centering
\begin{tabular}[t]{|\X{1}{4}|\X{1}{4}|\X{1}{4}|\X{1}{4}|}
\sphinxtoprule
\sphinxstyletheadfamily 
\sphinxAtStartPar
ライト
&\sphinxstyletheadfamily 
\sphinxAtStartPar
スタンダード
&\sphinxstyletheadfamily 
\sphinxAtStartPar
プロ
&\sphinxstyletheadfamily 
\sphinxAtStartPar
マスター
\\
\sphinxmidrule
\sphinxtableatstartofbodyhook
\sphinxAtStartPar
x
&
\sphinxAtStartPar
x
&
\sphinxAtStartPar
◯
&
\sphinxAtStartPar
◯
\\
\sphinxbottomrule
\end{tabular}
\sphinxtableafterendhook\par
\sphinxattableend\end{savenotes}

\sphinxAtStartPar
\sphinxstylestrong{トリガー:} 直接

\sphinxAtStartPar
\sphinxstylestrong{スピード:} 通常

\sphinxAtStartPar
\sphinxstylestrong{タイミング:} クイック

\sphinxAtStartPar
\sphinxstylestrong{キーカード:} {\normalsize $\spadesuit$} A〜10 を 2枚

\sphinxAtStartPar
\sphinxstylestrong{対象:}

\sphinxAtStartPar
兵士1体

\sphinxAtStartPar
\sphinxstylestrong{効果:}

\sphinxAtStartPar
対象とした兵士を墓地に移す。

\sphinxAtStartPar
\sphinxstylestrong{導入バージョン:}  v1.0

\sphinxAtStartPar
\sphinxstylestrong{更新バージョン:}  v8.0,v8.1


\bigskip\hrule\bigskip



\paragraph{停戦 (速攻魔法)}
\label{\detokenize{auto/actionlist:act-truce}}\label{\detokenize{auto/actionlist:id26}}
\sphinxAtStartPar
\sphinxstylestrong{フォーマット:}


\begin{savenotes}\sphinxattablestart
\sphinxthistablewithglobalstyle
\centering
\begin{tabular}[t]{|\X{1}{4}|\X{1}{4}|\X{1}{4}|\X{1}{4}|}
\sphinxtoprule
\sphinxstyletheadfamily 
\sphinxAtStartPar
ライト
&\sphinxstyletheadfamily 
\sphinxAtStartPar
スタンダード
&\sphinxstyletheadfamily 
\sphinxAtStartPar
プロ
&\sphinxstyletheadfamily 
\sphinxAtStartPar
マスター
\\
\sphinxmidrule
\sphinxtableatstartofbodyhook
\sphinxAtStartPar
x
&
\sphinxAtStartPar
x
&
\sphinxAtStartPar
◯
&
\sphinxAtStartPar
◯
\\
\sphinxbottomrule
\end{tabular}
\sphinxtableafterendhook\par
\sphinxattableend\end{savenotes}

\sphinxAtStartPar
\sphinxstylestrong{トリガー:} 直接

\sphinxAtStartPar
\sphinxstylestrong{スピード:} 通常

\sphinxAtStartPar
\sphinxstylestrong{タイミング:} クイック

\sphinxAtStartPar
\sphinxstylestrong{キーカード:} {\normalsize $\diamondsuit$} A〜10 を 2枚

\sphinxAtStartPar
\sphinxstylestrong{対象:}

\sphinxAtStartPar
ダメージ判定アクションのリクエスト

\sphinxAtStartPar
\sphinxstylestrong{効果:}

\sphinxAtStartPar
対象のリクエストを無効にする。対象のリクエストをステージから取り除き、対象のリクエストのキーカードを墓地に移す。

\sphinxAtStartPar
\sphinxstylestrong{導入バージョン:}  v1.0

\sphinxAtStartPar
\sphinxstylestrong{更新バージョン:}  v8.0,v8.1


\bigskip\hrule\bigskip



\paragraph{対象変更 (速攻魔法)}
\label{\detokenize{auto/actionlist:act-changetarget}}\label{\detokenize{auto/actionlist:id27}}
\sphinxAtStartPar
\sphinxstylestrong{フォーマット:}


\begin{savenotes}\sphinxattablestart
\sphinxthistablewithglobalstyle
\centering
\begin{tabular}[t]{|\X{1}{4}|\X{1}{4}|\X{1}{4}|\X{1}{4}|}
\sphinxtoprule
\sphinxstyletheadfamily 
\sphinxAtStartPar
ライト
&\sphinxstyletheadfamily 
\sphinxAtStartPar
スタンダード
&\sphinxstyletheadfamily 
\sphinxAtStartPar
プロ
&\sphinxstyletheadfamily 
\sphinxAtStartPar
マスター
\\
\sphinxmidrule
\sphinxtableatstartofbodyhook
\sphinxAtStartPar
x
&
\sphinxAtStartPar
x
&
\sphinxAtStartPar
◯
&
\sphinxAtStartPar
◯
\\
\sphinxbottomrule
\end{tabular}
\sphinxtableafterendhook\par
\sphinxattableend\end{savenotes}

\sphinxAtStartPar
\sphinxstylestrong{トリガー:} 直接

\sphinxAtStartPar
\sphinxstylestrong{スピード:} 通常

\sphinxAtStartPar
\sphinxstylestrong{タイミング:} クイック

\sphinxAtStartPar
\sphinxstylestrong{キーカード:} {\normalsize $\clubsuit$} A〜10 を 2枚

\sphinxAtStartPar
\sphinxstylestrong{対象:}

\sphinxAtStartPar
対象が指定されているリクエスト

\sphinxAtStartPar
\sphinxstylestrong{効果:}

\sphinxAtStartPar
対象のリクエストで指定されている対象をそのアクションが指定できる範囲で変更する。リクエストの対象を対象変更に変更することは可能、アクションの対象にできないものへの対象の変更は不可能とする。

\sphinxAtStartPar
\sphinxstylestrong{導入バージョン:}  v1.0

\sphinxAtStartPar
\sphinxstylestrong{更新バージョン:}  v8.0,v8.1


\bigskip\hrule\bigskip



\paragraph{サーチ (速攻魔法)}
\label{\detokenize{auto/actionlist:act-search}}\label{\detokenize{auto/actionlist:id28}}
\sphinxAtStartPar
\sphinxstylestrong{フォーマット:}


\begin{savenotes}\sphinxattablestart
\sphinxthistablewithglobalstyle
\centering
\begin{tabular}[t]{|\X{1}{4}|\X{1}{4}|\X{1}{4}|\X{1}{4}|}
\sphinxtoprule
\sphinxstyletheadfamily 
\sphinxAtStartPar
ライト
&\sphinxstyletheadfamily 
\sphinxAtStartPar
スタンダード
&\sphinxstyletheadfamily 
\sphinxAtStartPar
プロ
&\sphinxstyletheadfamily 
\sphinxAtStartPar
マスター
\\
\sphinxmidrule
\sphinxtableatstartofbodyhook
\sphinxAtStartPar
◯
&
\sphinxAtStartPar
◯
&
\sphinxAtStartPar
◯
&
\sphinxAtStartPar
◯
\\
\sphinxbottomrule
\end{tabular}
\sphinxtableafterendhook\par
\sphinxattableend\end{savenotes}

\sphinxAtStartPar
\sphinxstylestrong{トリガー:} 直接

\sphinxAtStartPar
\sphinxstylestrong{スピード:} 即時

\sphinxAtStartPar
\sphinxstylestrong{タイミング:} クイック

\sphinxAtStartPar
\sphinxstylestrong{キーカード:} Joker

\sphinxAtStartPar
\sphinxstylestrong{効果:}

\sphinxAtStartPar
ライフから好きなカードを1枚選び対戦相手に見せ手札に加える。その後ライフを切りなおす。

\sphinxAtStartPar
\sphinxstylestrong{導入バージョン:}  v0.1

\sphinxAtStartPar
\sphinxstylestrong{更新バージョン:}  v8.0


\bigskip\hrule\bigskip



\paragraph{B・J (速攻魔法)}
\label{\detokenize{auto/actionlist:bj}}\label{\detokenize{auto/actionlist:act-bj}}
\sphinxAtStartPar
\sphinxstylestrong{フォーマット:}


\begin{savenotes}\sphinxattablestart
\sphinxthistablewithglobalstyle
\centering
\begin{tabular}[t]{|\X{1}{4}|\X{1}{4}|\X{1}{4}|\X{1}{4}|}
\sphinxtoprule
\sphinxstyletheadfamily 
\sphinxAtStartPar
ライト
&\sphinxstyletheadfamily 
\sphinxAtStartPar
スタンダード
&\sphinxstyletheadfamily 
\sphinxAtStartPar
プロ
&\sphinxstyletheadfamily 
\sphinxAtStartPar
マスター
\\
\sphinxmidrule
\sphinxtableatstartofbodyhook
\sphinxAtStartPar
x
&
\sphinxAtStartPar
x
&
\sphinxAtStartPar
x
&
\sphinxAtStartPar
◯
\\
\sphinxbottomrule
\end{tabular}
\sphinxtableafterendhook\par
\sphinxattableend\end{savenotes}

\sphinxAtStartPar
\sphinxstylestrong{トリガー:} 直接

\sphinxAtStartPar
\sphinxstylestrong{スピード:} 通常

\sphinxAtStartPar
\sphinxstylestrong{タイミング:} クイック

\sphinxAtStartPar
\sphinxstylestrong{キーカード:} A と J

\sphinxAtStartPar
\sphinxstylestrong{効果:}
\begin{enumerate}
\sphinxsetlistlabels{\arabic}{enumi}{enumii}{}{.}%
\item {} 
\sphinxAtStartPar
合計点 を 0 とし、自分のライフの山札から1枚めくり表向きにする。公開したカードをまとめて 公開カード と呼ぶ。

\item {} 
\sphinxAtStartPar
めくったカードの点数を合計点に加える。

\sphinxAtStartPar
・A は 11 点として扱う。ただし合計点が 22 点以上になる場合は 1 点として扱う。

\sphinxAtStartPar
・J,Q,K は 10 点、Joker は 0 点、2〜10 は数字どおりとして扱う。

\item {} 
\sphinxAtStartPar
合計点が 22 点以上になった場合、公開カードをすべて墓地に置き、効果を終了する。

\item {} 
\sphinxAtStartPar
合計点が 21 点以下の場合、ヒット もしくは、スタンドを行う。

\sphinxAtStartPar
・ヒット: さらに 1 枚めくり、手順 2 に戻る。

\sphinxAtStartPar
・スタンド: 公開カードすべてを手札に加え、効果を終了する。

\end{enumerate}

\sphinxAtStartPar
\sphinxstylestrong{ストーリー:}

\sphinxAtStartPar
若きジャックはエースの力をかり革命的な行動を起こした。自国の膿を排除し新しい風を送り込んだ。

\sphinxAtStartPar
\sphinxstylestrong{フレーバー:}

\sphinxAtStartPar
\sphinxstyleemphasis{「悪しき習慣を捨て去り、新しい風を吹き込む」}

\sphinxAtStartPar
\textasciitilde{} 革命者より \textasciitilde{}

\sphinxAtStartPar
\sphinxstylestrong{導入バージョン:}  v1.0

\sphinxAtStartPar
\sphinxstylestrong{更新バージョン:}  v8.0,v8.1


\bigskip\hrule\bigskip



\paragraph{R・S・F (速攻魔法)}
\label{\detokenize{auto/actionlist:rsf}}\label{\detokenize{auto/actionlist:act-rsf}}
\sphinxAtStartPar
\sphinxstylestrong{フォーマット:}


\begin{savenotes}\sphinxattablestart
\sphinxthistablewithglobalstyle
\centering
\begin{tabular}[t]{|\X{1}{4}|\X{1}{4}|\X{1}{4}|\X{1}{4}|}
\sphinxtoprule
\sphinxstyletheadfamily 
\sphinxAtStartPar
ライト
&\sphinxstyletheadfamily 
\sphinxAtStartPar
スタンダード
&\sphinxstyletheadfamily 
\sphinxAtStartPar
プロ
&\sphinxstyletheadfamily 
\sphinxAtStartPar
マスター
\\
\sphinxmidrule
\sphinxtableatstartofbodyhook
\sphinxAtStartPar
x
&
\sphinxAtStartPar
x
&
\sphinxAtStartPar
x
&
\sphinxAtStartPar
◯
\\
\sphinxbottomrule
\end{tabular}
\sphinxtableafterendhook\par
\sphinxattableend\end{savenotes}

\sphinxAtStartPar
\sphinxstylestrong{トリガー:} 直接

\sphinxAtStartPar
\sphinxstylestrong{スピード:} 通常

\sphinxAtStartPar
\sphinxstylestrong{タイミング:} クイック

\sphinxAtStartPar
\sphinxstylestrong{キーカード:} 同じスートのA,10〜K を 5枚

\sphinxAtStartPar
\sphinxstylestrong{コスト:} BB

\sphinxAtStartPar
\sphinxstylestrong{対象:}

\sphinxAtStartPar
プレイヤー1人

\sphinxAtStartPar
\sphinxstylestrong{効果:}

\sphinxAtStartPar
対象のプレイヤーに40点のダメージを与える。

\sphinxAtStartPar
\sphinxstylestrong{導入バージョン:}  v1.0

\sphinxAtStartPar
\sphinxstylestrong{更新バージョン:}  v8.0,v8.1


\bigskip\hrule\bigskip



\paragraph{リバース (速攻魔法)}
\label{\detokenize{auto/actionlist:act-reverse}}\label{\detokenize{auto/actionlist:id29}}
\sphinxAtStartPar
\sphinxstylestrong{フォーマット:}


\begin{savenotes}\sphinxattablestart
\sphinxthistablewithglobalstyle
\centering
\begin{tabular}[t]{|\X{1}{4}|\X{1}{4}|\X{1}{4}|\X{1}{4}|}
\sphinxtoprule
\sphinxstyletheadfamily 
\sphinxAtStartPar
ライト
&\sphinxstyletheadfamily 
\sphinxAtStartPar
スタンダード
&\sphinxstyletheadfamily 
\sphinxAtStartPar
プロ
&\sphinxstyletheadfamily 
\sphinxAtStartPar
マスター
\\
\sphinxmidrule
\sphinxtableatstartofbodyhook
\sphinxAtStartPar
x
&
\sphinxAtStartPar
x
&
\sphinxAtStartPar
◯
&
\sphinxAtStartPar
◯
\\
\sphinxbottomrule
\end{tabular}
\sphinxtableafterendhook\par
\sphinxattableend\end{savenotes}

\sphinxAtStartPar
\sphinxstylestrong{トリガー:} 直接

\sphinxAtStartPar
\sphinxstylestrong{スピード:} 通常

\sphinxAtStartPar
\sphinxstylestrong{タイミング:} クイック

\sphinxAtStartPar
\sphinxstylestrong{キーカード:} 同じ数字を２枚

\sphinxAtStartPar
\sphinxstylestrong{対象:}

\sphinxAtStartPar
キャラクター1体

\sphinxAtStartPar
\sphinxstylestrong{効果:}
\begin{enumerate}
\sphinxsetlistlabels{\arabic}{enumi}{enumii}{}{.}%
\item {} 
\sphinxAtStartPar
必要であれば、対象のキャラクターをチャージ状態またはドライブ状態にする。

\item {} 
\sphinxAtStartPar
対象が兵士の場合、兵士を防壁にする。兵士の時に受けた効果、能力は無くなる。兵士が複数のカードから成る場合、1枚ずつ防壁にする。防壁の置き方は「防壁の置き方」参照

\item {} 
\sphinxAtStartPar
対象が防壁の場合、防壁を兵士にする。防壁の時に受けた効果、能力は無くなる。このターンに場に出た防壁を兵士にする場合、その兵士はこのターンに出た兵士と同様に扱う。

\item {} 
\sphinxAtStartPar
対象がアタッカーもしくは、ブロッカーの場合、それを解除する。

\end{enumerate}

\sphinxAtStartPar
\sphinxstylestrong{導入バージョン:}  v0.1

\sphinxAtStartPar
\sphinxstylestrong{更新バージョン:}  v8.0,v8.1


\bigskip\hrule\bigskip



\paragraph{帰還 (速攻魔法)}
\label{\detokenize{auto/actionlist:act-unsummons}}\label{\detokenize{auto/actionlist:id30}}
\sphinxAtStartPar
\sphinxstylestrong{フォーマット:}


\begin{savenotes}\sphinxattablestart
\sphinxthistablewithglobalstyle
\centering
\begin{tabular}[t]{|\X{1}{4}|\X{1}{4}|\X{1}{4}|\X{1}{4}|}
\sphinxtoprule
\sphinxstyletheadfamily 
\sphinxAtStartPar
ライト
&\sphinxstyletheadfamily 
\sphinxAtStartPar
スタンダード
&\sphinxstyletheadfamily 
\sphinxAtStartPar
プロ
&\sphinxstyletheadfamily 
\sphinxAtStartPar
マスター
\\
\sphinxmidrule
\sphinxtableatstartofbodyhook
\sphinxAtStartPar
x
&
\sphinxAtStartPar
◯
&
\sphinxAtStartPar
◯
&
\sphinxAtStartPar
◯
\\
\sphinxbottomrule
\end{tabular}
\sphinxtableafterendhook\par
\sphinxattableend\end{savenotes}

\sphinxAtStartPar
\sphinxstylestrong{トリガー:} 直接

\sphinxAtStartPar
\sphinxstylestrong{スピード:} 通常

\sphinxAtStartPar
\sphinxstylestrong{タイミング:} クイック

\sphinxAtStartPar
\sphinxstylestrong{キーカード:} 同じスートを2枚

\sphinxAtStartPar
\sphinxstylestrong{コスト:} B

\sphinxAtStartPar
\sphinxstylestrong{対象:}

\sphinxAtStartPar
自分のキャラクター1体

\sphinxAtStartPar
\sphinxstylestrong{効果:}
\begin{enumerate}
\sphinxsetlistlabels{\arabic}{enumi}{enumii}{}{.}%
\item {} 
\sphinxAtStartPar
対象のキャラクターがチャージ状態の場合、対象のキャラクターを手札に戻す。

\item {} 
\sphinxAtStartPar
キーカードを手札に戻す。

\end{enumerate}

\sphinxAtStartPar
\sphinxstylestrong{導入バージョン:}  v1.0

\sphinxAtStartPar
\sphinxstylestrong{更新バージョン:}  v8.0,v8.1


\bigskip\hrule\bigskip



\subsubsection{通常魔法}
\label{\detokenize{auto/actionlist:id31}}

\paragraph{防壁破壊 (通常魔法)}
\label{\detokenize{auto/actionlist:act-destroybulwark}}\label{\detokenize{auto/actionlist:id32}}
\sphinxAtStartPar
\sphinxstylestrong{フォーマット:}


\begin{savenotes}\sphinxattablestart
\sphinxthistablewithglobalstyle
\centering
\begin{tabular}[t]{|\X{1}{4}|\X{1}{4}|\X{1}{4}|\X{1}{4}|}
\sphinxtoprule
\sphinxstyletheadfamily 
\sphinxAtStartPar
ライト
&\sphinxstyletheadfamily 
\sphinxAtStartPar
スタンダード
&\sphinxstyletheadfamily 
\sphinxAtStartPar
プロ
&\sphinxstyletheadfamily 
\sphinxAtStartPar
マスター
\\
\sphinxmidrule
\sphinxtableatstartofbodyhook
\sphinxAtStartPar
◯
&
\sphinxAtStartPar
◯
&
\sphinxAtStartPar
◯
&
\sphinxAtStartPar
◯
\\
\sphinxbottomrule
\end{tabular}
\sphinxtableafterendhook\par
\sphinxattableend\end{savenotes}

\sphinxAtStartPar
\sphinxstylestrong{トリガー:} 直接

\sphinxAtStartPar
\sphinxstylestrong{スピード:} 通常

\sphinxAtStartPar
\sphinxstylestrong{タイミング:} メイン

\sphinxAtStartPar
\sphinxstylestrong{キーカード:} {\normalsize $\heartsuit$} A〜K と {\normalsize $\diamondsuit$} A〜K

\sphinxAtStartPar
\sphinxstylestrong{対象:}

\sphinxAtStartPar
防壁1体

\sphinxAtStartPar
\sphinxstylestrong{効果:}

\sphinxAtStartPar
対象の防壁を場から墓地に移す。

\sphinxAtStartPar
\sphinxstylestrong{ストーリー:}

\sphinxAtStartPar
防壁はとても後ろ向きで心が弱い。愛{\normalsize $\heartsuit$} と金{\normalsize $\diamondsuit$} に目移りして簡単に籠絡してしまうことも。

\sphinxAtStartPar
\sphinxstylestrong{フレーバー:}

\sphinxAtStartPar
\sphinxstyleemphasis{「あれ？いつの間にかなくなってるよ？」}

\sphinxAtStartPar
\textasciitilde{} かくれんぼしていた子供より \textasciitilde{}

\sphinxAtStartPar
\sphinxstylestrong{導入バージョン:}  v0.2

\sphinxAtStartPar
\sphinxstylestrong{更新バージョン:}  v8.0,v8.1


\bigskip\hrule\bigskip



\paragraph{投擲 (通常魔法)}
\label{\detokenize{auto/actionlist:act-throwing}}\label{\detokenize{auto/actionlist:id33}}
\sphinxAtStartPar
\sphinxstylestrong{フォーマット:}


\begin{savenotes}\sphinxattablestart
\sphinxthistablewithglobalstyle
\centering
\begin{tabular}[t]{|\X{1}{4}|\X{1}{4}|\X{1}{4}|\X{1}{4}|}
\sphinxtoprule
\sphinxstyletheadfamily 
\sphinxAtStartPar
ライト
&\sphinxstyletheadfamily 
\sphinxAtStartPar
スタンダード
&\sphinxstyletheadfamily 
\sphinxAtStartPar
プロ
&\sphinxstyletheadfamily 
\sphinxAtStartPar
マスター
\\
\sphinxmidrule
\sphinxtableatstartofbodyhook
\sphinxAtStartPar
◯
&
\sphinxAtStartPar
◯
&
\sphinxAtStartPar
◯
&
\sphinxAtStartPar
◯
\\
\sphinxbottomrule
\end{tabular}
\sphinxtableafterendhook\par
\sphinxattableend\end{savenotes}

\sphinxAtStartPar
\sphinxstylestrong{トリガー:} 直接

\sphinxAtStartPar
\sphinxstylestrong{スピード:} 通常

\sphinxAtStartPar
\sphinxstylestrong{タイミング:} メイン

\sphinxAtStartPar
\sphinxstylestrong{キーカード:} {\normalsize $\spadesuit$} A〜K と {\normalsize $\clubsuit$} A〜K

\sphinxAtStartPar
\sphinxstylestrong{対象:}

\sphinxAtStartPar
対戦相手1人

\sphinxAtStartPar
\sphinxstylestrong{効果:}

\sphinxAtStartPar
対象の対戦相手にX点のダメージを与える。Xはキーカードの{\normalsize $\spadesuit$} カードの数字に等しい。

\sphinxAtStartPar
\sphinxstylestrong{ストーリー:}

\sphinxAtStartPar
知恵{\normalsize $\clubsuit$} と殺意{\normalsize $\spadesuit$} をもってすれば兵士を出さずとも飛び道具で相手にダメージを負わすことができるでしょう。とにかく投げればいいのです。

\sphinxAtStartPar
\sphinxstylestrong{フレーバー:}

\sphinxAtStartPar
\sphinxstyleemphasis{「この戦法は戦い方を一変させた。なぜなら兵がいなくても勝てるからだ。」}

\sphinxAtStartPar
\textasciitilde{} 再現者の心得第十一章より \textasciitilde{}

\sphinxAtStartPar
\sphinxstylestrong{導入バージョン:}  v0.1

\sphinxAtStartPar
\sphinxstylestrong{更新バージョン:}  v8.0,v8.1


\bigskip\hrule\bigskip



\paragraph{死の槍 (通常魔法)}
\label{\detokenize{auto/actionlist:act-deathlance}}\label{\detokenize{auto/actionlist:id34}}
\sphinxAtStartPar
\sphinxstylestrong{フォーマット:}


\begin{savenotes}\sphinxattablestart
\sphinxthistablewithglobalstyle
\centering
\begin{tabular}[t]{|\X{1}{4}|\X{1}{4}|\X{1}{4}|\X{1}{4}|}
\sphinxtoprule
\sphinxstyletheadfamily 
\sphinxAtStartPar
ライト
&\sphinxstyletheadfamily 
\sphinxAtStartPar
スタンダード
&\sphinxstyletheadfamily 
\sphinxAtStartPar
プロ
&\sphinxstyletheadfamily 
\sphinxAtStartPar
マスター
\\
\sphinxmidrule
\sphinxtableatstartofbodyhook
\sphinxAtStartPar
x
&
\sphinxAtStartPar
◯
&
\sphinxAtStartPar
◯
&
\sphinxAtStartPar
◯
\\
\sphinxbottomrule
\end{tabular}
\sphinxtableafterendhook\par
\sphinxattableend\end{savenotes}

\sphinxAtStartPar
\sphinxstylestrong{トリガー:} 直接

\sphinxAtStartPar
\sphinxstylestrong{スピード:} 通常

\sphinxAtStartPar
\sphinxstylestrong{タイミング:} メイン

\sphinxAtStartPar
\sphinxstylestrong{キーカード:} {\normalsize $\spadesuit$} A〜K と {\normalsize $\diamondsuit$} A〜K

\sphinxAtStartPar
\sphinxstylestrong{対象:}

\sphinxAtStartPar
兵士1体

\sphinxAtStartPar
\sphinxstylestrong{効果:}

\sphinxAtStartPar
対象の兵士のサイズが0以外かつ、キーカードの{\normalsize $\diamondsuit$} カードの数字で割切れる場合、次を行う。
\begin{enumerate}
\sphinxsetlistlabels{\arabic}{enumi}{enumii}{}{.}%
\item {} 
\sphinxAtStartPar
対象の兵士をオーナーのライフの一番上に裏向きで移す。兵士が複数のカードから成る場合、任意の順でライフの一番上に裏向きで移す。

\item {} 
\sphinxAtStartPar
対象の兵士のオーナーにX点のダメージを与える。Xはキーカードの{\normalsize $\spadesuit$} カードの数字に等しい。

\end{enumerate}

\sphinxAtStartPar
\sphinxstylestrong{ストーリー:}

\sphinxAtStartPar
死{\normalsize $\spadesuit$} の槍{\normalsize $\diamondsuit$} 。そのままだな。

\sphinxAtStartPar
\sphinxstylestrong{フレーバー:}

\sphinxAtStartPar
\sphinxstyleemphasis{「この槍は非常に厄介な槍だ。刺さったら急に自国に戻り力尽きるまで暴れ回る。迷惑にも程がある。」}

\sphinxAtStartPar
\textasciitilde{} 再現者の心得第十二章より \textasciitilde{}

\sphinxAtStartPar
\sphinxstylestrong{導入バージョン:}  v2.0

\sphinxAtStartPar
\sphinxstylestrong{更新バージョン:}  v3.0,8.0,v8.1


\bigskip\hrule\bigskip



\paragraph{防壁補充 (通常魔法)}
\label{\detokenize{auto/actionlist:act-addbulwark}}\label{\detokenize{auto/actionlist:id35}}
\sphinxAtStartPar
\sphinxstylestrong{フォーマット:}


\begin{savenotes}\sphinxattablestart
\sphinxthistablewithglobalstyle
\centering
\begin{tabular}[t]{|\X{1}{4}|\X{1}{4}|\X{1}{4}|\X{1}{4}|}
\sphinxtoprule
\sphinxstyletheadfamily 
\sphinxAtStartPar
ライト
&\sphinxstyletheadfamily 
\sphinxAtStartPar
スタンダード
&\sphinxstyletheadfamily 
\sphinxAtStartPar
プロ
&\sphinxstyletheadfamily 
\sphinxAtStartPar
マスター
\\
\sphinxmidrule
\sphinxtableatstartofbodyhook
\sphinxAtStartPar
x
&
\sphinxAtStartPar
◯
&
\sphinxAtStartPar
◯
&
\sphinxAtStartPar
◯
\\
\sphinxbottomrule
\end{tabular}
\sphinxtableafterendhook\par
\sphinxattableend\end{savenotes}

\sphinxAtStartPar
\sphinxstylestrong{トリガー:} 直接

\sphinxAtStartPar
\sphinxstylestrong{スピード:} 通常

\sphinxAtStartPar
\sphinxstylestrong{タイミング:} メイン

\sphinxAtStartPar
\sphinxstylestrong{キーカード:} {\normalsize $\heartsuit$} A〜K と {\normalsize $\clubsuit$} A〜K

\sphinxAtStartPar
\sphinxstylestrong{効果:}

\sphinxAtStartPar
自分のライフの一番上から1枚を防壁として裏向きかつチャージ状態で場に出す。もしくは、自分のライフの一番上から2枚を防壁として裏向きかつドライブ状態で場に出す。防壁の能力はキャラクターリスト参照。

\sphinxAtStartPar
防壁の置き方は「防壁の置き方」参照

\sphinxAtStartPar
\sphinxstylestrong{ストーリー:}

\sphinxAtStartPar
愛する人{\normalsize $\heartsuit$} を守るため知恵{\normalsize $\clubsuit$} を絞った結果守りを固めることにした。

\sphinxAtStartPar
\sphinxstylestrong{フレーバー:}

\sphinxAtStartPar
\sphinxstyleemphasis{「守れ！守れ！ひたすら守れ！防御は最大の防御？なり！」}

\sphinxAtStartPar
\textasciitilde{} 無能な指揮官より \textasciitilde{}

\sphinxAtStartPar
\sphinxstylestrong{導入バージョン:}  v3.0

\sphinxAtStartPar
\sphinxstylestrong{更新バージョン:}  v8.0


\bigskip\hrule\bigskip



\paragraph{リアニメイト (通常魔法)}
\label{\detokenize{auto/actionlist:act-reanimate}}\label{\detokenize{auto/actionlist:id36}}
\sphinxAtStartPar
\sphinxstylestrong{フォーマット:}


\begin{savenotes}\sphinxattablestart
\sphinxthistablewithglobalstyle
\centering
\begin{tabular}[t]{|\X{1}{4}|\X{1}{4}|\X{1}{4}|\X{1}{4}|}
\sphinxtoprule
\sphinxstyletheadfamily 
\sphinxAtStartPar
ライト
&\sphinxstyletheadfamily 
\sphinxAtStartPar
スタンダード
&\sphinxstyletheadfamily 
\sphinxAtStartPar
プロ
&\sphinxstyletheadfamily 
\sphinxAtStartPar
マスター
\\
\sphinxmidrule
\sphinxtableatstartofbodyhook
\sphinxAtStartPar
x
&
\sphinxAtStartPar
◯
&
\sphinxAtStartPar
◯
&
\sphinxAtStartPar
◯
\\
\sphinxbottomrule
\end{tabular}
\sphinxtableafterendhook\par
\sphinxattableend\end{savenotes}

\sphinxAtStartPar
\sphinxstylestrong{トリガー:} 直接

\sphinxAtStartPar
\sphinxstylestrong{スピード:} 通常

\sphinxAtStartPar
\sphinxstylestrong{タイミング:} メイン

\sphinxAtStartPar
\sphinxstylestrong{キーカード:} {\normalsize $\spadesuit$} A〜K と {\normalsize $\heartsuit$} A〜K

\sphinxAtStartPar
\sphinxstylestrong{対象:}

\sphinxAtStartPar
自分のキャラクター1体

\sphinxAtStartPar
\sphinxstylestrong{効果:}

\sphinxAtStartPar
自分の墓地にあるカード1枚を選ぶ。対象のキャラクターを墓地に移せた場合、選んだカードを兵士として表向きかつチャージ状態で場に出す。移せない場合、選んだカードを墓地に戻す。

\sphinxAtStartPar
\sphinxstylestrong{導入バージョン:}  v3.0

\sphinxAtStartPar
\sphinxstylestrong{更新バージョン:}  v8.0,v8.1


\bigskip\hrule\bigskip



\paragraph{ハンデス (通常魔法)}
\label{\detokenize{auto/actionlist:act-handeth}}\label{\detokenize{auto/actionlist:id37}}
\sphinxAtStartPar
\sphinxstylestrong{フォーマット:}


\begin{savenotes}\sphinxattablestart
\sphinxthistablewithglobalstyle
\centering
\begin{tabular}[t]{|\X{1}{4}|\X{1}{4}|\X{1}{4}|\X{1}{4}|}
\sphinxtoprule
\sphinxstyletheadfamily 
\sphinxAtStartPar
ライト
&\sphinxstyletheadfamily 
\sphinxAtStartPar
スタンダード
&\sphinxstyletheadfamily 
\sphinxAtStartPar
プロ
&\sphinxstyletheadfamily 
\sphinxAtStartPar
マスター
\\
\sphinxmidrule
\sphinxtableatstartofbodyhook
\sphinxAtStartPar
x
&
\sphinxAtStartPar
◯
&
\sphinxAtStartPar
◯
&
\sphinxAtStartPar
◯
\\
\sphinxbottomrule
\end{tabular}
\sphinxtableafterendhook\par
\sphinxattableend\end{savenotes}

\sphinxAtStartPar
\sphinxstylestrong{トリガー:} 直接

\sphinxAtStartPar
\sphinxstylestrong{スピード:} 通常

\sphinxAtStartPar
\sphinxstylestrong{タイミング:} メイン

\sphinxAtStartPar
\sphinxstylestrong{キーカード:} {\normalsize $\diamondsuit$} A〜K と {\normalsize $\clubsuit$} A〜K

\sphinxAtStartPar
\sphinxstylestrong{対象:}

\sphinxAtStartPar
対戦相手1人

\sphinxAtStartPar
\sphinxstylestrong{効果:}

\sphinxAtStartPar
対戦相手の手札を見て1枚カードを指定する。対戦相手は指定されたカードを手札から捨てる。

\sphinxAtStartPar
\sphinxstylestrong{導入バージョン:}  v3.0

\sphinxAtStartPar
\sphinxstylestrong{更新バージョン:}  v8.0,v8.1


\bigskip\hrule\bigskip



\paragraph{フォース (通常魔法)}
\label{\detokenize{auto/actionlist:act-force}}\label{\detokenize{auto/actionlist:id38}}
\sphinxAtStartPar
\sphinxstylestrong{フォーマット:}


\begin{savenotes}\sphinxattablestart
\sphinxthistablewithglobalstyle
\centering
\begin{tabular}[t]{|\X{1}{4}|\X{1}{4}|\X{1}{4}|\X{1}{4}|}
\sphinxtoprule
\sphinxstyletheadfamily 
\sphinxAtStartPar
ライト
&\sphinxstyletheadfamily 
\sphinxAtStartPar
スタンダード
&\sphinxstyletheadfamily 
\sphinxAtStartPar
プロ
&\sphinxstyletheadfamily 
\sphinxAtStartPar
マスター
\\
\sphinxmidrule
\sphinxtableatstartofbodyhook
\sphinxAtStartPar
x
&
\sphinxAtStartPar
x
&
\sphinxAtStartPar
x
&
\sphinxAtStartPar
◯
\\
\sphinxbottomrule
\end{tabular}
\sphinxtableafterendhook\par
\sphinxattableend\end{savenotes}

\sphinxAtStartPar
\sphinxstylestrong{トリガー:} 直接

\sphinxAtStartPar
\sphinxstylestrong{スピード:} 通常

\sphinxAtStartPar
\sphinxstylestrong{タイミング:} メイン

\sphinxAtStartPar
\sphinxstylestrong{キーカード:} {\normalsize $\heartsuit$} A〜10 を 2枚

\sphinxAtStartPar
\sphinxstylestrong{コスト:} BB

\sphinxAtStartPar
\sphinxstylestrong{特記事項:} ※このアクションはカウンターアクションの対象にならない。

\sphinxAtStartPar
\sphinxstylestrong{効果:}

\sphinxAtStartPar
このターンが終わるまで自分の兵士全ての数字は、キーカードの合計値分加算される。

\sphinxAtStartPar
その印としてフォグにキーカードを置き、フォースとする。

\sphinxAtStartPar
\sphinxstylestrong{導入バージョン:}  v3.0

\sphinxAtStartPar
\sphinxstylestrong{更新バージョン:}  v8.0


\bigskip\hrule\bigskip



\paragraph{剣の雨 (通常魔法)}
\label{\detokenize{auto/actionlist:act-swordrain}}\label{\detokenize{auto/actionlist:id39}}
\sphinxAtStartPar
\sphinxstylestrong{フォーマット:}


\begin{savenotes}\sphinxattablestart
\sphinxthistablewithglobalstyle
\centering
\begin{tabular}[t]{|\X{1}{4}|\X{1}{4}|\X{1}{4}|\X{1}{4}|}
\sphinxtoprule
\sphinxstyletheadfamily 
\sphinxAtStartPar
ライト
&\sphinxstyletheadfamily 
\sphinxAtStartPar
スタンダード
&\sphinxstyletheadfamily 
\sphinxAtStartPar
プロ
&\sphinxstyletheadfamily 
\sphinxAtStartPar
マスター
\\
\sphinxmidrule
\sphinxtableatstartofbodyhook
\sphinxAtStartPar
x
&
\sphinxAtStartPar
x
&
\sphinxAtStartPar
x
&
\sphinxAtStartPar
◯
\\
\sphinxbottomrule
\end{tabular}
\sphinxtableafterendhook\par
\sphinxattableend\end{savenotes}

\sphinxAtStartPar
\sphinxstylestrong{トリガー:} 直接

\sphinxAtStartPar
\sphinxstylestrong{スピード:} 通常

\sphinxAtStartPar
\sphinxstylestrong{タイミング:} メイン

\sphinxAtStartPar
\sphinxstylestrong{キーカード:} {\normalsize $\spadesuit$} A〜10 を 2枚

\sphinxAtStartPar
\sphinxstylestrong{コスト:} BB

\sphinxAtStartPar
\sphinxstylestrong{特記事項:} ※このアクションはカウンターアクションの対象にならない。

\sphinxAtStartPar
\sphinxstylestrong{効果:}

\sphinxAtStartPar
サイズがキーカードの合計値以下の全ての兵士を墓地に移す。

\sphinxAtStartPar
\sphinxstylestrong{ストーリー:}

\sphinxAtStartPar
死{\normalsize $\spadesuit$} と剣{\normalsize $\spadesuit$} が合わさり雨のように戦場に降り注いだ。

\sphinxAtStartPar
\sphinxstylestrong{フレーバー:}

\sphinxAtStartPar
\sphinxstyleemphasis{「え！助かる訳ないじゃん！！」}

\sphinxAtStartPar
\textasciitilde{} 平和主義の{\normalsize $\clubsuit$} 3一般兵より \textasciitilde{}

\sphinxAtStartPar
\sphinxstylestrong{導入バージョン:}  v3.0

\sphinxAtStartPar
\sphinxstylestrong{更新バージョン:}  v8.0


\bigskip\hrule\bigskip



\paragraph{徴募 (通常魔法)}
\label{\detokenize{auto/actionlist:act-recruit}}\label{\detokenize{auto/actionlist:id40}}
\sphinxAtStartPar
\sphinxstylestrong{フォーマット:}


\begin{savenotes}\sphinxattablestart
\sphinxthistablewithglobalstyle
\centering
\begin{tabular}[t]{|\X{1}{4}|\X{1}{4}|\X{1}{4}|\X{1}{4}|}
\sphinxtoprule
\sphinxstyletheadfamily 
\sphinxAtStartPar
ライト
&\sphinxstyletheadfamily 
\sphinxAtStartPar
スタンダード
&\sphinxstyletheadfamily 
\sphinxAtStartPar
プロ
&\sphinxstyletheadfamily 
\sphinxAtStartPar
マスター
\\
\sphinxmidrule
\sphinxtableatstartofbodyhook
\sphinxAtStartPar
x
&
\sphinxAtStartPar
x
&
\sphinxAtStartPar
x
&
\sphinxAtStartPar
◯
\\
\sphinxbottomrule
\end{tabular}
\sphinxtableafterendhook\par
\sphinxattableend\end{savenotes}

\sphinxAtStartPar
\sphinxstylestrong{トリガー:} 直接

\sphinxAtStartPar
\sphinxstylestrong{スピード:} 通常

\sphinxAtStartPar
\sphinxstylestrong{タイミング:} メイン

\sphinxAtStartPar
\sphinxstylestrong{キーカード:} {\normalsize $\diamondsuit$} A〜10 を 2枚

\sphinxAtStartPar
\sphinxstylestrong{コスト:} BB

\sphinxAtStartPar
\sphinxstylestrong{特記事項:} ※このアクションはカウンターアクションの対象にならない。

\sphinxAtStartPar
\sphinxstylestrong{効果:}
\begin{enumerate}
\sphinxsetlistlabels{\arabic}{enumi}{enumii}{}{.}%
\item {} 
\sphinxAtStartPar
ライフの一番上から4枚めくり、キーカードの合計値以下のカードを兵士としてドライブ状態で場に出す。

\item {} 
\sphinxAtStartPar
残りのカードをライフの一番下に好きな順で移す。

\end{enumerate}

\sphinxAtStartPar
\sphinxstylestrong{ストーリー:}

\sphinxAtStartPar
金{\normalsize $\diamondsuit$} が手元にたくさんあれば、戦力強化が定石だろう。

\sphinxAtStartPar
\sphinxstylestrong{フレーバー:}

\sphinxAtStartPar
\sphinxstyleemphasis{「君はどう？兵士になる？...あっそう。」
「君はどう？兵士になる？...あっそう。」
「はぁ。金がなけれりゃ誰も来やしねぇ」}

\sphinxAtStartPar
\textasciitilde{} 孤独な徴募兵より \textasciitilde{}

\sphinxAtStartPar
\sphinxstylestrong{導入バージョン:}  v3.0

\sphinxAtStartPar
\sphinxstylestrong{更新バージョン:}  v8.0


\bigskip\hrule\bigskip



\paragraph{奇襲 (通常魔法)}
\label{\detokenize{auto/actionlist:act-surprise}}\label{\detokenize{auto/actionlist:id41}}
\sphinxAtStartPar
\sphinxstylestrong{フォーマット:}


\begin{savenotes}\sphinxattablestart
\sphinxthistablewithglobalstyle
\centering
\begin{tabular}[t]{|\X{1}{4}|\X{1}{4}|\X{1}{4}|\X{1}{4}|}
\sphinxtoprule
\sphinxstyletheadfamily 
\sphinxAtStartPar
ライト
&\sphinxstyletheadfamily 
\sphinxAtStartPar
スタンダード
&\sphinxstyletheadfamily 
\sphinxAtStartPar
プロ
&\sphinxstyletheadfamily 
\sphinxAtStartPar
マスター
\\
\sphinxmidrule
\sphinxtableatstartofbodyhook
\sphinxAtStartPar
x
&
\sphinxAtStartPar
x
&
\sphinxAtStartPar
x
&
\sphinxAtStartPar
◯
\\
\sphinxbottomrule
\end{tabular}
\sphinxtableafterendhook\par
\sphinxattableend\end{savenotes}

\sphinxAtStartPar
\sphinxstylestrong{トリガー:} 直接

\sphinxAtStartPar
\sphinxstylestrong{スピード:} 通常

\sphinxAtStartPar
\sphinxstylestrong{タイミング:} メイン

\sphinxAtStartPar
\sphinxstylestrong{キーカード:} {\normalsize $\clubsuit$} A〜10 を 2枚

\sphinxAtStartPar
\sphinxstylestrong{コスト:} BB

\sphinxAtStartPar
\sphinxstylestrong{特記事項:} ※このアクションはカウンターアクションの対象にならない。

\sphinxAtStartPar
\sphinxstylestrong{効果:}
\begin{enumerate}
\sphinxsetlistlabels{\arabic}{enumi}{enumii}{}{.}%
\item {} 
\sphinxAtStartPar
自分の場にいる防壁を全て兵士にする。このターンに場に出た防壁を兵士にする場合、その兵士はこのターンに出た兵士と同様に扱う。

\item {} 
\sphinxAtStartPar
自分の場にいる全ての兵士をチャージ状態にする。

\end{enumerate}

\sphinxAtStartPar
\sphinxstylestrong{導入バージョン:}  v3.0

\sphinxAtStartPar
\sphinxstylestrong{更新バージョン:}  v8.0


\subsection{キャラクターリスト}
\label{\detokenize{auto/actionlist:characterlist-act}}\label{\detokenize{auto/actionlist:id42}}
\begin{sphinxShadowBox}
\sphinxstyletopictitle{目次}
\begin{itemize}
\item {} 
\sphinxAtStartPar
\phantomsection\label{\detokenize{auto/actionlist:id99}}{\hyperref[\detokenize{auto/actionlist:id44}]{\sphinxcrossref{兵士}}}
\begin{itemize}
\item {} 
\sphinxAtStartPar
\phantomsection\label{\detokenize{auto/actionlist:id100}}{\hyperref[\detokenize{auto/actionlist:char-soldier}]{\sphinxcrossref{一般兵 (兵士)}}}

\item {} 
\sphinxAtStartPar
\phantomsection\label{\detokenize{auto/actionlist:id101}}{\hyperref[\detokenize{auto/actionlist:char-hero}]{\sphinxcrossref{英雄 (兵士)}}}

\item {} 
\sphinxAtStartPar
\phantomsection\label{\detokenize{auto/actionlist:id102}}{\hyperref[\detokenize{auto/actionlist:char-ace}]{\sphinxcrossref{エース (兵士)}}}

\item {} 
\sphinxAtStartPar
\phantomsection\label{\detokenize{auto/actionlist:id103}}{\hyperref[\detokenize{auto/actionlist:char-magician}]{\sphinxcrossref{魔術士 (兵士)}}}

\item {} 
\sphinxAtStartPar
\phantomsection\label{\detokenize{auto/actionlist:id104}}{\hyperref[\detokenize{auto/actionlist:char-armedsoldier}]{\sphinxcrossref{装備兵 (兵士)}}}

\end{itemize}

\item {} 
\sphinxAtStartPar
\phantomsection\label{\detokenize{auto/actionlist:id105}}{\hyperref[\detokenize{auto/actionlist:id50}]{\sphinxcrossref{防壁}}}
\begin{itemize}
\item {} 
\sphinxAtStartPar
\phantomsection\label{\detokenize{auto/actionlist:id106}}{\hyperref[\detokenize{auto/actionlist:char-bulwark}]{\sphinxcrossref{防壁 (防壁)}}}

\end{itemize}

\end{itemize}
\end{sphinxShadowBox}


\bigskip\hrule\bigskip



\subsubsection{兵士}
\label{\detokenize{auto/actionlist:id44}}

\paragraph{一般兵 (兵士)}
\label{\detokenize{auto/actionlist:char-soldier}}\label{\detokenize{auto/actionlist:id45}}
\sphinxAtStartPar
\sphinxstylestrong{フォーマット:}


\begin{savenotes}\sphinxattablestart
\sphinxthistablewithglobalstyle
\centering
\begin{tabular}[t]{|\X{1}{4}|\X{1}{4}|\X{1}{4}|\X{1}{4}|}
\sphinxtoprule
\sphinxstyletheadfamily 
\sphinxAtStartPar
ライト
&\sphinxstyletheadfamily 
\sphinxAtStartPar
スタンダード
&\sphinxstyletheadfamily 
\sphinxAtStartPar
プロ
&\sphinxstyletheadfamily 
\sphinxAtStartPar
マスター
\\
\sphinxmidrule
\sphinxtableatstartofbodyhook
\sphinxAtStartPar
◯
&
\sphinxAtStartPar
◯
&
\sphinxAtStartPar
◯
&
\sphinxAtStartPar
◯
\\
\sphinxbottomrule
\end{tabular}
\sphinxtableafterendhook\par
\sphinxattableend\end{savenotes}

\sphinxAtStartPar
\sphinxstylestrong{キャラクターキーカード:} 2〜10

\sphinxAtStartPar
\sphinxstylestrong{サイズ:} ユニットカードの数字

\sphinxAtStartPar
\sphinxstylestrong{ラベル:} アタッカー, ブロッカー

\sphinxAtStartPar
\sphinxstylestrong{導入バージョン:}  v1.0

\sphinxAtStartPar
\sphinxstylestrong{更新バージョン:}  v8.0


\bigskip\hrule\bigskip



\paragraph{英雄 (兵士)}
\label{\detokenize{auto/actionlist:char-hero}}\label{\detokenize{auto/actionlist:id46}}
\sphinxAtStartPar
\sphinxstylestrong{フォーマット:}


\begin{savenotes}\sphinxattablestart
\sphinxthistablewithglobalstyle
\centering
\begin{tabular}[t]{|\X{1}{4}|\X{1}{4}|\X{1}{4}|\X{1}{4}|}
\sphinxtoprule
\sphinxstyletheadfamily 
\sphinxAtStartPar
ライト
&\sphinxstyletheadfamily 
\sphinxAtStartPar
スタンダード
&\sphinxstyletheadfamily 
\sphinxAtStartPar
プロ
&\sphinxstyletheadfamily 
\sphinxAtStartPar
マスター
\\
\sphinxmidrule
\sphinxtableatstartofbodyhook
\sphinxAtStartPar
◯
&
\sphinxAtStartPar
◯
&
\sphinxAtStartPar
◯
&
\sphinxAtStartPar
◯
\\
\sphinxbottomrule
\end{tabular}
\sphinxtableafterendhook\par
\sphinxattableend\end{savenotes}

\sphinxAtStartPar
\sphinxstylestrong{キャラクターキーカード:} J〜K

\sphinxAtStartPar
\sphinxstylestrong{サイズ:} ユニットカードがJなら11,Qなら12,Kなら13

\sphinxAtStartPar
\sphinxstylestrong{ラベル:} アタッカー, ブロッカー

\sphinxAtStartPar
\sphinxstylestrong{導入バージョン:}  v1.0

\sphinxAtStartPar
\sphinxstylestrong{更新バージョン:}  v8.0


\bigskip\hrule\bigskip



\paragraph{エース (兵士)}
\label{\detokenize{auto/actionlist:char-ace}}\label{\detokenize{auto/actionlist:id47}}
\sphinxAtStartPar
\sphinxstylestrong{フォーマット:}


\begin{savenotes}\sphinxattablestart
\sphinxthistablewithglobalstyle
\centering
\begin{tabular}[t]{|\X{1}{4}|\X{1}{4}|\X{1}{4}|\X{1}{4}|}
\sphinxtoprule
\sphinxstyletheadfamily 
\sphinxAtStartPar
ライト
&\sphinxstyletheadfamily 
\sphinxAtStartPar
スタンダード
&\sphinxstyletheadfamily 
\sphinxAtStartPar
プロ
&\sphinxstyletheadfamily 
\sphinxAtStartPar
マスター
\\
\sphinxmidrule
\sphinxtableatstartofbodyhook
\sphinxAtStartPar
◯
&
\sphinxAtStartPar
◯
&
\sphinxAtStartPar
◯
&
\sphinxAtStartPar
◯
\\
\sphinxbottomrule
\end{tabular}
\sphinxtableafterendhook\par
\sphinxattableend\end{savenotes}

\sphinxAtStartPar
\sphinxstylestrong{キャラクターキーカード:} A

\sphinxAtStartPar
\sphinxstylestrong{サイズ:} 1

\sphinxAtStartPar
\sphinxstylestrong{ラベル:} アタッカー, ブロッカー, 速攻

\sphinxAtStartPar
\sphinxstylestrong{ストーリー:}

\sphinxAtStartPar
エースは先駆者であり皆の先頭に立ち切り込んで行く力があります。一人では大きなダメージにはなりませんが、大きな力を持っています。あなたもエースがまさにエースであることに気づくでしょう。

\sphinxAtStartPar
\sphinxstylestrong{フレーバー:}

\sphinxAtStartPar
\sphinxstyleemphasis{「俺には、ちびですばしっこいだけにしか見えませんが？」}

\sphinxAtStartPar
\textasciitilde{} 再現者見習いより \textasciitilde{}

\sphinxAtStartPar
\sphinxstylestrong{導入バージョン:}  v1.0

\sphinxAtStartPar
\sphinxstylestrong{更新バージョン:}  v8.0


\bigskip\hrule\bigskip



\paragraph{魔術士 (兵士)}
\label{\detokenize{auto/actionlist:char-magician}}\label{\detokenize{auto/actionlist:id48}}
\sphinxAtStartPar
\sphinxstylestrong{フォーマット:}


\begin{savenotes}\sphinxattablestart
\sphinxthistablewithglobalstyle
\centering
\begin{tabular}[t]{|\X{1}{4}|\X{1}{4}|\X{1}{4}|\X{1}{4}|}
\sphinxtoprule
\sphinxstyletheadfamily 
\sphinxAtStartPar
ライト
&\sphinxstyletheadfamily 
\sphinxAtStartPar
スタンダード
&\sphinxstyletheadfamily 
\sphinxAtStartPar
プロ
&\sphinxstyletheadfamily 
\sphinxAtStartPar
マスター
\\
\sphinxmidrule
\sphinxtableatstartofbodyhook
\sphinxAtStartPar
x
&
\sphinxAtStartPar
◯
&
\sphinxAtStartPar
◯
&
\sphinxAtStartPar
◯
\\
\sphinxbottomrule
\end{tabular}
\sphinxtableafterendhook\par
\sphinxattableend\end{savenotes}

\sphinxAtStartPar
\sphinxstylestrong{キャラクターキーカード:} Joker

\sphinxAtStartPar
\sphinxstylestrong{サイズ:} 0

\sphinxAtStartPar
\sphinxstylestrong{ラベル:} アタッカー, ブロッカー, 速攻

\sphinxAtStartPar
\sphinxstylestrong{能力:}

\sphinxAtStartPar
・魔力増加（このキャラクターが場にいる間、タイプ：速攻魔法のコストDが無しになる。）

\sphinxAtStartPar
\sphinxstylestrong{導入バージョン:}  v3.0

\sphinxAtStartPar
\sphinxstylestrong{更新バージョン:}  v8.0


\bigskip\hrule\bigskip



\paragraph{装備兵 (兵士)}
\label{\detokenize{auto/actionlist:char-armedsoldier}}\label{\detokenize{auto/actionlist:id49}}
\sphinxAtStartPar
\sphinxstylestrong{フォーマット:}


\begin{savenotes}\sphinxattablestart
\sphinxthistablewithglobalstyle
\centering
\begin{tabular}[t]{|\X{1}{4}|\X{1}{4}|\X{1}{4}|\X{1}{4}|}
\sphinxtoprule
\sphinxstyletheadfamily 
\sphinxAtStartPar
ライト
&\sphinxstyletheadfamily 
\sphinxAtStartPar
スタンダード
&\sphinxstyletheadfamily 
\sphinxAtStartPar
プロ
&\sphinxstyletheadfamily 
\sphinxAtStartPar
マスター
\\
\sphinxmidrule
\sphinxtableatstartofbodyhook
\sphinxAtStartPar
◯
&
\sphinxAtStartPar
◯
&
\sphinxAtStartPar
◯
&
\sphinxAtStartPar
◯
\\
\sphinxbottomrule
\end{tabular}
\sphinxtableafterendhook\par
\sphinxattableend\end{savenotes}

\sphinxAtStartPar
\sphinxstylestrong{キャラクターキーカード:} 同じスートを2枚以上

\sphinxAtStartPar
\sphinxstylestrong{サイズ:} ユニットカードの数字の合計

\sphinxAtStartPar
\sphinxstylestrong{ラベル:} アタッカー, ブロッカー, 速攻(キーカードにAが含まれる場合のみ)

\sphinxAtStartPar
\sphinxstylestrong{導入バージョン:}  v5.0

\sphinxAtStartPar
\sphinxstylestrong{更新バージョン:}  v8.0


\bigskip\hrule\bigskip



\subsubsection{防壁}
\label{\detokenize{auto/actionlist:id50}}

\paragraph{防壁 (防壁)}
\label{\detokenize{auto/actionlist:char-bulwark}}\label{\detokenize{auto/actionlist:id51}}
\sphinxAtStartPar
\sphinxstylestrong{フォーマット:}


\begin{savenotes}\sphinxattablestart
\sphinxthistablewithglobalstyle
\centering
\begin{tabular}[t]{|\X{1}{4}|\X{1}{4}|\X{1}{4}|\X{1}{4}|}
\sphinxtoprule
\sphinxstyletheadfamily 
\sphinxAtStartPar
ライト
&\sphinxstyletheadfamily 
\sphinxAtStartPar
スタンダード
&\sphinxstyletheadfamily 
\sphinxAtStartPar
プロ
&\sphinxstyletheadfamily 
\sphinxAtStartPar
マスター
\\
\sphinxmidrule
\sphinxtableatstartofbodyhook
\sphinxAtStartPar
◯
&
\sphinxAtStartPar
◯
&
\sphinxAtStartPar
◯
&
\sphinxAtStartPar
◯
\\
\sphinxbottomrule
\end{tabular}
\sphinxtableafterendhook\par
\sphinxattableend\end{savenotes}

\sphinxAtStartPar
\sphinxstylestrong{キャラクターキーカード:} 全て(裏向き)

\sphinxAtStartPar
\sphinxstylestrong{ラベル:} ブロッカー

\sphinxAtStartPar
\sphinxstylestrong{導入バージョン:}  v1.0

\sphinxAtStartPar
\sphinxstylestrong{更新バージョン:}  v8.0


\subsection{フォグリスト}
\label{\detokenize{auto/actionlist:foglist-act}}\label{\detokenize{auto/actionlist:id52}}
\begin{sphinxShadowBox}
\sphinxstyletopictitle{目次}
\begin{itemize}
\item {} 
\sphinxAtStartPar
\phantomsection\label{\detokenize{auto/actionlist:id107}}{\hyperref[\detokenize{auto/actionlist:id54}]{\sphinxcrossref{滞留魔法}}}
\begin{itemize}
\item {} 
\sphinxAtStartPar
\phantomsection\label{\detokenize{auto/actionlist:id108}}{\hyperref[\detokenize{auto/actionlist:fog-upfog}]{\sphinxcrossref{アップ (滞留魔法)}}}

\item {} 
\sphinxAtStartPar
\phantomsection\label{\detokenize{auto/actionlist:id109}}{\hyperref[\detokenize{auto/actionlist:fog-downfog}]{\sphinxcrossref{ダウン (滞留魔法)}}}

\item {} 
\sphinxAtStartPar
\phantomsection\label{\detokenize{auto/actionlist:id110}}{\hyperref[\detokenize{auto/actionlist:fog-forcefog}]{\sphinxcrossref{フォース (滞留魔法)}}}

\end{itemize}

\end{itemize}
\end{sphinxShadowBox}


\bigskip\hrule\bigskip



\subsubsection{滞留魔法}
\label{\detokenize{auto/actionlist:id54}}

\paragraph{アップ (滞留魔法)}
\label{\detokenize{auto/actionlist:fog-upfog}}\label{\detokenize{auto/actionlist:id55}}
\sphinxAtStartPar
\sphinxstylestrong{フォグキーカード:} {\normalsize $\heartsuit$} A〜10

\sphinxAtStartPar
\sphinxstylestrong{能力:}

\sphinxAtStartPar
対象とした兵士のサイズは、このターンが終わるまでキーカードの数字分加算される。

\sphinxAtStartPar
\sphinxstylestrong{導入バージョン:}  v8.0


\bigskip\hrule\bigskip



\paragraph{ダウン (滞留魔法)}
\label{\detokenize{auto/actionlist:fog-downfog}}\label{\detokenize{auto/actionlist:id56}}
\sphinxAtStartPar
\sphinxstylestrong{フォグキーカード:} {\normalsize $\spadesuit$} A〜10

\sphinxAtStartPar
\sphinxstylestrong{能力:}

\sphinxAtStartPar
対象とした兵士のサイズは、このターンが終わるまでキーカードの数字分減算される。

\sphinxAtStartPar
\sphinxstylestrong{導入バージョン:}  v8.0


\bigskip\hrule\bigskip



\paragraph{フォース (滞留魔法)}
\label{\detokenize{auto/actionlist:fog-forcefog}}\label{\detokenize{auto/actionlist:id57}}
\sphinxAtStartPar
\sphinxstylestrong{フォグキーカード:} {\normalsize $\heartsuit$} A〜10 を 2枚

\sphinxAtStartPar
\sphinxstylestrong{能力:}

\sphinxAtStartPar
このターンが終わるまで自分の兵士全ての数字は、キーカードの合計値分加算される。

\sphinxAtStartPar
\sphinxstylestrong{導入バージョン:}  v8.0

\sphinxstepscope


\chapter{フレーム}
\label{\detokenize{frame/frame:frame-rst}}\label{\detokenize{frame/frame:id1}}\label{\detokenize{frame/frame::doc}}
\index{フレーム@\spxentry{フレーム}}\ignorespaces 

\section{フレームとは}
\label{\detokenize{frame/frame:index-0}}\label{\detokenize{frame/frame:id2}}
\sphinxAtStartPar
フレームとはフォーマットとは違った角度で対戦方法を決める枠組みです。

\sphinxAtStartPar
フレームはフォーマットと異なり、アクション以外にも領域(ゾーン)、ゲームの始め方などについても定義しています。


\section{定義項目}
\label{\detokenize{frame/frame:id3}}
\sphinxAtStartPar
フレームには次の項目が定義されています。
\begin{description}
\sphinxlineitem{配置図}
\sphinxAtStartPar
共通ルールにて定義されていない領域を使う場合やゲームの始め方の補足説明で必要な場合、定義されます。
共通ルールで定義された領域：ライフ、墓地、場、手札、フォグ

\sphinxlineitem{フォーマット}
\sphinxAtStartPar
対応しているフォーマット

\sphinxlineitem{デッキ条件}
\sphinxAtStartPar
デッキのカード構成が54枚（2枚のJokerを含む）でない場合、定義されます。

\sphinxlineitem{ゲームの始め方}
\sphinxAtStartPar
共通ルールに記載されているゲームの始め方（カードの初期配置）を変更する場合、定義されます。

\sphinxlineitem{フレームアクションリスト}
\sphinxAtStartPar
追加で起こせるアクション（フレーム依存アクション）がある場合、定義されます。

\sphinxlineitem{フレームフォグリスト}
\sphinxAtStartPar
追加で起こせるフォグ（フレーム依存フォグ）がある場合、定義されます。

\sphinxlineitem{フレーム切札リスト}
\sphinxAtStartPar
追加で起こせる切札（フレーム依存切札）がある場合、定義されます。

\end{description}


\section{フレーム定義}
\label{\detokenize{frame/frame:id4}}
\sphinxAtStartPar
各フレームの詳細は次の公式フレームを参照してください。

\sphinxstepscope


\subsection{公式フレーム}
\label{\detokenize{auto/framelist:id1}}\label{\detokenize{auto/framelist::doc}}
\sphinxAtStartPar
本ドキュメントでは、BlackPokerのすべてのフレームについて記載します。

\begin{sphinxShadowBox}
\sphinxstyletopictitle{目次}
\begin{itemize}
\item {} 
\sphinxAtStartPar
\phantomsection\label{\detokenize{auto/framelist:id13}}{\hyperref[\detokenize{auto/framelist:frame-entry20}]{\sphinxcrossref{エントリー20}}}

\item {} 
\sphinxAtStartPar
\phantomsection\label{\detokenize{auto/framelist:id14}}{\hyperref[\detokenize{auto/framelist:frame-pack}]{\sphinxcrossref{パック}}}

\item {} 
\sphinxAtStartPar
\phantomsection\label{\detokenize{auto/framelist:id15}}{\hyperref[\detokenize{auto/framelist:frame-rarepack}]{\sphinxcrossref{レアパック}}}

\item {} 
\sphinxAtStartPar
\phantomsection\label{\detokenize{auto/framelist:id16}}{\hyperref[\detokenize{auto/framelist:frame-strategy}]{\sphinxcrossref{ストラテジー}}}

\item {} 
\sphinxAtStartPar
\phantomsection\label{\detokenize{auto/framelist:id17}}{\hyperref[\detokenize{auto/framelist:frame-extra}]{\sphinxcrossref{エクストラ}}}

\end{itemize}
\end{sphinxShadowBox}


\subsubsection{エントリー20}
\label{\detokenize{auto/framelist:frame-entry20}}\label{\detokenize{auto/framelist:id3}}
\sphinxAtStartPar
\sphinxstylestrong{配置図:}

\begin{figure}[htbp]
\centering
\capstart

\noindent\sphinxincludegraphics{{entry20}.pdf}
\caption{配置図\sphinxhyphen{}エントリー20}\label{\detokenize{auto/framelist:id8}}\label{\detokenize{auto/framelist:frame-entry20-image}}\end{figure}

\sphinxAtStartPar
\sphinxstylestrong{フォーマット:}


\begin{savenotes}\sphinxattablestart
\sphinxthistablewithglobalstyle
\centering
\begin{tabular}[t]{|\X{1}{4}|\X{1}{4}|\X{1}{4}|\X{1}{4}|}
\sphinxtoprule
\sphinxstyletheadfamily 
\sphinxAtStartPar
ライト
&\sphinxstyletheadfamily 
\sphinxAtStartPar
スタンダード
&\sphinxstyletheadfamily 
\sphinxAtStartPar
プロ
&\sphinxstyletheadfamily 
\sphinxAtStartPar
マスター
\\
\sphinxmidrule
\sphinxtableatstartofbodyhook
\sphinxAtStartPar
◯
&
\sphinxAtStartPar
x
&
\sphinxAtStartPar
x
&
\sphinxAtStartPar
x
\\
\sphinxbottomrule
\end{tabular}
\sphinxtableafterendhook\par
\sphinxattableend\end{savenotes}

\sphinxAtStartPar
\sphinxstylestrong{デッキ条件:}

\sphinxAtStartPar
次のカードのみで構成された20枚のデッキを使う
\begin{itemize}
\item {} 
\sphinxAtStartPar
{\normalsize $\spadesuit$} A, 2, 3, 4, 5

\item {} 
\sphinxAtStartPar
{\normalsize $\heartsuit$} A, 8, 9, 10, J

\item {} 
\sphinxAtStartPar
{\normalsize $\diamondsuit$} A, 3, 7, 10, Q

\item {} 
\sphinxAtStartPar
{\normalsize $\clubsuit$} A, 5, 6, 10, K

\end{itemize}

\sphinxAtStartPar
\sphinxstylestrong{ゲームの始め方:}
\begin{enumerate}
\sphinxsetlistlabels{\arabic}{enumi}{enumii}{}{.}%
\item {} 
\sphinxAtStartPar
デッキをシャッフルする。

\item {} 
\sphinxAtStartPar
デッキをライフとし、配置図の場所に伏せておく。

\item {} 
\sphinxAtStartPar
ライフの一番上のカード7枚引き手札とする。

\item {} 
\sphinxAtStartPar
\hyperref[\detokenize{common/common:common-gamestart-preset}]{\ref{\detokenize{common/common:common-gamestart-preset}} \nameref{\detokenize{common/common:common-gamestart-preset}}} に従ってプリセットする。

\item {} 
\sphinxAtStartPar
\hyperref[\detokenize{common/common:common-gamestart-first}]{\ref{\detokenize{common/common:common-gamestart-first}} \nameref{\detokenize{common/common:common-gamestart-first}}} に従って先攻プレイヤーを決定する。

\item {} 
\sphinxAtStartPar
\hyperref[\detokenize{common/common:common-gamestart-start}]{\ref{\detokenize{common/common:common-gamestart-start}} \nameref{\detokenize{common/common:common-gamestart-start}}} に従ってゲームを開始する。

\end{enumerate}


\bigskip\hrule\bigskip



\subsubsection{パック}
\label{\detokenize{auto/framelist:frame-pack}}\label{\detokenize{auto/framelist:id4}}
\sphinxAtStartPar
\sphinxstylestrong{配置図:}

\begin{figure}[htbp]
\centering
\capstart

\noindent\sphinxincludegraphics{{pack}.pdf}
\caption{配置図\sphinxhyphen{}パック}\label{\detokenize{auto/framelist:id9}}\label{\detokenize{auto/framelist:frame-pack-image}}\end{figure}

\sphinxAtStartPar
\sphinxstylestrong{フォーマット:}


\begin{savenotes}\sphinxattablestart
\sphinxthistablewithglobalstyle
\centering
\begin{tabular}[t]{|\X{1}{4}|\X{1}{4}|\X{1}{4}|\X{1}{4}|}
\sphinxtoprule
\sphinxstyletheadfamily 
\sphinxAtStartPar
ライト
&\sphinxstyletheadfamily 
\sphinxAtStartPar
スタンダード
&\sphinxstyletheadfamily 
\sphinxAtStartPar
プロ
&\sphinxstyletheadfamily 
\sphinxAtStartPar
マスター
\\
\sphinxmidrule
\sphinxtableatstartofbodyhook
\sphinxAtStartPar
x
&
\sphinxAtStartPar
◯
&
\sphinxAtStartPar
◯
&
\sphinxAtStartPar
◯
\\
\sphinxbottomrule
\end{tabular}
\sphinxtableafterendhook\par
\sphinxattableend\end{savenotes}

\sphinxAtStartPar
\sphinxstylestrong{デッキ条件:}

\sphinxAtStartPar
40枚以上

\sphinxAtStartPar
\sphinxstylestrong{ゲームの始め方:}
\begin{enumerate}
\sphinxsetlistlabels{\arabic}{enumi}{enumii}{}{.}%
\item {} 
\sphinxAtStartPar
デッキをシャッフルする。

\item {} 
\sphinxAtStartPar
デッキの上から14枚を取り除いてパックとし、配置図の場所に伏せておく。

\item {} 
\sphinxAtStartPar
残りのデッキをライフとし配置図の場所に伏せておく。

\item {} 
\sphinxAtStartPar
ライフの一番上のカード7枚引き手札とする。

\item {} 
\sphinxAtStartPar
\hyperref[\detokenize{common/common:common-gamestart-preset}]{\ref{\detokenize{common/common:common-gamestart-preset}} \nameref{\detokenize{common/common:common-gamestart-preset}}} に従ってプリセットする。

\item {} 
\sphinxAtStartPar
\hyperref[\detokenize{common/common:common-gamestart-first}]{\ref{\detokenize{common/common:common-gamestart-first}} \nameref{\detokenize{common/common:common-gamestart-first}}} に従って先攻プレイヤーを決定する。

\item {} 
\sphinxAtStartPar
\hyperref[\detokenize{common/common:common-gamestart-start}]{\ref{\detokenize{common/common:common-gamestart-start}} \nameref{\detokenize{common/common:common-gamestart-start}}} に従ってゲームを開始する。

\end{enumerate}

\sphinxAtStartPar
\sphinxstylestrong{アクションリスト:}

\sphinxAtStartPar
次のアクションが追加で利用可能です。
\begin{itemize}
\item {} 
\sphinxAtStartPar
パック操作アクション

\end{itemize}

\sphinxAtStartPar
詳細は次を参照してください。

\sphinxAtStartPar
\hyperref[\detokenize{auto/frameActionlist:act-act-frame}]{\ref{\detokenize{auto/frameActionlist:act-act-frame}} \nameref{\detokenize{auto/frameActionlist:act-act-frame}}}

\sphinxAtStartPar
\sphinxstylestrong{公開レベル:}

\sphinxAtStartPar
このフレームで定義されている領域の公開レベルは次の通りです。
\begin{description}
\sphinxlineitem{パック}
\begin{DUlineblock}{0em}
\item[] 完全公開：開封済み・未開封
\item[] 個人公開：開封後のパックの中身
\item[] 非公開：未開封のパックの中身
\end{DUlineblock}

\end{description}


\bigskip\hrule\bigskip



\subsubsection{レアパック}
\label{\detokenize{auto/framelist:frame-rarepack}}\label{\detokenize{auto/framelist:id5}}
\sphinxAtStartPar
\sphinxstylestrong{配置図:}

\begin{figure}[htbp]
\centering
\capstart

\noindent\sphinxincludegraphics{{rarePack}.pdf}
\caption{配置図\sphinxhyphen{}レアパック}\label{\detokenize{auto/framelist:id10}}\label{\detokenize{auto/framelist:frame-rarepack-image}}\end{figure}

\sphinxAtStartPar
\sphinxstylestrong{フォーマット:}


\begin{savenotes}\sphinxattablestart
\sphinxthistablewithglobalstyle
\centering
\begin{tabular}[t]{|\X{1}{4}|\X{1}{4}|\X{1}{4}|\X{1}{4}|}
\sphinxtoprule
\sphinxstyletheadfamily 
\sphinxAtStartPar
ライト
&\sphinxstyletheadfamily 
\sphinxAtStartPar
スタンダード
&\sphinxstyletheadfamily 
\sphinxAtStartPar
プロ
&\sphinxstyletheadfamily 
\sphinxAtStartPar
マスター
\\
\sphinxmidrule
\sphinxtableatstartofbodyhook
\sphinxAtStartPar
x
&
\sphinxAtStartPar
◯
&
\sphinxAtStartPar
◯
&
\sphinxAtStartPar
◯
\\
\sphinxbottomrule
\end{tabular}
\sphinxtableafterendhook\par
\sphinxattableend\end{savenotes}

\sphinxAtStartPar
\sphinxstylestrong{デッキ条件:}

\sphinxAtStartPar
40枚以上

\sphinxAtStartPar
\sphinxstylestrong{ゲームの始め方:}
\begin{enumerate}
\sphinxsetlistlabels{\arabic}{enumi}{enumii}{}{.}%
\item {} 
\sphinxAtStartPar
デッキからレアカード1枚を選び配置図の場所に伏せておく。

\item {} 
\sphinxAtStartPar
デッキをシャッフルする。

\item {} 
\sphinxAtStartPar
デッキの上から14枚を取り除いてパックとし、配置図の場所に伏せておく。

\item {} 
\sphinxAtStartPar
残りのデッキをライフとし配置図の場所に伏せておく。

\item {} 
\sphinxAtStartPar
ライフの一番上のカード7枚引き手札とする。

\item {} 
\sphinxAtStartPar
\hyperref[\detokenize{common/common:common-gamestart-preset}]{\ref{\detokenize{common/common:common-gamestart-preset}} \nameref{\detokenize{common/common:common-gamestart-preset}}} に従ってプリセットする。

\item {} 
\sphinxAtStartPar
\hyperref[\detokenize{common/common:common-gamestart-first}]{\ref{\detokenize{common/common:common-gamestart-first}} \nameref{\detokenize{common/common:common-gamestart-first}}} に従って先攻プレイヤーを決定する。

\item {} 
\sphinxAtStartPar
\hyperref[\detokenize{common/common:common-gamestart-start}]{\ref{\detokenize{common/common:common-gamestart-start}} \nameref{\detokenize{common/common:common-gamestart-start}}} に従ってゲームを開始する。

\end{enumerate}

\sphinxAtStartPar
\sphinxstylestrong{アクションリスト:}

\sphinxAtStartPar
次のアクションが追加で利用可能です。
\begin{itemize}
\item {} 
\sphinxAtStartPar
パック操作アクション

\item {} 
\sphinxAtStartPar
レアカード操作アクション

\end{itemize}

\sphinxAtStartPar
詳細は次を参照してください。

\sphinxAtStartPar
\hyperref[\detokenize{auto/frameActionlist:act-act-frame}]{\ref{\detokenize{auto/frameActionlist:act-act-frame}} \nameref{\detokenize{auto/frameActionlist:act-act-frame}}}

\sphinxAtStartPar
\sphinxstylestrong{公開レベル:}

\sphinxAtStartPar
このフレームで定義されている領域の公開レベルは次の通りです。
\begin{description}
\sphinxlineitem{パック}
\begin{DUlineblock}{0em}
\item[] 完全公開：開封済み・未開封
\item[] 個人公開：開封後のパックの中身
\item[] 非公開：未開封のパックの中身
\end{DUlineblock}

\sphinxlineitem{レアカード}
\begin{DUlineblock}{0em}
\item[] 完全公開：レアカードの枚数
\item[] 個人公開：レアカードの中身
\item[] 非公開：なし
\end{DUlineblock}

\end{description}


\bigskip\hrule\bigskip



\subsubsection{ストラテジー}
\label{\detokenize{auto/framelist:frame-strategy}}\label{\detokenize{auto/framelist:id6}}
\sphinxAtStartPar
\sphinxstylestrong{配置図:}

\begin{figure}[htbp]
\centering
\capstart

\noindent\sphinxincludegraphics{{rarePack}.pdf}
\caption{配置図\sphinxhyphen{}ストラテジー}\label{\detokenize{auto/framelist:id11}}\label{\detokenize{auto/framelist:frame-strategy-image}}\end{figure}

\sphinxAtStartPar
\sphinxstylestrong{フォーマット:}


\begin{savenotes}\sphinxattablestart
\sphinxthistablewithglobalstyle
\centering
\begin{tabular}[t]{|\X{1}{4}|\X{1}{4}|\X{1}{4}|\X{1}{4}|}
\sphinxtoprule
\sphinxstyletheadfamily 
\sphinxAtStartPar
ライト
&\sphinxstyletheadfamily 
\sphinxAtStartPar
スタンダード
&\sphinxstyletheadfamily 
\sphinxAtStartPar
プロ
&\sphinxstyletheadfamily 
\sphinxAtStartPar
マスター
\\
\sphinxmidrule
\sphinxtableatstartofbodyhook
\sphinxAtStartPar
x
&
\sphinxAtStartPar
x
&
\sphinxAtStartPar
◯
&
\sphinxAtStartPar
◯
\\
\sphinxbottomrule
\end{tabular}
\sphinxtableafterendhook\par
\sphinxattableend\end{savenotes}

\sphinxAtStartPar
\sphinxstylestrong{デッキ条件:}

\sphinxAtStartPar
40枚以上

\sphinxAtStartPar
\sphinxstylestrong{ゲームの始め方:}
\begin{enumerate}
\sphinxsetlistlabels{\arabic}{enumi}{enumii}{}{.}%
\item {} 
\sphinxAtStartPar
デッキから最初の手札に入れる3枚を選ぶ。（ストラテジー）

\item {} 
\sphinxAtStartPar
デッキからレアカード1枚を選び配置図の場所に伏せておく。

\item {} 
\sphinxAtStartPar
デッキをシャッフルする。

\item {} 
\sphinxAtStartPar
デッキの上から14枚を取り除いてパックとし、配置図の場所に伏せておく。

\item {} 
\sphinxAtStartPar
残りのデッキをライフとし配置図の場所に伏せておく。

\item {} 
\sphinxAtStartPar
ライフの一番上のカードを4枚引いて手札を7枚とする。

\item {} 
\sphinxAtStartPar
\hyperref[\detokenize{common/common:common-gamestart-preset}]{\ref{\detokenize{common/common:common-gamestart-preset}} \nameref{\detokenize{common/common:common-gamestart-preset}}} に従ってプリセットする。

\item {} 
\sphinxAtStartPar
\hyperref[\detokenize{common/common:common-gamestart-first}]{\ref{\detokenize{common/common:common-gamestart-first}} \nameref{\detokenize{common/common:common-gamestart-first}}} に従って先攻プレイヤーを決定する。

\item {} 
\sphinxAtStartPar
\hyperref[\detokenize{common/common:common-gamestart-start}]{\ref{\detokenize{common/common:common-gamestart-start}} \nameref{\detokenize{common/common:common-gamestart-start}}} に従ってゲームを開始する。

\end{enumerate}

\sphinxAtStartPar
\sphinxstylestrong{アクションリスト:}

\sphinxAtStartPar
次のアクションが追加で利用可能です。
\begin{itemize}
\item {} 
\sphinxAtStartPar
パック操作アクション

\item {} 
\sphinxAtStartPar
レアカード操作アクション

\end{itemize}

\sphinxAtStartPar
詳細は次を参照してください。

\sphinxAtStartPar
\hyperref[\detokenize{auto/frameActionlist:act-act-frame}]{\ref{\detokenize{auto/frameActionlist:act-act-frame}} \nameref{\detokenize{auto/frameActionlist:act-act-frame}}}

\sphinxAtStartPar
\sphinxstylestrong{公開レベル:}

\sphinxAtStartPar
このフレームで定義されている領域の公開レベルは次の通りです。
\begin{description}
\sphinxlineitem{パック}
\begin{DUlineblock}{0em}
\item[] 完全公開：開封済み・未開封
\item[] 個人公開：開封後のパックの中身
\item[] 非公開：未開封のパックの中身
\end{DUlineblock}

\sphinxlineitem{レアカード}
\begin{DUlineblock}{0em}
\item[] 完全公開：レアカードの枚数
\item[] 個人公開：レアカードの中身
\item[] 非公開：なし
\end{DUlineblock}

\end{description}


\bigskip\hrule\bigskip



\subsubsection{エクストラ}
\label{\detokenize{auto/framelist:frame-extra}}\label{\detokenize{auto/framelist:id7}}
\sphinxAtStartPar
\sphinxstylestrong{配置図:}

\begin{figure}[htbp]
\centering
\capstart

\noindent\sphinxincludegraphics{{extra}.pdf}
\caption{配置図\sphinxhyphen{}エクストラ}\label{\detokenize{auto/framelist:id12}}\label{\detokenize{auto/framelist:frame-extra-image}}\end{figure}

\sphinxAtStartPar
\sphinxstylestrong{フォーマット:}


\begin{savenotes}\sphinxattablestart
\sphinxthistablewithglobalstyle
\centering
\begin{tabular}[t]{|\X{1}{4}|\X{1}{4}|\X{1}{4}|\X{1}{4}|}
\sphinxtoprule
\sphinxstyletheadfamily 
\sphinxAtStartPar
ライト
&\sphinxstyletheadfamily 
\sphinxAtStartPar
スタンダード
&\sphinxstyletheadfamily 
\sphinxAtStartPar
プロ
&\sphinxstyletheadfamily 
\sphinxAtStartPar
マスター
\\
\sphinxmidrule
\sphinxtableatstartofbodyhook
\sphinxAtStartPar
x
&
\sphinxAtStartPar
x
&
\sphinxAtStartPar
x
&
\sphinxAtStartPar
◯
\\
\sphinxbottomrule
\end{tabular}
\sphinxtableafterendhook\par
\sphinxattableend\end{savenotes}

\sphinxAtStartPar
\sphinxstylestrong{デッキ条件:}

\sphinxAtStartPar
40枚以上

\sphinxAtStartPar
\sphinxstylestrong{ゲームの始め方:}
\begin{enumerate}
\sphinxsetlistlabels{\arabic}{enumi}{enumii}{}{.}%
\item {} 
\sphinxAtStartPar
デッキから2枚を切札として選び、2枚とも配置図の切札の場所に伏せておく。

\item {} 
\sphinxAtStartPar
デッキから最初の手札に入れる3枚を選ぶ。（ストラテジー)

\item {} 
\sphinxAtStartPar
デッキからレアカード1枚を選び配置図の場所に伏せておく。

\item {} 
\sphinxAtStartPar
残りのデッキをシャッフルし、上から14枚を取り除いてパックとし、配置図の場所に伏せておく。

\item {} 
\sphinxAtStartPar
残りのデッキをライフとし配置図の場所に伏せておく。

\item {} 
\sphinxAtStartPar
ライフの一番上のカードを4枚引いて手札を7枚とする。

\item {} 
\sphinxAtStartPar
ライフの一番上のカードを防壁として場に出し、次のカードを兵士の場所に裏向きで出す。（プリセット）

\item {} 
\sphinxAtStartPar
対戦相手も手札を7枚用意でき、準備ができてから \hyperref[\detokenize{common/common:common-gamestart-first}]{\ref{\detokenize{common/common:common-gamestart-first}} \nameref{\detokenize{common/common:common-gamestart-first}}} に従って先攻プレイヤーを決定する。

\item {} 
\sphinxAtStartPar
先攻プレイヤーは、プリセットで兵士の場所に裏向きで出したカードを表にし、切札の場所の2枚のうち1枚を表向きにする。

\item {} 
\sphinxAtStartPar
後攻プレイヤーも、プリセットで兵士の場所に裏向きで出したカードを表にし、切札の場所の2枚のうち1枚を表向きにする。

\item {} 
\sphinxAtStartPar
\hyperref[\detokenize{common/common:common-gamestart-start}]{\ref{\detokenize{common/common:common-gamestart-start}} \nameref{\detokenize{common/common:common-gamestart-start}}} に従ってゲームを開始する。

\end{enumerate}

\sphinxAtStartPar
\sphinxstylestrong{アクションリスト:}

\sphinxAtStartPar
次のアクションが追加で利用可能です。
\begin{itemize}
\item {} 
\sphinxAtStartPar
パック操作アクション

\item {} 
\sphinxAtStartPar
レアカード操作アクション

\item {} 
\sphinxAtStartPar
切札操作アクション

\item {} 
\sphinxAtStartPar
切札関連のアクション

\end{itemize}

\sphinxAtStartPar
詳細は次を参照してください。

\sphinxAtStartPar
\hyperref[\detokenize{auto/frameActionlist:act-act-frame}]{\ref{\detokenize{auto/frameActionlist:act-act-frame}} \nameref{\detokenize{auto/frameActionlist:act-act-frame}}}

\sphinxAtStartPar
\sphinxstylestrong{公開レベル:}

\sphinxAtStartPar
このフレームで定義されている領域の公開レベルは次の通りです。
\begin{description}
\sphinxlineitem{パック}
\begin{DUlineblock}{0em}
\item[] 完全公開：開封済み・未開封
\item[] 個人公開：開封後のパックの中身
\item[] 非公開：未開封のパックの中身
\end{DUlineblock}

\sphinxlineitem{レアカード}
\begin{DUlineblock}{0em}
\item[] 完全公開：レアカードの枚数
\item[] 個人公開：レアカードの中身
\item[] 非公開：なし
\end{DUlineblock}

\sphinxlineitem{切札}
\begin{DUlineblock}{0em}
\item[] 完全公開：表裏を変えずに見えるカード
\item[] 個人公開：伏せてあるカード
\item[] 非公開：なし
\end{DUlineblock}

\end{description}

\sphinxAtStartPar
各フレームで追加で使えるアクション、フォグ、切札については次を参照してください。

\sphinxstepscope


\subsection{フレームアクションリスト}
\label{\detokenize{auto/frameActionlist:actionlist-act-frame}}\label{\detokenize{auto/frameActionlist:act-act-frame}}\label{\detokenize{auto/frameActionlist:id1}}\label{\detokenize{auto/frameActionlist::doc}}
\begin{sphinxShadowBox}
\sphinxstyletopictitle{目次}
\begin{itemize}
\item {} 
\sphinxAtStartPar
\phantomsection\label{\detokenize{auto/frameActionlist:id171}}{\hyperref[\detokenize{auto/frameActionlist:id3}]{\sphinxcrossref{パック操作}}}
\begin{itemize}
\item {} 
\sphinxAtStartPar
\phantomsection\label{\detokenize{auto/frameActionlist:id172}}{\hyperref[\detokenize{auto/frameActionlist:act-packopen}]{\sphinxcrossref{パック開封 (パック操作)}}}

\end{itemize}

\item {} 
\sphinxAtStartPar
\phantomsection\label{\detokenize{auto/frameActionlist:id173}}{\hyperref[\detokenize{auto/frameActionlist:id5}]{\sphinxcrossref{レアカード操作}}}
\begin{itemize}
\item {} 
\sphinxAtStartPar
\phantomsection\label{\detokenize{auto/frameActionlist:id174}}{\hyperref[\detokenize{auto/frameActionlist:act-raredraw}]{\sphinxcrossref{レアドロー (レアカード操作)}}}

\item {} 
\sphinxAtStartPar
\phantomsection\label{\detokenize{auto/frameActionlist:id175}}{\hyperref[\detokenize{auto/frameActionlist:act-summonsrare}]{\sphinxcrossref{レア召喚 (召喚)}}}

\item {} 
\sphinxAtStartPar
\phantomsection\label{\detokenize{auto/frameActionlist:id176}}{\hyperref[\detokenize{auto/frameActionlist:act-trapcounter}]{\sphinxcrossref{罠カウンター (レアカード操作)}}}

\end{itemize}

\item {} 
\sphinxAtStartPar
\phantomsection\label{\detokenize{auto/frameActionlist:id177}}{\hyperref[\detokenize{auto/frameActionlist:id9}]{\sphinxcrossref{切札操作}}}
\begin{itemize}
\item {} 
\sphinxAtStartPar
\phantomsection\label{\detokenize{auto/frameActionlist:id178}}{\hyperref[\detokenize{auto/frameActionlist:act-quickopen}]{\sphinxcrossref{クイックオープン (切札操作)}}}

\item {} 
\sphinxAtStartPar
\phantomsection\label{\detokenize{auto/frameActionlist:id179}}{\hyperref[\detokenize{auto/frameActionlist:act-open}]{\sphinxcrossref{メインオープン (切札操作)}}}

\item {} 
\sphinxAtStartPar
\phantomsection\label{\detokenize{auto/frameActionlist:id180}}{\hyperref[\detokenize{auto/frameActionlist:act-close}]{\sphinxcrossref{クローズ (切札操作)}}}

\item {} 
\sphinxAtStartPar
\phantomsection\label{\detokenize{auto/frameActionlist:id181}}{\hyperref[\detokenize{auto/frameActionlist:act-trumpdestroy}]{\sphinxcrossref{切札破壊 (切札操作)}}}

\end{itemize}

\item {} 
\sphinxAtStartPar
\phantomsection\label{\detokenize{auto/frameActionlist:id182}}{\hyperref[\detokenize{auto/frameActionlist:a}]{\sphinxcrossref{切札{\normalsize $\spadesuit$} A}}}
\begin{itemize}
\item {} 
\sphinxAtStartPar
\phantomsection\label{\detokenize{auto/frameActionlist:id183}}{\hyperref[\detokenize{auto/frameActionlist:act-revolutiondraw}]{\sphinxcrossref{革命ドロー (誘発魔法)}}}

\end{itemize}

\item {} 
\sphinxAtStartPar
\phantomsection\label{\detokenize{auto/frameActionlist:id184}}{\hyperref[\detokenize{auto/frameActionlist:id15}]{\sphinxcrossref{切札{\normalsize $\spadesuit$} 2}}}
\begin{itemize}
\item {} 
\sphinxAtStartPar
\phantomsection\label{\detokenize{auto/frameActionlist:id185}}{\hyperref[\detokenize{auto/frameActionlist:act-slash}]{\sphinxcrossref{スラッシュ (兵士起因)}}}

\end{itemize}

\item {} 
\sphinxAtStartPar
\phantomsection\label{\detokenize{auto/frameActionlist:id186}}{\hyperref[\detokenize{auto/frameActionlist:id17}]{\sphinxcrossref{切札{\normalsize $\spadesuit$} 3}}}
\begin{itemize}
\item {} 
\sphinxAtStartPar
\phantomsection\label{\detokenize{auto/frameActionlist:id187}}{\hyperref[\detokenize{auto/frameActionlist:act-drain}]{\sphinxcrossref{ドレイン (通常魔法)}}}

\end{itemize}

\item {} 
\sphinxAtStartPar
\phantomsection\label{\detokenize{auto/frameActionlist:id188}}{\hyperref[\detokenize{auto/frameActionlist:id19}]{\sphinxcrossref{切札{\normalsize $\spadesuit$} 4}}}
\begin{itemize}
\item {} 
\sphinxAtStartPar
\phantomsection\label{\detokenize{auto/frameActionlist:id189}}{\hyperref[\detokenize{auto/frameActionlist:act-intimidate}]{\sphinxcrossref{威圧 (通常魔法)}}}

\end{itemize}

\item {} 
\sphinxAtStartPar
\phantomsection\label{\detokenize{auto/frameActionlist:id190}}{\hyperref[\detokenize{auto/frameActionlist:id21}]{\sphinxcrossref{切札{\normalsize $\spadesuit$} 5}}}
\begin{itemize}
\item {} 
\sphinxAtStartPar
\phantomsection\label{\detokenize{auto/frameActionlist:id191}}{\hyperref[\detokenize{auto/frameActionlist:act-resurrection}]{\sphinxcrossref{蘇生 (通常魔法)}}}

\end{itemize}

\item {} 
\sphinxAtStartPar
\phantomsection\label{\detokenize{auto/frameActionlist:id192}}{\hyperref[\detokenize{auto/frameActionlist:id23}]{\sphinxcrossref{切札{\normalsize $\spadesuit$} 7}}}
\begin{itemize}
\item {} 
\sphinxAtStartPar
\phantomsection\label{\detokenize{auto/frameActionlist:id193}}{\hyperref[\detokenize{auto/frameActionlist:act-instantkill}]{\sphinxcrossref{瞬殺 (速攻魔法)}}}

\end{itemize}

\item {} 
\sphinxAtStartPar
\phantomsection\label{\detokenize{auto/frameActionlist:id194}}{\hyperref[\detokenize{auto/frameActionlist:j}]{\sphinxcrossref{切札{\normalsize $\spadesuit$} J}}}
\begin{itemize}
\item {} 
\sphinxAtStartPar
\phantomsection\label{\detokenize{auto/frameActionlist:id195}}{\hyperref[\detokenize{auto/frameActionlist:act-summonreanimator}]{\sphinxcrossref{リアニメーター召喚 (召喚)}}}

\item {} 
\sphinxAtStartPar
\phantomsection\label{\detokenize{auto/frameActionlist:id196}}{\hyperref[\detokenize{auto/frameActionlist:act-quicksummonreanimator}]{\sphinxcrossref{クイックリアニメーター召喚 (召喚)}}}

\item {} 
\sphinxAtStartPar
\phantomsection\label{\detokenize{auto/frameActionlist:id197}}{\hyperref[\detokenize{auto/frameActionlist:act-reanimatorreanimate}]{\sphinxcrossref{リアニメーターリアニメイト (兵士起因)}}}

\item {} 
\sphinxAtStartPar
\phantomsection\label{\detokenize{auto/frameActionlist:id198}}{\hyperref[\detokenize{auto/frameActionlist:act-reanimatorrevive}]{\sphinxcrossref{リアニメーター再生 (切札操作)}}}

\end{itemize}

\item {} 
\sphinxAtStartPar
\phantomsection\label{\detokenize{auto/frameActionlist:id199}}{\hyperref[\detokenize{auto/frameActionlist:k}]{\sphinxcrossref{切札{\normalsize $\spadesuit$} K}}}
\begin{itemize}
\item {} 
\sphinxAtStartPar
\phantomsection\label{\detokenize{auto/frameActionlist:id200}}{\hyperref[\detokenize{auto/frameActionlist:act-summondarklord}]{\sphinxcrossref{魔王召喚 (召喚)}}}

\item {} 
\sphinxAtStartPar
\phantomsection\label{\detokenize{auto/frameActionlist:id201}}{\hyperref[\detokenize{auto/frameActionlist:act-darkwhisper}]{\sphinxcrossref{魔王の囁き (誘発魔法)}}}

\end{itemize}

\item {} 
\sphinxAtStartPar
\phantomsection\label{\detokenize{auto/frameActionlist:id202}}{\hyperref[\detokenize{auto/frameActionlist:id31}]{\sphinxcrossref{切札{\normalsize $\heartsuit$} 2}}}
\begin{itemize}
\item {} 
\sphinxAtStartPar
\phantomsection\label{\detokenize{auto/frameActionlist:id203}}{\hyperref[\detokenize{auto/frameActionlist:act-summonreinforcement}]{\sphinxcrossref{援軍 (通常魔法)}}}

\end{itemize}

\item {} 
\sphinxAtStartPar
\phantomsection\label{\detokenize{auto/frameActionlist:id204}}{\hyperref[\detokenize{auto/frameActionlist:id33}]{\sphinxcrossref{切札{\normalsize $\heartsuit$} 3}}}
\begin{itemize}
\item {} 
\sphinxAtStartPar
\phantomsection\label{\detokenize{auto/frameActionlist:id205}}{\hyperref[\detokenize{auto/frameActionlist:act-rally}]{\sphinxcrossref{激励 (兵士起因)}}}

\end{itemize}

\item {} 
\sphinxAtStartPar
\phantomsection\label{\detokenize{auto/frameActionlist:id206}}{\hyperref[\detokenize{auto/frameActionlist:id35}]{\sphinxcrossref{切札{\normalsize $\heartsuit$} 4}}}
\begin{itemize}
\item {} 
\sphinxAtStartPar
\phantomsection\label{\detokenize{auto/frameActionlist:id207}}{\hyperref[\detokenize{auto/frameActionlist:act-retreat}]{\sphinxcrossref{休戦 (速攻魔法)}}}

\end{itemize}

\item {} 
\sphinxAtStartPar
\phantomsection\label{\detokenize{auto/frameActionlist:id208}}{\hyperref[\detokenize{auto/frameActionlist:id37}]{\sphinxcrossref{切札{\normalsize $\heartsuit$} 5}}}
\begin{itemize}
\item {} 
\sphinxAtStartPar
\phantomsection\label{\detokenize{auto/frameActionlist:id209}}{\hyperref[\detokenize{auto/frameActionlist:act-pursuit}]{\sphinxcrossref{追撃 (速攻魔法)}}}

\end{itemize}

\item {} 
\sphinxAtStartPar
\phantomsection\label{\detokenize{auto/frameActionlist:id210}}{\hyperref[\detokenize{auto/frameActionlist:id39}]{\sphinxcrossref{切札{\normalsize $\heartsuit$} 7}}}
\begin{itemize}
\item {} 
\sphinxAtStartPar
\phantomsection\label{\detokenize{auto/frameActionlist:id211}}{\hyperref[\detokenize{auto/frameActionlist:act-phantomreverse}]{\sphinxcrossref{ファントムリバース (通常魔法)}}}

\item {} 
\sphinxAtStartPar
\phantomsection\label{\detokenize{auto/frameActionlist:id212}}{\hyperref[\detokenize{auto/frameActionlist:act-illusion}]{\sphinxcrossref{幻影 (誘発魔法)}}}

\end{itemize}

\item {} 
\sphinxAtStartPar
\phantomsection\label{\detokenize{auto/frameActionlist:id213}}{\hyperref[\detokenize{auto/frameActionlist:id42}]{\sphinxcrossref{切札{\normalsize $\heartsuit$} 8}}}
\begin{itemize}
\item {} 
\sphinxAtStartPar
\phantomsection\label{\detokenize{auto/frameActionlist:id214}}{\hyperref[\detokenize{auto/frameActionlist:act-salvation}]{\sphinxcrossref{救済 (通常魔法)}}}

\end{itemize}

\item {} 
\sphinxAtStartPar
\phantomsection\label{\detokenize{auto/frameActionlist:id215}}{\hyperref[\detokenize{auto/frameActionlist:id44}]{\sphinxcrossref{切札{\normalsize $\heartsuit$} J}}}
\begin{itemize}
\item {} 
\sphinxAtStartPar
\phantomsection\label{\detokenize{auto/frameActionlist:id216}}{\hyperref[\detokenize{auto/frameActionlist:act-summonchariot}]{\sphinxcrossref{チャリオット召喚 (召喚)}}}

\item {} 
\sphinxAtStartPar
\phantomsection\label{\detokenize{auto/frameActionlist:id217}}{\hyperref[\detokenize{auto/frameActionlist:act-quicksummonchariot}]{\sphinxcrossref{クイックチャリオット召喚 (召喚)}}}

\item {} 
\sphinxAtStartPar
\phantomsection\label{\detokenize{auto/frameActionlist:id218}}{\hyperref[\detokenize{auto/frameActionlist:act-chariotdestroy}]{\sphinxcrossref{チャリオット防壁破壊 (兵士起因)}}}

\item {} 
\sphinxAtStartPar
\phantomsection\label{\detokenize{auto/frameActionlist:id219}}{\hyperref[\detokenize{auto/frameActionlist:act-chariotrevive}]{\sphinxcrossref{チャリオット再生 (切札操作)}}}

\end{itemize}

\item {} 
\sphinxAtStartPar
\phantomsection\label{\detokenize{auto/frameActionlist:id220}}{\hyperref[\detokenize{auto/frameActionlist:id49}]{\sphinxcrossref{切札{\normalsize $\heartsuit$} K}}}
\begin{itemize}
\item {} 
\sphinxAtStartPar
\phantomsection\label{\detokenize{auto/frameActionlist:id221}}{\hyperref[\detokenize{auto/frameActionlist:act-summongiant}]{\sphinxcrossref{巨人召喚 (召喚)}}}

\end{itemize}

\item {} 
\sphinxAtStartPar
\phantomsection\label{\detokenize{auto/frameActionlist:id222}}{\hyperref[\detokenize{auto/frameActionlist:id51}]{\sphinxcrossref{切札{\normalsize $\diamondsuit$} A}}}
\begin{itemize}
\item {} 
\sphinxAtStartPar
\phantomsection\label{\detokenize{auto/frameActionlist:id223}}{\hyperref[\detokenize{auto/frameActionlist:act-putbulwark}]{\sphinxcrossref{防壁追加 (通常魔法)}}}

\end{itemize}

\item {} 
\sphinxAtStartPar
\phantomsection\label{\detokenize{auto/frameActionlist:id224}}{\hyperref[\detokenize{auto/frameActionlist:id53}]{\sphinxcrossref{切札{\normalsize $\diamondsuit$} 3}}}
\begin{itemize}
\item {} 
\sphinxAtStartPar
\phantomsection\label{\detokenize{auto/frameActionlist:id225}}{\hyperref[\detokenize{auto/frameActionlist:act-negotiate}]{\sphinxcrossref{交渉 (兵士起因)}}}

\end{itemize}

\item {} 
\sphinxAtStartPar
\phantomsection\label{\detokenize{auto/frameActionlist:id226}}{\hyperref[\detokenize{auto/frameActionlist:id55}]{\sphinxcrossref{切札{\normalsize $\diamondsuit$} 4}}}
\begin{itemize}
\item {} 
\sphinxAtStartPar
\phantomsection\label{\detokenize{auto/frameActionlist:id227}}{\hyperref[\detokenize{auto/frameActionlist:act-reckoning}]{\sphinxcrossref{相殺 (速攻魔法)}}}

\end{itemize}

\item {} 
\sphinxAtStartPar
\phantomsection\label{\detokenize{auto/frameActionlist:id228}}{\hyperref[\detokenize{auto/frameActionlist:id57}]{\sphinxcrossref{切札{\normalsize $\diamondsuit$} 5}}}
\begin{itemize}
\item {} 
\sphinxAtStartPar
\phantomsection\label{\detokenize{auto/frameActionlist:id229}}{\hyperref[\detokenize{auto/frameActionlist:act-exhaust}]{\sphinxcrossref{疲弊 (誘発魔法)}}}

\end{itemize}

\item {} 
\sphinxAtStartPar
\phantomsection\label{\detokenize{auto/frameActionlist:id230}}{\hyperref[\detokenize{auto/frameActionlist:id59}]{\sphinxcrossref{切札{\normalsize $\diamondsuit$} 8}}}
\begin{itemize}
\item {} 
\sphinxAtStartPar
\phantomsection\label{\detokenize{auto/frameActionlist:id231}}{\hyperref[\detokenize{auto/frameActionlist:act-repair}]{\sphinxcrossref{修繕 (通常魔法)}}}

\end{itemize}

\item {} 
\sphinxAtStartPar
\phantomsection\label{\detokenize{auto/frameActionlist:id232}}{\hyperref[\detokenize{auto/frameActionlist:id61}]{\sphinxcrossref{切札{\normalsize $\diamondsuit$} 9}}}
\begin{itemize}
\item {} 
\sphinxAtStartPar
\phantomsection\label{\detokenize{auto/frameActionlist:id233}}{\hyperref[\detokenize{auto/frameActionlist:act-rest}]{\sphinxcrossref{休息 (通常魔法)}}}

\item {} 
\sphinxAtStartPar
\phantomsection\label{\detokenize{auto/frameActionlist:id234}}{\hyperref[\detokenize{auto/frameActionlist:act-satisfaction}]{\sphinxcrossref{充足 (誘発魔法)}}}

\end{itemize}

\item {} 
\sphinxAtStartPar
\phantomsection\label{\detokenize{auto/frameActionlist:id235}}{\hyperref[\detokenize{auto/frameActionlist:id64}]{\sphinxcrossref{切札{\normalsize $\diamondsuit$} J}}}
\begin{itemize}
\item {} 
\sphinxAtStartPar
\phantomsection\label{\detokenize{auto/frameActionlist:id236}}{\hyperref[\detokenize{auto/frameActionlist:act-summonstrategist}]{\sphinxcrossref{策士召喚 (召喚)}}}

\item {} 
\sphinxAtStartPar
\phantomsection\label{\detokenize{auto/frameActionlist:id237}}{\hyperref[\detokenize{auto/frameActionlist:act-quicksummonstrategist}]{\sphinxcrossref{クイック策士召喚 (召喚)}}}

\item {} 
\sphinxAtStartPar
\phantomsection\label{\detokenize{auto/frameActionlist:id238}}{\hyperref[\detokenize{auto/frameActionlist:act-strategistpredict}]{\sphinxcrossref{先読み (兵士起因)}}}

\item {} 
\sphinxAtStartPar
\phantomsection\label{\detokenize{auto/frameActionlist:id239}}{\hyperref[\detokenize{auto/frameActionlist:act-strategistrevive}]{\sphinxcrossref{策士再生 (切札操作)}}}

\end{itemize}

\item {} 
\sphinxAtStartPar
\phantomsection\label{\detokenize{auto/frameActionlist:id240}}{\hyperref[\detokenize{auto/frameActionlist:id69}]{\sphinxcrossref{切札{\normalsize $\clubsuit$} 4}}}
\begin{itemize}
\item {} 
\sphinxAtStartPar
\phantomsection\label{\detokenize{auto/frameActionlist:id241}}{\hyperref[\detokenize{auto/frameActionlist:act-switch}]{\sphinxcrossref{スイッチ (兵士起因)}}}

\end{itemize}

\item {} 
\sphinxAtStartPar
\phantomsection\label{\detokenize{auto/frameActionlist:id242}}{\hyperref[\detokenize{auto/frameActionlist:id71}]{\sphinxcrossref{切札{\normalsize $\clubsuit$} 5}}}
\begin{itemize}
\item {} 
\sphinxAtStartPar
\phantomsection\label{\detokenize{auto/frameActionlist:id243}}{\hyperref[\detokenize{auto/frameActionlist:act-civilwar}]{\sphinxcrossref{内乱 (通常魔法)}}}

\item {} 
\sphinxAtStartPar
\phantomsection\label{\detokenize{auto/frameActionlist:id244}}{\hyperref[\detokenize{auto/frameActionlist:act-riot}]{\sphinxcrossref{暴動 (誘発魔法)}}}

\end{itemize}

\item {} 
\sphinxAtStartPar
\phantomsection\label{\detokenize{auto/frameActionlist:id245}}{\hyperref[\detokenize{auto/frameActionlist:id74}]{\sphinxcrossref{切札{\normalsize $\clubsuit$} 7}}}
\begin{itemize}
\item {} 
\sphinxAtStartPar
\phantomsection\label{\detokenize{auto/frameActionlist:id246}}{\hyperref[\detokenize{auto/frameActionlist:act-handrefill}]{\sphinxcrossref{手札補充 (誘発魔法)}}}

\end{itemize}

\item {} 
\sphinxAtStartPar
\phantomsection\label{\detokenize{auto/frameActionlist:id247}}{\hyperref[\detokenize{auto/frameActionlist:id76}]{\sphinxcrossref{切札{\normalsize $\clubsuit$} 8}}}
\begin{itemize}
\item {} 
\sphinxAtStartPar
\phantomsection\label{\detokenize{auto/frameActionlist:id248}}{\hyperref[\detokenize{auto/frameActionlist:act-reset}]{\sphinxcrossref{リセット (通常魔法)}}}

\end{itemize}

\item {} 
\sphinxAtStartPar
\phantomsection\label{\detokenize{auto/frameActionlist:id249}}{\hyperref[\detokenize{auto/frameActionlist:id78}]{\sphinxcrossref{切札{\normalsize $\clubsuit$} 10}}}
\begin{itemize}
\item {} 
\sphinxAtStartPar
\phantomsection\label{\detokenize{auto/frameActionlist:id250}}{\hyperref[\detokenize{auto/frameActionlist:act-lastditch}]{\sphinxcrossref{悪あがき (速攻魔法)}}}

\end{itemize}

\item {} 
\sphinxAtStartPar
\phantomsection\label{\detokenize{auto/frameActionlist:id251}}{\hyperref[\detokenize{auto/frameActionlist:id80}]{\sphinxcrossref{切札{\normalsize $\clubsuit$} J}}}
\begin{itemize}
\item {} 
\sphinxAtStartPar
\phantomsection\label{\detokenize{auto/frameActionlist:id252}}{\hyperref[\detokenize{auto/frameActionlist:act-summonknight}]{\sphinxcrossref{騎士召喚 (召喚)}}}

\item {} 
\sphinxAtStartPar
\phantomsection\label{\detokenize{auto/frameActionlist:id253}}{\hyperref[\detokenize{auto/frameActionlist:act-quicksummonknight}]{\sphinxcrossref{クイック騎士召喚 (召喚)}}}

\item {} 
\sphinxAtStartPar
\phantomsection\label{\detokenize{auto/frameActionlist:id254}}{\hyperref[\detokenize{auto/frameActionlist:act-knightsweep}]{\sphinxcrossref{薙払 (兵士起因)}}}

\item {} 
\sphinxAtStartPar
\phantomsection\label{\detokenize{auto/frameActionlist:id255}}{\hyperref[\detokenize{auto/frameActionlist:act-knightrevive}]{\sphinxcrossref{騎士再生 (切札操作)}}}

\end{itemize}

\end{itemize}
\end{sphinxShadowBox}


\bigskip\hrule\bigskip



\subsubsection{パック操作}
\label{\detokenize{auto/frameActionlist:id3}}

\paragraph{パック開封 (パック操作)}
\label{\detokenize{auto/frameActionlist:act-packopen}}\label{\detokenize{auto/frameActionlist:id4}}
\sphinxAtStartPar
\sphinxstylestrong{フレーム:}

\sphinxAtStartPar
・パック

\sphinxAtStartPar
・レアパック

\sphinxAtStartPar
・ストラテジー

\sphinxAtStartPar
・エクストラ

\sphinxAtStartPar
\sphinxstylestrong{トリガー:} 直接

\sphinxAtStartPar
\sphinxstylestrong{スピード:} 即時

\sphinxAtStartPar
\sphinxstylestrong{タイミング:} クイック

\sphinxAtStartPar
\sphinxstylestrong{起動条件:}

\sphinxAtStartPar
パックが開封済みでない場合のみこのアクションを起こすことができる。

\sphinxAtStartPar
\sphinxstylestrong{効果:}
\begin{enumerate}
\sphinxsetlistlabels{\arabic}{enumi}{enumii}{}{.}%
\item {} 
\sphinxAtStartPar
パックの中から好きなカードを1枚選び対戦相手に見せ手札に加える。レアカードがある場合、代わりにそれを対戦相手に見せ手札に加えてもよい。

\item {} 
\sphinxAtStartPar
パックを表向きにし、開封済みとする。

\end{enumerate}


\bigskip\hrule\bigskip



\subsubsection{レアカード操作}
\label{\detokenize{auto/frameActionlist:id5}}

\paragraph{レアドロー (レアカード操作)}
\label{\detokenize{auto/frameActionlist:act-raredraw}}\label{\detokenize{auto/frameActionlist:id6}}
\sphinxAtStartPar
\sphinxstylestrong{フレーム:}

\sphinxAtStartPar
・レアパック

\sphinxAtStartPar
・ストラテジー

\sphinxAtStartPar
・エクストラ

\sphinxAtStartPar
\sphinxstylestrong{トリガー:} 直接

\sphinxAtStartPar
\sphinxstylestrong{スピード:} 即時

\sphinxAtStartPar
\sphinxstylestrong{タイミング:} クイック

\sphinxAtStartPar
\sphinxstylestrong{起動条件:}

\sphinxAtStartPar
プレイヤーのライフが9以下の場合にしかこのアクションを起こすことができない。

\sphinxAtStartPar
\sphinxstylestrong{効果:}

\sphinxAtStartPar
レアカードを1枚手札に加える。


\bigskip\hrule\bigskip



\paragraph{レア召喚 (召喚)}
\label{\detokenize{auto/frameActionlist:act-summonsrare}}\label{\detokenize{auto/frameActionlist:id7}}
\sphinxAtStartPar
\sphinxstylestrong{フレーム:}

\sphinxAtStartPar
・レアパック

\sphinxAtStartPar
・ストラテジー

\sphinxAtStartPar
・エクストラ

\sphinxAtStartPar
\sphinxstylestrong{トリガー:} 直接

\sphinxAtStartPar
\sphinxstylestrong{スピード:} 通常

\sphinxAtStartPar
\sphinxstylestrong{タイミング:} クイック

\sphinxAtStartPar
\sphinxstylestrong{キーカード:} レアカード1枚

\sphinxAtStartPar
\sphinxstylestrong{コスト:} S

\sphinxAtStartPar
\sphinxstylestrong{特記事項:} ※このアクションが起こされた時、キーカード(レアカード)はステージ上に置かれ完全公開となる。

\sphinxAtStartPar
\sphinxstylestrong{効果:}

\sphinxAtStartPar
キーカードを兵士として表向きかつチャージ状態で場に出す。


\bigskip\hrule\bigskip



\paragraph{罠カウンター (レアカード操作)}
\label{\detokenize{auto/frameActionlist:act-trapcounter}}\label{\detokenize{auto/frameActionlist:id8}}
\sphinxAtStartPar
\sphinxstylestrong{フレーム:}

\sphinxAtStartPar
・レアパック

\sphinxAtStartPar
・ストラテジー

\sphinxAtStartPar
・エクストラ

\sphinxAtStartPar
\sphinxstylestrong{トリガー:} 誘発

\sphinxAtStartPar
\sphinxstylestrong{スピード:} 通常

\sphinxAtStartPar
\sphinxstylestrong{タイミング:} クイック

\sphinxAtStartPar
\sphinxstylestrong{キーカード:} レアカード1枚

\sphinxAtStartPar
\sphinxstylestrong{特記事項:} ※このアクションが起こされた時、キーカード(レアカード)はステージ上に置かれ完全公開となる。

\sphinxAtStartPar
\sphinxstylestrong{対象:}

\sphinxAtStartPar
キーカードが1枚または2枚のリクエスト

\sphinxAtStartPar
\sphinxstylestrong{効果:}

\sphinxAtStartPar
対象リクエストのキーカードのいずれかがこのアクションのキーカードと同じ場合、対象リクエストをステージから取り除き、対象リクエストのキーカードを墓地に移す。


\bigskip\hrule\bigskip



\subsubsection{切札操作}
\label{\detokenize{auto/frameActionlist:id9}}

\paragraph{クイックオープン (切札操作)}
\label{\detokenize{auto/frameActionlist:act-quickopen}}\label{\detokenize{auto/frameActionlist:id10}}
\sphinxAtStartPar
\sphinxstylestrong{フレーム:}

\sphinxAtStartPar
・エクストラ

\sphinxAtStartPar
\sphinxstylestrong{トリガー:} 直接

\sphinxAtStartPar
\sphinxstylestrong{スピード:} 即時

\sphinxAtStartPar
\sphinxstylestrong{タイミング:} クイック

\sphinxAtStartPar
\sphinxstylestrong{コスト:} DD

\sphinxAtStartPar
\sphinxstylestrong{効果:}

\sphinxAtStartPar
自分の裏切札を1枚表にし表切札とする。


\bigskip\hrule\bigskip



\paragraph{メインオープン (切札操作)}
\label{\detokenize{auto/frameActionlist:act-open}}\label{\detokenize{auto/frameActionlist:id11}}
\sphinxAtStartPar
\sphinxstylestrong{フレーム:}

\sphinxAtStartPar
・エクストラ

\sphinxAtStartPar
\sphinxstylestrong{トリガー:} 直接

\sphinxAtStartPar
\sphinxstylestrong{スピード:} 即時

\sphinxAtStartPar
\sphinxstylestrong{タイミング:} メイン

\sphinxAtStartPar
\sphinxstylestrong{コスト:} BBL

\sphinxAtStartPar
\sphinxstylestrong{効果:}

\sphinxAtStartPar
自分の裏切札を1枚表にし表切札とする。


\bigskip\hrule\bigskip



\paragraph{クローズ (切札操作)}
\label{\detokenize{auto/frameActionlist:act-close}}\label{\detokenize{auto/frameActionlist:id12}}
\sphinxAtStartPar
\sphinxstylestrong{フレーム:}

\sphinxAtStartPar
・エクストラ

\sphinxAtStartPar
\sphinxstylestrong{トリガー:} 直接

\sphinxAtStartPar
\sphinxstylestrong{スピード:} 通常

\sphinxAtStartPar
\sphinxstylestrong{タイミング:} クイック

\sphinxAtStartPar
\sphinxstylestrong{キーカード:} Q

\sphinxAtStartPar
\sphinxstylestrong{コスト:} D

\sphinxAtStartPar
\sphinxstylestrong{対象:}

\sphinxAtStartPar
表切札1枚

\sphinxAtStartPar
\sphinxstylestrong{効果:}

\sphinxAtStartPar
対象の表切札を裏向きにする。このターンが終わるまで、それは表切札にならない。その印としてフォグにキーカードを置く。


\bigskip\hrule\bigskip



\paragraph{切札破壊 (切札操作)}
\label{\detokenize{auto/frameActionlist:act-trumpdestroy}}\label{\detokenize{auto/frameActionlist:id13}}
\sphinxAtStartPar
\sphinxstylestrong{フレーム:}

\sphinxAtStartPar
・エクストラ

\sphinxAtStartPar
\sphinxstylestrong{トリガー:} 直接

\sphinxAtStartPar
\sphinxstylestrong{スピード:} 通常

\sphinxAtStartPar
\sphinxstylestrong{タイミング:} クイック

\sphinxAtStartPar
\sphinxstylestrong{キーカード:} A〜K

\sphinxAtStartPar
\sphinxstylestrong{コスト:} B

\sphinxAtStartPar
\sphinxstylestrong{対象:}

\sphinxAtStartPar
表切札1枚

\sphinxAtStartPar
\sphinxstylestrong{効果:}

\sphinxAtStartPar
対象の表切札とキーカードのスートと数字が同じ場合、対象の表切札を墓地に移す。


\bigskip\hrule\bigskip



\subsubsection{切札{\normalsize $\spadesuit$} A}
\label{\detokenize{auto/frameActionlist:a}}

\paragraph{革命ドロー (誘発魔法)}
\label{\detokenize{auto/frameActionlist:act-revolutiondraw}}\label{\detokenize{auto/frameActionlist:id14}}
\sphinxAtStartPar
\sphinxstylestrong{フレーム:}

\sphinxAtStartPar
・エクストラ

\sphinxAtStartPar
\sphinxstylestrong{トリガー:} 誘発

\sphinxAtStartPar
\sphinxstylestrong{スピード:} 即時

\sphinxAtStartPar
\sphinxstylestrong{タイミング:} クイック

\sphinxAtStartPar
\sphinxstylestrong{誘発条件:}

\sphinxAtStartPar
自分の表切札に「革命」があるかつ、自分の \textless{}革命ドロー\textgreater{} を持つ兵士がアタッカーかつ、ダメージ判定アクションにてブロックされていない時に誘発する。

\sphinxAtStartPar
\sphinxstylestrong{効果:}

\sphinxAtStartPar
ライフの一番上からカードを1枚引き手札に加える。

\sphinxAtStartPar
\sphinxstylestrong{導入バージョン:}  ex4.14.0

\sphinxAtStartPar
\sphinxstylestrong{更新バージョン:}  v8.2


\bigskip\hrule\bigskip



\subsubsection{切札{\normalsize $\spadesuit$} 2}
\label{\detokenize{auto/frameActionlist:id15}}

\paragraph{スラッシュ (兵士起因)}
\label{\detokenize{auto/frameActionlist:act-slash}}\label{\detokenize{auto/frameActionlist:id16}}
\sphinxAtStartPar
\sphinxstylestrong{フレーム:}

\sphinxAtStartPar
・エクストラ

\sphinxAtStartPar
\sphinxstylestrong{トリガー:} 直接

\sphinxAtStartPar
\sphinxstylestrong{スピード:} 即時

\sphinxAtStartPar
\sphinxstylestrong{タイミング:} メイン

\sphinxAtStartPar
\sphinxstylestrong{キーカード:} \textless{}スラッシュ\textgreater{} を持つ兵士1体

\sphinxAtStartPar
\sphinxstylestrong{コスト:} K

\sphinxAtStartPar
\sphinxstylestrong{起動条件:}

\sphinxAtStartPar
自分の表切札に「スラッシュ」がある

\sphinxAtStartPar
\sphinxstylestrong{対象:}

\sphinxAtStartPar
兵士1体

\sphinxAtStartPar
\sphinxstylestrong{効果:}

\sphinxAtStartPar
対象の兵士のサイズが X 以下なら、その兵士を墓地に移す。X はキーカードのサイズに等しい。

\sphinxAtStartPar
\sphinxstylestrong{導入バージョン:}  ex3.4.0

\sphinxAtStartPar
\sphinxstylestrong{更新バージョン:}  v8.1


\bigskip\hrule\bigskip



\subsubsection{切札{\normalsize $\spadesuit$} 3}
\label{\detokenize{auto/frameActionlist:id17}}

\paragraph{ドレイン (通常魔法)}
\label{\detokenize{auto/frameActionlist:act-drain}}\label{\detokenize{auto/frameActionlist:id18}}
\sphinxAtStartPar
\sphinxstylestrong{フレーム:}

\sphinxAtStartPar
・エクストラ

\sphinxAtStartPar
\sphinxstylestrong{トリガー:} 直接

\sphinxAtStartPar
\sphinxstylestrong{スピード:} 通常

\sphinxAtStartPar
\sphinxstylestrong{タイミング:} メイン

\sphinxAtStartPar
\sphinxstylestrong{起動条件:}

\sphinxAtStartPar
自分の表切札に「ドレイン」がある

\sphinxAtStartPar
\sphinxstylestrong{対象:}

\sphinxAtStartPar
対戦相手1人

\sphinxAtStartPar
\sphinxstylestrong{効果:}

\sphinxAtStartPar
自分の場にいるキャラクター1体をライフの一番下に移す。移せた場合、対象の対戦相手に2点のダメージを与える。

\sphinxAtStartPar
\sphinxstylestrong{導入バージョン:}  ex5.30.0

\sphinxAtStartPar
\sphinxstylestrong{更新バージョン:}  v8.2


\bigskip\hrule\bigskip



\subsubsection{切札{\normalsize $\spadesuit$} 4}
\label{\detokenize{auto/frameActionlist:id19}}

\paragraph{威圧 (通常魔法)}
\label{\detokenize{auto/frameActionlist:act-intimidate}}\label{\detokenize{auto/frameActionlist:id20}}
\sphinxAtStartPar
\sphinxstylestrong{フレーム:}

\sphinxAtStartPar
・エクストラ

\sphinxAtStartPar
\sphinxstylestrong{トリガー:} 直接

\sphinxAtStartPar
\sphinxstylestrong{スピード:} 通常

\sphinxAtStartPar
\sphinxstylestrong{タイミング:} メイン

\sphinxAtStartPar
\sphinxstylestrong{コスト:} L

\sphinxAtStartPar
\sphinxstylestrong{起動条件:}

\sphinxAtStartPar
自分の表切札に「威圧」がある

\sphinxAtStartPar
\sphinxstylestrong{効果:}

\sphinxAtStartPar
サイズが10以下の兵士を全て墓地に移す。

\sphinxAtStartPar
\sphinxstylestrong{導入バージョン:}  ex4.18.0

\sphinxAtStartPar
\sphinxstylestrong{更新バージョン:}  v8.2


\bigskip\hrule\bigskip



\subsubsection{切札{\normalsize $\spadesuit$} 5}
\label{\detokenize{auto/frameActionlist:id21}}

\paragraph{蘇生 (通常魔法)}
\label{\detokenize{auto/frameActionlist:act-resurrection}}\label{\detokenize{auto/frameActionlist:id22}}
\sphinxAtStartPar
\sphinxstylestrong{フレーム:}

\sphinxAtStartPar
・エクストラ

\sphinxAtStartPar
\sphinxstylestrong{トリガー:} 直接

\sphinxAtStartPar
\sphinxstylestrong{スピード:} 通常

\sphinxAtStartPar
\sphinxstylestrong{タイミング:} メイン

\sphinxAtStartPar
\sphinxstylestrong{コスト:} BL

\sphinxAtStartPar
\sphinxstylestrong{起動条件:}

\sphinxAtStartPar
自分の表切札に「蘇生」がある

\sphinxAtStartPar
\sphinxstylestrong{効果:}

\sphinxAtStartPar
自分の墓地にある{\normalsize $\spadesuit$} 2〜10のカードを1枚選ぶ。選んだカードを兵士としてチャージ状態で場に出す。

\sphinxAtStartPar
\sphinxstylestrong{導入バージョン:}  ex3.10.0

\sphinxAtStartPar
\sphinxstylestrong{更新バージョン:}  v8.2


\bigskip\hrule\bigskip



\subsubsection{切札{\normalsize $\spadesuit$} 7}
\label{\detokenize{auto/frameActionlist:id23}}

\paragraph{瞬殺 (速攻魔法)}
\label{\detokenize{auto/frameActionlist:act-instantkill}}\label{\detokenize{auto/frameActionlist:id24}}
\sphinxAtStartPar
\sphinxstylestrong{フレーム:}

\sphinxAtStartPar
・エクストラ

\sphinxAtStartPar
\sphinxstylestrong{トリガー:} 直接

\sphinxAtStartPar
\sphinxstylestrong{スピード:} 即時

\sphinxAtStartPar
\sphinxstylestrong{タイミング:} クイック

\sphinxAtStartPar
\sphinxstylestrong{起動条件:}

\sphinxAtStartPar
自分の表切札に「瞬殺」がある

\sphinxAtStartPar
\sphinxstylestrong{対象:}

\sphinxAtStartPar
兵士1体

\sphinxAtStartPar
\sphinxstylestrong{効果:}

\sphinxAtStartPar
自分の場にいる兵士を1体以上任意の数だけ墓地に移す。この方法で墓地に移した兵士のサイズの合計値が対象の兵士のサイズ以上の場合、対象の兵士を墓地に移す。

\sphinxAtStartPar
\sphinxstylestrong{導入バージョン:}  ex5.30.0

\sphinxAtStartPar
\sphinxstylestrong{更新バージョン:}  v8.2


\bigskip\hrule\bigskip



\subsubsection{切札{\normalsize $\spadesuit$} J}
\label{\detokenize{auto/frameActionlist:j}}

\paragraph{リアニメーター召喚 (召喚)}
\label{\detokenize{auto/frameActionlist:act-summonreanimator}}\label{\detokenize{auto/frameActionlist:id25}}
\sphinxAtStartPar
\sphinxstylestrong{フレーム:}

\sphinxAtStartPar
・エクストラ

\sphinxAtStartPar
\sphinxstylestrong{トリガー:} 直接

\sphinxAtStartPar
\sphinxstylestrong{スピード:} 通常

\sphinxAtStartPar
\sphinxstylestrong{タイミング:} メイン

\sphinxAtStartPar
\sphinxstylestrong{キーカード:} {\normalsize $\heartsuit$} 2〜10

\sphinxAtStartPar
\sphinxstylestrong{コスト:} BL

\sphinxAtStartPar
\sphinxstylestrong{起動条件:}

\sphinxAtStartPar
自分の表切札に「リアニメーター」がある

\sphinxAtStartPar
\sphinxstylestrong{効果:}

\sphinxAtStartPar
切札の{\normalsize $\spadesuit$} Jが表の場合、次を行う。
\begin{enumerate}
\sphinxsetlistlabels{\arabic}{enumi}{enumii}{}{.}%
\item {} 
\sphinxAtStartPar
墓地にあるカードを1枚選び、兵士として表向きかつチャージ状態で場に出す。

\item {} 
\sphinxAtStartPar
キーカードと切札の{\normalsize $\spadesuit$} Jをあわせてリアニメーターとしてチャージ状態で場に出す。

\end{enumerate}

\sphinxAtStartPar
\sphinxstylestrong{導入バージョン:}  ex4.22.0

\sphinxAtStartPar
\sphinxstylestrong{更新バージョン:}  v8.2


\bigskip\hrule\bigskip



\paragraph{クイックリアニメーター召喚 (召喚)}
\label{\detokenize{auto/frameActionlist:act-quicksummonreanimator}}\label{\detokenize{auto/frameActionlist:id26}}
\sphinxAtStartPar
\sphinxstylestrong{フレーム:}

\sphinxAtStartPar
・エクストラ

\sphinxAtStartPar
\sphinxstylestrong{トリガー:} 直接

\sphinxAtStartPar
\sphinxstylestrong{スピード:} 通常

\sphinxAtStartPar
\sphinxstylestrong{タイミング:} クイック

\sphinxAtStartPar
\sphinxstylestrong{キーカード:} {\normalsize $\heartsuit$} 2〜10

\sphinxAtStartPar
\sphinxstylestrong{コスト:} D

\sphinxAtStartPar
\sphinxstylestrong{起動条件:}

\sphinxAtStartPar
自分の表切札に「リアニメーター」がある

\sphinxAtStartPar
\sphinxstylestrong{効果:}

\sphinxAtStartPar
効果の内容は「リアニメーター召喚」を参照。

\sphinxAtStartPar
\sphinxstylestrong{導入バージョン:}  ex4.22.0

\sphinxAtStartPar
\sphinxstylestrong{更新バージョン:}  v8.2


\bigskip\hrule\bigskip



\paragraph{リアニメーターリアニメイト (兵士起因)}
\label{\detokenize{auto/frameActionlist:act-reanimatorreanimate}}\label{\detokenize{auto/frameActionlist:id27}}
\sphinxAtStartPar
\sphinxstylestrong{フレーム:}

\sphinxAtStartPar
・エクストラ

\sphinxAtStartPar
\sphinxstylestrong{トリガー:} 直接

\sphinxAtStartPar
\sphinxstylestrong{スピード:} 通常

\sphinxAtStartPar
\sphinxstylestrong{タイミング:} メイン

\sphinxAtStartPar
\sphinxstylestrong{キーカード:} 「リアニメーター」1体

\sphinxAtStartPar
\sphinxstylestrong{コスト:} K

\sphinxAtStartPar
\sphinxstylestrong{対象:}

\sphinxAtStartPar
自分のキャラクター1体

\sphinxAtStartPar
\sphinxstylestrong{効果:}

\sphinxAtStartPar
効果の内容は「リアニメイト」を参照。

\sphinxAtStartPar
\sphinxstylestrong{導入バージョン:}  ex4.22.0

\sphinxAtStartPar
\sphinxstylestrong{更新バージョン:}  v8.2


\bigskip\hrule\bigskip



\paragraph{リアニメーター再生 (切札操作)}
\label{\detokenize{auto/frameActionlist:act-reanimatorrevive}}\label{\detokenize{auto/frameActionlist:id28}}
\sphinxAtStartPar
\sphinxstylestrong{フレーム:}

\sphinxAtStartPar
・エクストラ

\sphinxAtStartPar
\sphinxstylestrong{トリガー:} 誘発

\sphinxAtStartPar
\sphinxstylestrong{スピード:} 即時

\sphinxAtStartPar
\sphinxstylestrong{タイミング:} クイック

\sphinxAtStartPar
\sphinxstylestrong{誘発条件:}

\sphinxAtStartPar
自分の 「リアニメーター」 が場から墓地に移る場合

\sphinxAtStartPar
\sphinxstylestrong{効果:}

\sphinxAtStartPar
墓地から{\normalsize $\spadesuit$} Jを1枚選び、自分の表切札として置く。


\bigskip\hrule\bigskip



\subsubsection{切札{\normalsize $\spadesuit$} K}
\label{\detokenize{auto/frameActionlist:k}}

\paragraph{魔王召喚 (召喚)}
\label{\detokenize{auto/frameActionlist:act-summondarklord}}\label{\detokenize{auto/frameActionlist:id29}}
\sphinxAtStartPar
\sphinxstylestrong{フレーム:}

\sphinxAtStartPar
・エクストラ

\sphinxAtStartPar
\sphinxstylestrong{トリガー:} 直接

\sphinxAtStartPar
\sphinxstylestrong{スピード:} 通常

\sphinxAtStartPar
\sphinxstylestrong{タイミング:} メイン

\sphinxAtStartPar
\sphinxstylestrong{コスト:} SS

\sphinxAtStartPar
\sphinxstylestrong{起動条件:}

\sphinxAtStartPar
自分の表切札に「魔王」がある

\sphinxAtStartPar
\sphinxstylestrong{効果:}

\sphinxAtStartPar
切札の{\normalsize $\spadesuit$} Kが表の場合、次を行う。
\begin{enumerate}
\sphinxsetlistlabels{\arabic}{enumi}{enumii}{}{.}%
\item {} 
\sphinxAtStartPar
全ての兵士を墓地に移す。

\item {} 
\sphinxAtStartPar
自分の全ての領域から{\normalsize $\clubsuit$} Kを探す。非公開の領域を見た場合、その領域をシャッフルする。

\item {} 
\sphinxAtStartPar
{\normalsize $\clubsuit$} Kを見つけた場合、切札の{\normalsize $\spadesuit$} Kと{\normalsize $\clubsuit$} Kを合わせて魔王としてチャージ状態で場に出す。

\item {} 
\sphinxAtStartPar
{\normalsize $\clubsuit$} Kが無い場合、切札の{\normalsize $\spadesuit$} Kを墓地に移す。

\end{enumerate}

\sphinxAtStartPar
\sphinxstylestrong{導入バージョン:}  ex3.10.0

\sphinxAtStartPar
\sphinxstylestrong{更新バージョン:}  v8.2


\bigskip\hrule\bigskip



\paragraph{魔王の囁き (誘発魔法)}
\label{\detokenize{auto/frameActionlist:act-darkwhisper}}\label{\detokenize{auto/frameActionlist:id30}}
\sphinxAtStartPar
\sphinxstylestrong{フレーム:}

\sphinxAtStartPar
・エクストラ

\sphinxAtStartPar
\sphinxstylestrong{トリガー:} 誘発

\sphinxAtStartPar
\sphinxstylestrong{スピード:} 即時

\sphinxAtStartPar
\sphinxstylestrong{タイミング:} クイック

\sphinxAtStartPar
\sphinxstylestrong{誘発条件:}

\sphinxAtStartPar
自分の表切札に「魔王」があるかつ、エンドアクションをリクエストした時に誘発する。

\sphinxAtStartPar
\sphinxstylestrong{効果:}

\sphinxAtStartPar
ライフの一番上からJoker,A,J,Q,Kのいずれかが出るまで墓地にカードを移動し、出たら手札に加える。

\sphinxAtStartPar
\sphinxstylestrong{導入バージョン:}  ex3.10.0

\sphinxAtStartPar
\sphinxstylestrong{更新バージョン:}  v8.2


\bigskip\hrule\bigskip



\subsubsection{切札{\normalsize $\heartsuit$} 2}
\label{\detokenize{auto/frameActionlist:id31}}

\paragraph{援軍 (通常魔法)}
\label{\detokenize{auto/frameActionlist:act-summonreinforcement}}\label{\detokenize{auto/frameActionlist:id32}}
\sphinxAtStartPar
\sphinxstylestrong{フレーム:}

\sphinxAtStartPar
・エクストラ

\sphinxAtStartPar
\sphinxstylestrong{トリガー:} 直接

\sphinxAtStartPar
\sphinxstylestrong{スピード:} 通常

\sphinxAtStartPar
\sphinxstylestrong{タイミング:} メイン

\sphinxAtStartPar
\sphinxstylestrong{コスト:} BL

\sphinxAtStartPar
\sphinxstylestrong{起動条件:}

\sphinxAtStartPar
・自分の表切札に「援軍」がある
・プレイヤーは、1ターンに1回までこのアクションを起こすことができる。

\sphinxAtStartPar
\sphinxstylestrong{効果:}

\sphinxAtStartPar
ライフの一番上から2枚カードをめくり1枚を兵士としてチャージ状態で場に出し、もう1枚をライフの一番下に移す。

\sphinxAtStartPar
\sphinxstylestrong{導入バージョン:}  ex3.10.0

\sphinxAtStartPar
\sphinxstylestrong{更新バージョン:}  v8.2


\bigskip\hrule\bigskip



\subsubsection{切札{\normalsize $\heartsuit$} 3}
\label{\detokenize{auto/frameActionlist:id33}}

\paragraph{激励 (兵士起因)}
\label{\detokenize{auto/frameActionlist:act-rally}}\label{\detokenize{auto/frameActionlist:id34}}
\sphinxAtStartPar
\sphinxstylestrong{フレーム:}

\sphinxAtStartPar
・エクストラ

\sphinxAtStartPar
\sphinxstylestrong{トリガー:} 直接

\sphinxAtStartPar
\sphinxstylestrong{スピード:} 即時

\sphinxAtStartPar
\sphinxstylestrong{タイミング:} クイック

\sphinxAtStartPar
\sphinxstylestrong{キーカード:} \textless{}激励\textgreater{} を持つ兵士1体

\sphinxAtStartPar
\sphinxstylestrong{コスト:} K

\sphinxAtStartPar
\sphinxstylestrong{起動条件:}

\sphinxAtStartPar
自分の表切札に「激励」がある

\sphinxAtStartPar
\sphinxstylestrong{対象:}

\sphinxAtStartPar
兵士1体

\sphinxAtStartPar
\sphinxstylestrong{効果:}

\sphinxAtStartPar
対象の兵士のサイズをこのターンが終わるまで X 加算する。X はキーカードのサイズに等しい。

\sphinxAtStartPar
\sphinxstylestrong{導入バージョン:}  ex3.4.0

\sphinxAtStartPar
\sphinxstylestrong{更新バージョン:}  v8.1


\bigskip\hrule\bigskip



\subsubsection{切札{\normalsize $\heartsuit$} 4}
\label{\detokenize{auto/frameActionlist:id35}}

\paragraph{休戦 (速攻魔法)}
\label{\detokenize{auto/frameActionlist:act-retreat}}\label{\detokenize{auto/frameActionlist:id36}}
\sphinxAtStartPar
\sphinxstylestrong{フレーム:}

\sphinxAtStartPar
・エクストラ

\sphinxAtStartPar
\sphinxstylestrong{トリガー:} 直接

\sphinxAtStartPar
\sphinxstylestrong{スピード:} 通常

\sphinxAtStartPar
\sphinxstylestrong{タイミング:} クイック

\sphinxAtStartPar
\sphinxstylestrong{キーカード:} \textless{}激励\textgreater{} を持つ兵士1体

\sphinxAtStartPar
\sphinxstylestrong{コスト:} BL

\sphinxAtStartPar
\sphinxstylestrong{起動条件:}

\sphinxAtStartPar
自分の表切札に「休戦」がある

\sphinxAtStartPar
\sphinxstylestrong{対象:}

\sphinxAtStartPar
アタックアクションのリクエスト

\sphinxAtStartPar
\sphinxstylestrong{効果:}
\begin{enumerate}
\sphinxsetlistlabels{\arabic}{enumi}{enumii}{}{.}%
\item {} 
\sphinxAtStartPar
対象のリクエストを無効し、対象のリクエストをステージから取り除く。

\item {} 
\sphinxAtStartPar
自分の切札にある{\normalsize $\heartsuit$} 4を裏向きにする。

\item {} 
\sphinxAtStartPar
自分の場にいるキャラクターを全てチャージ状態にする。

\end{enumerate}

\sphinxAtStartPar
\sphinxstylestrong{導入バージョン:}  ex5.30.0

\sphinxAtStartPar
\sphinxstylestrong{更新バージョン:}  v8.2


\bigskip\hrule\bigskip



\subsubsection{切札{\normalsize $\heartsuit$} 5}
\label{\detokenize{auto/frameActionlist:id37}}

\paragraph{追撃 (速攻魔法)}
\label{\detokenize{auto/frameActionlist:act-pursuit}}\label{\detokenize{auto/frameActionlist:id38}}
\sphinxAtStartPar
\sphinxstylestrong{フレーム:}

\sphinxAtStartPar
・エクストラ

\sphinxAtStartPar
\sphinxstylestrong{トリガー:} 直接

\sphinxAtStartPar
\sphinxstylestrong{スピード:} 通常

\sphinxAtStartPar
\sphinxstylestrong{タイミング:} メイン

\sphinxAtStartPar
\sphinxstylestrong{コスト:} BDD

\sphinxAtStartPar
\sphinxstylestrong{起動条件:}

\sphinxAtStartPar
・自分の表切札に「追撃」がある
・プレイヤーは、1ターンに1回までこのアクションを起こすことができる。

\sphinxAtStartPar
\sphinxstylestrong{効果:}
\begin{enumerate}
\sphinxsetlistlabels{\arabic}{enumi}{enumii}{}{.}%
\item {} 
\sphinxAtStartPar
自分の場にいる兵士を全てチャージ状態にする。

\item {} 
\sphinxAtStartPar
アタックアクションをリクエストする。

\end{enumerate}

\sphinxAtStartPar
\sphinxstylestrong{導入バージョン:}  ex5.30.0

\sphinxAtStartPar
\sphinxstylestrong{更新バージョン:}  v8.2


\bigskip\hrule\bigskip



\subsubsection{切札{\normalsize $\heartsuit$} 7}
\label{\detokenize{auto/frameActionlist:id39}}

\paragraph{ファントムリバース (通常魔法)}
\label{\detokenize{auto/frameActionlist:act-phantomreverse}}\label{\detokenize{auto/frameActionlist:id40}}
\sphinxAtStartPar
\sphinxstylestrong{フレーム:}

\sphinxAtStartPar
・エクストラ

\sphinxAtStartPar
\sphinxstylestrong{トリガー:} 直接

\sphinxAtStartPar
\sphinxstylestrong{スピード:} 通常

\sphinxAtStartPar
\sphinxstylestrong{タイミング:} メイン

\sphinxAtStartPar
\sphinxstylestrong{キーカード:} \textless{}ファントム\textgreater{} を持つ防壁1体

\sphinxAtStartPar
\sphinxstylestrong{コスト:} K

\sphinxAtStartPar
\sphinxstylestrong{起動条件:}

\sphinxAtStartPar
自分の表切札に「ファントム」がある

\sphinxAtStartPar
\sphinxstylestrong{効果:}

\sphinxAtStartPar
キーカードの防壁をチャージ状態の兵士にする。

\sphinxAtStartPar
\sphinxstylestrong{導入バージョン:}  ex5.30.0

\sphinxAtStartPar
\sphinxstylestrong{更新バージョン:}  v8.2


\bigskip\hrule\bigskip



\paragraph{幻影 (誘発魔法)}
\label{\detokenize{auto/frameActionlist:act-illusion}}\label{\detokenize{auto/frameActionlist:id41}}
\sphinxAtStartPar
\sphinxstylestrong{フレーム:}

\sphinxAtStartPar
・エクストラ

\sphinxAtStartPar
\sphinxstylestrong{トリガー:} 誘発

\sphinxAtStartPar
\sphinxstylestrong{スピード:} 通常

\sphinxAtStartPar
\sphinxstylestrong{タイミング:} クイック

\sphinxAtStartPar
\sphinxstylestrong{コスト:} K

\sphinxAtStartPar
\sphinxstylestrong{誘発条件:}

\sphinxAtStartPar
自分の表切札に「ファントム」があるかつ、自分のエンドアクションのリクエストがステージに追加された時に誘発する。

\sphinxAtStartPar
\sphinxstylestrong{効果:}

\sphinxAtStartPar
1．自分の場にいる兵士を好きな数だけ選び、チャージ状態の防壁にする。兵士の時に受けた効果、能力は無くなる。兵士が複数のカードから成る場合、1枚ずつ防壁にする。

\sphinxAtStartPar
2．自分の場にいる防壁の順番を任意の順番に並び替えてもよい。

\sphinxAtStartPar
\sphinxstylestrong{導入バージョン:}  ex5.30.0

\sphinxAtStartPar
\sphinxstylestrong{更新バージョン:}  v8.2


\bigskip\hrule\bigskip



\subsubsection{切札{\normalsize $\heartsuit$} 8}
\label{\detokenize{auto/frameActionlist:id42}}

\paragraph{救済 (通常魔法)}
\label{\detokenize{auto/frameActionlist:act-salvation}}\label{\detokenize{auto/frameActionlist:id43}}
\sphinxAtStartPar
\sphinxstylestrong{フレーム:}

\sphinxAtStartPar
・エクストラ

\sphinxAtStartPar
\sphinxstylestrong{トリガー:} 直接

\sphinxAtStartPar
\sphinxstylestrong{スピード:} 通常

\sphinxAtStartPar
\sphinxstylestrong{タイミング:} メイン

\sphinxAtStartPar
\sphinxstylestrong{キーカード:} {\normalsize $\heartsuit$} A〜K

\sphinxAtStartPar
\sphinxstylestrong{コスト:} B

\sphinxAtStartPar
\sphinxstylestrong{起動条件:}

\sphinxAtStartPar
自分の表切札に「救済」がある

\sphinxAtStartPar
\sphinxstylestrong{効果:}

\sphinxAtStartPar
キーカードの数字以上の全ての兵士をオーナーのライフの一番上にオーナーの好きな順で移す。

\sphinxAtStartPar
\sphinxstylestrong{導入バージョン:}  ex5.30.0

\sphinxAtStartPar
\sphinxstylestrong{更新バージョン:}  v8.2


\bigskip\hrule\bigskip



\subsubsection{切札{\normalsize $\heartsuit$} J}
\label{\detokenize{auto/frameActionlist:id44}}

\paragraph{チャリオット召喚 (召喚)}
\label{\detokenize{auto/frameActionlist:act-summonchariot}}\label{\detokenize{auto/frameActionlist:id45}}
\sphinxAtStartPar
\sphinxstylestrong{フレーム:}

\sphinxAtStartPar
・エクストラ

\sphinxAtStartPar
\sphinxstylestrong{トリガー:} 直接

\sphinxAtStartPar
\sphinxstylestrong{スピード:} 通常

\sphinxAtStartPar
\sphinxstylestrong{タイミング:} メイン

\sphinxAtStartPar
\sphinxstylestrong{キーカード:} {\normalsize $\diamondsuit$} 2〜10

\sphinxAtStartPar
\sphinxstylestrong{コスト:} BL

\sphinxAtStartPar
\sphinxstylestrong{起動条件:}

\sphinxAtStartPar
自分の表切札に「チャリオット」がある

\sphinxAtStartPar
\sphinxstylestrong{効果:}

\sphinxAtStartPar
切札の{\normalsize $\heartsuit$} Jが表の場合、次を行う。
\begin{enumerate}
\sphinxsetlistlabels{\arabic}{enumi}{enumii}{}{.}%
\item {} 
\sphinxAtStartPar
防壁1体を選び墓地に移す。

\item {} 
\sphinxAtStartPar
キーカードと切札の{\normalsize $\heartsuit$} Jをあわせてチャリオットとしてチャージ状態で場に出す。

\end{enumerate}

\sphinxAtStartPar
\sphinxstylestrong{導入バージョン:}  ex4.22.0

\sphinxAtStartPar
\sphinxstylestrong{更新バージョン:}  v8.2


\bigskip\hrule\bigskip



\paragraph{クイックチャリオット召喚 (召喚)}
\label{\detokenize{auto/frameActionlist:act-quicksummonchariot}}\label{\detokenize{auto/frameActionlist:id46}}
\sphinxAtStartPar
\sphinxstylestrong{フレーム:}

\sphinxAtStartPar
・エクストラ

\sphinxAtStartPar
\sphinxstylestrong{トリガー:} 直接

\sphinxAtStartPar
\sphinxstylestrong{スピード:} 通常

\sphinxAtStartPar
\sphinxstylestrong{タイミング:} クイック

\sphinxAtStartPar
\sphinxstylestrong{キーカード:} {\normalsize $\diamondsuit$} 2〜10

\sphinxAtStartPar
\sphinxstylestrong{コスト:} D

\sphinxAtStartPar
\sphinxstylestrong{起動条件:}

\sphinxAtStartPar
自分の表切札に「チャリオット」がある

\sphinxAtStartPar
\sphinxstylestrong{効果:}

\sphinxAtStartPar
効果の内容は「チャリオット召喚」を参照。

\sphinxAtStartPar
\sphinxstylestrong{導入バージョン:}  ex4.22.0

\sphinxAtStartPar
\sphinxstylestrong{更新バージョン:}  v8.2


\bigskip\hrule\bigskip



\paragraph{チャリオット防壁破壊 (兵士起因)}
\label{\detokenize{auto/frameActionlist:act-chariotdestroy}}\label{\detokenize{auto/frameActionlist:id47}}
\sphinxAtStartPar
\sphinxstylestrong{フレーム:}

\sphinxAtStartPar
・エクストラ

\sphinxAtStartPar
\sphinxstylestrong{トリガー:} 直接

\sphinxAtStartPar
\sphinxstylestrong{スピード:} 通常

\sphinxAtStartPar
\sphinxstylestrong{タイミング:} メイン

\sphinxAtStartPar
\sphinxstylestrong{キーカード:} 「チャリオット」1体

\sphinxAtStartPar
\sphinxstylestrong{コスト:} K

\sphinxAtStartPar
\sphinxstylestrong{対象:}

\sphinxAtStartPar
防壁1体

\sphinxAtStartPar
\sphinxstylestrong{効果:}

\sphinxAtStartPar
効果の内容は「防壁破壊」を参照。

\sphinxAtStartPar
\sphinxstylestrong{導入バージョン:}  ex4.22.0

\sphinxAtStartPar
\sphinxstylestrong{更新バージョン:}  v8.2


\bigskip\hrule\bigskip



\paragraph{チャリオット再生 (切札操作)}
\label{\detokenize{auto/frameActionlist:act-chariotrevive}}\label{\detokenize{auto/frameActionlist:id48}}
\sphinxAtStartPar
\sphinxstylestrong{フレーム:}

\sphinxAtStartPar
・エクストラ

\sphinxAtStartPar
\sphinxstylestrong{トリガー:} 誘発

\sphinxAtStartPar
\sphinxstylestrong{スピード:} 即時

\sphinxAtStartPar
\sphinxstylestrong{タイミング:} クイック

\sphinxAtStartPar
\sphinxstylestrong{誘発条件:}

\sphinxAtStartPar
自分の 「チャリオット」 が場から墓地に移る場合

\sphinxAtStartPar
\sphinxstylestrong{効果:}

\sphinxAtStartPar
墓地から{\normalsize $\heartsuit$} Jを1枚選び、自分の表切札として置く。


\bigskip\hrule\bigskip



\subsubsection{切札{\normalsize $\heartsuit$} K}
\label{\detokenize{auto/frameActionlist:id49}}

\paragraph{巨人召喚 (召喚)}
\label{\detokenize{auto/frameActionlist:act-summongiant}}\label{\detokenize{auto/frameActionlist:id50}}
\sphinxAtStartPar
\sphinxstylestrong{フレーム:}

\sphinxAtStartPar
・エクストラ

\sphinxAtStartPar
\sphinxstylestrong{トリガー:} 直接

\sphinxAtStartPar
\sphinxstylestrong{スピード:} 通常

\sphinxAtStartPar
\sphinxstylestrong{タイミング:} メイン

\sphinxAtStartPar
\sphinxstylestrong{コスト:} SS

\sphinxAtStartPar
\sphinxstylestrong{起動条件:}

\sphinxAtStartPar
自分の表切札に「巨人」がある

\sphinxAtStartPar
\sphinxstylestrong{効果:}

\sphinxAtStartPar
切札の{\normalsize $\heartsuit$} Kが表の場合、次を行う。
\begin{enumerate}
\sphinxsetlistlabels{\arabic}{enumi}{enumii}{}{.}%
\item {} 
\sphinxAtStartPar
全ての防壁を墓地に移す。

\item {} 
\sphinxAtStartPar
自分の全ての領域から{\normalsize $\diamondsuit$} Kを探す。非公開の領域を見た場合、その領域をシャッフルする。

\item {} 
\sphinxAtStartPar
{\normalsize $\diamondsuit$} Kを見つけた場合、切札の{\normalsize $\heartsuit$} Kと{\normalsize $\diamondsuit$} Kを合わせて巨人としてチャージ状態で場に出す。

\item {} 
\sphinxAtStartPar
{\normalsize $\diamondsuit$} Kが無い場合、切札の{\normalsize $\heartsuit$} Kを墓地に移す。

\end{enumerate}

\sphinxAtStartPar
\sphinxstylestrong{導入バージョン:}  ex3.10.0

\sphinxAtStartPar
\sphinxstylestrong{更新バージョン:}  v8.2


\bigskip\hrule\bigskip



\subsubsection{切札{\normalsize $\diamondsuit$} A}
\label{\detokenize{auto/frameActionlist:id51}}

\paragraph{防壁追加 (通常魔法)}
\label{\detokenize{auto/frameActionlist:act-putbulwark}}\label{\detokenize{auto/frameActionlist:id52}}
\sphinxAtStartPar
\sphinxstylestrong{フレーム:}

\sphinxAtStartPar
・エクストラ

\sphinxAtStartPar
\sphinxstylestrong{トリガー:} 直接

\sphinxAtStartPar
\sphinxstylestrong{スピード:} 即時

\sphinxAtStartPar
\sphinxstylestrong{タイミング:} メイン

\sphinxAtStartPar
\sphinxstylestrong{コスト:} L

\sphinxAtStartPar
\sphinxstylestrong{起動条件:}

\sphinxAtStartPar
・自分の表切札に「防壁追加」がある
・プレイヤーは、1ターンに1回までこのアクションを起こすことができる。

\sphinxAtStartPar
\sphinxstylestrong{効果:}

\sphinxAtStartPar
対戦相手の場にいる防壁の数が自分の場にいる防壁の数以上の場合、自分のライフの一番上から1枚を防壁として裏向きかつチャージ状態で場に出す。

\sphinxAtStartPar
\sphinxstylestrong{導入バージョン:}  ex3.10.0

\sphinxAtStartPar
\sphinxstylestrong{更新バージョン:}  v8.2


\bigskip\hrule\bigskip



\subsubsection{切札{\normalsize $\diamondsuit$} 3}
\label{\detokenize{auto/frameActionlist:id53}}

\paragraph{交渉 (兵士起因)}
\label{\detokenize{auto/frameActionlist:act-negotiate}}\label{\detokenize{auto/frameActionlist:id54}}
\sphinxAtStartPar
\sphinxstylestrong{フレーム:}

\sphinxAtStartPar
・エクストラ

\sphinxAtStartPar
\sphinxstylestrong{トリガー:} 直接

\sphinxAtStartPar
\sphinxstylestrong{スピード:} 即時

\sphinxAtStartPar
\sphinxstylestrong{タイミング:} クイック

\sphinxAtStartPar
\sphinxstylestrong{キーカード:} \textless{}交渉\textgreater{} を持つ兵士1体

\sphinxAtStartPar
\sphinxstylestrong{コスト:} K

\sphinxAtStartPar
\sphinxstylestrong{起動条件:}

\sphinxAtStartPar
自分の表切札に「交渉」がある

\sphinxAtStartPar
\sphinxstylestrong{対象:}

\sphinxAtStartPar
キャラクター1体

\sphinxAtStartPar
\sphinxstylestrong{効果:}

\sphinxAtStartPar
対象のキャラクターをドライブする。

\sphinxAtStartPar
\sphinxstylestrong{導入バージョン:}  ex3.4.0

\sphinxAtStartPar
\sphinxstylestrong{更新バージョン:}  v8.1


\bigskip\hrule\bigskip



\subsubsection{切札{\normalsize $\diamondsuit$} 4}
\label{\detokenize{auto/frameActionlist:id55}}

\paragraph{相殺 (速攻魔法)}
\label{\detokenize{auto/frameActionlist:act-reckoning}}\label{\detokenize{auto/frameActionlist:id56}}
\sphinxAtStartPar
\sphinxstylestrong{フレーム:}

\sphinxAtStartPar
・エクストラ

\sphinxAtStartPar
\sphinxstylestrong{トリガー:} 直接

\sphinxAtStartPar
\sphinxstylestrong{スピード:} 即時

\sphinxAtStartPar
\sphinxstylestrong{タイミング:} クイック

\sphinxAtStartPar
\sphinxstylestrong{コスト:} D

\sphinxAtStartPar
\sphinxstylestrong{起動条件:}

\sphinxAtStartPar
自分の表切札に「防壁追加」がある

\sphinxAtStartPar
\sphinxstylestrong{対象:}

\sphinxAtStartPar
キーカードがあるリクエスト

\sphinxAtStartPar
\sphinxstylestrong{効果:}

\sphinxAtStartPar
対象リクエストのキーカードの数字の合計値が5以上の場合、次を行う。

\sphinxAtStartPar
1.対象リクエストのキーカードの数字の合計値分ダメージを受ける。

\sphinxAtStartPar
2.対象リクエストをステージから取り除き、対象リクエストのキーカードを墓地に移す。

\sphinxAtStartPar
\sphinxstylestrong{導入バージョン:}  ex5.30.0

\sphinxAtStartPar
\sphinxstylestrong{更新バージョン:}  v8.2


\bigskip\hrule\bigskip



\subsubsection{切札{\normalsize $\diamondsuit$} 5}
\label{\detokenize{auto/frameActionlist:id57}}

\paragraph{疲弊 (誘発魔法)}
\label{\detokenize{auto/frameActionlist:act-exhaust}}\label{\detokenize{auto/frameActionlist:id58}}
\sphinxAtStartPar
\sphinxstylestrong{フレーム:}

\sphinxAtStartPar
・エクストラ

\sphinxAtStartPar
\sphinxstylestrong{トリガー:} 誘発

\sphinxAtStartPar
\sphinxstylestrong{スピード:} 通常

\sphinxAtStartPar
\sphinxstylestrong{タイミング:} クイック

\sphinxAtStartPar
\sphinxstylestrong{起動条件:}

\sphinxAtStartPar
自分の表切札に「喪失」があるかつ、エンドアクションをリクエストした時に誘発する。

\sphinxAtStartPar
\sphinxstylestrong{効果:}

\sphinxAtStartPar
自分の表切札にある「喪失」を裏切札にする。そうしたくない場合、2点のダメージを受ける。

\sphinxAtStartPar
\sphinxstylestrong{導入バージョン:}  ex5.30.0

\sphinxAtStartPar
\sphinxstylestrong{更新バージョン:}  v8.2


\bigskip\hrule\bigskip



\subsubsection{切札{\normalsize $\diamondsuit$} 8}
\label{\detokenize{auto/frameActionlist:id59}}

\paragraph{修繕 (通常魔法)}
\label{\detokenize{auto/frameActionlist:act-repair}}\label{\detokenize{auto/frameActionlist:id60}}
\sphinxAtStartPar
\sphinxstylestrong{フレーム:}

\sphinxAtStartPar
・エクストラ

\sphinxAtStartPar
\sphinxstylestrong{トリガー:} 直接

\sphinxAtStartPar
\sphinxstylestrong{スピード:} 即時

\sphinxAtStartPar
\sphinxstylestrong{タイミング:} メイン

\sphinxAtStartPar
\sphinxstylestrong{コスト:} L

\sphinxAtStartPar
\sphinxstylestrong{起動条件:}

\sphinxAtStartPar
自分の表切札に「修繕」がある

\sphinxAtStartPar
\sphinxstylestrong{対象:}

\sphinxAtStartPar
自分の場にいる{\normalsize $\diamondsuit$} 2〜10のカードを含んだ兵士1体

\sphinxAtStartPar
\sphinxstylestrong{効果:}

\sphinxAtStartPar
対象の兵士が装備アクションの対象にできる場合、次を行う。
\begin{enumerate}
\sphinxsetlistlabels{\arabic}{enumi}{enumii}{}{.}%
\item {} 
\sphinxAtStartPar
自分の墓地から{\normalsize $\diamondsuit$} J〜Kを1枚選ぶ。

\item {} 
\sphinxAtStartPar
選べた場合、「装備」アクションをコストを支払わずにリクエストする。対象は対象の兵士、キーカードは選んだカードとする。

\item {} 
\sphinxAtStartPar
選べなかった場合、対象の兵士を墓地に移す。

\end{enumerate}

\sphinxAtStartPar
\sphinxstylestrong{導入バージョン:}  ex4.14.0

\sphinxAtStartPar
\sphinxstylestrong{更新バージョン:}  v8.2


\bigskip\hrule\bigskip



\subsubsection{切札{\normalsize $\diamondsuit$} 9}
\label{\detokenize{auto/frameActionlist:id61}}

\paragraph{休息 (通常魔法)}
\label{\detokenize{auto/frameActionlist:act-rest}}\label{\detokenize{auto/frameActionlist:id62}}
\sphinxAtStartPar
\sphinxstylestrong{フレーム:}

\sphinxAtStartPar
・エクストラ

\sphinxAtStartPar
\sphinxstylestrong{トリガー:} 直接

\sphinxAtStartPar
\sphinxstylestrong{スピード:} 通常

\sphinxAtStartPar
\sphinxstylestrong{タイミング:} メイン

\sphinxAtStartPar
\sphinxstylestrong{コスト:} B

\sphinxAtStartPar
\sphinxstylestrong{起動条件:}

\sphinxAtStartPar
自分の表切札に「充足」がある

\sphinxAtStartPar
\sphinxstylestrong{効果:}

\sphinxAtStartPar
自分の場にいる全ての兵士をチャージ状態にする。

\sphinxAtStartPar
\sphinxstylestrong{導入バージョン:}  ex4.18.0

\sphinxAtStartPar
\sphinxstylestrong{更新バージョン:}  v8.2


\bigskip\hrule\bigskip



\paragraph{充足 (誘発魔法)}
\label{\detokenize{auto/frameActionlist:act-satisfaction}}\label{\detokenize{auto/frameActionlist:id63}}
\sphinxAtStartPar
\sphinxstylestrong{フレーム:}

\sphinxAtStartPar
・エクストラ

\sphinxAtStartPar
\sphinxstylestrong{トリガー:} 誘発

\sphinxAtStartPar
\sphinxstylestrong{スピード:} 通常

\sphinxAtStartPar
\sphinxstylestrong{タイミング:} クイック

\sphinxAtStartPar
\sphinxstylestrong{誘発条件:}

\sphinxAtStartPar
自分の表切札に「充足」があるかつ、エンドアクションをリクエストした時に誘発する。

\sphinxAtStartPar
\sphinxstylestrong{効果:}

\sphinxAtStartPar
自分の場にいる全てのキャラクターをチャージ状態にする。

\sphinxAtStartPar
\sphinxstylestrong{導入バージョン:}  ex4.18.0

\sphinxAtStartPar
\sphinxstylestrong{更新バージョン:}  v8.2


\bigskip\hrule\bigskip



\subsubsection{切札{\normalsize $\diamondsuit$} J}
\label{\detokenize{auto/frameActionlist:id64}}

\paragraph{策士召喚 (召喚)}
\label{\detokenize{auto/frameActionlist:act-summonstrategist}}\label{\detokenize{auto/frameActionlist:id65}}
\sphinxAtStartPar
\sphinxstylestrong{フレーム:}

\sphinxAtStartPar
・エクストラ

\sphinxAtStartPar
\sphinxstylestrong{トリガー:} 直接

\sphinxAtStartPar
\sphinxstylestrong{スピード:} 通常

\sphinxAtStartPar
\sphinxstylestrong{タイミング:} メイン

\sphinxAtStartPar
\sphinxstylestrong{キーカード:} {\normalsize $\clubsuit$} 2〜10

\sphinxAtStartPar
\sphinxstylestrong{コスト:} BL

\sphinxAtStartPar
\sphinxstylestrong{起動条件:}

\sphinxAtStartPar
自分の表切札に「策士」がある

\sphinxAtStartPar
\sphinxstylestrong{効果:}

\sphinxAtStartPar
切札の{\normalsize $\diamondsuit$} Jが表の場合、次を行う。
\begin{enumerate}
\sphinxsetlistlabels{\arabic}{enumi}{enumii}{}{.}%
\item {} 
\sphinxAtStartPar
ステージ上にあるリクエストを1つ選び無効化しステージから取り除く。そのリクエストは次のいずれかに該当するものとする。Xはこのリクエストのキーカードの数字とする。

\sphinxAtStartPar
・キーカードが1枚かつそのキーカードの数字がX以下

\sphinxAtStartPar
・キーカードが2枚

\item {} 
\sphinxAtStartPar
キーカードと切札の{\normalsize $\diamondsuit$} Jをあわせて策士としてチャージ状態で場に出す。

\end{enumerate}

\sphinxAtStartPar
\sphinxstylestrong{導入バージョン:}  ex4.22.0

\sphinxAtStartPar
\sphinxstylestrong{更新バージョン:}  v8.2


\bigskip\hrule\bigskip



\paragraph{クイック策士召喚 (召喚)}
\label{\detokenize{auto/frameActionlist:act-quicksummonstrategist}}\label{\detokenize{auto/frameActionlist:id66}}
\sphinxAtStartPar
\sphinxstylestrong{フレーム:}

\sphinxAtStartPar
・エクストラ

\sphinxAtStartPar
\sphinxstylestrong{トリガー:} 直接

\sphinxAtStartPar
\sphinxstylestrong{スピード:} 通常

\sphinxAtStartPar
\sphinxstylestrong{タイミング:} クイック

\sphinxAtStartPar
\sphinxstylestrong{キーカード:} {\normalsize $\clubsuit$} 2〜10

\sphinxAtStartPar
\sphinxstylestrong{コスト:} D

\sphinxAtStartPar
\sphinxstylestrong{起動条件:}

\sphinxAtStartPar
自分の表切札に「策士」がある

\sphinxAtStartPar
\sphinxstylestrong{効果:}

\sphinxAtStartPar
効果の内容は「策士召喚」を参照。

\sphinxAtStartPar
\sphinxstylestrong{導入バージョン:}  ex4.22.0

\sphinxAtStartPar
\sphinxstylestrong{更新バージョン:}  v8.2


\bigskip\hrule\bigskip



\paragraph{先読み (兵士起因)}
\label{\detokenize{auto/frameActionlist:act-strategistpredict}}\label{\detokenize{auto/frameActionlist:id67}}
\sphinxAtStartPar
\sphinxstylestrong{フレーム:}

\sphinxAtStartPar
・エクストラ

\sphinxAtStartPar
\sphinxstylestrong{トリガー:} 直接

\sphinxAtStartPar
\sphinxstylestrong{スピード:} 即時

\sphinxAtStartPar
\sphinxstylestrong{タイミング:} クイック

\sphinxAtStartPar
\sphinxstylestrong{キーカード:} 「策士」1体

\sphinxAtStartPar
\sphinxstylestrong{コスト:} K

\sphinxAtStartPar
\sphinxstylestrong{効果:}
\begin{enumerate}
\sphinxsetlistlabels{\arabic}{enumi}{enumii}{}{.}%
\item {} 
\sphinxAtStartPar
ライフの一番上からカードを2枚引き手札に加える。

\item {} 
\sphinxAtStartPar
手札から1枚選びライフの一番上もしくは一番下に移す。

\end{enumerate}

\sphinxAtStartPar
\sphinxstylestrong{導入バージョン:}  ex4.22.0

\sphinxAtStartPar
\sphinxstylestrong{更新バージョン:}  v8.2


\bigskip\hrule\bigskip



\paragraph{策士再生 (切札操作)}
\label{\detokenize{auto/frameActionlist:act-strategistrevive}}\label{\detokenize{auto/frameActionlist:id68}}
\sphinxAtStartPar
\sphinxstylestrong{フレーム:}

\sphinxAtStartPar
・エクストラ

\sphinxAtStartPar
\sphinxstylestrong{トリガー:} 誘発

\sphinxAtStartPar
\sphinxstylestrong{スピード:} 即時

\sphinxAtStartPar
\sphinxstylestrong{タイミング:} クイック

\sphinxAtStartPar
\sphinxstylestrong{誘発条件:}

\sphinxAtStartPar
自分の 「策士」 が場から墓地に移る場合

\sphinxAtStartPar
\sphinxstylestrong{効果:}

\sphinxAtStartPar
墓地から{\normalsize $\diamondsuit$} Jを1枚選び、自分の表切札として置く。


\bigskip\hrule\bigskip



\subsubsection{切札{\normalsize $\clubsuit$} 4}
\label{\detokenize{auto/frameActionlist:id69}}

\paragraph{スイッチ (兵士起因)}
\label{\detokenize{auto/frameActionlist:act-switch}}\label{\detokenize{auto/frameActionlist:id70}}
\sphinxAtStartPar
\sphinxstylestrong{フレーム:}

\sphinxAtStartPar
・エクストラ

\sphinxAtStartPar
\sphinxstylestrong{トリガー:} 直接

\sphinxAtStartPar
\sphinxstylestrong{スピード:} 即時

\sphinxAtStartPar
\sphinxstylestrong{タイミング:} クイック

\sphinxAtStartPar
\sphinxstylestrong{キーカード:} \textless{}スイッチ\textgreater{} を持つ兵士1体

\sphinxAtStartPar
\sphinxstylestrong{コスト:} K

\sphinxAtStartPar
\sphinxstylestrong{起動条件:}

\sphinxAtStartPar
自分の表切札に「スイッチ」がある

\sphinxAtStartPar
\sphinxstylestrong{効果:}
\begin{enumerate}
\sphinxsetlistlabels{\arabic}{enumi}{enumii}{}{.}%
\item {} 
\sphinxAtStartPar
手札からカードを1枚選び、兵士としてチャージ状態で場に出す。

\item {} 
\sphinxAtStartPar
キーカードを手札に戻す。

\end{enumerate}

\sphinxAtStartPar
\sphinxstylestrong{導入バージョン:}  ex3.4.0

\sphinxAtStartPar
\sphinxstylestrong{更新バージョン:}  v8.1


\bigskip\hrule\bigskip



\subsubsection{切札{\normalsize $\clubsuit$} 5}
\label{\detokenize{auto/frameActionlist:id71}}

\paragraph{内乱 (通常魔法)}
\label{\detokenize{auto/frameActionlist:act-civilwar}}\label{\detokenize{auto/frameActionlist:id72}}
\sphinxAtStartPar
\sphinxstylestrong{フレーム:}

\sphinxAtStartPar
・エクストラ

\sphinxAtStartPar
\sphinxstylestrong{トリガー:} 直接

\sphinxAtStartPar
\sphinxstylestrong{スピード:} 通常

\sphinxAtStartPar
\sphinxstylestrong{タイミング:} メイン

\sphinxAtStartPar
\sphinxstylestrong{キーカード:} {\normalsize $\clubsuit$} A〜K

\sphinxAtStartPar
\sphinxstylestrong{コスト:} B

\sphinxAtStartPar
\sphinxstylestrong{起動条件:}

\sphinxAtStartPar
自分の表切札に「内乱」がある

\sphinxAtStartPar
\sphinxstylestrong{対象:}

\sphinxAtStartPar
対戦相手1人

\sphinxAtStartPar
\sphinxstylestrong{効果:}

\sphinxAtStartPar
対象の対戦相手の手札の枚数が自分の手札の枚数以上の場合、次を行う。

\sphinxAtStartPar
・対象の対戦相手は手札を2枚捨てる。

\sphinxAtStartPar
\sphinxstylestrong{導入バージョン:}  ex4.14.0

\sphinxAtStartPar
\sphinxstylestrong{更新バージョン:}  v8.2


\bigskip\hrule\bigskip



\paragraph{暴動 (誘発魔法)}
\label{\detokenize{auto/frameActionlist:act-riot}}\label{\detokenize{auto/frameActionlist:id73}}
\sphinxAtStartPar
\sphinxstylestrong{フレーム:}

\sphinxAtStartPar
・エクストラ

\sphinxAtStartPar
\sphinxstylestrong{トリガー:} 誘発

\sphinxAtStartPar
\sphinxstylestrong{スピード:} 即時

\sphinxAtStartPar
\sphinxstylestrong{タイミング:} クイック

\sphinxAtStartPar
\sphinxstylestrong{誘発条件:}

\sphinxAtStartPar
自分の表切札に「内乱」があるかつ、対戦相手が手札を捨てた時に誘発する。

\sphinxAtStartPar
\sphinxstylestrong{効果:}

\sphinxAtStartPar
このアクションを誘発した時に捨てられた手札の枚数×2点のダメージを、手札を捨てたプレイヤーに与える。

\sphinxAtStartPar
\sphinxstylestrong{導入バージョン:}  ex4.14.0

\sphinxAtStartPar
\sphinxstylestrong{更新バージョン:}  v8.2


\bigskip\hrule\bigskip



\subsubsection{切札{\normalsize $\clubsuit$} 7}
\label{\detokenize{auto/frameActionlist:id74}}

\paragraph{手札補充 (誘発魔法)}
\label{\detokenize{auto/frameActionlist:act-handrefill}}\label{\detokenize{auto/frameActionlist:id75}}
\sphinxAtStartPar
\sphinxstylestrong{フレーム:}

\sphinxAtStartPar
・エクストラ

\sphinxAtStartPar
\sphinxstylestrong{トリガー:} 誘発

\sphinxAtStartPar
\sphinxstylestrong{スピード:} 通常

\sphinxAtStartPar
\sphinxstylestrong{タイミング:} クイック

\sphinxAtStartPar
\sphinxstylestrong{誘発条件:}

\sphinxAtStartPar
自分の表切札に「手札補充」があるかつ、エンドアクションをリクエストした時に誘発する。

\sphinxAtStartPar
\sphinxstylestrong{効果:}

\sphinxAtStartPar
自分の手札が2枚以下の場合、2枚カードを引き、3枚の場合1枚カードを引く。

\sphinxAtStartPar
\sphinxstylestrong{導入バージョン:}  ex3.10.0

\sphinxAtStartPar
\sphinxstylestrong{更新バージョン:}  v8.2


\bigskip\hrule\bigskip



\subsubsection{切札{\normalsize $\clubsuit$} 8}
\label{\detokenize{auto/frameActionlist:id76}}

\paragraph{リセット (通常魔法)}
\label{\detokenize{auto/frameActionlist:act-reset}}\label{\detokenize{auto/frameActionlist:id77}}
\sphinxAtStartPar
\sphinxstylestrong{フレーム:}

\sphinxAtStartPar
・エクストラ

\sphinxAtStartPar
\sphinxstylestrong{トリガー:} 直接

\sphinxAtStartPar
\sphinxstylestrong{スピード:} 通常

\sphinxAtStartPar
\sphinxstylestrong{タイミング:} メイン

\sphinxAtStartPar
\sphinxstylestrong{キーカード:} {\normalsize $\clubsuit$} A〜K

\sphinxAtStartPar
\sphinxstylestrong{コスト:} L

\sphinxAtStartPar
\sphinxstylestrong{起動条件:}

\sphinxAtStartPar
自分の表切札に「リセット」がある

\sphinxAtStartPar
\sphinxstylestrong{効果:}
\begin{enumerate}
\sphinxsetlistlabels{\arabic}{enumi}{enumii}{}{.}%
\item {} 
\sphinxAtStartPar
全ての防壁と兵士を墓地に移す。

\item {} 
\sphinxAtStartPar
全ての切札を裏切札にする。

\end{enumerate}

\sphinxAtStartPar
\sphinxstylestrong{導入バージョン:}  ex5.30.0

\sphinxAtStartPar
\sphinxstylestrong{更新バージョン:}  v8.2


\bigskip\hrule\bigskip



\subsubsection{切札{\normalsize $\clubsuit$} 10}
\label{\detokenize{auto/frameActionlist:id78}}

\paragraph{悪あがき (速攻魔法)}
\label{\detokenize{auto/frameActionlist:act-lastditch}}\label{\detokenize{auto/frameActionlist:id79}}
\sphinxAtStartPar
\sphinxstylestrong{フレーム:}

\sphinxAtStartPar
・エクストラ

\sphinxAtStartPar
\sphinxstylestrong{トリガー:} 直接

\sphinxAtStartPar
\sphinxstylestrong{スピード:} 即時

\sphinxAtStartPar
\sphinxstylestrong{タイミング:} クイック

\sphinxAtStartPar
\sphinxstylestrong{起動条件:}

\sphinxAtStartPar
自分の表切札に「悪あがき」がある

\sphinxAtStartPar
\sphinxstylestrong{対象:}

\sphinxAtStartPar
キーカードがあるリクエストを対象とする。

\sphinxAtStartPar
\sphinxstylestrong{効果:}
\begin{enumerate}
\sphinxsetlistlabels{\arabic}{enumi}{enumii}{}{.}%
\item {} 
\sphinxAtStartPar
対象リクエストをステージから取り除き、対象リクエストのキーカードを墓地に移す。

\item {} 
\sphinxAtStartPar
手札を全て捨てる。手札がない場合、自分の切札{\normalsize $\clubsuit$} 10を墓地に移す。

\end{enumerate}

\sphinxAtStartPar
\sphinxstylestrong{導入バージョン:}  ex5.30.0

\sphinxAtStartPar
\sphinxstylestrong{更新バージョン:}  v8.2


\bigskip\hrule\bigskip



\subsubsection{切札{\normalsize $\clubsuit$} J}
\label{\detokenize{auto/frameActionlist:id80}}

\paragraph{騎士召喚 (召喚)}
\label{\detokenize{auto/frameActionlist:act-summonknight}}\label{\detokenize{auto/frameActionlist:id81}}
\sphinxAtStartPar
\sphinxstylestrong{フレーム:}

\sphinxAtStartPar
・エクストラ

\sphinxAtStartPar
\sphinxstylestrong{トリガー:} 直接

\sphinxAtStartPar
\sphinxstylestrong{スピード:} 通常

\sphinxAtStartPar
\sphinxstylestrong{タイミング:} メイン

\sphinxAtStartPar
\sphinxstylestrong{キーカード:} {\normalsize $\spadesuit$} 2〜10

\sphinxAtStartPar
\sphinxstylestrong{コスト:} BL

\sphinxAtStartPar
\sphinxstylestrong{起動条件:}

\sphinxAtStartPar
自分の表切札に「騎士」がある

\sphinxAtStartPar
\sphinxstylestrong{効果:}

\sphinxAtStartPar
切札の{\normalsize $\clubsuit$} Jが表の場合、次を行う。
\begin{enumerate}
\sphinxsetlistlabels{\arabic}{enumi}{enumii}{}{.}%
\item {} 
\sphinxAtStartPar
対戦相手は自分の場にある兵士を1枚選び墓地に移す。

\item {} 
\sphinxAtStartPar
キーカードと切札の{\normalsize $\spadesuit$} Jをあわせて騎士としてチャージ状態で場に出す。

\end{enumerate}

\sphinxAtStartPar
\sphinxstylestrong{導入バージョン:}  ex4.22.0

\sphinxAtStartPar
\sphinxstylestrong{更新バージョン:}  v8.2


\bigskip\hrule\bigskip



\paragraph{クイック騎士召喚 (召喚)}
\label{\detokenize{auto/frameActionlist:act-quicksummonknight}}\label{\detokenize{auto/frameActionlist:id82}}
\sphinxAtStartPar
\sphinxstylestrong{フレーム:}

\sphinxAtStartPar
・エクストラ

\sphinxAtStartPar
\sphinxstylestrong{トリガー:} 直接

\sphinxAtStartPar
\sphinxstylestrong{スピード:} 通常

\sphinxAtStartPar
\sphinxstylestrong{タイミング:} クイック

\sphinxAtStartPar
\sphinxstylestrong{キーカード:} {\normalsize $\spadesuit$} 2〜10

\sphinxAtStartPar
\sphinxstylestrong{コスト:} D

\sphinxAtStartPar
\sphinxstylestrong{起動条件:}

\sphinxAtStartPar
自分の表切札に「騎士」がある

\sphinxAtStartPar
\sphinxstylestrong{効果:}

\sphinxAtStartPar
効果の内容は「騎士召喚」を参照。

\sphinxAtStartPar
\sphinxstylestrong{導入バージョン:}  ex4.22.0

\sphinxAtStartPar
\sphinxstylestrong{更新バージョン:}  v8.2


\bigskip\hrule\bigskip



\paragraph{薙払 (兵士起因)}
\label{\detokenize{auto/frameActionlist:act-knightsweep}}\label{\detokenize{auto/frameActionlist:id83}}
\sphinxAtStartPar
\sphinxstylestrong{フレーム:}

\sphinxAtStartPar
・エクストラ

\sphinxAtStartPar
\sphinxstylestrong{トリガー:} 直接

\sphinxAtStartPar
\sphinxstylestrong{スピード:} 即時

\sphinxAtStartPar
\sphinxstylestrong{タイミング:} クイック

\sphinxAtStartPar
\sphinxstylestrong{キーカード:} 「騎士」1体

\sphinxAtStartPar
\sphinxstylestrong{コスト:} K

\sphinxAtStartPar
\sphinxstylestrong{対象:}

\sphinxAtStartPar
対戦相手1人

\sphinxAtStartPar
\sphinxstylestrong{効果:}

\sphinxAtStartPar
対象の対戦相手は自分の場にいるキャラクターを1体墓地に移す。

\sphinxAtStartPar
\sphinxstylestrong{導入バージョン:}  ex4.22.0

\sphinxAtStartPar
\sphinxstylestrong{更新バージョン:}  v8.2


\bigskip\hrule\bigskip



\paragraph{騎士再生 (切札操作)}
\label{\detokenize{auto/frameActionlist:act-knightrevive}}\label{\detokenize{auto/frameActionlist:id84}}
\sphinxAtStartPar
\sphinxstylestrong{フレーム:}

\sphinxAtStartPar
・エクストラ

\sphinxAtStartPar
\sphinxstylestrong{トリガー:} 誘発

\sphinxAtStartPar
\sphinxstylestrong{スピード:} 即時

\sphinxAtStartPar
\sphinxstylestrong{タイミング:} クイック

\sphinxAtStartPar
\sphinxstylestrong{誘発条件:}

\sphinxAtStartPar
自分の 「騎士」 が場から墓地に移る場合

\sphinxAtStartPar
\sphinxstylestrong{効果:}

\sphinxAtStartPar
墓地から{\normalsize $\clubsuit$} Jを1枚選び、自分の表切札として置く。


\subsection{フレームキャラクターリスト}
\label{\detokenize{auto/frameActionlist:characterlist-act-frame}}\label{\detokenize{auto/frameActionlist:id85}}
\begin{sphinxShadowBox}
\sphinxstyletopictitle{目次}
\begin{itemize}
\item {} 
\sphinxAtStartPar
\phantomsection\label{\detokenize{auto/frameActionlist:id256}}{\hyperref[\detokenize{auto/frameActionlist:id87}]{\sphinxcrossref{切札{\normalsize $\spadesuit$} J}}}
\begin{itemize}
\item {} 
\sphinxAtStartPar
\phantomsection\label{\detokenize{auto/frameActionlist:id257}}{\hyperref[\detokenize{auto/frameActionlist:char-reanimator}]{\sphinxcrossref{リアニメーター (兵士)}}}

\end{itemize}

\item {} 
\sphinxAtStartPar
\phantomsection\label{\detokenize{auto/frameActionlist:id258}}{\hyperref[\detokenize{auto/frameActionlist:id89}]{\sphinxcrossref{切札{\normalsize $\spadesuit$} K}}}
\begin{itemize}
\item {} 
\sphinxAtStartPar
\phantomsection\label{\detokenize{auto/frameActionlist:id259}}{\hyperref[\detokenize{auto/frameActionlist:char-darklord}]{\sphinxcrossref{魔王 (兵士)}}}

\end{itemize}

\item {} 
\sphinxAtStartPar
\phantomsection\label{\detokenize{auto/frameActionlist:id260}}{\hyperref[\detokenize{auto/frameActionlist:id91}]{\sphinxcrossref{切札{\normalsize $\heartsuit$} J}}}
\begin{itemize}
\item {} 
\sphinxAtStartPar
\phantomsection\label{\detokenize{auto/frameActionlist:id261}}{\hyperref[\detokenize{auto/frameActionlist:char-chariot}]{\sphinxcrossref{チャリオット (兵士)}}}

\end{itemize}

\item {} 
\sphinxAtStartPar
\phantomsection\label{\detokenize{auto/frameActionlist:id262}}{\hyperref[\detokenize{auto/frameActionlist:id93}]{\sphinxcrossref{切札{\normalsize $\heartsuit$} K}}}
\begin{itemize}
\item {} 
\sphinxAtStartPar
\phantomsection\label{\detokenize{auto/frameActionlist:id263}}{\hyperref[\detokenize{auto/frameActionlist:char-giant}]{\sphinxcrossref{巨人 (兵士)}}}

\end{itemize}

\item {} 
\sphinxAtStartPar
\phantomsection\label{\detokenize{auto/frameActionlist:id264}}{\hyperref[\detokenize{auto/frameActionlist:id95}]{\sphinxcrossref{切札{\normalsize $\diamondsuit$} J}}}
\begin{itemize}
\item {} 
\sphinxAtStartPar
\phantomsection\label{\detokenize{auto/frameActionlist:id265}}{\hyperref[\detokenize{auto/frameActionlist:char-strategist}]{\sphinxcrossref{策士 (兵士)}}}

\end{itemize}

\item {} 
\sphinxAtStartPar
\phantomsection\label{\detokenize{auto/frameActionlist:id266}}{\hyperref[\detokenize{auto/frameActionlist:id97}]{\sphinxcrossref{切札{\normalsize $\clubsuit$} J}}}
\begin{itemize}
\item {} 
\sphinxAtStartPar
\phantomsection\label{\detokenize{auto/frameActionlist:id267}}{\hyperref[\detokenize{auto/frameActionlist:char-knight}]{\sphinxcrossref{騎士 (兵士)}}}

\end{itemize}

\end{itemize}
\end{sphinxShadowBox}


\bigskip\hrule\bigskip



\subsubsection{切札{\normalsize $\spadesuit$} J}
\label{\detokenize{auto/frameActionlist:id87}}

\paragraph{リアニメーター (兵士)}
\label{\detokenize{auto/frameActionlist:char-reanimator}}\label{\detokenize{auto/frameActionlist:id88}}
\sphinxAtStartPar
\sphinxstylestrong{フレーム:}

\sphinxAtStartPar
・エクストラ

\sphinxAtStartPar
\sphinxstylestrong{キャラクターキーカード:} {\normalsize $\spadesuit$} J と {\normalsize $\heartsuit$} 2〜10

\sphinxAtStartPar
\sphinxstylestrong{サイズ:} 11

\sphinxAtStartPar
\sphinxstylestrong{ラベル:} アタッカー, ブロッカー

\sphinxAtStartPar
\sphinxstylestrong{能力:}

\sphinxAtStartPar
・装備アクションの対象にできない。

\sphinxAtStartPar
・{\normalsize $\spadesuit$} 兵士としても{\normalsize $\heartsuit$} 兵士としても扱う。

\sphinxAtStartPar
\sphinxstylestrong{導入バージョン:}  ex4.22.0

\sphinxAtStartPar
\sphinxstylestrong{更新バージョン:}  v8.2


\bigskip\hrule\bigskip



\subsubsection{切札{\normalsize $\spadesuit$} K}
\label{\detokenize{auto/frameActionlist:id89}}

\paragraph{魔王 (兵士)}
\label{\detokenize{auto/frameActionlist:char-darklord}}\label{\detokenize{auto/frameActionlist:id90}}
\sphinxAtStartPar
\sphinxstylestrong{フレーム:}

\sphinxAtStartPar
・エクストラ

\sphinxAtStartPar
\sphinxstylestrong{キャラクターキーカード:} {\normalsize $\spadesuit$} K と {\normalsize $\clubsuit$} K

\sphinxAtStartPar
\sphinxstylestrong{サイズ:} 26

\sphinxAtStartPar
\sphinxstylestrong{ラベル:} アタッカー, ブロッカー

\sphinxAtStartPar
\sphinxstylestrong{能力:}

\sphinxAtStartPar
・アクションの対象にできない。

\sphinxAtStartPar
・{\normalsize $\spadesuit$} 兵士としても{\normalsize $\clubsuit$} 兵士としても扱う。

\sphinxAtStartPar
\sphinxstylestrong{導入バージョン:}  ex3.10.0

\sphinxAtStartPar
\sphinxstylestrong{更新バージョン:}  v8.2


\bigskip\hrule\bigskip



\subsubsection{切札{\normalsize $\heartsuit$} J}
\label{\detokenize{auto/frameActionlist:id91}}

\paragraph{チャリオット (兵士)}
\label{\detokenize{auto/frameActionlist:char-chariot}}\label{\detokenize{auto/frameActionlist:id92}}
\sphinxAtStartPar
\sphinxstylestrong{フレーム:}

\sphinxAtStartPar
・エクストラ

\sphinxAtStartPar
\sphinxstylestrong{キャラクターキーカード:} {\normalsize $\heartsuit$} J と {\normalsize $\diamondsuit$} 2〜10

\sphinxAtStartPar
\sphinxstylestrong{サイズ:} 11

\sphinxAtStartPar
\sphinxstylestrong{ラベル:} アタッカー, ブロッカー

\sphinxAtStartPar
\sphinxstylestrong{能力:}

\sphinxAtStartPar
・装備アクションの対象にできない。

\sphinxAtStartPar
・{\normalsize $\heartsuit$} 兵士としても{\normalsize $\diamondsuit$} 兵士としても扱う。

\sphinxAtStartPar
\sphinxstylestrong{導入バージョン:}  ex4.22.0

\sphinxAtStartPar
\sphinxstylestrong{更新バージョン:}  v8.2


\bigskip\hrule\bigskip



\subsubsection{切札{\normalsize $\heartsuit$} K}
\label{\detokenize{auto/frameActionlist:id93}}

\paragraph{巨人 (兵士)}
\label{\detokenize{auto/frameActionlist:char-giant}}\label{\detokenize{auto/frameActionlist:id94}}
\sphinxAtStartPar
\sphinxstylestrong{フレーム:}

\sphinxAtStartPar
・エクストラ

\sphinxAtStartPar
\sphinxstylestrong{キャラクターキーカード:} {\normalsize $\heartsuit$} K と {\normalsize $\diamondsuit$} K

\sphinxAtStartPar
\sphinxstylestrong{サイズ:} 26

\sphinxAtStartPar
\sphinxstylestrong{ラベル:} アタッカー, ブロッカー

\sphinxAtStartPar
\sphinxstylestrong{能力:}

\sphinxAtStartPar
・アクションの対象にできない。

\sphinxAtStartPar
・{\normalsize $\heartsuit$} 兵士としても{\normalsize $\diamondsuit$} 兵士としても扱う。

\sphinxAtStartPar
・自分のターンに可能な限りアタックアクションをリクエストし、巨人がチャージ状態の場合、アタッカーに指定する。

\sphinxAtStartPar
・「ダメージ判定」にてブロッカーが防壁の場合、条件に該当してもこの巨人は墓地に移されない。

\sphinxAtStartPar
\sphinxstylestrong{導入バージョン:}  ex3.10.0

\sphinxAtStartPar
\sphinxstylestrong{更新バージョン:}  v8.2


\bigskip\hrule\bigskip



\subsubsection{切札{\normalsize $\diamondsuit$} J}
\label{\detokenize{auto/frameActionlist:id95}}

\paragraph{策士 (兵士)}
\label{\detokenize{auto/frameActionlist:char-strategist}}\label{\detokenize{auto/frameActionlist:id96}}
\sphinxAtStartPar
\sphinxstylestrong{フレーム:}

\sphinxAtStartPar
・エクストラ

\sphinxAtStartPar
\sphinxstylestrong{キャラクターキーカード:} {\normalsize $\diamondsuit$} J と {\normalsize $\clubsuit$} 2〜10

\sphinxAtStartPar
\sphinxstylestrong{サイズ:} 11

\sphinxAtStartPar
\sphinxstylestrong{ラベル:} アタッカー, ブロッカー

\sphinxAtStartPar
\sphinxstylestrong{能力:}

\sphinxAtStartPar
・装備アクションの対象にできない。

\sphinxAtStartPar
・{\normalsize $\diamondsuit$} 兵士としても{\normalsize $\clubsuit$} 兵士としても扱う。

\sphinxAtStartPar
\sphinxstylestrong{導入バージョン:}  ex4.22.0

\sphinxAtStartPar
\sphinxstylestrong{更新バージョン:}  v8.2


\bigskip\hrule\bigskip



\subsubsection{切札{\normalsize $\clubsuit$} J}
\label{\detokenize{auto/frameActionlist:id97}}

\paragraph{騎士 (兵士)}
\label{\detokenize{auto/frameActionlist:char-knight}}\label{\detokenize{auto/frameActionlist:id98}}
\sphinxAtStartPar
\sphinxstylestrong{フレーム:}

\sphinxAtStartPar
・エクストラ

\sphinxAtStartPar
\sphinxstylestrong{キャラクターキーカード:} {\normalsize $\clubsuit$} J と {\normalsize $\spadesuit$} 2〜10

\sphinxAtStartPar
\sphinxstylestrong{サイズ:} 11

\sphinxAtStartPar
\sphinxstylestrong{ラベル:} アタッカー, ブロッカー

\sphinxAtStartPar
\sphinxstylestrong{能力:}

\sphinxAtStartPar
・装備アクションの対象にできない。

\sphinxAtStartPar
・{\normalsize $\clubsuit$} 兵士としても{\normalsize $\spadesuit$} 兵士としても扱う。

\sphinxAtStartPar
\sphinxstylestrong{導入バージョン:}  ex4.22.0

\sphinxAtStartPar
\sphinxstylestrong{更新バージョン:}  v8.2


\subsection{フレームフォグリスト}
\label{\detokenize{auto/frameActionlist:foglist-act-frame}}\label{\detokenize{auto/frameActionlist:id99}}
\begin{sphinxShadowBox}
\sphinxstyletopictitle{目次}
\begin{itemize}
\item {} 
\sphinxAtStartPar
\phantomsection\label{\detokenize{auto/frameActionlist:id268}}{\hyperref[\detokenize{auto/frameActionlist:id101}]{\sphinxcrossref{滞留魔法}}}
\begin{itemize}
\item {} 
\sphinxAtStartPar
\phantomsection\label{\detokenize{auto/frameActionlist:id269}}{\hyperref[\detokenize{auto/frameActionlist:fog-closefog}]{\sphinxcrossref{クローズ (滞留魔法)}}}

\end{itemize}

\end{itemize}
\end{sphinxShadowBox}


\bigskip\hrule\bigskip



\subsubsection{滞留魔法}
\label{\detokenize{auto/frameActionlist:id101}}

\paragraph{クローズ (滞留魔法)}
\label{\detokenize{auto/frameActionlist:fog-closefog}}\label{\detokenize{auto/frameActionlist:id102}}
\sphinxAtStartPar
\sphinxstylestrong{フレーム:}

\sphinxAtStartPar
・エクストラ

\sphinxAtStartPar
\sphinxstylestrong{フォグキーカード:} Q

\sphinxAtStartPar
\sphinxstylestrong{能力:}

\sphinxAtStartPar
対象とした切札はこのターンが終わるまで表向きにできない。

\sphinxAtStartPar
\sphinxstylestrong{導入バージョン:}  v8.1


\subsection{フレーム切札リスト}
\label{\detokenize{auto/frameActionlist:trumplist-act-frame}}\label{\detokenize{auto/frameActionlist:id103}}
\sphinxAtStartPar
切札は、ゲームにさらなる戦略性と構築の自由度をもたらす拡張要素です。

\sphinxAtStartPar
それぞれの切札にはスートと数字に応じた固有の能力が設定されており、多彩なコンボや戦術を可能にします。

\sphinxAtStartPar
また、トランプとの関連性が深いタロットの象徴イメージを採り入れ、ビジュアル面からも覚えやすさを追求しています。

\begin{sphinxShadowBox}
\sphinxstyletopictitle{目次}
\begin{itemize}
\item {} 
\sphinxAtStartPar
\phantomsection\label{\detokenize{auto/frameActionlist:id270}}{\hyperref[\detokenize{auto/frameActionlist:id105}]{\sphinxcrossref{{\normalsize $\spadesuit$} スペード}}}
\begin{itemize}
\item {} 
\sphinxAtStartPar
\phantomsection\label{\detokenize{auto/frameActionlist:id271}}{\hyperref[\detokenize{auto/frameActionlist:trump-sa}]{\sphinxcrossref{革命 (表切札)}}}

\item {} 
\sphinxAtStartPar
\phantomsection\label{\detokenize{auto/frameActionlist:id272}}{\hyperref[\detokenize{auto/frameActionlist:trump-s2}]{\sphinxcrossref{スラッシュ (表切札)}}}

\item {} 
\sphinxAtStartPar
\phantomsection\label{\detokenize{auto/frameActionlist:id273}}{\hyperref[\detokenize{auto/frameActionlist:trump-s3}]{\sphinxcrossref{ドレイン (表切札)}}}

\item {} 
\sphinxAtStartPar
\phantomsection\label{\detokenize{auto/frameActionlist:id274}}{\hyperref[\detokenize{auto/frameActionlist:trump-s4}]{\sphinxcrossref{威圧 (表切札)}}}

\item {} 
\sphinxAtStartPar
\phantomsection\label{\detokenize{auto/frameActionlist:id275}}{\hyperref[\detokenize{auto/frameActionlist:trump-s5}]{\sphinxcrossref{蘇生 (表切札)}}}

\item {} 
\sphinxAtStartPar
\phantomsection\label{\detokenize{auto/frameActionlist:id276}}{\hyperref[\detokenize{auto/frameActionlist:trump-s7}]{\sphinxcrossref{瞬殺 (表切札)}}}

\item {} 
\sphinxAtStartPar
\phantomsection\label{\detokenize{auto/frameActionlist:id277}}{\hyperref[\detokenize{auto/frameActionlist:trump-sj}]{\sphinxcrossref{リアニメーター (表切札)}}}

\item {} 
\sphinxAtStartPar
\phantomsection\label{\detokenize{auto/frameActionlist:id278}}{\hyperref[\detokenize{auto/frameActionlist:trump-sk}]{\sphinxcrossref{魔王 (表切札)}}}

\end{itemize}

\item {} 
\sphinxAtStartPar
\phantomsection\label{\detokenize{auto/frameActionlist:id279}}{\hyperref[\detokenize{auto/frameActionlist:id114}]{\sphinxcrossref{{\normalsize $\heartsuit$} ハート}}}
\begin{itemize}
\item {} 
\sphinxAtStartPar
\phantomsection\label{\detokenize{auto/frameActionlist:id280}}{\hyperref[\detokenize{auto/frameActionlist:trump-h2}]{\sphinxcrossref{援軍 (表切札)}}}

\item {} 
\sphinxAtStartPar
\phantomsection\label{\detokenize{auto/frameActionlist:id281}}{\hyperref[\detokenize{auto/frameActionlist:trump-h3}]{\sphinxcrossref{激励 (表切札)}}}

\item {} 
\sphinxAtStartPar
\phantomsection\label{\detokenize{auto/frameActionlist:id282}}{\hyperref[\detokenize{auto/frameActionlist:trump-h4}]{\sphinxcrossref{休戦 (表切札)}}}

\item {} 
\sphinxAtStartPar
\phantomsection\label{\detokenize{auto/frameActionlist:id283}}{\hyperref[\detokenize{auto/frameActionlist:trump-h5}]{\sphinxcrossref{追撃 (表切札)}}}

\item {} 
\sphinxAtStartPar
\phantomsection\label{\detokenize{auto/frameActionlist:id284}}{\hyperref[\detokenize{auto/frameActionlist:trump-h7}]{\sphinxcrossref{ファントム (表切札)}}}

\item {} 
\sphinxAtStartPar
\phantomsection\label{\detokenize{auto/frameActionlist:id285}}{\hyperref[\detokenize{auto/frameActionlist:trump-h8}]{\sphinxcrossref{救済 (表切札)}}}

\item {} 
\sphinxAtStartPar
\phantomsection\label{\detokenize{auto/frameActionlist:id286}}{\hyperref[\detokenize{auto/frameActionlist:trump-hj}]{\sphinxcrossref{チャリオット (表切札)}}}

\item {} 
\sphinxAtStartPar
\phantomsection\label{\detokenize{auto/frameActionlist:id287}}{\hyperref[\detokenize{auto/frameActionlist:trump-hk}]{\sphinxcrossref{巨人 (表切札)}}}

\end{itemize}

\item {} 
\sphinxAtStartPar
\phantomsection\label{\detokenize{auto/frameActionlist:id288}}{\hyperref[\detokenize{auto/frameActionlist:id123}]{\sphinxcrossref{{\normalsize $\diamondsuit$} ダイヤ}}}
\begin{itemize}
\item {} 
\sphinxAtStartPar
\phantomsection\label{\detokenize{auto/frameActionlist:id289}}{\hyperref[\detokenize{auto/frameActionlist:trump-da}]{\sphinxcrossref{防壁追加 (表切札)}}}

\item {} 
\sphinxAtStartPar
\phantomsection\label{\detokenize{auto/frameActionlist:id290}}{\hyperref[\detokenize{auto/frameActionlist:trump-d3}]{\sphinxcrossref{交渉 (表切札)}}}

\item {} 
\sphinxAtStartPar
\phantomsection\label{\detokenize{auto/frameActionlist:id291}}{\hyperref[\detokenize{auto/frameActionlist:trump-d4}]{\sphinxcrossref{相殺 (表切札)}}}

\item {} 
\sphinxAtStartPar
\phantomsection\label{\detokenize{auto/frameActionlist:id292}}{\hyperref[\detokenize{auto/frameActionlist:trump-d5}]{\sphinxcrossref{喪失 (表切札)}}}

\item {} 
\sphinxAtStartPar
\phantomsection\label{\detokenize{auto/frameActionlist:id293}}{\hyperref[\detokenize{auto/frameActionlist:trump-d8}]{\sphinxcrossref{修繕 (表切札)}}}

\item {} 
\sphinxAtStartPar
\phantomsection\label{\detokenize{auto/frameActionlist:id294}}{\hyperref[\detokenize{auto/frameActionlist:trump-d9}]{\sphinxcrossref{充足 (表切札)}}}

\item {} 
\sphinxAtStartPar
\phantomsection\label{\detokenize{auto/frameActionlist:id295}}{\hyperref[\detokenize{auto/frameActionlist:trump-dj}]{\sphinxcrossref{策士 (表切札)}}}

\end{itemize}

\item {} 
\sphinxAtStartPar
\phantomsection\label{\detokenize{auto/frameActionlist:id296}}{\hyperref[\detokenize{auto/frameActionlist:id131}]{\sphinxcrossref{{\normalsize $\clubsuit$} クラブ}}}
\begin{itemize}
\item {} 
\sphinxAtStartPar
\phantomsection\label{\detokenize{auto/frameActionlist:id297}}{\hyperref[\detokenize{auto/frameActionlist:trump-c4}]{\sphinxcrossref{スイッチ (表切札)}}}

\item {} 
\sphinxAtStartPar
\phantomsection\label{\detokenize{auto/frameActionlist:id298}}{\hyperref[\detokenize{auto/frameActionlist:trump-c5}]{\sphinxcrossref{内乱 (表切札)}}}

\item {} 
\sphinxAtStartPar
\phantomsection\label{\detokenize{auto/frameActionlist:id299}}{\hyperref[\detokenize{auto/frameActionlist:trump-c7}]{\sphinxcrossref{手札補充 (表切札)}}}

\item {} 
\sphinxAtStartPar
\phantomsection\label{\detokenize{auto/frameActionlist:id300}}{\hyperref[\detokenize{auto/frameActionlist:trump-c8}]{\sphinxcrossref{手札補充 (表切札)}}}

\item {} 
\sphinxAtStartPar
\phantomsection\label{\detokenize{auto/frameActionlist:id301}}{\hyperref[\detokenize{auto/frameActionlist:trump-c9}]{\sphinxcrossref{要塞 (表切札)}}}

\item {} 
\sphinxAtStartPar
\phantomsection\label{\detokenize{auto/frameActionlist:id302}}{\hyperref[\detokenize{auto/frameActionlist:trump-c10}]{\sphinxcrossref{悪あがき (表切札)}}}

\item {} 
\sphinxAtStartPar
\phantomsection\label{\detokenize{auto/frameActionlist:id303}}{\hyperref[\detokenize{auto/frameActionlist:trump-cj}]{\sphinxcrossref{騎士 (表切札)}}}

\end{itemize}

\item {} 
\sphinxAtStartPar
\phantomsection\label{\detokenize{auto/frameActionlist:id304}}{\hyperref[\detokenize{auto/frameActionlist:id139}]{\sphinxcrossref{裏切札}}}
\begin{itemize}
\item {} 
\sphinxAtStartPar
\phantomsection\label{\detokenize{auto/frameActionlist:id305}}{\hyperref[\detokenize{auto/frameActionlist:trump-secrettrump}]{\sphinxcrossref{裏切札 (裏切札)}}}

\end{itemize}

\end{itemize}
\end{sphinxShadowBox}


\bigskip\hrule\bigskip



\subsubsection{{\normalsize $\spadesuit$} スペード}
\label{\detokenize{auto/frameActionlist:id105}}

\paragraph{革命 (表切札)}
\label{\detokenize{auto/frameActionlist:trump-sa}}\label{\detokenize{auto/frameActionlist:id106}}
\sphinxAtStartPar
\sphinxstylestrong{フレーム:}

\sphinxAtStartPar
・エクストラ

\sphinxAtStartPar
\sphinxstylestrong{切札キーカード:} {\normalsize $\spadesuit$} A

\sphinxAtStartPar
\sphinxstylestrong{能力:}

\sphinxAtStartPar
・自分の兵士は \textless{}革命ドロー\textgreater{} を持つ。

\sphinxAtStartPar
・この能力が有効で場合、対戦相手も含めた全てのダメージ判定アクションの「1.兵士（アタッカー）と兵士（ブロッカー）の場合」を以下のように変更する。

\sphinxAtStartPar
1.兵士（アタッカー）と兵士（ブロッカー）の場合、アタッカーとブロッカーで数字を比較し、大きい方を墓地に移動する。同じ場合は両方を墓地に移動する。アタッカーとブロッカーを比較した数字の差をダメージとして兵士を墓地に移した方のプレイヤーに与える。1アタッカーに対して複数ブロッカーがいる場合、ブロッカーの合計数字と比較する。

\sphinxAtStartPar
\sphinxstylestrong{イメージ:}
革命のイメージ (\hyperref[\detokenize{auto/frameActionlist:frame-sa-image}]{Fig.\@ \ref{\detokenize{auto/frameActionlist:frame-sa-image}}})

\begin{figure}[htbp]
\centering
\capstart

\noindent\sphinxincludegraphics[height=200\sphinxpxdimen]{{Pictorial_Key_to_the_Tarot_Swords_01}.jpg}
\caption{革命のイメージ}\label{\detokenize{auto/frameActionlist:id141}}\label{\detokenize{auto/frameActionlist:frame-sa-image}}\end{figure}

\sphinxAtStartPar
\sphinxstylestrong{導入バージョン:}  ex4.14.0

\sphinxAtStartPar
\sphinxstylestrong{更新バージョン:}  v8.2


\bigskip\hrule\bigskip



\paragraph{スラッシュ (表切札)}
\label{\detokenize{auto/frameActionlist:trump-s2}}\label{\detokenize{auto/frameActionlist:id107}}
\sphinxAtStartPar
\sphinxstylestrong{フレーム:}

\sphinxAtStartPar
・エクストラ

\sphinxAtStartPar
\sphinxstylestrong{切札キーカード:} {\normalsize $\spadesuit$} 2

\sphinxAtStartPar
\sphinxstylestrong{能力:}

\sphinxAtStartPar
自分の場にいる{\normalsize $\spadesuit$} の兵士は \textless{}スラッシュ\textgreater{} を持つ。

\sphinxAtStartPar
\sphinxstylestrong{イメージ:}
スラッシュのイメージ (\hyperref[\detokenize{auto/frameActionlist:frame-s2-image}]{Fig.\@ \ref{\detokenize{auto/frameActionlist:frame-s2-image}}})

\begin{figure}[htbp]
\centering
\capstart

\noindent\sphinxincludegraphics[height=200\sphinxpxdimen]{{Pictorial_Key_to_the_Tarot_Swords_02}.jpg}
\caption{スラッシュのイメージ}\label{\detokenize{auto/frameActionlist:id142}}\label{\detokenize{auto/frameActionlist:frame-s2-image}}\end{figure}

\sphinxAtStartPar
\sphinxstylestrong{導入バージョン:}  ex3.4.0

\sphinxAtStartPar
\sphinxstylestrong{更新バージョン:}  v8.1


\bigskip\hrule\bigskip



\paragraph{ドレイン (表切札)}
\label{\detokenize{auto/frameActionlist:trump-s3}}\label{\detokenize{auto/frameActionlist:id108}}
\sphinxAtStartPar
\sphinxstylestrong{フレーム:}

\sphinxAtStartPar
・エクストラ

\sphinxAtStartPar
\sphinxstylestrong{切札キーカード:} {\normalsize $\spadesuit$} 3

\sphinxAtStartPar
\sphinxstylestrong{能力:}

\sphinxAtStartPar
なし

\sphinxAtStartPar
\sphinxstylestrong{イメージ:}
ドレインのイメージ (\hyperref[\detokenize{auto/frameActionlist:frame-s3-image}]{Fig.\@ \ref{\detokenize{auto/frameActionlist:frame-s3-image}}})

\begin{figure}[htbp]
\centering
\capstart

\noindent\sphinxincludegraphics[height=200\sphinxpxdimen]{{Pictorial_Key_to_the_Tarot_Swords_03}.jpg}
\caption{ドレインのイメージ}\label{\detokenize{auto/frameActionlist:id143}}\label{\detokenize{auto/frameActionlist:frame-s3-image}}\end{figure}

\sphinxAtStartPar
\sphinxstylestrong{導入バージョン:}  ex5.30.0

\sphinxAtStartPar
\sphinxstylestrong{更新バージョン:}  v8.2


\bigskip\hrule\bigskip



\paragraph{威圧 (表切札)}
\label{\detokenize{auto/frameActionlist:trump-s4}}\label{\detokenize{auto/frameActionlist:id109}}
\sphinxAtStartPar
\sphinxstylestrong{フレーム:}

\sphinxAtStartPar
・エクストラ

\sphinxAtStartPar
\sphinxstylestrong{切札キーカード:} {\normalsize $\spadesuit$} 4

\sphinxAtStartPar
\sphinxstylestrong{能力:}

\sphinxAtStartPar
なし

\sphinxAtStartPar
\sphinxstylestrong{イメージ:}
威圧のイメージ (\hyperref[\detokenize{auto/frameActionlist:frame-s4-image}]{Fig.\@ \ref{\detokenize{auto/frameActionlist:frame-s4-image}}})

\begin{figure}[htbp]
\centering
\capstart

\noindent\sphinxincludegraphics[height=200\sphinxpxdimen]{{Pictorial_Key_to_the_Tarot_Swords_04}.jpg}
\caption{威圧のイメージ}\label{\detokenize{auto/frameActionlist:id144}}\label{\detokenize{auto/frameActionlist:frame-s4-image}}\end{figure}

\sphinxAtStartPar
\sphinxstylestrong{導入バージョン:}  ex4.18.0

\sphinxAtStartPar
\sphinxstylestrong{更新バージョン:}  v8.2


\bigskip\hrule\bigskip



\paragraph{蘇生 (表切札)}
\label{\detokenize{auto/frameActionlist:trump-s5}}\label{\detokenize{auto/frameActionlist:id110}}
\sphinxAtStartPar
\sphinxstylestrong{フレーム:}

\sphinxAtStartPar
・エクストラ

\sphinxAtStartPar
\sphinxstylestrong{切札キーカード:} {\normalsize $\spadesuit$} 5

\sphinxAtStartPar
\sphinxstylestrong{能力:}

\sphinxAtStartPar
なし

\sphinxAtStartPar
\sphinxstylestrong{イメージ:}
蘇生のイメージ (\hyperref[\detokenize{auto/frameActionlist:frame-s5-image}]{Fig.\@ \ref{\detokenize{auto/frameActionlist:frame-s5-image}}})

\begin{figure}[htbp]
\centering
\capstart

\noindent\sphinxincludegraphics[height=200\sphinxpxdimen]{{Pictorial_Key_to_the_Tarot_Swords_05}.jpg}
\caption{蘇生のイメージ}\label{\detokenize{auto/frameActionlist:id145}}\label{\detokenize{auto/frameActionlist:frame-s5-image}}\end{figure}

\sphinxAtStartPar
\sphinxstylestrong{導入バージョン:}  ex3.10.0

\sphinxAtStartPar
\sphinxstylestrong{更新バージョン:}  v8.2


\bigskip\hrule\bigskip



\paragraph{瞬殺 (表切札)}
\label{\detokenize{auto/frameActionlist:trump-s7}}\label{\detokenize{auto/frameActionlist:id111}}
\sphinxAtStartPar
\sphinxstylestrong{フレーム:}

\sphinxAtStartPar
・エクストラ

\sphinxAtStartPar
\sphinxstylestrong{切札キーカード:} {\normalsize $\spadesuit$} 7

\sphinxAtStartPar
\sphinxstylestrong{能力:}

\sphinxAtStartPar
なし

\sphinxAtStartPar
\sphinxstylestrong{イメージ:}
瞬殺のイメージ (\hyperref[\detokenize{auto/frameActionlist:frame-s7-image}]{Fig.\@ \ref{\detokenize{auto/frameActionlist:frame-s7-image}}})

\begin{figure}[htbp]
\centering
\capstart

\noindent\sphinxincludegraphics[height=200\sphinxpxdimen]{{Pictorial_Key_to_the_Tarot_Swords_07}.jpg}
\caption{瞬殺のイメージ}\label{\detokenize{auto/frameActionlist:id146}}\label{\detokenize{auto/frameActionlist:frame-s7-image}}\end{figure}

\sphinxAtStartPar
\sphinxstylestrong{導入バージョン:}  ex5.30.0

\sphinxAtStartPar
\sphinxstylestrong{更新バージョン:}  v8.2


\bigskip\hrule\bigskip



\paragraph{リアニメーター (表切札)}
\label{\detokenize{auto/frameActionlist:trump-sj}}\label{\detokenize{auto/frameActionlist:id112}}
\sphinxAtStartPar
\sphinxstylestrong{フレーム:}

\sphinxAtStartPar
・エクストラ

\sphinxAtStartPar
\sphinxstylestrong{切札キーカード:} {\normalsize $\spadesuit$} J

\sphinxAtStartPar
\sphinxstylestrong{能力:}

\sphinxAtStartPar
なし

\sphinxAtStartPar
\sphinxstylestrong{イメージ:}
リアニメーターのイメージ (\hyperref[\detokenize{auto/frameActionlist:frame-sj-image}]{Fig.\@ \ref{\detokenize{auto/frameActionlist:frame-sj-image}}})

\begin{figure}[htbp]
\centering
\capstart

\noindent\sphinxincludegraphics[height=200\sphinxpxdimen]{{Pictorial_Key_to_the_Tarot_20_Judgement}.jpg}
\caption{リアニメーターのイメージ}\label{\detokenize{auto/frameActionlist:id147}}\label{\detokenize{auto/frameActionlist:frame-sj-image}}\end{figure}

\sphinxAtStartPar
\sphinxstylestrong{導入バージョン:}  ex4.22.0

\sphinxAtStartPar
\sphinxstylestrong{更新バージョン:}  v8.2


\bigskip\hrule\bigskip



\paragraph{魔王 (表切札)}
\label{\detokenize{auto/frameActionlist:trump-sk}}\label{\detokenize{auto/frameActionlist:id113}}
\sphinxAtStartPar
\sphinxstylestrong{フレーム:}

\sphinxAtStartPar
・エクストラ

\sphinxAtStartPar
\sphinxstylestrong{切札キーカード:} {\normalsize $\spadesuit$} K

\sphinxAtStartPar
\sphinxstylestrong{能力:}

\sphinxAtStartPar
なし

\sphinxAtStartPar
\sphinxstylestrong{イメージ:}
魔王のイメージ (\hyperref[\detokenize{auto/frameActionlist:frame-sk-image}]{Fig.\@ \ref{\detokenize{auto/frameActionlist:frame-sk-image}}})

\begin{figure}[htbp]
\centering
\capstart

\noindent\sphinxincludegraphics[height=200\sphinxpxdimen]{{Pictorial_Key_to_the_Tarot_04_The_Emperor}.jpg}
\caption{魔王のイメージ}\label{\detokenize{auto/frameActionlist:id148}}\label{\detokenize{auto/frameActionlist:frame-sk-image}}\end{figure}

\sphinxAtStartPar
\sphinxstylestrong{導入バージョン:}  ex3.10.0

\sphinxAtStartPar
\sphinxstylestrong{更新バージョン:}  v8.2


\bigskip\hrule\bigskip



\subsubsection{{\normalsize $\heartsuit$} ハート}
\label{\detokenize{auto/frameActionlist:id114}}

\paragraph{援軍 (表切札)}
\label{\detokenize{auto/frameActionlist:trump-h2}}\label{\detokenize{auto/frameActionlist:id115}}
\sphinxAtStartPar
\sphinxstylestrong{フレーム:}

\sphinxAtStartPar
・エクストラ

\sphinxAtStartPar
\sphinxstylestrong{切札キーカード:} {\normalsize $\heartsuit$} 2

\sphinxAtStartPar
\sphinxstylestrong{能力:}

\sphinxAtStartPar
なし

\sphinxAtStartPar
\sphinxstylestrong{イメージ:}
援軍のイメージ (\hyperref[\detokenize{auto/frameActionlist:frame-h2-image}]{Fig.\@ \ref{\detokenize{auto/frameActionlist:frame-h2-image}}})

\begin{figure}[htbp]
\centering
\capstart

\noindent\sphinxincludegraphics[height=200\sphinxpxdimen]{{Pictorial_Key_to_the_Tarot_Cups_02}.jpg}
\caption{援軍のイメージ}\label{\detokenize{auto/frameActionlist:id149}}\label{\detokenize{auto/frameActionlist:frame-h2-image}}\end{figure}

\sphinxAtStartPar
\sphinxstylestrong{導入バージョン:}  ex3.10.0

\sphinxAtStartPar
\sphinxstylestrong{更新バージョン:}  v8.2


\bigskip\hrule\bigskip



\paragraph{激励 (表切札)}
\label{\detokenize{auto/frameActionlist:trump-h3}}\label{\detokenize{auto/frameActionlist:id116}}
\sphinxAtStartPar
\sphinxstylestrong{フレーム:}

\sphinxAtStartPar
・エクストラ

\sphinxAtStartPar
\sphinxstylestrong{切札キーカード:} {\normalsize $\heartsuit$} 3

\sphinxAtStartPar
\sphinxstylestrong{能力:}

\sphinxAtStartPar
自分の場にいる{\normalsize $\heartsuit$} の兵士は \textless{}激励\textgreater{} を持つ。

\sphinxAtStartPar
\sphinxstylestrong{イメージ:}
激励のイメージ (\hyperref[\detokenize{auto/frameActionlist:frame-h3-image}]{Fig.\@ \ref{\detokenize{auto/frameActionlist:frame-h3-image}}})

\begin{figure}[htbp]
\centering
\capstart

\noindent\sphinxincludegraphics[height=200\sphinxpxdimen]{{Pictorial_Key_to_the_Tarot_Cups_03}.jpg}
\caption{激励のイメージ}\label{\detokenize{auto/frameActionlist:id150}}\label{\detokenize{auto/frameActionlist:frame-h3-image}}\end{figure}

\sphinxAtStartPar
\sphinxstylestrong{導入バージョン:}  ex3.4.0

\sphinxAtStartPar
\sphinxstylestrong{更新バージョン:}  v8.1


\bigskip\hrule\bigskip



\paragraph{休戦 (表切札)}
\label{\detokenize{auto/frameActionlist:trump-h4}}\label{\detokenize{auto/frameActionlist:id117}}
\sphinxAtStartPar
\sphinxstylestrong{フレーム:}

\sphinxAtStartPar
・エクストラ

\sphinxAtStartPar
\sphinxstylestrong{切札キーカード:} {\normalsize $\heartsuit$} 4

\sphinxAtStartPar
\sphinxstylestrong{能力:}

\sphinxAtStartPar
なし

\sphinxAtStartPar
\sphinxstylestrong{イメージ:}
休戦のイメージ (\hyperref[\detokenize{auto/frameActionlist:frame-h4-image}]{Fig.\@ \ref{\detokenize{auto/frameActionlist:frame-h4-image}}})

\begin{figure}[htbp]
\centering
\capstart

\noindent\sphinxincludegraphics[height=200\sphinxpxdimen]{{Pictorial_Key_to_the_Tarot_Cups_04}.jpg}
\caption{休戦のイメージ}\label{\detokenize{auto/frameActionlist:id151}}\label{\detokenize{auto/frameActionlist:frame-h4-image}}\end{figure}

\sphinxAtStartPar
\sphinxstylestrong{導入バージョン:}  ex5.30.0

\sphinxAtStartPar
\sphinxstylestrong{更新バージョン:}  v8.2


\bigskip\hrule\bigskip



\paragraph{追撃 (表切札)}
\label{\detokenize{auto/frameActionlist:trump-h5}}\label{\detokenize{auto/frameActionlist:id118}}
\sphinxAtStartPar
\sphinxstylestrong{フレーム:}

\sphinxAtStartPar
・エクストラ

\sphinxAtStartPar
\sphinxstylestrong{切札キーカード:} {\normalsize $\heartsuit$} 5

\sphinxAtStartPar
\sphinxstylestrong{能力:}

\sphinxAtStartPar
なし

\sphinxAtStartPar
\sphinxstylestrong{イメージ:}
追撃のイメージ (\hyperref[\detokenize{auto/frameActionlist:frame-h5-image}]{Fig.\@ \ref{\detokenize{auto/frameActionlist:frame-h5-image}}})

\begin{figure}[htbp]
\centering
\capstart

\noindent\sphinxincludegraphics[height=200\sphinxpxdimen]{{Pictorial_Key_to_the_Tarot_Cups_05}.jpg}
\caption{追撃のイメージ}\label{\detokenize{auto/frameActionlist:id152}}\label{\detokenize{auto/frameActionlist:frame-h5-image}}\end{figure}

\sphinxAtStartPar
\sphinxstylestrong{導入バージョン:}  ex5.30.0

\sphinxAtStartPar
\sphinxstylestrong{更新バージョン:}  v8.2


\bigskip\hrule\bigskip



\paragraph{ファントム (表切札)}
\label{\detokenize{auto/frameActionlist:trump-h7}}\label{\detokenize{auto/frameActionlist:id119}}
\sphinxAtStartPar
\sphinxstylestrong{フレーム:}

\sphinxAtStartPar
・エクストラ

\sphinxAtStartPar
\sphinxstylestrong{切札キーカード:} {\normalsize $\heartsuit$} 7

\sphinxAtStartPar
\sphinxstylestrong{能力:}

\sphinxAtStartPar
自分の場にいる防壁は \textless{}ファントム\textgreater{} を持つ。

\sphinxAtStartPar
\sphinxstylestrong{イメージ:}
ファントムのイメージ (\hyperref[\detokenize{auto/frameActionlist:frame-h7-image}]{Fig.\@ \ref{\detokenize{auto/frameActionlist:frame-h7-image}}})

\begin{figure}[htbp]
\centering
\capstart

\noindent\sphinxincludegraphics[height=200\sphinxpxdimen]{{Pictorial_Key_to_the_Tarot_Cups_07}.jpg}
\caption{ファントムのイメージ}\label{\detokenize{auto/frameActionlist:id153}}\label{\detokenize{auto/frameActionlist:frame-h7-image}}\end{figure}

\sphinxAtStartPar
\sphinxstylestrong{導入バージョン:}  ex4.18.0

\sphinxAtStartPar
\sphinxstylestrong{更新バージョン:}  v8.2


\bigskip\hrule\bigskip



\paragraph{救済 (表切札)}
\label{\detokenize{auto/frameActionlist:trump-h8}}\label{\detokenize{auto/frameActionlist:id120}}
\sphinxAtStartPar
\sphinxstylestrong{フレーム:}

\sphinxAtStartPar
・エクストラ

\sphinxAtStartPar
\sphinxstylestrong{切札キーカード:} {\normalsize $\heartsuit$} 8

\sphinxAtStartPar
\sphinxstylestrong{能力:}

\sphinxAtStartPar
なし

\sphinxAtStartPar
\sphinxstylestrong{イメージ:}
救済のイメージ (\hyperref[\detokenize{auto/frameActionlist:frame-h8-image}]{Fig.\@ \ref{\detokenize{auto/frameActionlist:frame-h8-image}}})

\begin{figure}[htbp]
\centering
\capstart

\noindent\sphinxincludegraphics[height=200\sphinxpxdimen]{{Pictorial_Key_to_the_Tarot_Cups_08}.jpg}
\caption{救済のイメージ}\label{\detokenize{auto/frameActionlist:id154}}\label{\detokenize{auto/frameActionlist:frame-h8-image}}\end{figure}

\sphinxAtStartPar
\sphinxstylestrong{導入バージョン:}  ex4.14.0

\sphinxAtStartPar
\sphinxstylestrong{更新バージョン:}  v8.2


\bigskip\hrule\bigskip



\paragraph{チャリオット (表切札)}
\label{\detokenize{auto/frameActionlist:trump-hj}}\label{\detokenize{auto/frameActionlist:id121}}
\sphinxAtStartPar
\sphinxstylestrong{フレーム:}

\sphinxAtStartPar
・エクストラ

\sphinxAtStartPar
\sphinxstylestrong{切札キーカード:} {\normalsize $\heartsuit$} J

\sphinxAtStartPar
\sphinxstylestrong{能力:}

\sphinxAtStartPar
なし

\sphinxAtStartPar
\sphinxstylestrong{イメージ:}
チャリオットのイメージ (\hyperref[\detokenize{auto/frameActionlist:frame-hj-image}]{Fig.\@ \ref{\detokenize{auto/frameActionlist:frame-hj-image}}})

\begin{figure}[htbp]
\centering
\capstart

\noindent\sphinxincludegraphics[height=200\sphinxpxdimen]{{Pictorial_Key_to_the_Tarot_07_The_Chariot}.jpg}
\caption{チャリオットのイメージ}\label{\detokenize{auto/frameActionlist:id155}}\label{\detokenize{auto/frameActionlist:frame-hj-image}}\end{figure}

\sphinxAtStartPar
\sphinxstylestrong{導入バージョン:}  ex4.22.0

\sphinxAtStartPar
\sphinxstylestrong{更新バージョン:}  v8.2


\bigskip\hrule\bigskip



\paragraph{巨人 (表切札)}
\label{\detokenize{auto/frameActionlist:trump-hk}}\label{\detokenize{auto/frameActionlist:id122}}
\sphinxAtStartPar
\sphinxstylestrong{フレーム:}

\sphinxAtStartPar
・エクストラ

\sphinxAtStartPar
\sphinxstylestrong{切札キーカード:} {\normalsize $\heartsuit$} K

\sphinxAtStartPar
\sphinxstylestrong{能力:}

\sphinxAtStartPar
なし

\sphinxAtStartPar
\sphinxstylestrong{イメージ:}
巨人のイメージ (\hyperref[\detokenize{auto/frameActionlist:frame-hk-image}]{Fig.\@ \ref{\detokenize{auto/frameActionlist:frame-hk-image}}})

\begin{figure}[htbp]
\centering
\capstart

\noindent\sphinxincludegraphics[height=200\sphinxpxdimen]{{Pictorial_Key_to_the_Tarot_08_Strength}.jpg}
\caption{巨人のイメージ}\label{\detokenize{auto/frameActionlist:id156}}\label{\detokenize{auto/frameActionlist:frame-hk-image}}\end{figure}

\sphinxAtStartPar
\sphinxstylestrong{導入バージョン:}  ex3.10.0

\sphinxAtStartPar
\sphinxstylestrong{更新バージョン:}  v8.2


\bigskip\hrule\bigskip



\subsubsection{{\normalsize $\diamondsuit$} ダイヤ}
\label{\detokenize{auto/frameActionlist:id123}}

\paragraph{防壁追加 (表切札)}
\label{\detokenize{auto/frameActionlist:trump-da}}\label{\detokenize{auto/frameActionlist:id124}}
\sphinxAtStartPar
\sphinxstylestrong{フレーム:}

\sphinxAtStartPar
・エクストラ

\sphinxAtStartPar
\sphinxstylestrong{切札キーカード:} {\normalsize $\diamondsuit$} A

\sphinxAtStartPar
\sphinxstylestrong{能力:}

\sphinxAtStartPar
なし

\sphinxAtStartPar
\sphinxstylestrong{イメージ:}
防壁追加のイメージ (\hyperref[\detokenize{auto/frameActionlist:frame-da-image}]{Fig.\@ \ref{\detokenize{auto/frameActionlist:frame-da-image}}})

\begin{figure}[htbp]
\centering
\capstart

\noindent\sphinxincludegraphics[height=200\sphinxpxdimen]{{Pictorial_Key_to_the_Tarot_Pentacles_01}.jpg}
\caption{防壁追加のイメージ}\label{\detokenize{auto/frameActionlist:id157}}\label{\detokenize{auto/frameActionlist:frame-da-image}}\end{figure}

\sphinxAtStartPar
\sphinxstylestrong{導入バージョン:}  ex3.10.0

\sphinxAtStartPar
\sphinxstylestrong{更新バージョン:}  v8.2


\bigskip\hrule\bigskip



\paragraph{交渉 (表切札)}
\label{\detokenize{auto/frameActionlist:trump-d3}}\label{\detokenize{auto/frameActionlist:id125}}
\sphinxAtStartPar
\sphinxstylestrong{フレーム:}

\sphinxAtStartPar
・エクストラ

\sphinxAtStartPar
\sphinxstylestrong{切札キーカード:} {\normalsize $\diamondsuit$} 3

\sphinxAtStartPar
\sphinxstylestrong{能力:}

\sphinxAtStartPar
自分の場にいる{\normalsize $\diamondsuit$} の兵士は \textless{}交渉\textgreater{} を持つ。

\sphinxAtStartPar
\sphinxstylestrong{イメージ:}
交渉のイメージ (\hyperref[\detokenize{auto/frameActionlist:frame-d3-image}]{Fig.\@ \ref{\detokenize{auto/frameActionlist:frame-d3-image}}})

\begin{figure}[htbp]
\centering
\capstart

\noindent\sphinxincludegraphics[height=200\sphinxpxdimen]{{Pictorial_Key_to_the_Tarot_Pentacles_03}.jpg}
\caption{交渉のイメージ}\label{\detokenize{auto/frameActionlist:id158}}\label{\detokenize{auto/frameActionlist:frame-d3-image}}\end{figure}

\sphinxAtStartPar
\sphinxstylestrong{導入バージョン:}  ex3.4.0

\sphinxAtStartPar
\sphinxstylestrong{更新バージョン:}  v8.1


\bigskip\hrule\bigskip



\paragraph{相殺 (表切札)}
\label{\detokenize{auto/frameActionlist:trump-d4}}\label{\detokenize{auto/frameActionlist:id126}}
\sphinxAtStartPar
\sphinxstylestrong{フレーム:}

\sphinxAtStartPar
・エクストラ

\sphinxAtStartPar
\sphinxstylestrong{切札キーカード:} {\normalsize $\diamondsuit$} 4

\sphinxAtStartPar
\sphinxstylestrong{能力:}

\sphinxAtStartPar
なし

\sphinxAtStartPar
\sphinxstylestrong{イメージ:}
相殺のイメージ (\hyperref[\detokenize{auto/frameActionlist:frame-d4-image}]{Fig.\@ \ref{\detokenize{auto/frameActionlist:frame-d4-image}}})

\begin{figure}[htbp]
\centering
\capstart

\noindent\sphinxincludegraphics[height=200\sphinxpxdimen]{{Pictorial_Key_to_the_Tarot_Pentacles_04}.jpg}
\caption{相殺のイメージ}\label{\detokenize{auto/frameActionlist:id159}}\label{\detokenize{auto/frameActionlist:frame-d4-image}}\end{figure}

\sphinxAtStartPar
\sphinxstylestrong{導入バージョン:}  ex5.30.0

\sphinxAtStartPar
\sphinxstylestrong{更新バージョン:}  v8.2


\bigskip\hrule\bigskip



\paragraph{喪失 (表切札)}
\label{\detokenize{auto/frameActionlist:trump-d5}}\label{\detokenize{auto/frameActionlist:id127}}
\sphinxAtStartPar
\sphinxstylestrong{フレーム:}

\sphinxAtStartPar
・エクストラ

\sphinxAtStartPar
\sphinxstylestrong{切札キーカード:} {\normalsize $\diamondsuit$} 5

\sphinxAtStartPar
\sphinxstylestrong{能力:}

\sphinxAtStartPar
あなたの場にチャージ状態の防壁が2体以上存在する場合、全てのプレイヤーはチャージアクションの効果でキャラクターをチャージ状態にすることができない。

\sphinxAtStartPar
\sphinxstylestrong{イメージ:}
喪失のイメージ (\hyperref[\detokenize{auto/frameActionlist:frame-d5-image}]{Fig.\@ \ref{\detokenize{auto/frameActionlist:frame-d5-image}}})

\begin{figure}[htbp]
\centering
\capstart

\noindent\sphinxincludegraphics[height=200\sphinxpxdimen]{{Pictorial_Key_to_the_Tarot_Pentacles_05}.jpg}
\caption{喪失のイメージ}\label{\detokenize{auto/frameActionlist:id160}}\label{\detokenize{auto/frameActionlist:frame-d5-image}}\end{figure}

\sphinxAtStartPar
\sphinxstylestrong{導入バージョン:}  ex5.30.0

\sphinxAtStartPar
\sphinxstylestrong{更新バージョン:}  v8.2


\bigskip\hrule\bigskip



\paragraph{修繕 (表切札)}
\label{\detokenize{auto/frameActionlist:trump-d8}}\label{\detokenize{auto/frameActionlist:id128}}
\sphinxAtStartPar
\sphinxstylestrong{フレーム:}

\sphinxAtStartPar
・エクストラ

\sphinxAtStartPar
\sphinxstylestrong{切札キーカード:} {\normalsize $\diamondsuit$} 8

\sphinxAtStartPar
\sphinxstylestrong{能力:}

\sphinxAtStartPar
なし

\sphinxAtStartPar
\sphinxstylestrong{イメージ:}
修繕のイメージ (\hyperref[\detokenize{auto/frameActionlist:frame-d8-image}]{Fig.\@ \ref{\detokenize{auto/frameActionlist:frame-d8-image}}})

\begin{figure}[htbp]
\centering
\capstart

\noindent\sphinxincludegraphics[height=200\sphinxpxdimen]{{Pictorial_Key_to_the_Tarot_Pentacles_08}.jpg}
\caption{修繕のイメージ}\label{\detokenize{auto/frameActionlist:id161}}\label{\detokenize{auto/frameActionlist:frame-d8-image}}\end{figure}

\sphinxAtStartPar
\sphinxstylestrong{導入バージョン:}  ex4.14.0

\sphinxAtStartPar
\sphinxstylestrong{更新バージョン:}  v8.2


\bigskip\hrule\bigskip



\paragraph{充足 (表切札)}
\label{\detokenize{auto/frameActionlist:trump-d9}}\label{\detokenize{auto/frameActionlist:id129}}
\sphinxAtStartPar
\sphinxstylestrong{フレーム:}

\sphinxAtStartPar
・エクストラ

\sphinxAtStartPar
\sphinxstylestrong{切札キーカード:} {\normalsize $\diamondsuit$} 9

\sphinxAtStartPar
\sphinxstylestrong{能力:}

\sphinxAtStartPar
なし

\sphinxAtStartPar
\sphinxstylestrong{イメージ:}
充足のイメージ (\hyperref[\detokenize{auto/frameActionlist:frame-d9-image}]{Fig.\@ \ref{\detokenize{auto/frameActionlist:frame-d9-image}}})

\begin{figure}[htbp]
\centering
\capstart

\noindent\sphinxincludegraphics[height=200\sphinxpxdimen]{{Pictorial_Key_to_the_Tarot_Pentacles_09}.jpg}
\caption{充足のイメージ}\label{\detokenize{auto/frameActionlist:id162}}\label{\detokenize{auto/frameActionlist:frame-d9-image}}\end{figure}

\sphinxAtStartPar
\sphinxstylestrong{導入バージョン:}  ex4.18.0

\sphinxAtStartPar
\sphinxstylestrong{更新バージョン:}  v8.2


\bigskip\hrule\bigskip



\paragraph{策士 (表切札)}
\label{\detokenize{auto/frameActionlist:trump-dj}}\label{\detokenize{auto/frameActionlist:id130}}
\sphinxAtStartPar
\sphinxstylestrong{フレーム:}

\sphinxAtStartPar
・エクストラ

\sphinxAtStartPar
\sphinxstylestrong{切札キーカード:} {\normalsize $\diamondsuit$} J

\sphinxAtStartPar
\sphinxstylestrong{能力:}

\sphinxAtStartPar
なし

\sphinxAtStartPar
\sphinxstylestrong{イメージ:}
策士のイメージ (\hyperref[\detokenize{auto/frameActionlist:frame-dj-image}]{Fig.\@ \ref{\detokenize{auto/frameActionlist:frame-dj-image}}})

\begin{figure}[htbp]
\centering
\capstart

\noindent\sphinxincludegraphics[height=200\sphinxpxdimen]{{Pictorial_Key_to_the_Tarot_02_The_High_Priestess}.jpg}
\caption{策士のイメージ}\label{\detokenize{auto/frameActionlist:id163}}\label{\detokenize{auto/frameActionlist:frame-dj-image}}\end{figure}

\sphinxAtStartPar
\sphinxstylestrong{導入バージョン:}  ex4.22.0

\sphinxAtStartPar
\sphinxstylestrong{更新バージョン:}  v8.2


\bigskip\hrule\bigskip



\subsubsection{{\normalsize $\clubsuit$} クラブ}
\label{\detokenize{auto/frameActionlist:id131}}

\paragraph{スイッチ (表切札)}
\label{\detokenize{auto/frameActionlist:trump-c4}}\label{\detokenize{auto/frameActionlist:id132}}
\sphinxAtStartPar
\sphinxstylestrong{フレーム:}

\sphinxAtStartPar
・エクストラ

\sphinxAtStartPar
\sphinxstylestrong{切札キーカード:} {\normalsize $\clubsuit$} 4

\sphinxAtStartPar
\sphinxstylestrong{能力:}

\sphinxAtStartPar
自分の場にいる{\normalsize $\clubsuit$} の兵士は \textless{}スイッチ\textgreater{} を持つ。

\sphinxAtStartPar
\sphinxstylestrong{イメージ:}
スイッチのイメージ (\hyperref[\detokenize{auto/frameActionlist:frame-c4-image}]{Fig.\@ \ref{\detokenize{auto/frameActionlist:frame-c4-image}}})

\begin{figure}[htbp]
\centering
\capstart

\noindent\sphinxincludegraphics[height=200\sphinxpxdimen]{{Pictorial_Key_to_the_Tarot_Wands_04}.jpg}
\caption{スイッチのイメージ}\label{\detokenize{auto/frameActionlist:id164}}\label{\detokenize{auto/frameActionlist:frame-c4-image}}\end{figure}

\sphinxAtStartPar
\sphinxstylestrong{導入バージョン:}  ex3.4.0

\sphinxAtStartPar
\sphinxstylestrong{更新バージョン:}  v8.1


\bigskip\hrule\bigskip



\paragraph{内乱 (表切札)}
\label{\detokenize{auto/frameActionlist:trump-c5}}\label{\detokenize{auto/frameActionlist:id133}}
\sphinxAtStartPar
\sphinxstylestrong{フレーム:}

\sphinxAtStartPar
・エクストラ

\sphinxAtStartPar
\sphinxstylestrong{切札キーカード:} {\normalsize $\clubsuit$} 5

\sphinxAtStartPar
\sphinxstylestrong{能力:}

\sphinxAtStartPar
なし

\sphinxAtStartPar
\sphinxstylestrong{イメージ:}
内乱のイメージ (\hyperref[\detokenize{auto/frameActionlist:frame-c5-image}]{Fig.\@ \ref{\detokenize{auto/frameActionlist:frame-c5-image}}})

\begin{figure}[htbp]
\centering
\capstart

\noindent\sphinxincludegraphics[height=200\sphinxpxdimen]{{Pictorial_Key_to_the_Tarot_Wands_05}.jpg}
\caption{内乱のイメージ}\label{\detokenize{auto/frameActionlist:id165}}\label{\detokenize{auto/frameActionlist:frame-c5-image}}\end{figure}

\sphinxAtStartPar
\sphinxstylestrong{導入バージョン:}  ex4.14.0

\sphinxAtStartPar
\sphinxstylestrong{更新バージョン:}  v8.2


\bigskip\hrule\bigskip



\paragraph{手札補充 (表切札)}
\label{\detokenize{auto/frameActionlist:trump-c7}}\label{\detokenize{auto/frameActionlist:id134}}
\sphinxAtStartPar
\sphinxstylestrong{フレーム:}

\sphinxAtStartPar
・エクストラ

\sphinxAtStartPar
\sphinxstylestrong{切札キーカード:} {\normalsize $\clubsuit$} 7

\sphinxAtStartPar
\sphinxstylestrong{能力:}

\sphinxAtStartPar
なし

\sphinxAtStartPar
\sphinxstylestrong{イメージ:}
手札補充のイメージ (\hyperref[\detokenize{auto/frameActionlist:frame-c7-image}]{Fig.\@ \ref{\detokenize{auto/frameActionlist:frame-c7-image}}})

\begin{figure}[htbp]
\centering
\capstart

\noindent\sphinxincludegraphics[height=200\sphinxpxdimen]{{Pictorial_Key_to_the_Tarot_Wands_07}.jpg}
\caption{手札補充のイメージ}\label{\detokenize{auto/frameActionlist:id166}}\label{\detokenize{auto/frameActionlist:frame-c7-image}}\end{figure}

\sphinxAtStartPar
\sphinxstylestrong{導入バージョン:}  ex3.10.0

\sphinxAtStartPar
\sphinxstylestrong{更新バージョン:}  v8.2


\bigskip\hrule\bigskip



\paragraph{手札補充 (表切札)}
\label{\detokenize{auto/frameActionlist:trump-c8}}\label{\detokenize{auto/frameActionlist:id135}}
\sphinxAtStartPar
\sphinxstylestrong{フレーム:}

\sphinxAtStartPar
・エクストラ

\sphinxAtStartPar
\sphinxstylestrong{切札キーカード:} {\normalsize $\clubsuit$} 8

\sphinxAtStartPar
\sphinxstylestrong{能力:}

\sphinxAtStartPar
なし

\sphinxAtStartPar
\sphinxstylestrong{イメージ:}
手札補充のイメージ (\hyperref[\detokenize{auto/frameActionlist:frame-c8-image}]{Fig.\@ \ref{\detokenize{auto/frameActionlist:frame-c8-image}}})

\begin{figure}[htbp]
\centering
\capstart

\noindent\sphinxincludegraphics[height=200\sphinxpxdimen]{{Pictorial_Key_to_the_Tarot_Wands_08}.jpg}
\caption{手札補充のイメージ}\label{\detokenize{auto/frameActionlist:id167}}\label{\detokenize{auto/frameActionlist:frame-c8-image}}\end{figure}

\sphinxAtStartPar
\sphinxstylestrong{導入バージョン:}  ex5.30.0

\sphinxAtStartPar
\sphinxstylestrong{更新バージョン:}  v8.2


\bigskip\hrule\bigskip



\paragraph{要塞 (表切札)}
\label{\detokenize{auto/frameActionlist:trump-c9}}\label{\detokenize{auto/frameActionlist:id136}}
\sphinxAtStartPar
\sphinxstylestrong{フレーム:}

\sphinxAtStartPar
・エクストラ

\sphinxAtStartPar
\sphinxstylestrong{切札キーカード:} {\normalsize $\clubsuit$} 9

\sphinxAtStartPar
\sphinxstylestrong{能力:}

\sphinxAtStartPar
対戦相手がキーカードに{\normalsize $\spadesuit$} を含むアクションの効果を発揮する直前に自分の場にキャラクターがいる場合、その効果によってあなたはダメージを受けない。

\sphinxAtStartPar
\sphinxstylestrong{イメージ:}
要塞のイメージ (\hyperref[\detokenize{auto/frameActionlist:frame-c9-image}]{Fig.\@ \ref{\detokenize{auto/frameActionlist:frame-c9-image}}})

\begin{figure}[htbp]
\centering
\capstart

\noindent\sphinxincludegraphics[height=200\sphinxpxdimen]{{Pictorial_Key_to_the_Tarot_Wands_09}.jpg}
\caption{要塞のイメージ}\label{\detokenize{auto/frameActionlist:id168}}\label{\detokenize{auto/frameActionlist:frame-c9-image}}\end{figure}

\sphinxAtStartPar
\sphinxstylestrong{導入バージョン:}  ex4.18.0

\sphinxAtStartPar
\sphinxstylestrong{更新バージョン:}  v8.2


\bigskip\hrule\bigskip



\paragraph{悪あがき (表切札)}
\label{\detokenize{auto/frameActionlist:trump-c10}}\label{\detokenize{auto/frameActionlist:id137}}
\sphinxAtStartPar
\sphinxstylestrong{フレーム:}

\sphinxAtStartPar
・エクストラ

\sphinxAtStartPar
\sphinxstylestrong{切札キーカード:} {\normalsize $\clubsuit$} 10

\sphinxAtStartPar
\sphinxstylestrong{能力:}

\sphinxAtStartPar
なし

\sphinxAtStartPar
\sphinxstylestrong{イメージ:}
悪あがきのイメージ (\hyperref[\detokenize{auto/frameActionlist:frame-c10-image}]{Fig.\@ \ref{\detokenize{auto/frameActionlist:frame-c10-image}}})

\begin{figure}[htbp]
\centering
\capstart

\noindent\sphinxincludegraphics[height=200\sphinxpxdimen]{{Pictorial_Key_to_the_Tarot_Wands_10}.jpg}
\caption{悪あがきのイメージ}\label{\detokenize{auto/frameActionlist:id169}}\label{\detokenize{auto/frameActionlist:frame-c10-image}}\end{figure}

\sphinxAtStartPar
\sphinxstylestrong{導入バージョン:}  ex5.30.0

\sphinxAtStartPar
\sphinxstylestrong{更新バージョン:}  v8.2


\bigskip\hrule\bigskip



\paragraph{騎士 (表切札)}
\label{\detokenize{auto/frameActionlist:trump-cj}}\label{\detokenize{auto/frameActionlist:id138}}
\sphinxAtStartPar
\sphinxstylestrong{フレーム:}

\sphinxAtStartPar
・エクストラ

\sphinxAtStartPar
\sphinxstylestrong{切札キーカード:} {\normalsize $\clubsuit$} J

\sphinxAtStartPar
\sphinxstylestrong{能力:}

\sphinxAtStartPar
なし

\sphinxAtStartPar
\sphinxstylestrong{イメージ:}
騎士のイメージ (\hyperref[\detokenize{auto/frameActionlist:frame-cj-image}]{Fig.\@ \ref{\detokenize{auto/frameActionlist:frame-cj-image}}})

\begin{figure}[htbp]
\centering
\capstart

\noindent\sphinxincludegraphics[height=200\sphinxpxdimen]{{Pictorial_Key_to_the_Tarot_13_Death}.jpg}
\caption{騎士のイメージ}\label{\detokenize{auto/frameActionlist:id170}}\label{\detokenize{auto/frameActionlist:frame-cj-image}}\end{figure}

\sphinxAtStartPar
\sphinxstylestrong{導入バージョン:}  ex4.22.0

\sphinxAtStartPar
\sphinxstylestrong{更新バージョン:}  v8.2


\bigskip\hrule\bigskip



\subsubsection{裏切札}
\label{\detokenize{auto/frameActionlist:id139}}

\paragraph{裏切札 (裏切札)}
\label{\detokenize{auto/frameActionlist:trump-secrettrump}}\label{\detokenize{auto/frameActionlist:id140}}
\sphinxAtStartPar
\sphinxstylestrong{フレーム:}

\sphinxAtStartPar
・エクストラ

\sphinxAtStartPar
\sphinxstylestrong{切札キーカード:} 全て(裏向き)

\sphinxAtStartPar
\sphinxstylestrong{能力:}

\sphinxAtStartPar
なし

\sphinxAtStartPar
\sphinxstylestrong{導入バージョン:}  v8.1

\sphinxstepscope


\chapter{コアルール}
\label{\detokenize{core/core:core-rst}}\label{\detokenize{core/core:id1}}\label{\detokenize{core/core::doc}}
\sphinxAtStartPar
トランプのみでトレーディングカードゲームのように遊ぶために、
一般的なトレーディングカードゲームを抽象化し、トランプのみで遊べるように再構築しました。

\sphinxAtStartPar
コアルールは、割込みが可能なターン制ゲームの開始から勝敗が決まるまでを定義します。(\hyperref[\detokenize{core/core:abstract-core-rule}]{Fig.\@ \ref{\detokenize{core/core:abstract-core-rule}}})

\sphinxAtStartPar
BlackPokerはコアルールを実装しているため、ルールを変更する場合は、コアルールに従う必要があります。

\sphinxAtStartPar
このルールは、割込みが可能なターン制ゲームの標準化したモデルとして構築されており、他のゲームへの応用も可能です。

\begin{figure}[htbp]
\centering
\capstart

\noindent\sphinxincludegraphics[scale=0.5]{{plantuml-5869adc776d20b7021b88175e8de80f0f436330c}.pdf}
\caption{コアルールのイメージ}\label{\detokenize{core/core:id34}}\label{\detokenize{core/core:abstract-core-rule}}\end{figure}


\section{基礎概念}
\label{\detokenize{core/core:id2}}
\sphinxAtStartPar
割込み可能なターン制ゲームでは、現実のゲーム盤面に影響を与える前に、
どの行動を先に行うかを決める仮想的な場所があります。

\sphinxAtStartPar
割込み可能なターン制ゲームは、見方を変えると“許可制のゲーム”とも表現できます。

\sphinxAtStartPar
つまり、プレイヤーは相手に許可を得て、自分の行動を実行できるかどうか確認しながら進める形式のゲームです。

\sphinxAtStartPar
いくつかの基礎概念を紹介します。
その後、それらを組み合わせたイメージを説明します。

\index{ターン@\spxentry{ターン}}\ignorespaces 

\subsection{ターン}
\label{\detokenize{core/core:index-0}}\label{\detokenize{core/core:id3}}
\sphinxAtStartPar
プレイヤーはターンを“持つ”ことができます。
ターンを持っているプレイヤーは先に行動できます。
ターンを持っているプレイヤーをターンプレイヤーといいます。

\index{アクション(コア)@\spxentry{アクション(コア)}}\ignorespaces 

\subsection{アクション}
\label{\detokenize{core/core:index-1}}\label{\detokenize{core/core:id4}}
\sphinxAtStartPar
ターン制のゲームでは、プレイヤーは様々な行動を行います。
チェスであればコマを進めたり、ババ抜きであれば隣の人からカードを引いたりします。
また、ゲームのルールシステムが行う行動もあります。それらの行動をアクションと定義します。
プレイヤーが行う行動を直接アクション、ルールシステムが行うアクションを誘発アクションといいます。

\index{リクエスト@\spxentry{リクエスト}}\ignorespaces 

\subsection{リクエスト}
\label{\detokenize{core/core:index-2}}\label{\detokenize{core/core:id5}}
\sphinxAtStartPar
アクションを実行するために、ゲームシステム（コアフロー）に対して
どのようなアクションを実行するかを要求します。
この要求をリクエストと呼びます。正式名称はアクションリクエストです。

\sphinxAtStartPar
リクエストの種類によっては、ステージに蓄積されず、すぐに実行(解決)されるものもあります。


\subsection{解決}
\label{\detokenize{core/core:id6}}
\sphinxAtStartPar
リクエストを処理すること（要求されたアクションを実行すること）を解決と言います。

\index{チャンス@\spxentry{チャンス}}\ignorespaces 

\subsection{チャンス}
\label{\detokenize{core/core:index-3}}\label{\detokenize{core/core:id7}}
\sphinxAtStartPar
プレイヤーがリクエストできる機会をチャンスといいます。
チャンスを持っている間は何度でもリクエストすることができます。
正式名称はアクションチャンスです。

\index{ステージ@\spxentry{ステージ}}\ignorespaces 

\subsection{ステージ}
\label{\detokenize{core/core:index-4}}\label{\detokenize{core/core:id8}}
\sphinxAtStartPar
ステージはリクエストの解決順を整理するために使う領域です。
後入れ先出し方式で、最後に積まれたリクエストから順に解決されます。
リクエストの種類によっては、ステージに蓄積されず、すぐに解決されるものもあります。

\index{コンポーネント@\spxentry{コンポーネント}}\ignorespaces 

\subsection{コンポーネント}
\label{\detokenize{core/core:component}}\label{\detokenize{core/core:index-5}}\label{\detokenize{core/core:id9}}
\sphinxAtStartPar
コンポーネントとは、ゲーム盤面に配置されるモンスターや駒の定義です。
例えば、将棋の「歩」は駒としては複数存在しますが、「歩」の定義は1つです。
この定義をコンポーネントといいます。

\index{ユニット@\spxentry{ユニット}}\ignorespaces 

\subsection{ユニット}
\label{\detokenize{core/core:index-6}}\label{\detokenize{core/core:id10}}
\sphinxAtStartPar
コンポーネントとして定義されたものがゲーム盤面に配置されているとき、それをユニットといいます。
例えば、将棋の「歩」はコンポーネントとして定義されていますが、 実際の盤面に配置される「歩」は、コンポーネント定義という設計図から作られたユニットです。


\subsection{コアフロー(ルールシステム)}
\label{\detokenize{core/core:id11}}
\sphinxAtStartPar
リクエストはコアフロー(ルールシステム)によって整理され、処理（解決）されます。
リクエストは即時解決されるものと、ステージに蓄積されるものに分類され、順番に処理されます。

\sphinxAtStartPar
ゲームによって具体的に行う内容は異なりますが、処理する順番の制御はコアフローが担います。

\sphinxAtStartPar
コアフローはゲームの開始から勝敗が決まるまで動作し続けます。


\section{詳細}
\label{\detokenize{core/core:id12}}
\sphinxAtStartPar
基礎概念を図で表すと次のようになります。(\hyperref[\detokenize{core/core:abstract-core-image}]{Fig.\@ \ref{\detokenize{core/core:abstract-core-image}}})

\begin{figure}[htbp]
\centering
\capstart

\noindent\sphinxincludegraphics{{abstract}.pdf}
\caption{割込み可能なターン制ゲーム}\label{\detokenize{core/core:id35}}\label{\detokenize{core/core:abstract-core-image}}\end{figure}

\sphinxAtStartPar
仮想的な場所でリクエストを整理し、現実のゲーム盤面に承認された順で変更を反映します。

\sphinxAtStartPar
リクエストを整理することで割込みを実現しています。どのようにリクエストを処理するかは、コアフローに従います。

\sphinxAtStartPar
さらに、アクションとリクエスト、コンポーネントとユニットの関係は次のようになります。(\hyperref[\detokenize{core/core:action-request-image}]{Fig.\@ \ref{\detokenize{core/core:action-request-image}}})

\begin{figure}[htbp]
\centering
\capstart

\noindent\sphinxincludegraphics[scale=0.5]{{plantuml-b424b035f54a58f3126979e401d275cf680053a5}.pdf}
\caption{リクエストとユニットの関係}\label{\detokenize{core/core:id36}}\label{\detokenize{core/core:action-request-image}}\end{figure}

\sphinxAtStartPar
ゲーム盤面には複数のユニットが生成されます。

\sphinxAtStartPar
リクエストが解決されるたびに、ユニットが生成されたり、既存のユニットの状態が変化したりします。

\sphinxAtStartPar
ここからは、アクションとコンポーネントの各項目について説明します。


\subsection{アクションの定義項目}
\label{\detokenize{core/core:id13}}
\sphinxAtStartPar
アクションを定義する際には、次の項目を設定する必要があります。
その他の項目は、具体的なアクションに応じて追加してください。
\begin{itemize}
\item {} 
\sphinxAtStartPar
オーナー

\item {} 
\sphinxAtStartPar
トリガー

\item {} 
\sphinxAtStartPar
スピード

\item {} 
\sphinxAtStartPar
タイミング

\item {} 
\sphinxAtStartPar
起動条件

\item {} 
\sphinxAtStartPar
誘発条件

\item {} 
\sphinxAtStartPar
効果

\end{itemize}

\index{トリガー(アクション)@\spxentry{トリガー(アクション)}}\ignorespaces 

\subsubsection{オーナー}
\label{\detokenize{core/core:index-7}}\label{\detokenize{core/core:id14}}
\sphinxAtStartPar
アクションの所有者。
アクションの定義はプレイヤー毎に保持するため、同じ内容のアクションでもオーナーが異なります。

\index{トリガー(コア)@\spxentry{トリガー(コア)}}\ignorespaces 

\subsubsection{トリガー}
\label{\detokenize{core/core:trigger-core}}\label{\detokenize{core/core:index-8}}\label{\detokenize{core/core:id15}}
\sphinxAtStartPar
アクションが要求される方法は、大きく分けて次の2種類に分類されます。
\begin{description}
\sphinxlineitem{\sphinxstylestrong{直接}}
\sphinxAtStartPar
プレイヤーがコストを支払うなどの手続きを経てリクエストされるアクション。

\sphinxlineitem{\sphinxstylestrong{誘発}}
\sphinxAtStartPar
条件が満たされた場合、自動でリクエストされるアクション。

\end{description}

\sphinxAtStartPar
トリガー項目には「直接」または「誘発」のいずれかが設定されます。

\index{スピード(コア)@\spxentry{スピード(コア)}}\ignorespaces 

\subsubsection{スピード}
\label{\detokenize{core/core:speed-core}}\label{\detokenize{core/core:index-9}}\label{\detokenize{core/core:id16}}
\sphinxAtStartPar
リクエストが処理される速度は次の2種類に分類されます。
\begin{description}
\sphinxlineitem{\sphinxstylestrong{即時}}
\sphinxAtStartPar
リクエストはステージを用いず解決されます。

\sphinxlineitem{\sphinxstylestrong{通常}}
\sphinxAtStartPar
リクエストはステージを経由して解決されます。

\end{description}

\sphinxAtStartPar
スポード項目には「即時」または「通常」のいずれかが設定されます。

\index{タイミング(コア)@\spxentry{タイミング(コア)}}\ignorespaces 

\subsubsection{タイミング}
\label{\detokenize{core/core:timing}}\label{\detokenize{core/core:index-10}}\label{\detokenize{core/core:id17}}
\sphinxAtStartPar
タイミングとは、アクションがリクエストされるタイミングを示します。
タイミングには次の2種類があります。

\index{メイン@\spxentry{メイン}}\ignorespaces \begin{description}
\sphinxlineitem{\sphinxstylestrong{メイン}}
\sphinxAtStartPar
ターンプレイヤーかつステージが空のときに起こせるアクション。
実行条件：
\sphinxhyphen{} チャンスを持っている
\sphinxhyphen{} 自分のターンである
\sphinxhyphen{} ステージが空である

\end{description}

\index{クイック@\spxentry{クイック}}\ignorespaces \begin{description}
\sphinxlineitem{\sphinxstylestrong{クイック}}
\sphinxAtStartPar
いつでも起こすことができ、アクションをステージに積み重ねることが可能。
実行条件：
\sphinxhyphen{} チャンスを持っている

\end{description}

\begin{sphinxadmonition}{note}{注釈:}
\sphinxAtStartPar
エンドアクションの定義

\sphinxAtStartPar
最低1つはターンを別のプレイヤーに渡すアクションを定義してください。
これがないと、ターンが進行せずゲームが停止する可能性があります。
\end{sphinxadmonition}

\begin{sphinxadmonition}{note}{注釈:}
\sphinxAtStartPar
アクションのコントローラー

\sphinxAtStartPar
アクションを実行したプレイヤーを \sphinxstylestrong{アクションのコントローラー} と呼びます。
効果の解釈は、このコントローラーの視点で行われます。
\end{sphinxadmonition}

\index{き\textbar{}起動条件(コア)@\spxentry{き\textbar{}起動条件(コア)}}\ignorespaces 

\subsubsection{起動条件}
\label{\detokenize{core/core:index-13}}\label{\detokenize{core/core:id18}}
\sphinxAtStartPar
アクションを起こす（リクエストする）ための条件を示します。

\sphinxAtStartPar
トリガーが「直接」の場合、起動条件が定義されます。
コストの支払いや対象の指定など、様々な条件がアクションごとに設定されます。

\sphinxAtStartPar
BlackPokerでは、コストの支払いや対象の指定の記述が冗長にならないよう、省略されることが多いです。

\index{ゆ\textbar{}誘発条件(コア)@\spxentry{ゆ\textbar{}誘発条件(コア)}}\ignorespaces 

\subsubsection{誘発条件}
\label{\detokenize{core/core:index-14}}\label{\detokenize{core/core:id19}}
\sphinxAtStartPar
アクションが誘発される条件を示します。
条件が満たされると自動的にリクエストされます。

\sphinxAtStartPar
トリガーが「誘発」の場合、この項目が定義されます。

\sphinxAtStartPar
例：
\sphinxhyphen{} ダメージを受けたとき
\sphinxhyphen{} カードが墓地に移動したとき

\sphinxAtStartPar
これらの状況で誘発するアクションが設定されることがあります。

\index{こ\textbar{}効果(コア)@\spxentry{こ\textbar{}効果(コア)}}\ignorespaces 

\subsubsection{効果}
\label{\detokenize{core/core:index-15}}\label{\detokenize{core/core:id20}}
\sphinxAtStartPar
効果とは、アクションが解決された際に実行される処理を指します。


\subsection{リクエストの定義項目}
\label{\detokenize{core/core:id21}}\begin{description}
\sphinxlineitem{\sphinxstylestrong{アクション定義}}
\sphinxAtStartPar
リクエストが解決された際に実行されるアクションの内容。

\sphinxlineitem{\sphinxstylestrong{スピード}}
\sphinxAtStartPar
アクション定義のスピードとなります。

\sphinxlineitem{\sphinxstylestrong{タイミング}}
\sphinxAtStartPar
アクション定義のタイミングとなります。

\sphinxlineitem{\sphinxstylestrong{コントローラー}}
\sphinxAtStartPar
アクションを実行したプレイヤーがリクエストのコントローラーとなります。

\end{description}


\subsection{コンポーネントの定義項目}
\label{\detokenize{core/core:id22}}
\sphinxAtStartPar
コンポーネントには、次の項目が定義されます。
必要に応じて、ゲームに合わせた追加設定をしてください。
\begin{description}
\sphinxlineitem{\sphinxstylestrong{オーナー}}
\sphinxAtStartPar
コンポーネンの所有者。
コンポーネントの定義はプレイヤー毎に保持するため、同じ内容のコンポーネントでもオーナーが異なります。

\sphinxlineitem{\sphinxstylestrong{能力}}
\sphinxAtStartPar
能力の詳細については、後述のセクションを参照してください。(\hyperref[\detokenize{core/core:ability}]{\ref{\detokenize{core/core:ability}} \nameref{\detokenize{core/core:ability}}})

\end{description}


\subsection{ユニットの定義項目}
\label{\detokenize{core/core:id23}}
\sphinxAtStartPar
ユニットには、次の項目が設定されます。
必要に応じて、ゲームに合わせた追加設定をしてください。
\begin{description}
\sphinxlineitem{\sphinxstylestrong{コンポーネント定義}}
\sphinxAtStartPar
どのコンポーネント定義から生成されたのかを保持します。

\end{description}

\index{オーナー@\spxentry{オーナー}}\ignorespaces \begin{description}
\sphinxlineitem{\sphinxstylestrong{オーナー}}
\sphinxAtStartPar
ユニットの所有者。
一般的なトランプゲームでは無視されることが多いですが、TCGのようにデッキを個人所有するゲームでは必要な情報です。

\end{description}

\index{コントローラー@\spxentry{コントローラー}}\ignorespaces \begin{description}
\sphinxlineitem{\sphinxstylestrong{コントローラー}}
\sphinxAtStartPar
現在、そのユニットを操作しているプレイヤー。
通常はオーナーとコントローラーは同じですが、コントロールを奪うアクションがある場合、異なることがあります。

\end{description}

\begin{sphinxadmonition}{note}{注釈:}
\sphinxAtStartPar
ユニットとリクエストのコントローラー

\sphinxAtStartPar
コントローラーは制御している人という意味になるため、ユニットとリクエストのコントローラーは制御する対象が異なります。
ユニットとリクエストの属性を次の図に示します。(\hyperref[\detokenize{core/core:controller-attr}]{Fig.\@ \ref{\detokenize{core/core:controller-attr}}})
\end{sphinxadmonition}

\begin{figure}[htbp]
\centering
\capstart

\noindent\sphinxincludegraphics[scale=0.5]{{plantuml-e3598ea29fb852c45136cb979c397a63a8586247}.pdf}
\caption{コントローラー属性}\label{\detokenize{core/core:id37}}\label{\detokenize{core/core:controller-attr}}\end{figure}

\index{の\textbar{}能力(コア)@\spxentry{の\textbar{}能力(コア)}}\ignorespaces 

\section{能力}
\label{\detokenize{core/core:ability}}\label{\detokenize{core/core:index-18}}\label{\detokenize{core/core:id24}}
\sphinxAtStartPar
アクション、コンポーネントの定義項目を見てきました。
これらとは別の概念である \sphinxstylestrong{能力} について説明します。

\sphinxAtStartPar
能力とはアクションの効果とは異なる概念で、アクションを起こす際や効果を解釈する際に参照されます。

\sphinxAtStartPar
能力は解釈される際にコストは支払われず、ステージに置かれません。

\sphinxAtStartPar
能力を持つことができるのは、プレイヤーの他に駒やカードなどのゲームに登場するコンポーネントも含まれます。
(\hyperref[\detokenize{core/core:ability-image}]{Fig.\@ \ref{\detokenize{core/core:ability-image}}})

\begin{figure}[htbp]
\centering
\capstart

\noindent\sphinxincludegraphics[scale=0.5]{{plantuml-5c3b9d2c43d8740600145b89955601a615170288}.pdf}
\caption{能力のイメージ}\label{\detokenize{core/core:id38}}\label{\detokenize{core/core:ability-image}}\end{figure}

\begin{sphinxadmonition}{note}{注釈:}
\sphinxAtStartPar
7版までは、能力に誘発能力と常在型能力がありました。
8版からは、誘発型能力とアクションを起こせる能力をアクションの定義側に移動しました。
能力はそれ以外の常在型能力を示すものになりました。
\end{sphinxadmonition}

\index{コアフロー@\spxentry{コアフロー}}\ignorespaces 

\section{コアフロー}
\label{\detokenize{core/core:coreflowsec}}\label{\detokenize{core/core:index-19}}\label{\detokenize{core/core:id25}}
\sphinxAtStartPar
今まで説明してきた概念を用いて \sphinxstylestrong{コアフローの具体的な処理} を説明します。
この図は \sphinxstylestrong{ゲームの開始から勝敗が決まるまでの流れ(コアフロー)} を示しています。(\hyperref[\detokenize{core/core:coreflow-2}]{Fig.\@ \ref{\detokenize{core/core:coreflow-2}}})

\sphinxAtStartPar
BlackPokerはこのコアフローに則りリクエストが処理されます。

\sphinxAtStartPar
アナログゲーム用に作成したコアフローであるため、なるべく記憶する容量を減らすように設計しています。
デジタルゲームに応用する場合は、細部をゲームに合わせて変更してください。

\begin{figure}[htbp]
\centering
\capstart

\noindent\sphinxincludegraphics[scale=0.5]{{plantuml-56e0b8d8835bda254d3bb4ae5813ef3482975b0f}.pdf}
\caption{コアフロー}\label{\detokenize{core/core:id39}}\label{\detokenize{core/core:coreflow-2}}\end{figure}
\phantomsection\label{\detokenize{core/core:core-gamestart}}\begin{description}
\sphinxlineitem{\sphinxstylestrong{{[}1{]} ゲーム開始}}
\sphinxAtStartPar
先攻を決め、ゲームを始める準備を行います。

\sphinxlineitem{\sphinxstylestrong{{[}2{]} ターンプレイヤーにチャンスを移動}}
\sphinxAtStartPar
ターンを持っているプレイヤーにチャンスを移動します。

\sphinxlineitem{\sphinxstylestrong{{[}3{]} アクションをリクエストするか？}}
\sphinxAtStartPar
チャンスを持っているプレイヤーはアクションを起こすかを判断します。

\sphinxlineitem{\sphinxstylestrong{{[}4{]} パス記録のリセット}}
\sphinxAtStartPar
パスしたプレイヤーの記録をリセットします。

\sphinxlineitem{\sphinxstylestrong{{[}5{]} アクションをリクエストする}}
\sphinxAtStartPar
アクションをリクエストし、これからプレイヤーが行うことを宣言します。
ゲームによってアクションをリクエストする方法は異なります。
BlackPokerではアクション名を言い、コストの支払や対象を指定しアクションをリクエストします。
一方ババ抜きでは、隣のプレイヤーからカードを引く際に宣言せず暗黙にアクションがリクエストされている場合もあります。

\sphinxAtStartPar
このときに有効になっている能力を考慮します。
能力の適用順については \hyperref[\detokenize{core/core:ability-order}]{\ref{\detokenize{core/core:ability-order}} \nameref{\detokenize{core/core:ability-order}}} 参照してください。

\sphinxlineitem{\sphinxstylestrong{{[}6{]} 誘発チェック}}
\sphinxAtStartPar
ここに至るまでに誘発したアクションがないかチェックします。
誘発した場合、対象のアクションをリクエストします。
詳しいフローは \hyperref[\detokenize{core/core:trigger-check}]{\ref{\detokenize{core/core:trigger-check}} \nameref{\detokenize{core/core:trigger-check}}} を参照してください。

\sphinxlineitem{\sphinxstylestrong{{[}7{]} 即時か？}}
\sphinxAtStartPar
リクエストのスピードが即時か判定します。

\end{description}
\phantomsection\label{\detokenize{core/core:actresolve}}\begin{description}
\sphinxlineitem{\sphinxstylestrong{{[}8{]} リクエストの解決}}
\sphinxAtStartPar
アクションの効果に定義されている内容を実行します。
その他にコンポーネントを捨て山に移動するなどゲームによって決まった処理があれば行います。
アクションの解決の中でも効果に定義されている内容を実行することのみを指す場合「効果を発揮する」と言います。

\sphinxAtStartPar
このときに有効になっている能力を考慮します。
能力の適用順については \hyperref[\detokenize{core/core:ability-order}]{\ref{\detokenize{core/core:ability-order}} \nameref{\detokenize{core/core:ability-order}}} 参照してください。

\end{description}
\phantomsection\label{\detokenize{core/core:winlose}}\begin{description}
\sphinxlineitem{\sphinxstylestrong{{[}9{]} 勝敗判定}}
\sphinxAtStartPar
ゲームの勝敗を判定します。決着した場合ゲームが終了します。判定の方法はゲームにより異なります。

\sphinxlineitem{\sphinxstylestrong{{[}10{]} ステージに追加}}
\sphinxAtStartPar
リクエストをステージに追加します。

\sphinxlineitem{\sphinxstylestrong{{[}11{]} パス記録に登録}}
\sphinxAtStartPar
パスしたプレイヤーを記録します。パス記録がリセットされるため、同じプレイヤー名は2回登録されません。

\sphinxlineitem{\sphinxstylestrong{{[}12{]} 全員がパスしたか？}}
\sphinxAtStartPar
パス記録に全てのプレイヤー名が記録されているか判定します。

\sphinxlineitem{\sphinxstylestrong{{[}13{]} ルールシステムにチャンスを移動}}
\sphinxAtStartPar
ルールシステムにチャンスを移動します。

\sphinxlineitem{\sphinxstylestrong{{[}14{]} ステージにリクエストが存在するか？}}
\sphinxAtStartPar
ステージにリクエストが存在するか判定します。

\sphinxlineitem{\sphinxstylestrong{{[}15{]} 最後のリクエストを解決}}
\sphinxAtStartPar
最後にステージに追加されたリクエストを解決します。
解決方法は {\hyperref[\detokenize{core/core:actresolve}]{\sphinxcrossref{\DUrole{std,std-ref}{{[}8{]} リクエストの解決}}}} 参照してください。
ステージ上のリクエストを解決する場合、 {\hyperref[\detokenize{core/core:actresolve}]{\sphinxcrossref{\DUrole{std,std-ref}{{[}8{]} リクエストの解決}}}} を行った後、次の内容も合わせて行います。
\begin{itemize}
\item {} 
\sphinxAtStartPar
リクエストをステージから取り除く

\end{itemize}

\sphinxlineitem{\sphinxstylestrong{{[}16{]} チャンス移動}}
\sphinxAtStartPar
チャンスを持っているプレイヤーからチャンスを持っていないプレイヤーにチャンスを移動します。
チャンスを移動するルールはゲームによって異なります。

\end{description}


\subsection{誘発チェック}
\label{\detokenize{core/core:trigger-check}}\label{\detokenize{core/core:id26}}
\sphinxAtStartPar
アクションの中には誘発条件を持っているアクションがあります。
誘発条件に該当した場合、アクションからリクエストが誘発されます。

\sphinxAtStartPar
誘発チェックでは、誘発したリクエストを解決またはステージに追加します。
誘発したリクエストのコントローラーは起因となったアクションのオーナーがコントローラーとなります。
誘発チェックは次の図のように行います。(\hyperref[\detokenize{core/core:trigger-flow}]{Fig.\@ \ref{\detokenize{core/core:trigger-flow}}})

\begin{sphinxadmonition}{note}{注釈:}
\sphinxAtStartPar
バッファ

\sphinxAtStartPar
誘発したリクエストを一時的に溜めておくバッファという領域があります。正式名称はアクションバッファです。
\end{sphinxadmonition}

\begin{figure}[htbp]
\centering
\capstart

\noindent\sphinxincludegraphics[scale=0.5]{{plantuml-6f625fd3091913318d2a09c4f5dda21eb1e08d44}.pdf}
\caption{誘発チェック}\label{\detokenize{core/core:id40}}\label{\detokenize{core/core:trigger-flow}}\end{figure}
\phantomsection\label{\detokenize{core/core:trigger-act-gather}}\begin{description}
\sphinxlineitem{\sphinxstylestrong{{[}6\sphinxhyphen{}1{]} 誘発したリクエストを分類しバッファに追加}}
\sphinxAtStartPar
各プレイヤーが誘発させたリクエストを、スピードおよび
タイミングに基づいて分類し、一旦バッファに追加します。

\sphinxlineitem{\sphinxstylestrong{{[}6\sphinxhyphen{}2{]} バッファは空か？}}
\sphinxAtStartPar
バッファが空であるかどうかを判定します。
未処理のリクエストが残っている場合は、以降の処理ループを継続します。

\sphinxlineitem{\sphinxstylestrong{{[}6\sphinxhyphen{}3{]} バッファに即時はあるか？}}
\sphinxAtStartPar
バッファ内にスピードが即時のリクエストが存在するかを判定します。
存在する場合、
ターンプレイヤーから順に即時の処理グループへ進みます。

\sphinxlineitem{\sphinxstylestrong{{[}6\sphinxhyphen{}4{]} 該当プレイヤーに即時があるか？}}
\sphinxAtStartPar
現在処理対象となっているプレイヤーがコントローラーとなっているリクエストがバッファにあるかを確認します。
スピードが即時のリクエストが存在するかを判定し、
存在しない場合は、そのプレイヤーでの処理を終了し、次のプレイヤーへ移行します。

\sphinxlineitem{\sphinxstylestrong{{[}6\sphinxhyphen{}5{]} タイミング=メインの即時アクションを処理}}
\sphinxAtStartPar
該当プレイヤーについて、タイミングが「メイン」、スピードが「即時」のアクションを実行します。
詳細は \hyperref[\detokenize{core/core:trigger-act-s}]{\ref{\detokenize{core/core:trigger-act-s}} \nameref{\detokenize{core/core:trigger-act-s}}} をタイミング=メインとして参照してください。

\sphinxlineitem{\sphinxstylestrong{{[}6\sphinxhyphen{}6{]} タイミング=クイックの即時アクションを処理}}
\sphinxAtStartPar
同じプレイヤーについて、タイミングが「クイック」、スピードが「即時」のアクションを実行します。
詳細は \hyperref[\detokenize{core/core:trigger-act-s}]{\ref{\detokenize{core/core:trigger-act-s}} \nameref{\detokenize{core/core:trigger-act-s}}} をタイミング=クイックとして参照してください。

\sphinxlineitem{\sphinxstylestrong{{[}6\sphinxhyphen{}7{]} 次のプレイヤーへ}}
\sphinxAtStartPar
現在のプレイヤーでの処理が完了した後、
ターン順に次のプレイヤーへ処理を移行します。

\sphinxlineitem{\sphinxstylestrong{{[}6\sphinxhyphen{}8{]} バッファに通常はあるか？}}
\sphinxAtStartPar
バッファ内にスピードが通常のリクエストが存在するかを判定します。
存在する場合、
ターンプレイヤーから順に即時の処理グループへ進みます。

\sphinxlineitem{\sphinxstylestrong{{[}6\sphinxhyphen{}9{]} 該当プレイヤーに通常があるか？}}
\sphinxAtStartPar
現在処理対象となっているプレイヤーがコントローラーとなっているリクエストがバッファにあるかを確認します。
スピードが即時のリクエストが存在するかを判定し、
存在しない場合は、そのプレイヤーでの処理を終了し、次のプレイヤーへ移行します。

\sphinxlineitem{\sphinxstylestrong{{[}6\sphinxhyphen{}10{]} タイミング=メインの通常アクションを処理}}
\sphinxAtStartPar
該当プレイヤーについて、タイミングが「メイン」、スピードが「通常」のアクションを実行します。
詳細は \hyperref[\detokenize{core/core:trigger-act-n}]{\ref{\detokenize{core/core:trigger-act-n}} \nameref{\detokenize{core/core:trigger-act-n}}} をタイミング=メインとして参照してください。

\sphinxlineitem{\sphinxstylestrong{{[}6\sphinxhyphen{}11{]} タイミング=クイックの通常アクションを処理}}
\sphinxAtStartPar
同じプレイヤーについて、タイミングが「クイック」、スピードが「通常」のアクションを実行します。
詳細は \hyperref[\detokenize{core/core:trigger-act-n}]{\ref{\detokenize{core/core:trigger-act-n}} \nameref{\detokenize{core/core:trigger-act-n}}} をタイミング=クイックとして参照してください。

\sphinxlineitem{\sphinxstylestrong{{[}6\sphinxhyphen{}12{]} 次のプレイヤーへ}}
\sphinxAtStartPar
現在のプレイヤーでの処理が完了した後、
ターン順に次のプレイヤーへ処理を移行します。

\end{description}

\begin{sphinxadmonition}{note}{注釈:}
\sphinxAtStartPar
各処理グループ内では、必ずターンプレイヤーから始まり、
ターンが回る順に全プレイヤーに対して確認およびアクションの処理を実施します。
また、ループはバッファに未処理のアクションが存在する限り繰り返されます。
\end{sphinxadmonition}


\subsubsection{誘発即時解決}
\label{\detokenize{core/core:trigger-act-s}}\label{\detokenize{core/core:id27}}
\sphinxAtStartPar
誘発チェックで誘発したスピードが即時のリクエストを処理します。
呼び出し元で指定されたプレイヤーおよびタイミングに基づいて処理します。
誘発チェックは次の図のように行います。(\hyperref[\detokenize{core/core:trigger-flow-s}]{Fig.\@ \ref{\detokenize{core/core:trigger-flow-s}}})

\begin{figure}[htbp]
\centering
\capstart

\noindent\sphinxincludegraphics[scale=0.5]{{plantuml-caaca34586e95540adb102e9bc8309243f59b393}.pdf}
\caption{誘発チェック\sphinxhyphen{}即時処理}\label{\detokenize{core/core:id41}}\label{\detokenize{core/core:trigger-flow-s}}\end{figure}
\begin{description}
\sphinxlineitem{\sphinxstylestrong{{[}6\sphinxhyphen{}5\sphinxhyphen{}1{]} バッファから該当の即時はあるか？}}
\sphinxAtStartPar
バッファから対象プレイヤーかつ、該当するタイミングかつ、スピードが即時のリクエストが存在するかを判定します。
存在する場合、以降の処理へ進み、存在しなくなるまでこのループを継続します。

\sphinxlineitem{\sphinxstylestrong{{[}6\sphinxhyphen{}5\sphinxhyphen{}2{]} バッファからリクエストを取り出す}}
\sphinxAtStartPar
条件を満たしたリクエストを、バッファから1つ取り出します。
どのリクエストを取り出すかは対象プレイヤーが決定します。
取り出されたリクエストは、解決処理の対象となります。

\sphinxlineitem{\sphinxstylestrong{{[}6\sphinxhyphen{}5\sphinxhyphen{}3{]} 取り出したリクエストを解決する}}
\sphinxAtStartPar
取り出されたリクエストの効果を実行し、解決します。
詳しくは {\hyperref[\detokenize{core/core:actresolve}]{\sphinxcrossref{\DUrole{std,std-ref}{{[}8{]} リクエストの解決}}}} 参照してください。

\sphinxlineitem{\sphinxstylestrong{{[}6\sphinxhyphen{}5\sphinxhyphen{}4{]} 勝敗判定}}
\sphinxAtStartPar
勝敗を判定します。
詳しくは {\hyperref[\detokenize{core/core:winlose}]{\sphinxcrossref{\DUrole{std,std-ref}{{[}9{]} 勝敗判定}}}} 参照。

\sphinxlineitem{\sphinxstylestrong{{[}6\sphinxhyphen{}5\sphinxhyphen{}5{]} リクエストが新たな誘発を発生させたか？}}
\sphinxAtStartPar
効果の解決後、それが起因となり新たな誘発を発生させたかどうかを確認します。
発生している場合は、その誘発アクションを再度バッファに追加する必要があります。

\sphinxlineitem{\sphinxstylestrong{{[}6\sphinxhyphen{}5\sphinxhyphen{}6{]} 誘発したリクエストをバッファに追加}}
\sphinxAtStartPar
新たに誘発されたリクエストが存在する場合、該当アクションをバッファに追加します。
これにより、再帰的なアクション処理が可能となり、次のループで該当するリクエストの取り出し処理が実行されます。

\end{description}


\subsubsection{通常：アクション毎に処理}
\label{\detokenize{core/core:trigger-act-n}}\label{\detokenize{core/core:id28}}
\sphinxAtStartPar
誘発チェックで誘発したスピードが通常のリクエストを処理します。
呼び出し元で指定されたプレイヤーおよびタイミングに基づいて処理します。
誘発チェックは次の図のように行います。(\hyperref[\detokenize{core/core:trigger-flow-n}]{Fig.\@ \ref{\detokenize{core/core:trigger-flow-n}}})

\begin{figure}[htbp]
\centering
\capstart

\noindent\sphinxincludegraphics[scale=0.5]{{plantuml-5d2051ec7bb14a770950bcb23c25d203c53ec5f1}.pdf}
\caption{誘発チェック\sphinxhyphen{}通常処理}\label{\detokenize{core/core:id42}}\label{\detokenize{core/core:trigger-flow-n}}\end{figure}
\begin{description}
\sphinxlineitem{\sphinxstylestrong{{[}6\sphinxhyphen{}10\sphinxhyphen{}1{]} バッファから該当の通常はあるか？}}
\sphinxAtStartPar
バッファから対象プレイヤーかつ、該当するタイミングかつ、スピードが通常のリクエストが存在するかを判定します。
存在する場合、以降の処理へ進み、存在しなくなるまでこのループを継続します。

\sphinxlineitem{\sphinxstylestrong{{[}6\sphinxhyphen{}10\sphinxhyphen{}2{]} バッファからリクエストを取り出す}}
\sphinxAtStartPar
条件を満たしたリクエストを、バッファから1つ取り出します。
どのリクエストを取り出すかは対象プレイヤーが決定します。
取り出されたリクエストは、解決処理の対象となります。

\sphinxlineitem{\sphinxstylestrong{{[}6\sphinxhyphen{}10\sphinxhyphen{}3{]} 取り出したリクエストのタイミングがメインか判定}}
\sphinxAtStartPar
取り出したリクエストのタイミングが「メイン」であるかどうかを判定します。
「Yes」と判定された場合は、ステージへの追加前に空き状況の確認へ進みます。
「No」の場合は、クイックタイミングとして処理を行います。

\sphinxlineitem{\sphinxstylestrong{{[}6\sphinxhyphen{}10\sphinxhyphen{}4{]} ステージが空か判定}}
\sphinxAtStartPar
タイミングがメインの場合、そのリクエストをステージに追加できるかどうか、
すなわちステージに空きがあるかを判定します。

\sphinxlineitem{\sphinxstylestrong{{[}6\sphinxhyphen{}10\sphinxhyphen{}5{]} リクエストをステージに追加}}
\sphinxAtStartPar
ステージが空いている場合、取り出したリクエストをステージに追加します。

\sphinxlineitem{\sphinxstylestrong{{[}6\sphinxhyphen{}10\sphinxhyphen{}6{]} リクエストを破棄}}
\sphinxAtStartPar
ステージが埋まっている場合、取り出したタイミングがメインのリクエストを破棄します。

\sphinxlineitem{\sphinxstylestrong{{[}6\sphinxhyphen{}10\sphinxhyphen{}7{]} リクエストをステージに追加}}
\sphinxAtStartPar
取り出したリクエストのタイミングが「クイック」である場合、
ステージの空き状況にかかわらず無条件でリクエストをステージに追加します。

\sphinxlineitem{\sphinxstylestrong{{[}6\sphinxhyphen{}10\sphinxhyphen{}8{]} リクエストが新たな誘発を発生させたか？}}
\sphinxAtStartPar
リクエストをステージに追加した後、それが起因となり新たな誘発を発生させたかどうかを確認します。
誘発が発生している場合は、後続の処理でアクションバッファへの追加が行われます。

\sphinxlineitem{\sphinxstylestrong{{[}6\sphinxhyphen{}10\sphinxhyphen{}9{]} 誘発したリクエストをバッファに追加}}
\sphinxAtStartPar
新たに誘発されたリクエストが存在する場合、該当アクションをバッファに追加します。
これにより、誘発処理の再実行が可能となります。

\end{description}


\subsection{能力の適応順}
\label{\detokenize{core/core:ability-order}}\label{\detokenize{core/core:id29}}
\sphinxAtStartPar
複数の能力が同時に有効になった場合、効果の内容によっては矛盾を引き起こす可能性があります。
その場合は、次の \hyperref[\detokenize{core/core:ability-order-rule}]{\ref{\detokenize{core/core:ability-order-rule}} \nameref{\detokenize{core/core:ability-order-rule}}} に沿って能力を解釈してください。

\sphinxAtStartPar
なお、ゲームによって細部の裁定が異なる場合があります。ここでは BlackPoker を例に説明します。


\subsubsection{考え方の原則}
\label{\detokenize{core/core:id30}}
\sphinxAtStartPar
影響範囲が広い能力ほど優先順位を低く扱い、影響範囲が狭い（個別の対象を限定する）能力ほど優先順位を高く扱います。
つまり、「広く浅く影響する能力」は先に適用し、「狭く深く影響する能力」は後から適用します。


\subsubsection{能力の適用ルール}
\label{\detokenize{core/core:ability-order-rule}}\label{\detokenize{core/core:id31}}\begin{enumerate}
\sphinxsetlistlabels{\arabic}{enumi}{enumii}{}{.}%
\item {} 
\sphinxAtStartPar
能力を分類する
各能力を「影響範囲」に注目して、次の 3 つに分類します。
\begin{description}
\sphinxlineitem{\sphinxstylestrong{{[}1{]} 全ての (all)}}\begin{itemize}
\item {} 
\sphinxAtStartPar
ゲームの場全体、または全プレイヤーなど「すべて」を有効範囲に含む能力です。

\item {} 
\sphinxAtStartPar
例: 「全ての{\normalsize $\spadesuit$} の兵士はサイズを1加算する」など。

\end{itemize}

\sphinxlineitem{\sphinxstylestrong{{[}2{]} あなたの (your)}}\begin{itemize}
\item {} 
\sphinxAtStartPar
あなたの場や、あなたがコントロールしているコンポーネントだけを有効範囲にする能力です。

\item {} 
\sphinxAtStartPar
例: 「あなたの兵士はドライブ状態でも攻撃できる」など。

\end{itemize}

\sphinxlineitem{\sphinxstylestrong{{[}3{]} この (this)}}\begin{itemize}
\item {} 
\sphinxAtStartPar
対象のユニット 1 つを有効範囲にします。

\item {} 
\sphinxAtStartPar
例: 「この兵士はサイズを3加算する」「対象の兵士のサイズを2減算する」など。

\end{itemize}

\end{description}

\item {} 
\sphinxAtStartPar
能力を適用する
上記で分類した番号の小さい順、すなわち \sphinxstylestrong{{[}1{]} \(\rightarrow\) {[}2{]} \(\rightarrow\) {[}3{]}} の順で能力を適用します。
\begin{enumerate}
\sphinxsetlistlabels{\arabic}{enumii}{enumiii}{}{.}%
\item {} \begin{description}
\sphinxlineitem{\sphinxstylestrong{{[}1{]} 全ての (all)}}
\sphinxAtStartPar
まず「全ての ～」といった広範囲を対象とする能力を先に処理します。

\end{description}

\item {} \begin{description}
\sphinxlineitem{\sphinxstylestrong{{[}2{]} あなたの (your)}}
\sphinxAtStartPar
次に「あなたの ～」等、特定プレイヤーやその場を限定する能力を処理します。

\end{description}

\item {} \begin{description}
\sphinxlineitem{\sphinxstylestrong{{[}3{]} この (this)}}
\sphinxAtStartPar
最後に、個別のコンポーネント単位でより限定的な能力を処理し、必要に応じて上書きします。

\end{description}

\end{enumerate}

\sphinxAtStartPar
\sphinxstylestrong{同じ範囲で矛盾した能力がある場合}
それらの能力が発動（または有効化）された順番に応じて “タイムスタンプ” が古いものから適用し、
新しいものが後から上書きする形で処理します。

\sphinxAtStartPar
このように、同時に有効な効果が互いに矛盾する場合でも、最終的により限定的な効果が優先されます。

\end{enumerate}


\subsubsection{運用上の注意}
\label{\detokenize{core/core:id32}}\begin{itemize}
\item {} \begin{description}
\sphinxlineitem{\sphinxstylestrong{複数の「同じ範囲」の能力同士の扱い}}
\sphinxAtStartPar
たとえば「全ての兵士は攻撃できない」という能力と、「全ての兵士のサイズはこのターン中1加算される」という能力が同じ “{[}1{]} 全ての (all)” 範囲に当たる場合、どちらも基本的には同時に成立し共存します。ただし、まったく相反する内容で整合が取れない場合は、対戦ルールやジャッジの裁定に従ってください。

\end{description}

\item {} \begin{description}
\sphinxlineitem{\sphinxstylestrong{適用タイミング}}
\sphinxAtStartPar
能力の適用は、コアフロー内の特定のステップでまとめて行います。何度か再適用が行われる場面があっても、適用のたびに「{[}1{]} \(\rightarrow\) {[}2{]} \(\rightarrow\) {[}3{]}」の順でまっさらな状態から再計算します。そのため、兵士のサイズが重複して加算され続けるようなことはありません。

\end{description}

\item {} \begin{description}
\sphinxlineitem{\sphinxstylestrong{互いに矛盾しない場合}}
\sphinxAtStartPar
たとえば「全ての兵士のサイズを \sphinxhyphen{}1」「あなたの兵士はさらに +2」「この兵士はさらに +3」という 3 つの効果が同時に有効なとき、BlackPoker ではサイズ 0 以下でも兵士が場に存在し続けられるため、最終的にその兵士には合計 +5 の修正が加わります。特に抵触しないかぎりすべて重複適用されます。

\end{description}

\item {} \begin{description}
\sphinxlineitem{\sphinxstylestrong{互いに矛盾する場合}}
\sphinxAtStartPar
「全ての兵士は攻撃できない」と「この兵士は攻撃できる」というように真っ向から対立する効果が同時に発生した場合、最終的には範囲の狭い「この兵士は攻撃できる」の方が優先されます。

\end{description}

\end{itemize}

\sphinxAtStartPar
以上のルールが、複数の能力が同時に有効になった場合の処理基準となります。


\section{まとめ}
\label{\detokenize{core/core:id33}}
\sphinxAtStartPar
コアルールについて説明しました。
すでにあるターン制のゲームからアクションを洗い出し、能力を整理することで割込処理を可能としゲームの新しい遊び方が見つけられます。
また、新しく作成するゲームに関してもコアルールを意識して作成することで、ルール追加がしやすいゲームが考えやすいと思います。

\sphinxstepscope


\chapter{付録}
\label{\detokenize{appendix/appendix:appendix-rst}}\label{\detokenize{appendix/appendix:id1}}\label{\detokenize{appendix/appendix::doc}}

\section{アクションリスト}
\label{\detokenize{appendix/appendix:id2}}

\subsection{ライト}
\label{\detokenize{appendix/appendix:actionlist-lite}}\label{\detokenize{appendix/appendix:id3}}\begin{description}
\sphinxlineitem{URL}
\sphinxAtStartPar
\sphinxurl{https://blackpoker.github.io/BlackPoker/master/actionlist/html/all.html?format=\%E3\%83\%A9\%E3\%82\%A4\%E3\%83\%88}

\sphinxlineitem{PDF}
\sphinxAtStartPar
\sphinxurl{https://blackpoker.github.io/BlackPoker/master/actionlist/pdf/blackpoker-lite.pdf}

\sphinxAtStartPar
\sphinxurl{https://blackpoker.github.io/BlackPoker/master/actionlist/pdf/blackpoker-lite-2up.pdf}

\end{description}


\subsection{マスター}
\label{\detokenize{appendix/appendix:actionlist-master}}\label{\detokenize{appendix/appendix:id4}}\begin{description}
\sphinxlineitem{URL}
\sphinxAtStartPar
\sphinxurl{https://blackpoker.github.io/BlackPoker/master/actionlist/html/all.html?format=\%E3\%83\%9E\%E3\%82\%B9\%E3\%82\%BF\%E3\%83\%BC}

\sphinxlineitem{PDF}
\sphinxAtStartPar
\sphinxurl{https://blackpoker.github.io/BlackPoker/master/actionlist/pdf/blackpoker-mast.pdf}

\sphinxAtStartPar
\sphinxurl{https://blackpoker.github.io/BlackPoker/master/actionlist/pdf/blackpoker-mast-2up.pdf}

\end{description}


\subsection{エクストラフレーム}
\label{\detokenize{appendix/appendix:actionlist-extra}}\label{\detokenize{appendix/appendix:id5}}\begin{description}
\sphinxlineitem{URL}
\sphinxAtStartPar
\sphinxurl{https://blackpoker.github.io/BlackPoker/master/actionlist/html/all.html?format=\%E3\%83\%9E\%E3\%82\%B9\%E3\%82\%BF\%E3\%83\%BC\&frame=\%E3\%82\%A8\%E3\%82\%AF\%E3\%82\%B9\%E3\%83\%88\%E3\%83\%A9}

\end{description}
\phantomsection\label{\detokenize{appendix/appendix:extralist}}

\section{ガイドライン}
\label{\detokenize{appendix/appendix:id6}}
\sphinxAtStartPar
ガイドラインは、BlackPokerのアクションやルールを設計・記述する際の共通基準を示すものです。

\sphinxstepscope


\subsection{BlackPoker二次創作ガイドライン}
\label{\detokenize{appendix/guideline/fanworks:blackpoker}}\label{\detokenize{appendix/guideline/fanworks::doc}}
\begin{sphinxShadowBox}
\sphinxstyletopictitle{目次}
\begin{itemize}
\item {} 
\sphinxAtStartPar
\phantomsection\label{\detokenize{appendix/guideline/fanworks:id9}}{\hyperref[\detokenize{appendix/guideline/fanworks:id2}]{\sphinxcrossref{はじめに}}}

\item {} 
\sphinxAtStartPar
\phantomsection\label{\detokenize{appendix/guideline/fanworks:id10}}{\hyperref[\detokenize{appendix/guideline/fanworks:id3}]{\sphinxcrossref{二次創作の範囲}}}

\item {} 
\sphinxAtStartPar
\phantomsection\label{\detokenize{appendix/guideline/fanworks:id11}}{\hyperref[\detokenize{appendix/guideline/fanworks:id4}]{\sphinxcrossref{創作の自由と制限}}}

\item {} 
\sphinxAtStartPar
\phantomsection\label{\detokenize{appendix/guideline/fanworks:id12}}{\hyperref[\detokenize{appendix/guideline/fanworks:id5}]{\sphinxcrossref{クレジット表記について}}}

\item {} 
\sphinxAtStartPar
\phantomsection\label{\detokenize{appendix/guideline/fanworks:id13}}{\hyperref[\detokenize{appendix/guideline/fanworks:id6}]{\sphinxcrossref{禁止事項}}}

\item {} 
\sphinxAtStartPar
\phantomsection\label{\detokenize{appendix/guideline/fanworks:id14}}{\hyperref[\detokenize{appendix/guideline/fanworks:id7}]{\sphinxcrossref{商用利用について}}}

\item {} 
\sphinxAtStartPar
\phantomsection\label{\detokenize{appendix/guideline/fanworks:id15}}{\hyperref[\detokenize{appendix/guideline/fanworks:id8}]{\sphinxcrossref{ガイドラインにないケースについて}}}

\end{itemize}
\end{sphinxShadowBox}


\subsubsection{はじめに}
\label{\detokenize{appendix/guideline/fanworks:id2}}
\sphinxAtStartPar
BlackPokerの制作チームでは、ファンの皆さまによる創作活動を心から歓迎しています。BlackPokerを通じて生まれる新たなアイデアや作品は、私たちにとって大きな喜びであり、コミュニティをさらに盛り上げる原動力です。そこで、安心して二次創作を楽しんでいただけるように、本ガイドラインでは創作にあたってのルールとお願い事をまとめました。

\sphinxAtStartPar
このガイドラインを読むことで、BlackPokerという作品名に敬意を払いつつ、自由に新しい遊び方や作品を生み出せるようになることが目的です。皆さまの創造力でBlackPokerの世界がさらに豊かになることを、私たちはとても楽しみにしています。どうぞルールを守って、思う存分創作を楽しんでください！


\subsubsection{二次創作の範囲}
\label{\detokenize{appendix/guideline/fanworks:id3}}
\sphinxAtStartPar
BlackPokerでは、以下のようなジャンルの二次創作を行っていただけます。
\begin{itemize}
\item {} 
\sphinxAtStartPar
\sphinxstylestrong{ゲームルール・バリエーション}
BlackPokerのルールをアレンジした新しい遊び方やオリジナルのルール追加

\item {} 
\sphinxAtStartPar
\sphinxstylestrong{イラスト・漫画}
キャラクターや世界観を題材にしたイラストやマンガ

\item {} 
\sphinxAtStartPar
\sphinxstylestrong{小説・ストーリー}
BlackPokerの設定や物語を膨らませた小説やショートストーリー

\item {} 
\sphinxAtStartPar
\sphinxstylestrong{動画}
ゲームのプレイ動画、解説動画、ファンによるアニメーションやMAD動画など

\item {} 
\sphinxAtStartPar
\sphinxstylestrong{音楽}
イメージソング、BGMのアレンジ、リミックスなどの音楽作品

\item {} 
\sphinxAtStartPar
\sphinxstylestrong{グッズ}
手作りの同人グッズ（キーホルダー、カードケース、ファンアートグッズ等）

\end{itemize}

\sphinxAtStartPar
上記は一例です。この他にも常識の範囲内でBlackPokerに関連する創作であれば、基本的に自由にチャレンジしていただけます。


\subsubsection{創作の自由と制限}
\label{\detokenize{appendix/guideline/fanworks:id4}}
\sphinxAtStartPar
BlackPokerの二次創作は \sphinxstylestrong{誰でも自由に行うことができます} 。本ガイドラインの範囲内であれば、事前に個別の許可を取る必要はありませんので、安心して創作活動を行ってください。ファンの皆さまが考案する新ルールや描いてくださるイラストなど、私たちも心から楽しみにしています。

\sphinxAtStartPar
ただし、二次創作を許可するといっても \sphinxstylestrong{著作権や公式の権利を放棄したわけではない} 点にご注意ください。BlackPoker自体の著作権・商標などの権利は引き続き制作元に帰属しており、原作の無断転載や権利者を偽る行為は認められません。Apache License 2.0のように、作品の利用を許諾しても著作者はその権利を保持し続けるスタンスです。ガイドライン遵守のもとで、原作へのリスペクトを忘れずに創作をお楽しみください。


\subsubsection{クレジット表記について}
\label{\detokenize{appendix/guideline/fanworks:id5}}
\sphinxAtStartPar
創作した作品を公開する際は、 \sphinxstylestrong{「BlackPoker」の名前をクレジットに記載することを推奨します} 。必須ではありませんが、作品の説明文や動画のエンドロール、投稿時のコメント欄などに「これは \sphinxstylestrong{BlackPoker} の二次創作です」等の一文を入れていただけると幸いです。こうしたクレジット表記は原作への敬意を示すことになり、閲覧する方にもその作品がBlackPoker由来のファン創作であると伝わりやすくなります。

\begin{sphinxadmonition}{note}{注釈:}
\sphinxAtStartPar
作品中に「公式」や「公認」などの紛らわしい表現を入れることはお控えください。あくまでファンによる非公式の創作物であることが分かる形で公開していただければと思います。
\end{sphinxadmonition}


\subsubsection{禁止事項}
\label{\detokenize{appendix/guideline/fanworks:id6}}
\sphinxAtStartPar
BlackPokerの二次創作を楽しんでいただくにあたり、以下のような行為・表現はご遠慮ください（禁止事項）。
\begin{itemize}
\item {} 
\sphinxAtStartPar
\sphinxstylestrong{公式作品と誤認させる表現の禁止}
あたかも公式のコンテンツや公式発表であるかのような紛らわしい表現は禁止します。ファン創作であるにも関わらず公式と混同される恐れのある表現（公式ロゴの無断使用や、「公式が制作した」などの記載）は行わないでください。

\item {} 
\sphinxAtStartPar
\sphinxstylestrong{公序良俗に反する内容の禁止}
極端に暴力的・グロテスクな描写、過度な性的表現、差別的な言動やヘイト表現など、公序良俗に反する内容の作品は公開しないでください。見る人が不快になったり社会的に問題となり得る表現は避けてください。

\item {} 
\sphinxAtStartPar
\sphinxstylestrong{第三者の権利を侵害する内容の禁止}
他者の著作権を侵害する画像や音楽の無断使用、他人物のプライバシー侵害、誹謗中傷など、第三者の権利を侵害する内容は含めないでください。二次創作とはいえ、他の作品や人物の権利を尊重することが大切です。

\end{itemize}

\sphinxAtStartPar
以上の禁止事項に抵触する作品は、公開後であっても削除をお願いしたり、必要に応じて法的措置を検討する場合がありますのでご了承ください。


\subsubsection{商用利用について}
\label{\detokenize{appendix/guideline/fanworks:id7}}
\sphinxAtStartPar
BlackPokerの二次創作物の利用や頒布は、 \sphinxstylestrong{原則として非営利} で行ってくださいますようお願いいたします。ファン活動の範囲内（同人イベントでの頒布や、個人が少額の収益を得る程度の動画配信など）であれば収益が発生しても問題ありません。しかし、企業が関与する大規模な商用利用や、同人の範囲を超える大きな利益を伴う利用については、\sphinxstylestrong{事前に必ずご相談ください}。

\sphinxAtStartPar
例えば、BlackPokerの設定を用いた商品を大量生産して販売する、クラウドファンディングで多額の支援を募る、といったケースがこれに該当します。そのようなプロジェクトを計画されている場合は、事前に制作者までお問い合わせいただければ個別に協議させていただきます。ガイドラインの範囲内であれば同人活動は大歓迎ですので、常識の範囲内でお楽しみください。


\subsubsection{ガイドラインにないケースについて}
\label{\detokenize{appendix/guideline/fanworks:id8}}
\sphinxAtStartPar
本ガイドラインに記載されていないケースや、内容の解釈に迷う場合は、どうぞ遠慮なく \sphinxstylestrong{個別にお問い合わせ} ください。公式サイトやSNS等の連絡先を通じてご質問をいただければ、可能な範囲で回答し対応いたします。皆さまに安心して創作を続けていただけるよう、私たちもできる限りサポートさせていただきます。

\sphinxAtStartPar
BlackPokerの世界は、ファンの皆さまのクリエイティブな発想によって無限に広がります。ガイドラインを守りつつ、自由なアイデアでBlackPokerをもっと一緒に盛り上げていきましょう！

\sphinxstepscope


\subsection{BlackPoker アクション・ルールガイドライン}
\label{\detokenize{appendix/guideline/customRule:blackpoker}}\label{\detokenize{appendix/guideline/customRule::doc}}
\begin{sphinxShadowBox}
\sphinxstyletopictitle{目次}
\begin{itemize}
\item {} 
\sphinxAtStartPar
\phantomsection\label{\detokenize{appendix/guideline/customRule:id11}}{\hyperref[\detokenize{appendix/guideline/customRule:id2}]{\sphinxcrossref{はじめに}}}
\begin{itemize}
\item {} 
\sphinxAtStartPar
\phantomsection\label{\detokenize{appendix/guideline/customRule:id12}}{\hyperref[\detokenize{appendix/guideline/customRule:id3}]{\sphinxcrossref{アクションを追加する際のルール}}}

\end{itemize}

\item {} 
\sphinxAtStartPar
\phantomsection\label{\detokenize{appendix/guideline/customRule:id13}}{\hyperref[\detokenize{appendix/guideline/customRule:id4}]{\sphinxcrossref{アクションの対象は単対象 (必須)}}}

\item {} 
\sphinxAtStartPar
\phantomsection\label{\detokenize{appendix/guideline/customRule:id14}}{\hyperref[\detokenize{appendix/guideline/customRule:id5}]{\sphinxcrossref{キーカードなしの通常\sphinxhyphen{}クイックアクションを作らない (必須)}}}

\item {} 
\sphinxAtStartPar
\phantomsection\label{\detokenize{appendix/guideline/customRule:id15}}{\hyperref[\detokenize{appendix/guideline/customRule:id6}]{\sphinxcrossref{誘発アクションは即時 (推奨)}}}

\item {} 
\sphinxAtStartPar
\phantomsection\label{\detokenize{appendix/guideline/customRule:id16}}{\hyperref[\detokenize{appendix/guideline/customRule:id7}]{\sphinxcrossref{先行1ターン目に確実に勝利できる効果を作らない (必須)}}}

\item {} 
\sphinxAtStartPar
\phantomsection\label{\detokenize{appendix/guideline/customRule:id17}}{\hyperref[\detokenize{appendix/guideline/customRule:id8}]{\sphinxcrossref{誘発\textendash{}メインは同時に2つ以上誘発しないように制限 (必須)}}}
\begin{itemize}
\item {} 
\sphinxAtStartPar
\phantomsection\label{\detokenize{appendix/guideline/customRule:id18}}{\hyperref[\detokenize{appendix/guideline/customRule:id9}]{\sphinxcrossref{キャラクターを追加する際のルール}}}

\end{itemize}

\item {} 
\sphinxAtStartPar
\phantomsection\label{\detokenize{appendix/guideline/customRule:id19}}{\hyperref[\detokenize{appendix/guideline/customRule:id10}]{\sphinxcrossref{場のキャラクターは一目でわかるように定義 (必須)}}}

\end{itemize}
\end{sphinxShadowBox}


\subsubsection{はじめに}
\label{\detokenize{appendix/guideline/customRule:id2}}
\sphinxAtStartPar
本ガイドラインは、BlackPoker のアクションやルールを設計・記述する際の共通基準を示すものです。
各章を参照しながら、項目ごとに「内容」と「理由」を理解したうえで適用してください。

\sphinxAtStartPar
各項目には守るべきレベルを記載しています。
\sphinxhyphen{} \sphinxstylestrong{必須}: すべてのアクションで必ず守る必要があります。
\sphinxhyphen{} \sphinxstylestrong{推奨}: 守ることを強く推奨しますが、必須ではありません。


\paragraph{アクションを追加する際のルール}
\label{\detokenize{appendix/guideline/customRule:id3}}

\subsubsection{アクションの対象は単対象 (必須)}
\label{\detokenize{appendix/guideline/customRule:id4}}
\sphinxAtStartPar
\sphinxstylestrong{内容}
\begin{quote}

\sphinxAtStartPar
アクションの対象は必ず一つになるように設計する。
\end{quote}

\sphinxAtStartPar
\sphinxstylestrong{理由}
\begin{quote}

\sphinxAtStartPar
対象が複数存在すると同時解決や誘発判定の例外処理が膨大になり、ルール全体が煩雑化するからです。
\end{quote}


\subsubsection{キーカードなしの通常\sphinxhyphen{}クイックアクションを作らない (必須)}
\label{\detokenize{appendix/guideline/customRule:id5}}
\sphinxAtStartPar
\sphinxstylestrong{内容}
\begin{quote}

\sphinxAtStartPar
ステージに残る通常\sphinxhyphen{}クイックアクションは、必ず明示的なキーカードを指定します。
\end{quote}

\sphinxAtStartPar
\sphinxstylestrong{理由}
\begin{quote}

\sphinxAtStartPar
キーカードがないと、ステージ上で解決する順番が不明瞭になりルール運用が不安定になるためです。
\end{quote}


\subsubsection{誘発アクションは即時 (推奨)}
\label{\detokenize{appendix/guideline/customRule:id6}}
\sphinxAtStartPar
\sphinxstylestrong{内容}
\begin{quote}

\sphinxAtStartPar
トリガーが「誘発」のアクションはスピードを「即時」、タイミングを「クイック」にする。
\end{quote}

\sphinxAtStartPar
\sphinxstylestrong{理由}
\begin{quote}

\sphinxAtStartPar
誘発アクションは、通常のアクションと異なり、特定の条件が満たされたときに即座に発動することが期待されます。
そのため、即時に発動することで、プレイヤーの記憶領域を節約し、ゲームの進行がスムーズになります。
すでに公式のフォーマットで「ブロック」など即時でない誘発アクションがありますが、それは
ターンの大きな流れを作るために総設計しています。新たにオリジナルで誘発アクションを追加する場合は、即時\sphinxhyphen{}クイックにすることを推奨します。
\end{quote}


\subsubsection{先行1ターン目に確実に勝利できる効果を作らない (必須)}
\label{\detokenize{appendix/guideline/customRule:id7}}
\sphinxAtStartPar
\sphinxstylestrong{内容}
\begin{quote}

\sphinxAtStartPar
先攻プレイヤーが1ターン目に確実な勝利や致命的アドバンテージを得ることを目的とするアクションは禁止する。
\end{quote}

\sphinxAtStartPar
\sphinxstylestrong{理由}
\begin{quote}

\sphinxAtStartPar
ゲームの展開が一方的となり、プレイヤー間の駆け引きや戦略性が失われ、楽しさが損なわれるためです。
\end{quote}


\subsubsection{誘発\textendash{}メインは同時に2つ以上誘発しないように制限 (必須)}
\label{\detokenize{appendix/guideline/customRule:id8}}
\sphinxAtStartPar
\sphinxstylestrong{内容}
\begin{quote}

\sphinxAtStartPar
「誘発\textendash{}通常\textendash{}メイン」のアクションは、一度の解決で二重に誘発しないよう条件を設定します。
\end{quote}

\sphinxAtStartPar
\sphinxstylestrong{理由}
\begin{quote}

\sphinxAtStartPar
誘発\textendash{}メインは、同時に2つ以上の誘発が起こると、解決順序が不明瞭になり、ゲームの整合性が損なわれるためです。
公式のフォーマットでは「ブロック」「ダメージ判定」などターンの流れを作るアクションが 誘発\textendash{}メインで定義されています。
これらのアクションは、必ずステージの一番下になる前提で設計されており、同時に2つ以上の誘発が起こることを想定していません。
\end{quote}


\paragraph{キャラクターを追加する際のルール}
\label{\detokenize{appendix/guideline/customRule:id9}}

\subsubsection{場のキャラクターは一目でわかるように定義 (必須)}
\label{\detokenize{appendix/guideline/customRule:id10}}
\sphinxAtStartPar
\sphinxstylestrong{内容}
\begin{quote}

\sphinxAtStartPar
場に出るキャラクターは、スート、数字、枚数、裏表などを考慮して一意に定義します。
\end{quote}

\sphinxAtStartPar
\sphinxstylestrong{理由}
\begin{quote}

\sphinxAtStartPar
同じスート・数字や組み合わせでも能力や役割が異なる場合があるため、キャラクターを明確に特定できないと誤操作や解釈の混乱を招くためです。
\end{quote}

\sphinxstepscope


\chapter{変更履歴}
\label{\detokenize{revision-history/revision-history:id1}}\label{\detokenize{revision-history/revision-history::doc}}
\sphinxstepscope


\section{第8版の主な変更点}
\label{\detokenize{revision-history/8th:id1}}\label{\detokenize{revision-history/8th::doc}}
\sphinxAtStartPar
次はBlackPoker第8版における主要な変更点です。
8.1の変更点は (8.1) として記載しています。


\subsection{アクションの定義の変更}
\label{\detokenize{revision-history/8th:id2}}\begin{itemize}
\item {} 
\sphinxAtStartPar
B・Jのコスト、効果を見直し

\item {} 
\sphinxAtStartPar
クローズのコスト、効果、キーカードを見直し

\end{itemize}


\subsection{能力の解釈を変更}
\label{\detokenize{revision-history/8th:id3}}
\sphinxAtStartPar
第7版では、能力は「常在型能力」と「誘発型能力」の2種類が定義されていました。

\sphinxAtStartPar
第8版では、能力の解釈を整理し、アクションに「起動条件」と「誘発条件」を追加することで、能力は常在型能力のみを指すものとしました。

\sphinxAtStartPar
この変更により、エクストラリストに大きな影響が生じたため、第8.0版ではエクストラは一旦非表示としています。 なお、コアルール内において能力の影響範囲についての記載を追加しています。

\sphinxAtStartPar
第8.1版では、エクストラをフレーム化して切札を復活させました。(8.1)


\subsection{フォグ領域追加}
\label{\detokenize{revision-history/8th:id4}}\begin{itemize}
\item {} 
\sphinxAtStartPar
新たに「フォグ領域」を追加しました。
この領域は「アップ」「ダウン」などターン終了時まで持続する効果を配置するためのもので、効果が視認しやすくなります。

\item {} 
\sphinxAtStartPar
領域追加に伴い、新たにフォグリストも導入しました。

\end{itemize}


\subsection{用語追加}
\label{\detokenize{revision-history/8th:id5}}
\sphinxAtStartPar
ルールの分かりやすさ向上を目的として、以下の用語を新たに導入しました。
\begin{description}
\sphinxlineitem{【キャラクター関連】}\begin{itemize}
\item {} 
\sphinxAtStartPar
ラベル

\item {} 
\sphinxAtStartPar
サイズ

\end{itemize}

\sphinxlineitem{【アクション関連】}\begin{itemize}
\item {} 
\sphinxAtStartPar
リクエスト

\item {} 
\sphinxAtStartPar
トリガー

\item {} 
\sphinxAtStartPar
スピード

\item {} 
\sphinxAtStartPar
起動条件

\item {} 
\sphinxAtStartPar
誘発条件

\end{itemize}

\sphinxlineitem{【コンポーネント関連】 (8.1)}\begin{itemize}
\item {} 
\sphinxAtStartPar
コンポーネントキーカード

\item {} 
\sphinxAtStartPar
ユニットカード

\item {} 
\sphinxAtStartPar
キャラクターキーカード

\item {} 
\sphinxAtStartPar
フォグキーカード

\item {} 
\sphinxAtStartPar
切札キーカード

\end{itemize}

\end{description}


\subsection{オプション削除とフレームへの統合}
\label{\detokenize{revision-history/8th:id6}}
\sphinxAtStartPar
第7版では対戦レギュレーションにオプションが存在していましたが、
第8版ではオプションを削除し、各フレームに統合することで全体の遊び方の整理を図りました。

\sphinxAtStartPar
また、新たな構築要素として「ストラテジー」を導入しています。
結果として、フレームは以下の4種類のみとし、その他は廃止しました (8.1でエクストラを追加)。
\begin{itemize}
\item {} 
\sphinxAtStartPar
エントリー20

\item {} 
\sphinxAtStartPar
パック

\item {} 
\sphinxAtStartPar
レアパック

\item {} 
\sphinxAtStartPar
ストラテジー

\item {} 
\sphinxAtStartPar
エクストラ (8.1)

\end{itemize}

\sphinxAtStartPar
さらに、ハーフ系が廃止されたため、マッチに関する説明も削除しています。


\subsection{エクストラをフレーム化 (8.1)}
\label{\detokenize{revision-history/8th:id7}}
\sphinxAtStartPar
第8.0版で削除したエクストラ関連ルールを、第8.1版ではフレーム化し、切札を復活させました。
ゲーム開始の流れがこれまでと異なり、よりスピーディーな展開を楽しめます。

\sphinxAtStartPar
8.1版では「立ち上がる兵」の切札のみ復活し、ほかの切札は順次追加予定です。


\subsection{曖昧なルールの明確化}
\label{\detokenize{revision-history/8th:id8}}
\sphinxAtStartPar
以下の点について、ルールの曖昧な部分を明確化しました。
\begin{itemize}
\item {} 
\sphinxAtStartPar
能力の影響範囲をコアルールに明記

\item {} 
\sphinxAtStartPar
誘発チェックのフローを詳細化

\item {} 
\sphinxAtStartPar
コアフロー内で能力を確認するタイミングを追加

\item {} 
\sphinxAtStartPar
共通ルールおよびコアルールにおいて、リクエストの解決が完了するまでリクエストがステージ上に残ることを明記

\item {} 
\sphinxAtStartPar
フォグを配置する際のカードの向きを記載

\item {} 
\sphinxAtStartPar
ステージ上にあるキーカードの配置方法を記載

\item {} 
\sphinxAtStartPar
コンポーネントインスタンスの概念を導入。

\item {} 
\sphinxAtStartPar
コンポーネントインスタンスをユニットに改名。 (8.1)

\item {} 
\sphinxAtStartPar
共通ルールに記載されていた引き分けの記述を修正し、引き分けとならないよう変更

\item {} 
\sphinxAtStartPar
ゲームの始め方をフレーム単位で明確に記載できるよう修正 (8.1)

\item {} 
\sphinxAtStartPar
パック、切札などフレームを含めた全てのコンポーネントを説明する章を追加 (8.1)

\item {} 
\sphinxAtStartPar
デッキの構築条件を明確化 (8.1)

\item {} 
\sphinxAtStartPar
キーカードという単語はアクションだけのものになるよう修正 (8.1)

\end{itemize}


\subsection{ルール外変更点}
\label{\detokenize{revision-history/8th:id9}}
\sphinxAtStartPar
ルール本体に直接関係しない部分について、以下の変更を実施しました。
\begin{itemize}
\item {} 
\sphinxAtStartPar
コアルールの章を後ろに配置し、最初から読めばルール全体が理解できる章構成に変更

\item {} 
\sphinxAtStartPar
新たにフレームの章を設け、フォーマットや章内にアクションリストなどを含む構成に変更

\item {} 
\sphinxAtStartPar
HTML版のアクションリストで、フレーム依存のアクションも表示されるよう修正

\item {} 
\sphinxAtStartPar
各アクション、キャラクター、フォグの導入バージョンなど、細かな情報を記載

\item {} 
\sphinxAtStartPar
PDF版に索引を追加

\item {} 
\sphinxAtStartPar
PDF版の表紙を、新ロゴをあしらった新デザインに変更

\item {} 
\sphinxAtStartPar
アクションなどのyaml定義に対してjsonスキーマを作成

\item {} 
\sphinxAtStartPar
変更履歴を公式ルール内に記載

\item {} 
\sphinxAtStartPar
中綴じ印刷用のPDFを追加

\item {} 
\sphinxAtStartPar
対戦レギュレーションを第2章に移動し、ルール全体の流れを整理 (8.1)

\item {} 
\sphinxAtStartPar
切札リストにイラストを追加 (8.1)

\end{itemize}

\sphinxstepscope


\section{第7版の主な変更点}
\label{\detokenize{revision-history/7th:id1}}\label{\detokenize{revision-history/7th::doc}}
\sphinxAtStartPar
次はBlackPoker第7版における主要な変更点です。


\subsection{リバースのクイック化}
\label{\detokenize{revision-history/7th:id2}}\begin{itemize}
\item {} 
\sphinxAtStartPar
リバースがクイックタイミングで使用可能になりました。

\item {} 
\sphinxAtStartPar
リバースはスタンダードではなく、プロ向けに提供されるようになりました。

\end{itemize}


\subsection{ゲーム中のデッキ名称をライフに変更}
\label{\detokenize{revision-history/7th:id3}}
\sphinxAtStartPar
対戦レギュレーションで示されるデッキと、実際のゲーム中に使用するデッキでは意味が異なります。

\sphinxAtStartPar
その違いを明確にするため、ゲーム中のデッキ名称を「ライフ」に変更しました。


\subsection{対戦レギュレーションの定義方法変更}
\label{\detokenize{revision-history/7th:id4}}
\sphinxAtStartPar
対戦レギュレーションは、フレーム、オプション、フォーマットの3要素の決定により定まるようになりました。

\sphinxAtStartPar
この修正により、対戦レギュレーションの種類が大幅に増加し、決定事項が少なくてもルールが確定できるようになりました。


\subsection{コアフローをわかりやすく整理}
\label{\detokenize{revision-history/7th:id5}}
\sphinxAtStartPar
従来、複雑で分かりにくかったコアフローのフローチャートを整理し、より分かりやすいものに改訂しました。


\subsection{B・Jのクイック化}
\label{\detokenize{revision-history/7th:bj}}
\sphinxAtStartPar
B・Jがクイックタイミングで使用できるようになりました。


\subsection{引き分け廃止}
\label{\detokenize{revision-history/7th:id6}}
\sphinxAtStartPar
勝敗判定がターンプレイヤーから行う方式に変更されたため、引き分けが廃止されました。


\subsection{ルールの記載を微修正}
\label{\detokenize{revision-history/7th:id7}}\begin{itemize}
\item {} 
\sphinxAtStartPar
索引を追加しました。

\item {} 
\sphinxAtStartPar
PDF版において、大きすぎる図のサイズを縮小しました。

\item {} 
\sphinxAtStartPar
ストーリー部分を追加しました。

\item {} 
\sphinxAtStartPar
エクストラの対応版数に関する記載を削除しました。

\end{itemize}



\renewcommand{\indexname}{索引}
\printindex
\end{document}